% =====================================================================
%  H. Brøchner
%  Bidrag til Opfattelsen af Philosophiens historiske Udvikling.
%  Kjøbenhavn: P. G. Philipsens Forlag, 1869.   294 pp.
%  Printer: Bianco Lunos Bogtrykkeri ved F. S. Muhle.
%
%  DIPLOMATIC TRANSCRIPTION (Danish), from the Royal Danish Library scan
%  `bidrag-philosophiens-1869.pdf` (313 PDF pages, 300 ppi RGB JPEG, one
%  image per leaf, NO 1-bit mask).  KB shelfmark DA 1.-2.S 14 8°.
%
%  This header is the publication of record for the edition: the page map,
%  the errata, the emphasis and normalisation conventions, and the log of
%  printer's defects.  The per-book harness (pagemap.py, ocr.sh, check.py,
%  joints.py, splice.py, spacing.py, RECON.md, RESUME-NOTES.md) is
%  deliberately untracked, so an outside reader checking this edition
%  against the scan gets this header and note.md and nothing else.
% =====================================================================
%
%  ---------------------------------------------------------------
%  1. PAGE MAP  (verified by eye, recon pass, 4 September 2026)
%  ---------------------------------------------------------------
%  THE MAP IS NOT UNIFORM.  A leaf was scanned twice, so the offset steps
%  by two in the middle of the book:
%
%      printed   1--209   =   PDF page N +  9      (PDF  10--218)
%      printed 210--294   =   PDF page N + 11      (PDF 221--305)
%
%  PDF 219 and PDF 220 are a SECOND EXPOSURE of printed 208 and 209,
%  which are already at PDF 217 and PDF 218.  The duplication is not a
%  folio coincidence: PDF 217 and 219 carry the identical first line
%  ("bliver den dog, skjøndt et Ideelt i Subjectet, som en An-"), and
%  PDF 218 and 220 likewise.  Two exposures of the same setting.
%  Anyone re-deriving the offset from a folio read near PDF 219--220 will
%  corrupt everything from p. 210 to the end of the book by two pages.
%
%  Verified by reading the printed folio off the image at printed pp. 1,
%  2, 44, 137, 164, 207, 208, 209, 210, 211, 288, 289, 294, and at PDF 219
%  and 220, and corroborated by an ABBYY folio sweep of every leaf from
%  PDF 10 to PDF 305, which contradicts the map nowhere.  Six leaves carry
%  a folio the text layer garbles but the image reads plainly: PDF 53
%  (p. 44), 146 (p. 137), 173 (p. 164), 176 (p. 167), 302 (p. 291).
%  Printed p. 1 (PDF 10) is unnumbered by design.
%
%  Front matter: PDF 1 KB digitisation notice; PDF 2 front board;
%  PDF 3 pastedown with the KB label, barcode and the oval BIBLIOTHECA
%  REGIA HAFNIENSIS stamp; PDF 4--7 blank leaves; PDF 8 title page;
%  PDF 9 title verso, carrying the PRINTER'S IMPRINT (not blank).
%  THERE IS NO HALF-TITLE AND NO CONTENTS PAGE.
%
%  Back matter: printed 289--294 = PDF 300--305, the two-column Register;
%  PDF 306 the unnumbered "Trykfeil og Rettelser" leaf (handled by
%  pagemap.py as internal printed p. 295); PDF 307--313 blank binder's
%  leaves and the rear board.
%
%  ---------------------------------------------------------------
%  2. WHAT IS TRANSCRIBED, AND HOW
%  ---------------------------------------------------------------
%  Diplomatic.  19th-century orthography exactly as printed: Philosophie,
%  Christelige, Videnskaben, høiere, aa, Charakteer, capitalised nouns.
%  Printer's errors are transcribed AS PRINTED, logged in a % comment at
%  the site, and never silently corrected -- including the nine the author
%  himself lists on the Trykfeil leaf (§5).  Balance counts therefore
%  drift off zero on purpose; the running total is in RESUME-NOTES.md.
%
%  TYPEFACE.  Unlike its two Brøchner siblings in this repo (det-religioese
%  and svar-nielsen, both Fraktur), THIS BOOK IS SET IN ANTIQUA -- upright
%  roman throughout, no long s, no Fraktur letterforms anywhere, including
%  in the German (p. 255 "physische Anfangsgründe der Naturwissenschaft"
%  is in the same roman as the Danish, checked at 600 dpi).  No small caps
%  anywhere.  Sans-serif occurs only on the title page and in the Trykfeil
%  heading.
%    BOLD antiqua is used for display heads (levels L1, L2 and the run-in
%  heads such as p. 97 "a) Opfattelsen af Vorden.") but NOT for L3 -- see
%  §3, where the recon pass's claim that L3 is bold is corrected.
%    Bold does not otherwise occur inside a paragraph, WITH ONE ODDITY
%  worth recording as a fact about the printing house: a single lower-case
%  `i' set from a markedly heavier fount turns up mid-paragraph at p. 42
%  ("virker i Stoffet"), p. 89 and p. 95, found independently by two
%  batches.  Its stem and dot are visibly thicker than the neighbouring
%  `i' on the same line.  Three occurrences of the same sort is a WRONG
%  SORT -- the compositor's `i' box held stray heavy sorts -- not an
%  emphasis, and it is transcribed as an ordinary `i' with a % note at
%  each site.
%
%  QUOTATION MARKS.  The book has NO low-high „…“ pair at all.  Its
%  quotation mark is the GUILLEMET, used for Danish as well as foreign
%  matter -- i.e. it follows the sense, not the fount.  The normal pair
%  points INWARD, «…».  But the compositor is inconsistent: roughly 13 of
%  some 87 quotations OPEN with » and about 8 CLOSE with «, verified on
%  the image at 1200 dpi at pp. 23 and 78.  This is transcribed as printed
%  and NOT regularised, so check.py's guillemet balance is expected to be
%  nonzero.  No nested or doubled pairs occur.
%
%  EMPHASIS.  Letterspacing (Sperrsatz) is the book's only in-paragraph
%  emphatic device, and it is used more heavily here than in either
%  sibling -- on p. 209 twelve separate phrases are spaced, about a third
%  of the page.  All of it is rendered \emph{}.  It occurs inside
%  footnotes, inside the displayed schemata, and in the display heads.
%  Proper names are letterspaced, often, and never italicised.
%
%  ITALIC IS RESERVED FOR GREEK, with one exception: the algebraic
%  variables a, b, c, x on p. 5 are italic on the page, and are set
%  \textit{} here as printed.  That is the compositor italicising a
%  variable, a third use of the slant distinct from both the Greek fount
%  and the emphasis device.  Latin, German and French phrases are set
%  in the SAME upright roman as the Danish -- verified on the image at
%  pp. 70, 138, 147, 157, 160, 180, 253, 255 (tabula rasa, cogito ergo
%  sum, monas monadum, immaculata conceptio, res cogitans, physische
%  Anfangsgründe, Système de la nature).  Foreign matter therefore takes
%  \emph{} when it is letterspaced and nothing when it is not; it never
%  takes \textit{}.  This is the opposite of the practice in
%  det-religioese, where the Latin is antiqua against Fraktur.
%
%  GREEK.  There are an estimated 400--600 words of polytonic Greek, on 77
%  body pages plus the Register and the errata leaf, densest at pp. 94--136
%  (Aristotle, Plotinos) and 186--225.  NEITHER OCR WITNESS CONTAINS A
%  SINGLE GREEK CHARACTER in 313 pages -- both transliterate it into Latin
%  garbage ("Novs" for νοῦς, "Koauog" for κόσμος, "(pvGtg" for φύσις) --
%  so every Greek word in this edition was read off the page image at
%  1200 dpi and nowhere else.
%    The scan does resolve polytonic pointing at body size.  Calibrated on
%    words whose breathing is positively of the other kind: smooth on
%    οὐσία (pp. 97, 121) and ἐπέκεινα (p. 129) against rough on ὕλη
%    (p. 35) and ἡδονή (p. 193); grave on ἡδὺς against acute on ἡδονή
%    (p. 199); circumflex on Νοῦς (p. 88).  The distinction is positive
%    evidence, not inference.
%    THE ONE THING THE SCAN DOES NOT RESOLVE is the breathing-plus-acute
%    compound over a low vowel: in ἄπειρον (pp. 52, 81) the mark is a
%    single blob at 1200 dpi and ἄ cannot be separated from ἅ.  Where a
%    breathing or accent could not be resolved the letter is set
%    unaccented and a % comment records the undecided reading at the site.
%    Nothing is supplied from knowledge of Greek.
%    The Greek fount is a CURSIVE (italic).  In such a fount ϑ and ϱ are
%    the ordinary glyph variants of θ and ρ, not distinct sorts, so plain
%    θ and ρ are typed and the font renders the variant.  Greek is typed
%    as Unicode and wrapped \gk{...} -- see the macro's own comment in the
%    preamble for why a named macro and not a bare \textit{}.  Where a
%    Greek word is ALSO letterspaced it takes \emph{\gk{...}}.
%
%  RULES.  Two kinds, and they must not be confused.  The DIVISIONAL rule
%  is a short CENTRED hairline (~1.5--1.8 cm) with white space above and
%  below, standing above almost every display head; it is part of the text
%  and is transcribed \begin{center}\rule{2cm}{0.4pt}\end{center} with a
%  % note at the site.  The FOOTNOTE SEPARATOR is a short LEFT-FLUSH
%  hairline above the note area; it is LaTeX's \footnoterule job and is
%  NOT transcribed.  Sheet signatures at the foot (3*, 4*, ... 19*) are
%  not transcribed.  Running head: the folio alone, no words.
%
%  DASHES.  Measured at 1200 dpi, the parenthetical dashes run 128--156 px
%  and the numeral-range dashes 112--136 px against an em of 167 px; the
%  distributions OVERLAP and the samples are small, so per
%  TRANSCRIPTION-PLAYBOOK §2 rule 1 the scan does not establish a distinct
%  en-dash sort.  Every dash is transcribed ---.  An en would be 83 px and
%  nothing that short was measured.
%
%  THE ɔ: MARK (id est) occurs; confirmed on the image at p. 69.  Typed as
%  U+0254 followed by a colon and mapped in the preamble.  Neither witness
%  reads it (it comes out "o:", "9:", "d:"), so every page is read for it.
%  There is no § anywhere in the book; it cites by page ("S. 175 ff.").
%
%  MARGINALIA.  This copy carries MODERN PENCIL, and nothing that is not
%  printed type is transcribed.  Confirmed by eye: a two-line cursive note
%  in the left margin of p. 88 (which the ABBYY layer splices INTO the
%  text as "'pv^hj £"), margin ticks and brackets on p. 209, and -- inside
%  the type area -- `fra' struck through with `for' written above at
%  p. 18 l. 2, and `Forbrydelse' struck through at p. 110 l. 29.  Both of
%  those are a reader applying the printed errata; the PRINTED reading is
%  transcribed in each case.  Three mechanical sweeps for pencil failed
%  (it is only ~33 grey levels below the paper against ~106 for the ink),
%  so the true number of marked places is unknown; twelve further sites
%  are visible as OCR damage at pp. 26, 29, 31, 45, 67, 69, 88, 111, 177,
%  207, 286, 293.
%
%  ---------------------------------------------------------------
%  3. STRUCTURE  (read off the images; there is no contents page)
%  ---------------------------------------------------------------
%  The heading hierarchy, with the printed page of each head:
%
%    L0  large letterspaced roman, centred, alone on its line
%        Indledning. (1)
%        Hovedepocherne i Philosophiens Historie. (16--17, two lines)
%        Charakteristik af Philosophiens historiske Hovedformer. (216)
%        Slutning. (284)      Register. (289)
%    L1  large BOLD centred head with a short centred rule under it
%        Den græske Philosophies Udvikling. (31)
%    L2  BOLD centred head, preceded by a short centred rule
%        the four theses  I. (50)  II. (79)  III. (137)  IV. (170)
%        Første Hovedepoche: (219)  Anden Hovedepoche: *) (226)
%        Tredie Hovedepoche: (269)
%    L3  LETTERSPACED ROMAN at body size, set as a justified PARAGRAPH --
%        not bold and not centred.  (The recon pass called this level
%        "bold centred head"; that was WRONG, and four batch agents
%        refuted it independently by measuring letter gaps at 500--600
%        dpi: 10--18 px in these heads against 2--5 px in the certainly
%        bold L2 head on p. 79.  Checked at pp. 50, 54, 65, 80, 94.
%        The error came of judging by size and prominence rather than by
%        measurement, and it matters: a bold head must not be set \emph{},
%        and a letterspaced one must not be set \textbf{}.)
%        A. Den realistiske Philosophie. (231)
%        B. Den idealistiske Philosophie. (236)
%        a. Den nyere idealistiske Philosophie som Enhed. (237)
%        b. Den idealistiske Philosophies særegne Former. (240)
%        1. Den græske Atomlære. (186)  2. Den Aristippiske og
%        Epikureiske Lystlære. (192)  3. Sokrates. (201)
%        1./2./3. under each thesis: 50, 54, 65; 80, 94, 127; 137, 145, 150
%    L4  letterspaced (NOT bold) centred head, using GREEK letters
%        α. Den dogmatiske Idealisme. (240)
%        β. Den kritiske Idealisme. (250)
%    L5  BOLD centred head, under β
%        I.  Den kritiske Idealisme i snævrere Betydning. (252)
%        II. Den speculative Idealisme. (259)
%
%  α. and β. are GREEK letters used as outline levels, confirmed on the
%  image at pp. 240 and 250, and they sit BELOW a./b. in the hierarchy.
%  The I./II. at pp. 252/259 is a DIFFERENT sequence from the I.--IV. at
%  pp. 50/79/137/170.  Neither is renumbered or harmonised here.
%
%  Three DISPLAYED SCHEMATA are set apart from the running text and are
%  the hardest typesetting in the book -- pp. 29--30 (bold numbered heads
%  with a wide swash brace over two columns), p. 163 (a genealogical tree
%  built out of rules), and p. 180 (multi-line typographic braces with
%  connector rules).  What each contains, and what LaTeX construction was
%  chosen for it, is recorded in a % comment at the site.
%
%  There are no block quotations: no quotation is set smaller or indented,
%  and the long ones run into the paragraph inside guillemets.
%
%  ---------------------------------------------------------------
%  4. FOOTNOTES
%  ---------------------------------------------------------------
%  Marks are *) **) ***) then the DAGGER †) and the DOUBLE DAGGER ††),
%  superscripted, RESTARTING at *) on every printed page; hence the \ifcase
%  in the preamble plus \setcounter{footnote}{0} at the page marker of every
%  note-bearing page, and NOT footmisc[perpage], which resets on the typeset
%  page.  A heading can carry a mark (p. 226).  About 72 pages bear notes.
%    FIVE levels are used, not three.  The recon pass reported three and
%  found ***) on pp. 76 and 115 only.  p. 115 in fact carries FIVE notes --
%  *) **) ***) †) ††) -- read on the image at 1200 dpi in the text and again
%  at the head of each note by the pp. 113--126 batch.  The preamble's
%  \ifcase was extended accordingly; before that it would have set "4)" and
%  "5)" silently, which is the kind of defect that compiles cleanly and is
%  invisible until someone reads the PDF.
%    RUN-OVER NOTES EXIST.  The recon pass concluded they did not: all nine
%  pages whose mark counts looked asymmetric in its mechanical sweep (69,
%  77, 90, 104, 133, 135, 261, 272, 279) were read on the image and every
%  one was self-contained.  That conclusion was WRONG, and the error was
%  one of method -- the sweep could only test candidates an OCR asymmetry
%  made visible.  Two counter-examples were found in the first three
%  batches: p. 18's note fills two-thirds of p. 19, and p. 42's continues
%  onto p. 43; neither page was among the nine.  Every run-over note found
%  in this edition is set in full inside a single \footnote{} at its call,
%  with a % comment at the page break of the continuing page recording
%  that the note continues there in the print.  Further instances are
%  logged in RESUME-NOTES.md as they are found.
%
%  ---------------------------------------------------------------
%  5. TRYKFEIL OG RETTELSER  -- the author's own errata, NOT applied
%  ---------------------------------------------------------------
%  The unnumbered leaf at PDF 306 lists nine corrections.  Per the repo's
%  diplomatic stance the body keeps the PRINTED reading in every one of
%  them, with a % comment at the site, and the leaf is set separately at
%  the end of this edition.  Eight of the nine printed readings were
%  confirmed on the image during the recon pass; the ninth (p. 113 l. 35,
%  "Bostemmelse") was located during the batch that covers p. 113.
%
%  ---------------------------------------------------------------
%  6. TYPESETTING NOTE -- THIS FILE IS BUILT WITH LuaLaTeX
%  ---------------------------------------------------------------
%  The repo's default is pdfLaTeX with libertinus + textalpha, and that is
%  right for a book with a dozen Greek words.  It is not right for this
%  one.  Under pdfLaTeX the Greek is served by LGR, for which Libertinus
%  has no Type1 face, so it falls back to CB Greek -- a Computer Modern
%  Greek that clashes with the Libertinus text and sets upright, where
%  this book's Greek fount is a cursive.  Under LuaLaTeX the Greek comes
%  from Libertinus Serif itself, in italic, with the polytonic repertoire
%  and the iota subscript this book needs (ἐνεργείᾳ, p. 99 and the errata
%  leaf).  The Makefile already carries per-file LuaLaTeX overrides for
%  the Høffding texts; this file needs one too.
% =====================================================================

\documentclass[12pt,a4paper]{book}
\usepackage{fontspec}
\usepackage{libertinus}  % engine-aware: loads libertinus-otf under LuaLaTeX
\usepackage{graphicx}    % for \reflectbox (the old Danish »ɔ:« id-est mark)
\usepackage{newunicodechar}
\newunicodechar{ɔ}{\reflectbox{c}}
\usepackage[danish]{babel}
\usepackage[protrusion=true,expansion=true]{microtype}
\usepackage{setspace}
\usepackage[margin=1.2in]{geometry}
\usepackage{enumitem}
\usepackage{amsmath}     % \left\{ for the p. 180 brace schema
\usepackage{fancyhdr}
\setlength{\headheight}{28pt}
\addtolength{\topmargin}{-13.5pt}
\usepackage{hyperref}

\hypersetup{
  pdftitle={Bidrag til Opfattelsen af Philosophiens historiske Udvikling},
  pdfauthor={H. Brøchner},
  hidelinks
}

% The book prints *) **) ***) then the DAGGER and the DOUBLE DAGGER,
% restarting on every printed page.  The recon pass reported three levels
% only; p. 115 uses five, read on the image at 1200 dpi in the text and
% again at the head of each note.  It is the only five-note page found so
% far; p. 76 uses three.
\renewcommand{\thefootnote}{%
  \ifcase\value{footnote}\or *)\or **)\or ***)\or \dag)\or \dag\dag)\else
  \arabic{footnote})\fi}

% A short centred divisional rule, as printed above almost every display
% head.  NOT the footnote separator, which is left-flush and is LaTeX's job.
\newcommand{\divrule}{\begin{center}\rule{2cm}{0.4pt}\end{center}}

% The book's Greek fount is a CURSIVE, and Libertinus Serif's Greek is
% upright, so every Greek word is wrapped \gk{...} to get the slant the
% page shows.  One macro rather than a bare \textit{} for three reasons:
% it is greppable (the count is in check.py), it records the FOUNT as a
% fact about the print rather than as an emphasis, and if a future
% engine or font change makes the Greek cursive already, the whole book
% is corrected by editing this one line to {#1}.
% Where a Greek word is ALSO letterspaced in the print, write
% \emph{\gk{...}}; the slant and the emphasis are then independent.
\newcommand{\gk}[1]{\textit{#1}}

\pagestyle{fancy}
\fancyhf{}
\fancyhead[LE,RO]{\thepage}
\fancyhead[RE]{\textit{Bidrag til Opfattelsen af Philosophiens historiske Udvikling}}
\fancyhead[LO]{\textit{\leftmark}}

\setstretch{1.3}

\title{\textbf{Bidrag til Opfattelsen}\\[2pt]
  \large af\\[6pt]
  \textbf{\Large Philosophiens historiske Udvikling}}
\author{Dr.\ H.\ Brøchner\\[4pt]
  \small Professor i Philosophie}
\date{Kjøbenhavn: P.\ G.\ Philipsens Forlag, 1869}

\begin{document}
\maketitle
\thispagestyle{empty}
\newpage

% =====================================================================
%  BODY.  printed pp. 1--288.  One marker per batch; splice.py matches a
%  fragment .parts/pp<FIRST>-<LAST>.texfrag to the marker carrying the
%  same page pair, so these lines must not be edited.
%
%  The display heads listed in §3 above fall INSIDE batches, not at their
%  boundaries, and are written by the batch agent that meets them.
% =====================================================================

\chapter*{Bidrag til Opfattelsen af Philosophiens historiske Udvikling}
\markboth{Philosophiens historiske Udvikling}{Philosophiens historiske Udvikling}

% --- p. 1 ---
% HEAD, p. 1 (L0): "Indledning." --- large LETTERSPACED roman, centred, alone on
% its line.  NO divisional rule above it: this is the opening page of the book.
% The page carries no folio by design; the "1" at the foot is the sheet
% signature and is not transcribed.  The text opens with an oversized two-line
% initial capital D ("Den Opfattelse ..."), set here as an ordinary D.
% p. 1 carries no footnote and no Greek.
\begin{center}
{\Large\emph{Indledning.}}
\end{center}

Den Opfattelse af Philosophiens Historie, der tilstræbes i dette
Skrift, er en saadan, som til sin Forudsætning har den strænge
historisk-kritiske Behandling af Stoffet, der i de sidste Decen\-%
nier er bleven saagodtsom alle Philosophiens historiske Hoved\-%
former tildeel. Den omfattende og nøiagtige Kundskab til det
ved denne Behandling sigtede og opklarede Stof er Udgangs\-%
punktet og Grundlaget for den Behandling af Gjenstanden, der
ikke mere som den tidligere arbeider sig ad Inductionens Vei
opad fra det empiriske Enkelte til et sammenfattende Alment;
men fra den gjennem de forudgaaende Enkeltundersøgelser mu\-%
liggjorte Erkjendelse af det Almene stiger deducerende ned
til det Enkelte og giver dettes Opfattelse paa een Gang en
rigere Fylde og en strængere Nødvendighed ved at lade det
% p. 1, l. 15: a small stray dot stands between "Bestemthed" and "ved"
% (tesseract reads it as a full stop; ABBYY ignores it).  It is not printed
% type -- an ink speck or modern pencil -- and is not transcribed.
sees i dets Begrundethed og Bestemthed ved hint. Den Op\-%
fattelse af Philosophiens Historie, jeg forsøger, forholder sig
saaledes overalt \emph{til Totaliteten}, erkjendende dens Princip
og dens Værens Lov, og seer den historiske Udfoldelse som
bestemt ved denne Lov, der ikke er en blot empirisk funden,
% p. 1: "Værens Lov" (l. 18) and "denne Lov" (l. 19) are NOT letterspaced;
% only "tankegyldig" (l. 20) is.  Measured letter-gaps on the image: the
% spaced words run 12--18 px at 600 dpi, the unspaced ones 3--8 px.
men en \emph{tankegyldig} Lov, og som, idet den fremtræder i det
Concrete, aabenbarer sig gjennem en \emph{Flerhed} af væsensforbundne
Love. Mit Øiemed er altsaa at give et Bidrag til en \emph{Philosophi\-%
ens Histories Philosophie}, til en begribende Erkjendelse
af Philosophiens Historie i dens principbestemte Vorden, og
derved at forberede en \emph{Fremstilling af Philosophiens
Historie}, hvilende paa denne Erkjendelse.

% --- p. 2 ---
Idet dette Skrifts Opgave er saaledes bestemt, forudsætter
det hos Læseren et almindeligt Kjendskab til Stoffet i Philo\-%
sophiens Historie, dog ikke i et større Omfang end det, der
kan erhverves ved et paalideligt, ikke altfor kort Grundrids.

De historiske Partier, i hvilke Udviklingsloven i dens con\-%
crete Fremtræden vil blive paaviist, ville blive valgte, deels
med Hensyn til deres indre Betydenhed og deres Prægnants,
deels med Hensyn til, at der, under Behandlingen af dem,
gjennem Paaviisningen af Loven, og væsentlig ved den, kan
gives nye Bidrag til Opfattelsen af det historiske Stof i dets
Objectivitet, hvor dette i de foreliggende Bearbeidelser af Phi\-%
losophiens Historie er utilstrækkeligt og utilfredsstillende oplyst.

% DIVISIONAL RULE, p. 2, between "... utilfredsstillende oplyst." and "Naar
% Udviklingsloven ...": short, CENTRED, with white space above and below,
% about a third of the measure.  It stands alone -- no display head follows it
% -- and marks the division between the prefatory paragraphs and the argument.
% Not the footnote separator (which is left-flush); this page bears no note.
\divrule

Naar Udviklingsloven, der skal søges, bliver Loven, \emph{ikke}
for Philosophiens Udformning \emph{som System}, for dens indre,
% p. 2, l. 15--16: the emphasis resumes at "dens" on l. 16; the single letter
% "i" that ends l. 15 cannot itself show letterspacing and is left plain.
ligesom tidløse Organisationsproces, men for dens Udvikling i
\emph{dens Vorden som historisk Phænomen}, som Led i
Menneskeaandens totale Historie, --- en Udvikling, der dog
viser sig i væsentlig Sammenhæng med hiin allerede derigjen\-%
nem, at den historisk fremskridende Udvikling i Bestemmelsen
af Principet viser sig afgjørende for den systematiske Organi\-%
sations Fuldstændighed, for den gjennemførte Articularisation,
saa gaaer Opgaven ind under en mere omfattende. \emph{Philoso\-%
phiens Historie} er kun en Deel af den \emph{universelle Hi\-%
storie, Philosophiens Histories Philosophie}, den be\-%
gribende Erkjenden af Philosophien som vordende, maa derfor
gaae ind under \emph{Historiens Philosophie}, under den begri\-%
bende Erkjenden af det Historiskes Vorden, og Muligheden af
hiin maa være betinget ved Muligheden af denne som det
Særlige ved det Almene. Vilde man forestille sig den Mulig\-%
hed, der synes at forudsættes ved de theologiske Protester mod
«det Historiskes» Opfattelse gjennem Begrebet, at en enkelt
Side af den menneskelige Natur, f.\ Ex. Tanken, ifølge sin
% --- p. 3 ---
særegne Bestemthed, vel kunde være Nødvendigheden og Lov\-%
mæssigheden underkastet, uden at Væsenet i sin Heelhed, i
alle sine forskjellige Udviklingsformer, var det, saa vilde denne
Forestilling strax vise sig som uholdbar saasnart der reflecte\-%
redes paa Væsenets Enhed, der skulde omfatte baade det Lov\-%
bestemte og det Lovløse. I Væsenets Almindelighed, i dets
Totalitetsbestemmelse, er Loven at søge; den er netop kun
Udtrykket for den væsensbegrundede og derfor constante Virk\-%
somheds Form. Derfor maa Muligheden af en begribende Er\-%
kjenden af et hvilketsomhelst Særegent indenfor den historiske
Udvikling overhovedet være betinget ved Muligheden af den
begribende Erkjenden af det omfattende Almene. Dette ude\-%
lukker ikke, at det ene Særegne, ifølge sin Eiendommelighed,
kan ligge den begribende Erkjenden nærmere end et andet;
men Forskjellen bliver her kun quantitativ: det afgjørende
Qualitative er begrundet i Væsensheelheden. Alt, hvad der er
Yttring af Menneskevæsenet, er lovbestemt, og enhver Lov
maa som Lov kunne blive Gjenstand for Viden.

Muligheden af en Historiens Philosophie, som det bliver
denne omfattende Videnskabs Opgave at godtgjøre, maa den
da godtgjøre ved at gjøre det antydede Synspunkt gjældende:
at alt det, der i Menneskehistorien udvikler sig, som Yttring
af et Væsentligt og Substantielt, maa udvikle sig som bestemt
% p. 3, l. 24 and l. 32: the ɔ: id-est mark, confirmed on the image at 900 dpi
% (neither witness reads it: ABBYY gives "o:", tesseract "9;").
ved sit Væsen, ɔ: med Nødvendighed og Lovbundethed, og at
denne Nødvendighed og Lovmæssighed, som Aandens Bestemt\-%
hed, maa være tilgængelig for den erkjendende Aand. Hvad
der selvvirksomt virker, maa virke i Overensstemmelse med
hvad det er, det er altsaa i sin Virken bundet ved sin Væsens
Bestemthed. Virksomhed er kun den virkende Væren, nærmere:
den efter sin Bestemthed, sin Natur, virkende Væren. En lovløs
Virken vilde forudsætte en bestemmelsesløs Væren, men en ube\-%
stemt Væren ɔ: en Væren, der kun var en rastløs Vorden, af hvilken
Intet vordede, vilde være en selvmodsigende Væren, der opløste sig
i Intet. Denne almindelige Bestemmelse maa ogsaa gjælde for
det Menneskelige, hvor \emph{Frihedens} Grundbestemmelse, som
% --- p. 4 ---
% EMPHASIS, p. 4, l. 1: the compositor letterspaces only the FIRST element of
% each compound -- "V æ s e n s udtryk", "V æ s e n s bestemmelse" -- and sets
% the second element tight.  Measured at 600 dpi: the gaps inside "Væsens" run
% 15--20 px, those inside "udtryk" and "bestemmelse" 3--10 px.  The same hand
% shows again at p. 10 l. 9 ("T a n k e udviklings") and p. 14 l. 26
% ("V æ s e n s erkjendelse"), so it is a habit, not a slip.
\emph{Væsens}udtryk, \emph{Væsens}bestemmelse, som \emph{væsentlig Fri\-%
hed}, har \emph{Nødvendighedens} og \emph{Lovmæssighedens} Natur
i sig, og hvor den individuelle Selvbestemmen, forsaavidt den
faaer \emph{Vilkaarlighedens} Form, viser sig som dragen ind
under og som bestemt ved \emph{den universelle Magts Nød\-%
vendighed}, ved hvilken den i Vilkaarligheden virksomme
Frihed har sin Mulighed. Frihedens Lov er Erkjendelsens
Gjenstand ligesaafuldt som enhver anden Lov; den bliver, som
% p. 4, l. 9: the quotation OPENS and CLOSES with » -- »afspeilet» -- where the
% book's normal pair points inward, «…».  Read at 600 dpi; transcribed as
% printed and not regularised (transcription.tex header §2).
Menneskevæsenets Lov, ikke blot »afspeilet» i den menneske\-%
lige Viden, men begreben af denne.

Naar nu med Muligheden af en Historiens Philosophie
som det Almene Muligheden af det Særegne: af en Philoso\-%
phiens Histories Philosophie, er given, saa bliver, idet Mulig\-%
heden skal realiseres, først den historiske Udviklings \emph{almin\-%
delige} Charakteer at bestemme, og dernæst den \emph{Individua\-%
lisation} at paavise, som dette Almindelige modtager i An\-%
vendelsen paa \emph{Philosophien} som \emph{eiendommelig} Retning
af Aandsudviklingen.

De \emph{almindelige} Bestemmelser for det Historiskes lov\-%
mæssige Udvikling blive at tage som \emph{Laanesætninger fra
Historiens Philosophie}. De fremgaae af Begrebet om
\emph{Subjectet for den historiske Udvikling}, hvis Bestem\-%
melse atter viser tilbage til den hele Opfattelse af Tilværelsen.
Grundmodsætningen i Bestemmelsen fremgaaer af Modsætnin\-%
gen mellem en \emph{materialistisk} og en \emph{ideel} Opfattelse af
% p. 4, l. 26--27: "atomistisk" is letterspaced but "Individuelle," on the next
% line is NOT (measured 15 px against 4 px); the emphasis stops at the line end.
Tilværelsen. Hvor Virkeligheden sættes i det \emph{atomistisk}
Individuelle, hvor paa abstract nominalistisk Maade det Almene
opfattes som et blot Subjectivt, der bliver Menneskehistorien
\emph{Summen af Individernes Historie}, den faaer ikke eet
Subject, men en Utallighed af Subjecter, af hvis tilfældig\-%
nødvendige Vexelvirkning Historiens Naturbegivenheder blive
Resultanterne, de individuelle Aarsagers summerede Virkninger,
medens Historiens Love kun blive simple Consequentser af den
individuelle Menneskeaands Love. Denne Opfattelse forudsæt\-%
ter fra Begyndelsen af, at Individet kan bestemmes \emph{kun} som
% --- p. 5 ---
Individ, at Lovene for dets Virksomhed som \emph{menneskeligt},
ɔ: som selvbevidst og fornuftigt, aandeligt Individ, kunne op\-%
fattes uden at Forholdet til det \emph{Almene} og dettes Virksom\-%
hed bliver \emph{primitiv} Factor. Der gaaes ud fra den Opfattelse,
% ALGEBRA, p. 5, ll. 5--7.  Read off the image at 1200 dpi.  What is printed,
% exactly, is:  "at Individernes Enhed er = Individernes Sum, at, forat bruge
% en mathematisk Betegnelse, Enheden af a og b og c er = a + b + c, ikke
% = a + b + c + x, saaledes at dette x udtrykker det Høieste og
% Væsentligste."  (ABBYY's "a-\-b-f-c + a?" is its rendering of "a + b + c + x".)
% Set as TEXT, not math mode, per BATCH-AGENT §9.  The "=" is the book's
% slanted double-stroke equals sign; the "+" is an ordinary plus.
% The algebraic letters a, b, c and x are set in ITALIC on the page (a
% cursive roman, clearly slanted at 1200 dpi), and they are the ONLY
% non-Greek italic in the book.  Set with \textit{} here, as printed: this
% is the compositor's ordinary practice of italicising a variable, not the
% emphasis device (which is letterspacing) and not the Greek fount, so it
% is a third and separate use of the slant.  The header's rule "italic is
% reserved for Greek" is amended accordingly; BATCH-AGENT §3's ban on
% \textit{} for foreign matter is unaffected and still stands.
at Individernes Enhed er = Individernes Sum, at, forat bruge
en mathematisk Betegnelse, Enheden af \textit{a} og \textit{b} og \textit{c} er = \textit{a}
+ \textit{b} + \textit{c}, ikke = \textit{a} + \textit{b} + \textit{c} + \textit{x}, saaledes at dette \textit{x} ud\-%
trykker det Høieste og Væsentligste. Mod denne Opfattelse,
hvis Grundlags Uigjennemførlighed Metaphysiken paaviser,
idet den paaviser Grundmodsigelsen i Atombegrebet, ligesom
Erkjendelseslæren paaviser den, idet den viser Erkjendelsens
Tilintetgjørelse som dens Consequents, stiller sig den anden,
der i Historien ikke seer Individernes Aggregats Historie, men
\emph{en Individerne i sin Almeenhed sammenfattende
ideel Enheds Historie}. Denne Enhed bestemmer den nær\-%
mest, idet den bruger en Betegnelse, der overføres fra Natur\-%
sphæren, som \emph{Menneskeslægten}. Slægt forholder sig til
Individer som Udtryk for \emph{Væsensenheden}, der, som \emph{con\-%
cret} Enhed, sammenfatter disse, for det \emph{Almene}, der frem\-%
bringer og vedligeholder sig gjennem de vexlende Individer,
som den sætter. I Slægtbegrebet er det \emph{Almenes Realitet
sat}, men som \emph{Proces}, der har det \emph{Individuelle} til sit
Element. Slægten er ikke værende paa umiddelbar Maade;
dens Realitet er en sig frembringende. Slægten, det Almene,
er kun, idet den ved sin Energie resulterer af Individerne.
Dette er, i den yderste Abstraction, Forholdet paa Naturge\-%
betet i snævrere Forstand. Idet Slægtbegrebet overføres paa
det \emph{Menneskelige} og der tales om \emph{Menneskeslægten} i
dens Forhold til \emph{Menneskeindividerne}, nuanceres Begrebet
derved, at Individerne her ere \emph{selvbevidste} og staae i et
Bevidsthedsforhold til Slægten, hvilket kun bliver muligt der\-%
ved, at de \emph{i sig} have et Forhold \emph{af det Almene til det
Almene}. Menneskeindividet har ikke blot Slægtens almene
Bestemmelse som værende i sig, men har den i en Forsigvæ\-%
ren, og er bevidst virkende i Kraft af denne. I Slægtens og
% --- p. 6 ---
Individets Forhold, som et Forhold mellem det Almene og det
Individuelle, finder der her paa den ene af Siderne en \emph{For\-%
dobbling} Sted, idet det Individuelle bevidst bliver et \emph{Indi\-%
viduelt-Alment}. I denne Bestemmelse ligger \emph{Frihedens
Mulighed}, og med Friheden er Muligheden given til en in\-%
dividuel Handlen, der kan synes at emancipere sig fra det
Almene, at gjøre sig gjældende som en Magt ligeover\-%
for dette, og Udgangspunktet for en \emph{individuel psycho\-%
logisk og ethisk Historie}. Men fra denne sidste sondrer
sig \emph{Slægthistorien} som \emph{det Almenes Historie}, der som
saadan maa have en \emph{metaphysisk} Charakteer, bestemmes
ved \emph{metaphysiske Nødvendighedslove}, der beherske og\-%
saa hint tilsyneladende Emanciperede. Selve Menneskehedens
historiske Frihedsudvikling, dens Fremskriden i Frihed, maa
sees ogsaa fra den metaphysiske Nødvendigheds, den almene
Væsensnødvendigheds Synspunkt. Men idet Historiens \emph{almene}
Subject ikke er et i sin Almindelighed som færdigt værende,
men et af det Individuelle resulterende, gjennem det Indivi\-%
duelle sig frembringende Alment, saa maae de \emph{metaphysi\-%
ske} Nødvendighedsbestemmelser bestandig sees \emph{som resulte\-%
rende af den individuelle psychologisk-ethiske Fri\-%
hedsudvikling}, og bestandig vise hen paa Vexelforholdet
til denne. \emph{Slægthistoriens} Charakteer bliver da nærmere
at bestemme som \emph{metaphysisk med Reflex af det Psy\-%
chologiske og Ethiske}.

Fastholdes denne Anskuelse, for hvilken Historiens Subject
er det almene Subject: Menneskeslægten (Menneskeheden), der
% p. 6, l. 28: a small stray mark stands immediately before "de enkelte"
% (tesseract turns it into an opening guillemet, "«le enkelte"; ABBYY into a
% hyphen).  It is not printed type and is not transcribed.
omfatter og sammenslutter de enkelte virksomme Factorer, som
bestemmende deres Virken (der ogsaa faaer Formen af en Vir\-%
ken for det Almene), og som gjennemvirkende deres Bevidst\-%
hed, i hvilken det Almene forkynder sig, saa ligger heri, alle\-%
rede inden der sees hen til Spørgsmaalet om dette concrete
Almenes Forhold til \emph{det Absolute}, til det i sidste Instants
Bestemmende, en dobbelt Consequents for den historiske Ud\-%
vikling i Almindelighed.

% --- p. 7 ---
Denne bliver nemlig først at see \emph{som underkastet den
i det Almenes Magt begrundede Lovmæssighed}, der,
ifølge sit Udspring, faaer den \emph{metaphysiske Nødvendig\-%
hedscharakteer}. Og den bliver dernæst, idet den ved det
Almene bestemte Udvikling er \emph{selve det Almenes teleo\-%
logiske Selvvirkeliggjørelse gjennem den individu\-%
elle Virken}, at see som \emph{en Udvikling, i hvilken paa
ethvert Punkt et Vexelforhold maa gjøre sig gjæl\-%
dende mellem Individernes bevidste Stræben og det
gjennem dem og gjennem selve deres Bevidstheds\-%
liv for dem ubevidst virkende overgribende Almene},
der gjør sig gjældende som et objectivt Fornuftigt, som uni\-%
versel aandelig Magt paa en Naturmagts Basis.

I det antydede Forhold mellem \emph{det for Individet Be\-%
vidste} og \emph{det, der naaer ud over dets Bevidsthed} og
unddrager sig denne, gjør et for enhver Yttring af individuelt
Aandsliv Væsentligt sig gjældende, og det viser tilbage til Til\-%
værelsens dybeste Grundforhold. Psychologien viser, at alt
bevidst Liv ikke blot i en Tidsudvikling er fremgaaet af et
ubevidst, er anlagt igjennem det, men at det ogsaa, efterat
det er kommet til Virkelighed, i hvert Øieblik hviler paa det
ubevidste som paa sin Grund; at dets Indhold efter en nød\-%
vendig Lov gaaer tilbage i det Ubevidste for atter at hæve sig
op af dette til Bevidsthedens Lys; at Bevidstheds- og Bevidst\-%
løshedsperioder afvexle i Individet; at selve den enkelte Be\-%
vidsthedsact er Summationen af de uendeligt smaa i det Ube\-%
vidste blivende Elementer, og at Individet endelig, gjennem
det absolut og relativ ubevidste Liv staaer i dybere Sammen\-%
hæng med det universelle Livs Strømninger. Men denne over\-%
alt stedfindende Sammenhæng, der godtgjør Enheden mellem
Natursiden og Aandssiden i det Menneskelige, viser tilbage til,
\emph{at alt individuelt Bevidsthedsliv og Frihedsliv}, med
eet Ord: \emph{alt individuelt Fornuftliv}, er begrundet i et
omfattende og constituerende \emph{Universelt}, der i sin \emph{ideel\-%
reale} Bestemthed paa een Gang er \emph{Naturtilværelsens}
% --- p. 8 ---
teleologisk virkende Princip og Princip for den \emph{aandelige}
Tilværelse. Dette \emph{ideelt-reale} Aands-Princip, der teleolo\-%
gisk virkeliggjør det individuelt \emph{Aandelige} paa det \emph{ube\-%
vidste Naturliges} Grundlag, begrunder i denne Virken ud
fra det \emph{objectivt Fornuftige} den i Bevidsthedens og Fri\-%
hedens Form fremtrædende \emph{individuelt-subjective} Fornuft,
og bestemmer Jeget, der tænker, ved \emph{det, der tænker i
det}, ved det almene Fornuftprincips tankebegrundende Magt.
Og som det \emph{universelle} Aandsprincip sammenslutter det den
individuelle Bevidsthed og Frihed under sin overgribende Nød\-%
vendighed.

Idet de individuelle Frihedsacter sees saa omfattede og
bestemte ved det Almenes Nødvendighed, sees denne det Al\-%
menes Nødvendighed som Individernes \emph{Skjæbne}. Dette er
det første Udtryk for den; men det indeslutter et andet og
høiere i sig. Derved, at Skjæbnen er de \emph{selvbevidste, tæn\-%
kende} Væseners Skjæbne, \emph{Skjæbne for Aander}, under\-%
gaaer den en Transfiguration. Skjæbnen, den som Nødven\-%
dighed raadende Magt, forklares for Aanden, og udformer sig
for Bevidstheden til \emph{Naturens og Aandens fornuftbe\-%
stemte Lov}, til \emph{fornuftig styrende subjectiv Magt}.
Men denne Skjæbnens Transfiguration bliver dog bestandig
% GREEK, p. 8, l. 24: Ἀνάγκη.  Neither witness contains a Greek character here
% (ABBYY "’Avdyxri", tesseract "”Avoyxn"); read off the image at 1200 dpi.
% The breathing is set as a separate PRE-POSED mark before the capital Α, as
% the fount requires; at 1200 dpi it is a comma-shaped mark whose bowl lies to
% the right and whose tail runs down to the left, i.e. the SMOOTH breathing
% (rejected reading: the rough ῾, whose tail runs the other way).  The acute
% over the second α is unambiguous.  The κ is the cursive kappa (which is why
% both witnesses read it as x).  The word is not letterspaced.
\emph{kun relativ}, og Skjæbnen taber aldrig fuldstændig sin Cha\-%
rakteer og sit Præg som Skjæbne, som den \gk{Ἀνάγκη}, for hvil\-%
ken Grækerne ogsaa lod Guderne bøie sig; det Lys, til hvilket
den hæves op, har i hvert Øieblik Mørket bagved sig og straa\-%
ler ud fra dette.

Skjæbnen bevarer relativt sin Charakteer af Skjæbne fordi
selve den Bevidsthedsact, der forvandler den, og hæver den
fra den blinde Nødvendigheds Mørke til Bevidsthedens og Fri\-%
hedens Lys, som baaren og styret af et for Individet ubevidst
Virkende, ligesom enhver anden Bevidsthedsact, beholder et
ubevidst Nødvendighedsmoment i sig, gjennem hvilket Skjæbnen
hævder sig som Skjæbne, om end i en høiere Skikkelse som det
Bevidstes og Fries Skjæbne. Og at den gjør dette, viser sig ogsaa
% --- p. 9 ---
derigjennem, at det, der virker som Skjæbne, det de indivi\-%
duelle Aanders Virken styrende Universelle, er et Saadant, i
Forhold til hvilket den individuelle Bevidstheds Erkjenden skri\-%
der frem fra det Abstractere til det Concretere, men kun saa\-%
ledes skrider frem, at den vigende Grændse aldrig forsvinder,
men, skjøndt bevægelig Grændse, dog altid bliver en Bevidst\-%
hedsgrændse i Forhold til det Universelle.

Naar nu de \emph{almindelige} Bestemmelser for \emph{det Hi\-%
storiskes} Udvikling skulle anvendes paa \emph{Philosophiens
Historie}, maa den i Philosophien fremtrædende Aandsvirk\-%
somheds \emph{specifiske} Charakteer gjøre sig gjældende som \emph{in\-%
dividualiserende} hine. Det Eiendommelige ved den er, at
Vexelforholdet mellem Frihed og Nødvendighed, mellem det
Bevidste og det Ubevidste, gjennem Modsætningernes Potent\-%
sation, bliver mere dialektisk. Idet den philosophiske Aands\-%
virksomhed er den over sig selv reflecterende Tankes Virksom\-%
hed, træder Bevidstheds- og derigjennem Frihedsmomentet in\-%
tensivere frem end i de andre Former af den historiske Ud\-%
vikling; men med det Samme gjør, netop ifølge den Grund\-%
bestemmelse ved Tanken, at den har det Almene til sit Ele\-%
ment og gjennem det staaer i Sammenhæng med Nødvendig\-%
heden, dennes Bestemmelse sig ligesaa intensivt gjældende, og
med Jegets bevidste Tanke, der er rettet mod det Universelle,
% p. 9, l. 24: ABBYY reads a comma after "inderlig" and tesseract a full stop.
% At 600 dpi the mark is a light brown smudge standing clear of the word and
% above the baseline -- a stain, not a sort.  No punctuation is transcribed.
gaaer det i Jeget Tænkende, for Jeget Ubevidste, i inderlig
uadskillelig Enhed. Modsætningerne staae ikke i det Forhold
til hinanden, at den enes Stigen er den andens Aftagen, men
de stige med hinanden. Bevidsthedens Intensitet stiger med
den beherskende ubevidste Virkens Intensitet, Friheden bliver,
i sin prægnanteste Form, metaphysisk Nødvendighed.

Den samme Grundanskuelse af Udviklingen, der blev gjort
gjældende for det Historiske i Almindelighed, bliver altsaa og\-%
saa her at gjøre gjældende, men den maa i Gjennemførelsen
blive saaledes individuelt udformet, at Processens Bevægelse
bliver dybere, Momenterne skarpere udprægede, deres Spæn\-%
ding intensivere.

% --- p. 10 ---
% p. 10, ll. 2--8: ABBYY's layer indents these lines, but on the image they run
% the FULL measure like the rest of the paragraph.  They are not a displayed
% quotation; they are simply seven consecutive letterspaced lines.
Opgaven med Hensyn til Philosophiens Historie bliver da
\emph{at paavise den i den menneskelige Tankes Væsen
begrundede Lovmæssighed, der med sin Nødven\-%
dighed behersker Gangen i den historiske Tanke\-%
udvikling,} og \emph{at see denne Udvikling som det over\-%
gribende Almenes teleologiske, i Fornuftnødvendig\-%
hedens Form fremtrædende Selvvirkeliggjørelse
gjennem den bevidste individuelle Virken.}

% p. 10, l. 9: "T a n k e udviklings" -- the first element of the compound is
% letterspaced and "udviklings" is not (measured 12 px against 4--6 px at
% 600 dpi).  Same hand as p. 4 l. 1 and p. 14 l. 26.
Dette \emph{Almene}, den historiske \emph{Tanke}udviklings almene
Subject, bliver, efter det Foregaaende, først at bestemme som
\emph{Menneskeslægtens} (Menneskenaturens) \emph{almene Væsen i
Tankens Bestemthed}, og dette almene Væsens Nødven\-%
dighed bliver det, der behersker Tankens Tidsudvikling, dens
phænomenale Fremtræden. Men Begrebet om Menneskeslæg\-%
tens Væsen, om det i Forhold til de menneskelige Individer
\emph{concrete, subjective Almene}, viste tilbage til et Høiere
og \emph{Universellere}, til det \emph{absolute Princip}, der paa een
Gang er Princip for Menneskevæsenet, for Aanden, og Princip
for Naturen, den menneskelige Tankes Basis og Gjenstand.
Og idet Menneskeslægtens almene tænkende Væsen sees som
\emph{begrundet} i det absolute Princip, bliver den menneskelige
Tankes historiske Udvikling i sidste Instants seet som bestemt
\emph{ved det absolute ideelt-reale Princips Væsen}, og
som behersket af dettes Nødvendighed, der, som et \emph{subjectivt}
Bestemts Nødvendighed gjør sig gjældende \emph{gjennem Sub\-%
jecters frie Tankeudvikling}, beherskende denne idet den
bærer og hæver den.

Gjennem den individuelle Tankes Væsen føres saaledes
den menneskelige Tankeudviklings Lov tilbage til \emph{det Abso\-%
lute} som sin Kilde, og nærmere, til \emph{det Absolute, seet
under Synspunktet af Videns (Sandhedens) Idee}. I
denne Aabenbaringsform af den absolute Idee maa altsaa \emph{Ud\-%
viklingens nødvendige Lov} definitivt søges, og idet Ud\-%
viklingen, hvis Lovmæssighed skal erkjendes, ikke er den
logiske Ideeudfoldelse, men den ved Ideens metaphysiske Væ\-%
% --- p. 11 ---
sen betingede og bestemte \emph{Tidsudfoldelse} under den \emph{reale}
Tilværelses Vilkaar, --- en Udfoldelse, af hvilken metaphysiske
Væsensbestemmelser resultere for Bevidstheden som den indi\-%
viduelle Tankes substantielle Bestemtheder, --- saa bliver Videns
logiske Idee at see som med Nødvendighed bestemmende
\emph{gjennem psychologiske Forhold og Udviklingsbetin\-%
gelser}.

I \emph{Videns Idee} opfattes det \emph{Absolute} som det Totali\-%
teten af sine reale Bestemmelser (Momenter) til subjectiv
Enhed Idealiserende. Den \emph{absolute Viden}, som Videns
Idee betegner, er \emph{den sig i Alt vidende Viden}, saaledes,
at dette: \emph{Alt}, udtrykker en \emph{Realitetens} Bestemmelse, der
ikke blot som \emph{logisk} Moment er i \emph{Identitet} med den
\emph{ideelle} Vidende, men \emph{som væsensgyldigt i sin Reali\-%
tet}, under Identiteten er \emph{modsat} det Ideelle. \emph{Realite\-%
tens} Bestemmelse udtrykker en \emph{Andethedens} Bestemmelse,
men \emph{som begrundet i selve det ideelle Princip}. Andet\-%
heden hører saaledes med til dettes \emph{Væsen}: det er \emph{ideelt\-%
realt} Princip, Viden \emph{af sig} som Viden \emph{af sit Andet, Sub\-%
jectivitetsprincip}, idet det er Princip for det \emph{Objective}.
Den vidende Subjectivitets Bestemmelse er det \emph{ideelle Prius}
idet den er det \emph{Omfattende og Overgribende}; men under
denne det Subjectives ideelle Prioritet ere det Subjective og
det Objective \emph{lige nødvendige} og derfor \emph{lige oprindelige}
Væsensbestemmelser. Idet Subjectiviteten sees som \emph{Forkla\-%
ringsgrund} for sig selv og for det Objective, saa sees det
Objective, der forklares ved den, som \emph{dens egen Værens
Betingelse}.

Sees nu denne Videns Idee som bestemmende Princip i
den \emph{menneskelige} Viden, under \emph{Endelighedssondringen}
mellem Ideelt og Reelt, mellem Tanken og Tankens objective
Gjenstand, saa maa det i Videns Idee liggende Grundforhold:
\emph{Viden af sig selv som Viden af Totaliteten af sit
Andet}, gjøre sig gjældende i det efter sin Natur \emph{univer\-%
selle} menneskelige Væsens Tankeudvikling, men det maa der
% --- p. 12 ---
gjøre sig gjældende \emph{under psychologiske Udviklings\-%
betingelser}, under Tidsbetingelser, Endelighedsbetingelser.

Idet disse Betingelser nærmere bestemmes, træder en
\emph{Sondring} frem, der i det Foregaaende, hvor Grundsyns\-%
punkterne for Opfattelsen søgtes, ikke blev draget frem forat
% PRINTER'S ERROR, p. 12, l. 6: "skuide" is printed where "skulde" is meant.
% Transcribed as printed.  Both witnesses read "i", and the image confirms it:
% at 1200 dpi the fourth letter carries a detached mark over a short stem --
% the grey minimum in the break reaches 154 against a paper value of ~200 --
% whereas the "l" of "Forholdet" three words earlier on the same line is an
% unbroken stem whose grey minimum never rises above 50.  The vertical break is
% shallower than that of the certain "i" of "ikke" (182--200) at the head of
% the same line, so the sort is inked heavily, but the break is positive.
% This error is NOT on the author's Trykfeil leaf.
ikke Forholdet skuide compliceres: Sondringen mellem \emph{den}
Videns Form, der bliver at betegne som den \emph{philosophiske},
og Videns (Tænkningens) \emph{øvrige} Former. Det i Videns Idee
givne Forhold mellem Videns Subject og Videns Gjenstand gjør
sig, idet denne Idee virker som det under Tidsbetingelser og
i en Endelighedssondring sig udviklende Væsens Lov, gjældende
gjennem den nødvendighedsbestemte psychologiske Udvikling,
gjennem hvilken Individet udvikler sig fra umiddelbar (sand\-%
sende) til begribende Subjectivitet; den gjør sig gjældende i
Generationernes efter Progressionens Lov voxende Maximums\-%
punkters Gjennemløben af en i kategorisk Bestemmelse sti\-%
gende Scala; den gjør sig gjældende i det Forhold, i hvilket
den efter Gjenstandens objective Sondringer sig articulerende
reale Videns Indhold staaer til den ideelle omfattende Princip\-%
erkjendelse, og den gjør sig gjældende idet en saadan Erkjen\-%
delse teleologisk udvikler sig. Som \emph{teleologisk anlagt i}
og \emph{udviklet af} Menneskeaandens forudgaaende Tankeudvikling
er den \emph{philosophiske} Erkjendelse \emph{betinget} ved denne, ved
Individernes og Generationernes psychologiske Udvikling, ved
den reale Videns Omfang og Art, og den \emph{fuldstændige} Op\-%
fattelse af dens Love maa reflectere \emph{Forholdet} til disse
\emph{Betingelser}. Men netop idet den philosophiske Erkjendelse
sees som den menneskelige Tankeudviklings \emph{Maal}, afslutter
den sig i sig med en \emph{indre Selvstændighed}, ved hvilken
det bliver muligt, saalænge Bestemmelserne holdes \emph{abstrac\-%
tere}, at see den udfolde sig, ligesom \emph{ud af sig selv}, i
Overensstemmelse med det al menneskelig Viden constitue\-%
rende Princips Væsen.

Vi betragte altsaa her \emph{Philosophiens} historiske Udvik\-%
ling i denne \emph{relative Selvstændighed}, for i simplere
% --- p. 13 ---
Omrids at kunne fastholde Udviklingens Hovedformer, for lettere
at oversee den indre Bevægelse.

Den ovenfor angivne Opgave med Hensyn til Opfattelsen
% p. 13, l. 4: "denne" is printed with a word-space inside it, "den ne" (ABBYY
% reads the two halves as separate words); the compositor's spacing fault, on a
% loosely justified line.  Set here as the single word.
af denne \emph{særegne} menneskelige Tænkningsforms Udvikling
faaer nu, ifølge det i det Foregaaende Udviklede, en bestemtere
Form. Opgaven bliver \emph{først} at paavise den i det Universelles
(Principets) Væsensnødvendighed, i Ideeforholdene, begrundede
Lovmæssighed, der bestemmer den under psychologiske Be\-%
tingelser foregaaende Udvikling af Philosophiens store Epocher
i deres Rækkefølge, og af de Systemer, der fremtræde indenfor
de enkelte Epocher. Og den bliver \emph{dernæst}: at see det al\-%
mene Princips Herredømme aabenbare sig gjennem det i Phi\-%
losophiens teleologiske Udvikling fremtrædende \emph{Vexelforhold
mellem det for Individerne Ubevidste og Bevidste},
et Forhold, i hvilket Udviklingens ledende Princip, hvis meta\-%
physiske Nødvendighed i den teleologiske Udvikling gjennem\-%
virker den individuelle Virken, viser sig i hiin dialektiske
Dobbelthed som Skjæbne og som fornuftig styrende subjectiv
Magt. Under Individernes bevidste Stræben maa, snart som
fremskyndende, snart som tilbageholdende, den universelle Magt
paavises, der virker bag Individernes Bevidsthed, gjennem deres
bevidste Tænkning, og som netop mægtigst aabenbarer sig
hvor Individets productive Tanke er meest energisk. Og i
dette Vexelforhold maa Individernes Bestemmethed ved det
overgribende Universelle, gjennem de Grundtanker, i hvilke
dets sondrede Momenters vexlende Forhold fremtræde, paa
ethvert Punkt, selv der, hvor Bestemmetheden var en Hæmmen,
en Tilbageholden, erkjendes som fremmende for den Tankens
% p. 13, l. 30: the mark in "historisk-metaphysisk" is a short HYPHEN, not a
% dash (measured against the em-dash at the end of l. 31 on the same page);
% the white space on either side of it is the line's loose justification.
totale, teleologiske Udfoldelse, under hvilken Principets Be\-%
stemmelser sees resulterende som historisk-metaphysisk Re\-%
sultat af psychologiske Bevidsthedsacter. ---

Den \emph{første} Opgave: Paavisningen af den i Tankens Væ\-%
sen og Princip begrundede Lovmæssighed, der behersker hele
Udviklingsgangen i Philosophiens Historie, maa, ifølge det
Foregaaende, løses paa den Maade, at Loven, der findes for
% --- p. 14 ---
den menneskelige Bevidstheds (Tankes) Udvikling, viser sig
som reflecterende Ideens (Ideebevægelsens) Lov, at altsaa denne
dobbelte Lovmæssighed sees i Enhed.

Til Grundforholdene, Grundbestemmelserne i Videns Idee
% p. 14, l. 5: ABBYY reads a full stop after the second "i"; at 600 dpi the
% mark is a faint speck, not a sort, and is not transcribed.
svare i den i Endelighedssondring existerende menneskelige
Aand \emph{Selvbevidsthedens} og \emph{Bevidsthedens} Bestem\-%
melser. Selvbevidstheden: Jegets Forsigværen som Jeg, er det
\emph{abstracte} Udtryk for, den abstracte Reflex af, den absolute
concrete Subjectivitets Forsigværen i sin ideelle Enhed; Be\-%
vidstheden, som Bevidsthed om en fra Jeget forskjellig \emph{ydre}
Gjenstand, Reflexen af den i den absolute Videns Begreb lig\-%
gende Bestemmelse af Viden af et Andet som idealiseret
\emph{Væsensmoment} (ikke: Vidensmoment). I \emph{Selvbevidst\-%
hedens Grundforhold til Bevidstheden} ligger for Men\-%
neskeaanden \emph{Nødvendigheden} af en Stræben efter en \emph{uni\-%
versel} Erkjenden, der, som den endelige Tænknings Viden af
sig i sit Andet, den reale Tilværelse, med det Samme er en
Viden af sig og sit Andet gjennem en Viden af det Tanken
og Virkeligheden begrundende ideelt-reale Princip. \emph{Selv\-%
bevidstheden} er Menneskeaandens \emph{Grundbestemmelse},
det, der kategorisk sondrer den fra Livets underordnede Former
og i sig indeslutter dens væsentlige Eiendommeligheder, og
Selvbevidsthedens Grundforhold til Bevidstheden er det, at
den med Nødvendighed stræber henimod at forvandle ethvert
% p. 14, l. 26: "V æ s e n s erkjendelse" -- the first element of the compound
% is letterspaced, the second is not.  Same hand as p. 4 l. 1 and p. 10 l. 9.
\emph{Bevidsthedsobject til Selvbevidsthedsmoment}, stræ\-%
ber efter, gjennem en \emph{Væsens}erkjendelse af Gjenstanden, at
ophæve dens \emph{Fremmedhed}, at fjerne \emph{Modsigelsen} mellem
Selvbevidsthedens indre Universalitet og dens Begrændsethed
ved en Bevidsthed om et Andet, mellem Bevidsthedens sub\-%
jective Form og objective Indhold, mellem Jegets Viden af sig
som gjennemsigtig Enhed og som bestemt ved et i Jeget
Værende, fra Jeget Forskjelligt. Den Erkjendelsens Drift, der
fra første Færd rører sig i Menneskeaanden, kan derfor kun
tilfredsstilles ved en universel Erkjenden, d.\ v.\ s.\ ved en philo\-%
sophisk Erkjenden, bestemt efter dennes ideale Maal, og den
% --- p. 15 ---
% p. 15 opens mid-sentence: it continues the paragraph carried over from
% p. 14 ("... Den" / "philosophiske Erkjendelse ..."), so there must be NO
% paragraph break at this batch seam.  Source lines below follow the printed
% lines one-for-one; every printed line-break hyphen is set \-% .
philosophiske Erkjendelse fremgaaer saaledes med Nødvendighed
af Menneskenaturen, af den Drift, der betegner dennes væ\-%
sentlige Eiendommelighed, og den har sit Maal sat ved det,
der oprindelig ligger i Væsenets Bestemthed. Men som den
efter Begrebet universelle Erkjenden bliver Philosophien først
til naar Begrebet om \emph{en Tilværelsens Enhed i Principet}
kommer frem for Tanken; thi først dette Begreb indeslutter
den bevidste Fordring af en Erkjendelsens \emph{Universalitet}, og
det betinger Muligheden af Aandens Viden af sig gjennem det
reale Andets Totalitet. Idet \emph{Philosophien} saaledes ved sit
Væsen staaer i et \emph{nødvendigt} Forhold til det \emph{Absolutes}
% p. 15 l. 12: "Tanke," is NOT letterspaced -- measured at 300 dpi its
% inter-glyph gaps are 2--3 px against 8--9 px in "Absolutes" on the line
% above; the wide space before the comma is this loose line's justification.
% Only "Absolutes" is spaced, so the emphasis breaks mid-phrase as printed.
Tanke, saa maa, naar den historisk fremtræder i forskjellige
Hovedformer, \emph{enhver} af dem eie dette Begreb om det Abso\-%
lute, om Tilværelsens Enhedsprincip. Og da Principet er en
\emph{concret Modsætningsenhed} og i sin Concretion skal blive
til for den fra det Endelige udgaaende, ved psychologiske Ud\-%
viklingslove bundne Bevidsthed, saa maae de forskjellige Stil\-%
linger, Enhedens Grundmomenter under den fremskridende
Opfattelse indtage til hinanden, \emph{nødvendigviis} betinge
\emph{væsentlig forskjellige Opfattelser af Principet}, der
maae begrunde en Sondring af Epocher i Philosophiens,
den universelle Erkjendelses, historiske Udvikling og be\-%
stemme Epochernes Grundpræg. Epochernes Rækkefølge maa
blive bestemt ved Begrebsforholdet mellem det Absolutes
Momenter, \emph{saaledes som de reflectere sig i den men\-%
neskelige Bevidsthed}. Men i denne reflectere Subjectivi\-%
tetsmomentet og den objective Andetheds Moment i Videns
% p. 15 l. 28: the "og" between "Aand" and "Natur" is NOT spaced (its o--g
% gap measures 2 px against 7--10 px inside "Aand" and "Natur").
Idee sig som \emph{Aand} og \emph{Natur}.

Ud fra disse foreløbige Bestemmelser, der vise hen til
% p. 15 ll. 30--31: "og i" between "Hovedformer" and "lovbestemt" is NOT
% spaced (o--g gap 1 px); "i" is a single glyph and cannot be decided.
Nødvendigheden af en Udvikling gjennem \emph{en Flerhed af
Hovedformer} og i \emph{lovbestemt} Orden, bliver nu den
\emph{ene} Deel af den stillede Opgave først saavidt at løse, at de
historisk givne \emph{Hovedepocher} i Philosophien begribes i
deres principbestemte Lovmæssighed. Vi søge altsaa først
Lovmæssigheden i Rækkefølgen af \emph{Epocherne}.
% DIVISIONAL RULE, printed at the FOOT of p. 15 (short centred hairline,
% ~2 cm, white space above and below), standing over the L0 head that opens
% p. 16.  It is at the foot of 15, not the head of 16 -- p. 16 carries no
% rule.  This site is not in RECON.md's confirmed divisional-rule list.
\divrule

% --- p. 16 ---
% p. 16--17: L0 display head (transcription.tex §3), large roman, centred,
% set on TWO LINES exactly as broken here.  It is NOT letterspaced: the
% inter-glyph gaps measure 3--4 px, the same as the body fount, so
% spacing.py's RUN flag on it is the large-fount artefact of BATCH-AGENT §3
% and the head takes no \emph{}.  Both lines are centred on the measure
% (line 1 spans x341--1266, line 2 x681--918; measure centre x795).
\begin{center}
{\Large Hovedepocherne i Philosophiens\\
Historie.}
\end{center}

% p. 16: the paragraph opens with an enlarged initial "I" of "Ifølge"
% (about 1.5 x body height, sitting on the baseline).  Page-making
% furniture; set as ordinary text.
Ifølge det Foregaaende ville Hovedstadierne i Philosophiens
historiske Udvikling i deres lovmæssige Fremadskriden blive
at see som bestemte ved en \emph{metaphysisk} og en \emph{psycho\-%
logisk} Nødvendighed, der staae i indre Væsenssammenhæng.
Dette vil med andre Ord sige: de i deres Følge begrebsnød\-%
vendige Opfattelser af Forholdet mellem det Absolutes Grund\-%
momenter, der beherske de enkelte Epocher af Philosophiens
Historie, Philosophiens historiske Grundformer, maae paa een
Gang vise sig som Udtrykket for Grundforholdene mellem
Ideens Momenter, der bestemme Ideebevægelsen i Videns Idee,
og for de Bevidsthedens Grundforhold, der bestemmer den
% PENCIL, p. 16: a modern "+|" in the left margin beside this line, a
% two-word cursive note below it in the margin, and a "+|" written above
% the line just after "Deductionen" (which ABBYY renders "Deductionen'"
% and tesseract "Deductionen”").  None of it is printed type; not transcribed.
menneskelige Tænknings med psychologisk Nødvendighed fore\-%
gaaende Udvikling henimod dens Maal, og i Deductionen af
% p. 16 l. 14: "psychologiske" with ch, read on the image; ABBYY has
% "psykologiske" here and is wrong.
Epochernes Grundtanker (Principer) maa den psychologiske
nødvendige Tidsudvikling afspeile hiin Ideens Nødvendighed
og Ideebestemmelserne maae sees som viklende sig ud af hiin
som dens i Begrebet afklarede Udtryk.

Forat en saadan Resulteren, i hvilken det universelle
Princip, der i sin metaphysiske Nødvendighed bestemmer den
individuelle bevidste Tankeudvikling, for Bevidstheden frem\-%
% p. 16 l. 21: a small round ink speck stands between "skal" and "kunne"
% (both witnesses read it as an apostrophe or a guillemet).  Not a sort;
% not transcribed.
gaaer af denne i sin Nødvendighedsform, skal kunne blive
mulig, maa det Psychologiske have en Dobbeltsidighed: det
maa fra den ene Side sees fra det Individuelles Synspunkt
som Individets frie Bevidsthedsliv; paa den anden Side maa
det sees fra det Universelles Synspunkt som det, der i en
% --- p. 17 ---
nødvendig Væsensudfoldelse er baaret og omfattet af hvad der
bestemmer selve Individets Tanke- og Frihedsliv som dets
høiere Lov. Fordi det Psychologiske i sin Individualitet hviler
paa og bestemmes af et Alment, i sin Frihed af et Nødvendigt,
kan det metaphysisk Nødvendige resultere deraf i den Proces,
i hvilken et subjectivt Princip, i Virkeliggjørelsen af den hi\-%
storiske Tidsudvikling, virkeliggjør sig selv efter sin den in\-%
dividuelle Frihedsudvikling reflecterende Nødvendighed, og under
denne Virkeliggjøren, som det i det Individuelle Værende,
klarer sig i Individets Bevidsthed til den menneskelige Tanke\-%
helheds bestemmende, sammenfattende Idee-Princip.

Vi tale saaledes om det \emph{almene Subject}, der realiserer
sig gjennem den \emph{individuelle} Udvikling, og om \emph{dette
Subjects Væsenslov} til Forskjel fra den menneskelige Be\-%
vidstheds \emph{psychologiske Udviklingslov}, men saaledes, at
vi ikke glemme, at det almene Subject kun har en Væren for
os som værende i \emph{vor Bevidsthed}, og at det i denne Væren,
som Principet, i vor Bevidsthed maa fremgaae i metaphysisk
Bestemthed af den \emph{individuelle} Tankeudvikling. Idet vi
tale om Principet i dets Bestemmen af den individuelle, psy\-%
chologiske, Udvikling, tænke vi os det som det, der ifølge
Grundforholdet i Individet mellem det \emph{Psychologiske} og det
\emph{Ontologiske}, i Bevidstheden resulterer af de \emph{enkelte} Be\-%
vidsthedsacter \emph{som den høieste, Bevidsthedens substan\-%
tielle Grund udtrykkende Tanke}. Gjennem Forholdet
til Gjenstanden, gjennem Endelighedssondringen, virke som be\-%
stemmende Principets ontologiske Grundforhold, og de resultere
for Tanken af det psychologiske Forhold, og af de enkelte,
oprindelig med Endelighedsbestemmelser behæftede Tanker.
Historiens Princip realiserer sig i den menneskelige Tanke
saaledes, at denne i Realisationen reflecterer sig i hiin og i
denne Reflecteren optages i det Ontologiske.

Idet vi nu opkaste Spørgsmaalet om den \emph{første} Form,
i hvilken Philosophien maa fremtræde, om den \emph{første phi\-%
losophiske Hovedepoches} Eiendommelighed, saa bestemme
% Sheet signature "2" at the foot of p. 17; not transcribed.
% --- p. 18 ---
\setcounter{footnote}{0}%
% p. 18 l. 1: there is NO comma after "Udgangspunktet" -- read at 1200 dpi.
% What tesseract takes for one is the top of the modern pencil "for"
% written above the following line.
vi, idet vi tage Udgangspunktet fra det \emph{Psychologiske}, fra
% TRYKFEIL (author's own errata leaf, PDF 306, entry 1 of 9): the print
% reads "fra Philosophiens Tilbliven"; the leaf asks for "for".  Per the
% repo's diplomatic stance the PRINTED reading stands.
% PENCIL at the same site: a modern hand has struck "fra" through and
% written "for" above it.  Not transcribed (RECON.md; header §2 MARGINALIA).
den psychologiske Betingelse fra Philosophiens Tilbliven, den
første Forms Væsen gjennem en Reflexion over det Betingende.
Betingelsen for Philosophiens Existents er \emph{Aandens Frihed};
thi uden den, som \emph{indre Frigjørelse}, der indeslutter Aan\-%
dens instinctmæssige Tillid til sin Magt, vilde Tanken ikke
kunne stille sig den \emph{universelle} Opgave: gjennem den En\-%
hed, der behersker Phænomenets Totalitet i dens uendelige,
vexlende Mangfoldighed, at finde sit eget Væsen, at naae det
Forhold til sig selv i sin Gjenstand, der er Frihedens Forhold.
Af denne Grund maa Aandens Frihed være Philosophiens
Forudsætning, men Aandens Frihed bliver først til, hvor dens
Bevidsthed om at være \emph{høiere} end det Naturlige, at være
\emph{herskende} over dette, begynder at dæmre frem. Derfor maa
Tanken i Philosophiens første Hovedepoche udtrykke dette
\emph{det Aandeliges Principat}, men den maa tillige udtrykke
% p. 18 l. 17: measured letter by letter at 300 dpi, "det" (gaps 3, 4) and
% "Form," (gaps 3, 3, 3) are NOT spaced; only "den" (gaps 8, 10) is.
det i \emph{den} Form, der efter Begrebet er den \emph{første} for dette
Forhold. Og den første Form, i hvilken Aanden fremtræder
som fri og herskende, er \emph{Skjønhedens} Form; thi Skjønheden
er den \emph{umiddelbare}, endnu ikke gjennem Reflexion sikkrede
og befæstede Enhed af Naturen og den Naturen besjælende
og beherskende Aand: den er en Enhed, endnu i \emph{Naturens
Form, i Harmoniens Umiddelbarhed}. For den frie
Aand, der i Skjønhedens Bestemthed hviler i den umiddelbare
Harmonie af Tilværelsens Grundformer, maa den høieste, ab\-%
solute Form for \emph{Anskuelsen af Tilværelsens Heelhed}
blive Opfattelsen af den som et \emph{altomfattende Kunst\-%
værk}, Tanken om Universet som den ideebesjælede Kosmos,
Formens og Stoffets harmoniske Enhed. Denne Grundanskuelse
% RUN-OVER FOOTNOTE.  This note begins on p. 18 and CONTINUES on p. 19,
% where it fills the whole page below thirteen lines of text and opens
% WITHOUT a mark.  This contradicts transcription.tex §4 ("NO RUN-OVER NOTE
% EXISTS IN THIS BOOK"), whose nine test pages did not include 18/19 -- see
% the report.  The note is set here as one \footnote{}; the point at which
% the print carries it over to p. 19 is marked inside it.
% The footnote separator on both pages is the short LEFT-FLUSH hairline; not
% transcribed (BATCH-AGENT §6).
behersker den \emph{græske} Philosophie, Philosophiens \emph{første}\footnote{%
Idet vi sætte den \emph{græske} Philosophie som Philosophiens \emph{første}
Hovedform, som den egentlige Philosophies historiske Begyndelse,
udelukke vi fra Philosophien de Former, man pleier at betegne som
\emph{østerlandsk} Philosophie. Vi gjøre det, fordi disse Former, af
% [the note is carried over to p. 19 at this point]
hvilke kun de \emph{indiske} kunne komme i Betragtning, om de end
med den egentlige Philosophie dele en Stræben efter at finde Til\-%
værelsens Princip og ved Grundforholdet til dette Princip at be\-%
stemme Menneskelivets Opgave, dog kun gjøre det som afhængige
af Religionen og uden at løsne sig fra dens Grundform, eller kun i
Former, der, om end ikke identiske med Religionens, dog ikke ere
den \emph{philosophiske} Tankes. Aanden søger Uendeligheden som
Tilværelses- og Livsprincip, men den Uendelighed, den søger, er
ikke Tankens, men \emph{Phantasiens}, Tankens og dens Bestemthed
absorberende indholdsløse Uendelighed, fra hvilken ingen Udfoldelse
kan blive mulig. Og hvor der ved Siden af denne Stræben efter
\emph{Phantasieuendelighedens Hvile}, der søges naaet ad Veie og
ved Midler, der vise Maalets Natur, fremtræder en Reflexion, ikke
blot i Almindelighed over Hovedphænomener i den concrete Til\-%
værelse, men ogsaa særligt over Tankens Former, altsaa en Art af
logisk Theorie, der kommer denne Reflexion ikke frem som bevidst
bestemt ved et Princip, og Gjenstandens Art kan ikke, om end
% p. 19 note l. 18: a thin raised tick stands before "philosophisk" (both
% witnesses read it as an apostrophe).  It is a stray mark, not a sort --
% the book's quotation mark is the guillemet and no pair opens here.
Gjenstanden er Tanken selv, gjøre Erkjendelsen philosophisk, naar de
Former, i hvilke Gjenstanden bringes frem for Bevidstheden, kun ere
% p. 19 note l. 20: "den" is NOT spaced (gaps 2, 2); the run is
% "spredte Erfarings" (gaps 6--7).
den \emph{spredte Erfarings} og den \emph{principløse Reflexions}. --- Vi
lade derfor disse Former ude af Betragtningen, og sætte den græske
% p. 19 note ll. 22--23: the line-break hyphen falls inside "idet-høieste";
% resolved here as the single word "idethøieste", which is what the print
% gives.  Reported for review.
Philosophie, med hvis ældste Former de i deres Begyndelse idet\-%
høieste ere samtidige, som Philosophiens saavel \emph{efter Tiden} som
\emph{efter Begrebet} første Hovedform.}
% --- p. 19 ---
% (the note begun on p. 18 occupies the lower two thirds of this page; it is
% set above, inside the \footnote{}.  Thirteen printed lines of body text
% and a fourteenth follow here.)
historiske Hovedform, og det vil vise sig, at ud fra den dens
hele Udviklingsgang bliver forklarlig. Paa Basis af denne
Grundanskuelse, af dette Grundforhold mellem Momenterne,
maatte der udfolde sig en Videns Form, i hvilken Tanken,
Begrebet, som metaphysisk Begreb, \emph{umiddelbart} gjorde sit
Herredømme over det Naturligt-Materielle gjældende, saaledes,
at \emph{Andethedens} Moment i den \emph{reale} Enkelthed \emph{ikke}
kunde komme afgjørende frem, men det Enkelte, Reale kun
kunde gjøre sig gjældende indenfor samme Grændse som In\-%
dividualiteten i det plastiske Kunstværk. Der maatte udfolde
sig en Tankens Form, der drog det \emph{Physiske} og det \emph{Ethi\-%
ske} ind under det \emph{Metaphysiske}, der i Metaphysiken be\-%
% p. 19 l. 13: an ink speck sits above the second o of "Forhold" (ABBYY
% reads "Forhøld", tesseract "Forhoøld").  The sort is a plain o.
varede \emph{et umiddelbart Forhold til Anskuelsen}, der
lod Tankerne \emph{fæstne} sig til \emph{ideale Billeder} uden Tankens
% Sheet signature "2*" at the foot of p. 19; not transcribed.
% --- p. 20 ---
% GREEK, p. 20 (twice), read on the image at 1200 dpi -- neither OCR witness
% contains a Greek character.  νόησις: acute over ο, no breathing (none is
% printed and none is possible on ν).  τῆς: circumflex over η, drawn as a
% tilde in the first instance and as a flatter bar in the second; same sort.
% νοήσεως: acute over η.  Nothing is supplied from knowledge of Greek.
indre Bevægelse og culminerede i Begrebet om en formel
\gk{νόησις τῆς νοήσεως}.

I denne Bestemmelse af den philosophiske Tankes første
Form, der blev naaet idet Udgangspunktet toges fra de psy\-%
chologiske Betingelser, reflecterer sig den Bestemmelse i Videns
% p. 20 ll. 6--7: the letterspacing stops INSIDE the word: measured glyph by
% glyph, the gaps run 5--9 px through "Sub-jectivitets" and then drop to
% 2--3 px from "b" of "bestemmelsen" onward.  Set as printed.
Idee, der for Tanken maa blive dens første, netop fordi \emph{Sub\-%
jectivitets}bestemmelsen er Grundbestemmelsen i Ideen og
afgiver Grundsynspunktet for dens Opfattelse. Videns Idee, den
% p. 20 l. 9 and l. 30: the ɔ: id-est mark, confirmed on the image at
% 1200 dpi (a reversed c followed by a colon).  Both witnesses give "o:".
absolute ɔ: den sig i Alt vidende Viden, betegner først og \emph{umid\-%
delbarest} Subjectivitetens Idealiseren, og nærmere: \emph{umiddel\-%
bare} Idealiseren som overgribende over Objectivitetens Andet\-%
hed, der endnu ikke gjør sig gjældende som saadan i en Mod\-%
sætning til Idealitetsbestemmelsen. Det er den umiddelbare
Væren-for-sig, den umiddelbare Reflexion i sig fra sit ideali\-%
serede Andet, den umiddelbare Viden af Alt som sig, altsaa en
umiddelbar Viden af et Ideelt, der her træder frem. Dette
Moment i Videns Idee reflecterer sig i de ovenfor fra et psy\-%
chologisk Udgangspunkt vundne Bestemmelser og resulterer for
Tanken af den historiske Udvikling paa dens Culminationspunkt
i det ved Tankens Sondring af Verdenskunstværkets Momenter
naaede Begreb om en \gk{νόησις τῆς νοήσεως} som det Absolutes
Væsensbestemmelse.

I den græske Grundanskuelses Eiendommelighed ligger
Spiren til den første philosophiske Hovedforms Opløsning og
til Udviklingen af en ny Form af Aandens Verdensanskuelse.
I den \emph{umiddelbare} Forening af \emph{Modsætningerne}, af Form
og Stof, Aand og Natur, ligger \emph{Adskillelsen} skjult; Mod\-%
sætningen er ikke i Sandhed overvunden og forsonet, netop
fordi Foreningen kun er den \emph{umiddelbare} Forening af \emph{det
ikke ved hinanden Satte} ɔ: af det \emph{oprindelig Væ\-%
senssondrede}; den bryder derfor uundgaaelig frem saasnart
den begrebsbestemmende Tanke i sin klare Udvikling sondrer
% p. 20 l. 33: an ink speck sits over the o of the second "for" (tesseract
% reads "fór").  The sort is a plain o.
hvad der for \emph{Anskuelsen} og for den endnu i Anskuelsen
hvilende Tanke var uadskilt i Kunstværkets Enhed. Enheden
maa lade Sondringen fremtræde for Tanken, ligesom det frem\-%
% --- p. 21 ---
hævede Forhold i Videns Idee, ved at lade Andetheden glimte
frem, men ikke hævde sig \emph{som} Andethed, viser hen paa Nød\-%
vendigheden af dens Sætten \emph{som saadan}, saaledes at Ideali\-%
tetsbestemmelsen skiller sig fra Realitetsbestemmelsen, Viden
af sig fra Viden af det Andet.

Men idet for den græske Philosophie Adskillelsen mellem
de i dens Grundanskuelse under den tilsyneladende Enhed ufor\-%
sonede Modsætninger træder frem, saa kan Tanken, der har
sondret Modsætningerne, ikke igjen forene dem saalænge den
ikke fuldstændig forlader hiin Anskuelses Grund, og dette vil
den ikke kunne som \emph{græsk} Tanke, fordi den som saadan har
\emph{en naturbestemt Væsensbegrændsning}, der binder
den til hiin Enhedsform: \emph{Skjønhedens} umiddelbare Enhed,
der udtrykker den græske, til \emph{Anskuelsen bundne}, Tankes
eiendommelige Væsensbestemthed. Den maa derfor ende i
frugtesløse Forsøg paa, efterat Bevidstheden om Spliden, om
Modsætningernes Dualisme, er vakt, atter at finde en Enhed.
Det Maal, efter hvilket den stræber uden at kunne naae det,
maa blive: at hævde det af Tilværelsesmomenterne, der i den
græske Verdensanskuelse var opfattet som det \emph{høieste} og som
% PRINTER'S DEFECT, p. 21 l. 21: the comma after the second "Virkelighed"
% is set HIGH, riding at x-height instead of on the baseline (a damaged or
% wrongly mounted sort); ABBYY reads it ")" and tesseract "»".  The sense
% and the sort are a comma, and it is transcribed as one.
den \emph{sande} Virkelighed, som det \emph{eneste}, som al Virkelighed,
altsaa at uddanne en \emph{spiritualistisk} Anskuelse af Tilværel\-%
sen, som dog først en anden Tidsaand kunde give den Gjen\-%
nemførelse, der overhovedet er mulig.

I Modsætning til den eiendommelig græske Verdensan\-%
skuelse, der udtrykte \emph{Aandens umiddelbare Herre\-%
dømme over det af den besjælede Naturlige}, maa en
saadan spiritualistisk Verdensanskuelse opfatte det \emph{Aande\-%
lige}, i dets \emph{Sondring} fra det som \emph{væsenløst} bestemte
Naturlige, som det \emph{Enegyldige}. Denne Opfattelse bliver
Særkjendet for den Periode, der følger efter den classiske Old\-%
tid: for den \emph{christelige Middelalder}. Den finder sit Ud\-%
tryk i Christendommens Gudsforestilling, i dens Lære om Ska\-%
belsen, i dens hele Anskuelse af Guds overnaturlige Virken
paa Naturen, i Læren om Christus, om Menneskets sædelige
% --- p. 22 ---
og religiøse Opgave og om Salighedstilstanden. Idet vi søge
et \emph{Begrebsudtryk} for det i denne religiøse Anskuelses Form
Charakteristiske, betegne vi det som \emph{den abstracte (na\-%
turløse) Aandelighed, der med Bevidsthed sætter
Naturen som det Væsenløse}. I denne Anskuelse have vi
altsaa den \emph{anden} Stilling given, som Aand og Natur for den
frie menneskelige Bevidsthed kunne indtage til hinanden, og
vi have seet den udvikle sig med Nødvendighed af den førstes
Modsigelse. Men denne nye Opfattelse af Forholdet, i hvilken
den for Videns Idee nødvendige Gyldighed af Andetheden er
fornægtet, gjennem Fornægtelsen af det Naturliges væsentlige
Gyldighed, vil \emph{ikke} kunne gjennemføres philosophisk, ikke
% p. 22 l. 13: the "og" between the two spaced words is NOT spaced
% (o--g gap 2 px), so the emphasis is broken as printed.
danne Grundlaget for en \emph{selvstændig} og \emph{eiendommelig}
udformet Philosophie. Den christelige Bevidsthed, der opgiver
Naturen, taber med det Samme Tankens Basis og Udviklings\-%
incitament, det, i hvilket den har sin Bekræftelse; den faaer,
idet den drager sig tilbage i sin rene Aandelighed, og for\-%
nægter en væsentlig Side af det Menneskelige, kun en ind\-%
holdsløs, blot negativ Idealitet, af hvilken ingen Tankeudfol\-%
delse kan fremgaae, ingen indholdsrig, selvstændig Philosophie
skabes. Og den christelige Bevidsthed kan, consequent, end
ikke ville søge en virkelig \emph{Tankens} Udfoldelse, thi dens ne\-%
gative Forhold til det Naturlige indeslutter et negativt Forhold
til Menneskevæsenets naturlige Evner og Kræfter, altsaa ogsaa
til den af det Naturlige baarne og med Naturen overensstem\-%
% p. 22 l. 26: a short horizontal ink mark stands between "idet" and "den"
% (ABBYY "idet'den", tesseract "idet”den").  A speck; not transcribed.
mende Tanke, og den maa opfatte denne, idet den gaaer ind
under den skabte Naturs Bestemmelse: at være et Intet, og
endnu mere, idet den sees under den som en Virkelighed ube\-%
rettiget virkende Naturs Bestemmelse: at være et Syndigt, som
manglende Muligheden til den høiere Erkjendelse, der er Tro\-%
ens. Vel forvandles nu uundgaaeligt dette \emph{negative} Forhold
til Tanken og til det ved Tanken Frembragte, idet Christen\-%
dommen skal existere midt i den fremmede Verden, og under
Kampen og Forsvaret ikke er istand til, af sin abstracte Ne\-%
gativitet at frembringe et selvstændigt Tankeindhold, til \emph{en
% --- p. 23 ---
nødtvungen factisk Anerkjendelse og Optagelse af
det i Principet Fornægtede}, og ved de historisk udvik\-%
lede antike Tankeformers og det fra dem uadskillelige Tanke\-%
indholds Combination med det Christelige see vi det Phæno\-%
men opstaae, der er blevet betegnet som \emph{christelig Philo\-%
% GUILLEMETS, p. 23 l. 6: this quotation OPENS with » and CLOSES with »
% as well -- both marks point right, confirmed on the image at 1200 dpi.
% Set as printed; not regularised to the book's normal «…» (header §2).
sophie}. Men denne »christelige Philosophie» viser sig under
sin Udvikling i \emph{Patristikens} og \emph{Scholastikens} Tid som
en \emph{uselvstændig} og \emph{principløs} udvortes Sammenknytning
af det fra Oldtiden optagne, en fremmed, naturlig Verdensan\-%
skuelse tilhørende Tankeindhold, og af et christeligt, for Tan\-%
ken utilgængeligt, \emph{ubegribeligt} Troesindhold: to Elementer,
der, ifølge deres Væsensforskjellighed, aldrig kunne forenes til
en virkelig Enhed eller en Tankehelhed. Selve den christelige
% p. 23 l. 14: an ink speck stands between "indre" and "Strid" (ABBYY
% "indre»Strid", tesseract "indrøøStrid").  Not transcribed.
Philosophie bliver sig i Scholasticismen den indre Strid mel\-%
lem sine heterogene Bestanddele og deres heterogene Principer
bevidst, og udtaler gjennem Anerkjendelsen af deres Uensartet\-%
hed Dommen over sig selv som Philosophie.

Naar det nu var \emph{Fjernelsen fra Naturen}, Opfattelsen
af den som \emph{væsenløs}, og den deri begrundede \emph{Forkasten
af de naturlige Evner og Kræfter}, der for den middel\-%
alderlige Bevidsthed gjorde det umuligt at udvikle en eiendom\-%
melig og selvstændig Tænkning, saa viser det sig let hvad der
maatte være Betingelsen for, at en Philosophie atter kunde
opstaae, og gjennem denne Betingelse, og gjennem den nye
Tænknings Forhold til sine historiske Forudsætninger vil den
nye philosophiske Epoches Charakteer kunne bestemmes. \emph{Re\-%
habilitationen af Naturen og af det Naturlige i
Mennesket}, navnlig af \emph{den, sine egne Love følgende
menneskelige Tænkning}, maatte være den første Betingelse
for Philosophiens Gjenfødelse, og denne Rehabilitation fandt
Sted, da Scholasticismen havde opløst sig ved sin indre Mod\-%
sigelse, og da det fornyede Bekjendtskab med Oldtiden atter
tilnærmede den ved den historiske Udvikling dertil forberedte
Aand til den Livsanskuelse, der havde været den græske Old\-%
tid eiendommelig. I Stedet for Middelalderens Retning mod
% --- p. 24 ---
det Hiinsidige og dens Flugt fra Naturen og Virkeligheden; i
Stedet for dens ufrie Selvopgivelse og dens Miskjendelse af den
virkelighedsrige, paa en Naturgrund hvilende Sædelighed, traadte
i hint Tidsrum, i hvilket den nye Tid udviklede sig af Mid\-%
delalderen, en Retning mod Virkeligheden, mod det Nær\-%
værende, en Anerkjendelse af Naturen i dens Væsentlighed og
i dens Betydning for Aanden, en Hævdelse af Selvbevidsthedens
og Tankens Frihed, og en Stræben efter Udfoldelsen af et fyl\-%
digt humant Liv, Natur og Humanitet blev Løsenet for Livet
istedenfor Aandelighed og Gudelighed. Menneskeaandens Skue\-%
plads blev en anden: den var ikke længere hiinsides Virkelig\-%
heden, i Himlen, men atter indenfor det \emph{guddommelige
Universum}, i hvilket den \emph{mikrokosmiske Menneske\-%
aand}, i begeistret Tro paa Virkelighedens Sandhed, stræbte,
paa en Naturbasis at udfolde sit evige Indhold, i Kunst, i
Videnskab og i det sædelige Livs Former.

% PENCIL, p. 24: a small tick stands in the margin at the head of this
% paragraph, just left of "Betingelsen".  Not transcribed.
% p. 24 l. 18: "og" between "selvstændig" and "frit" is NOT spaced
% (o--g gap 1 px); "frit" IS (gaps 5, 7).
\emph{Betingelsen} for en Gjenoplivelse af Philosophien var
given i disse forandrede Forhold, og \emph{selvstændig} og \emph{frit
udpræget} fremtraadte Philosophiens \emph{anden Hovedform}
efterat Tænkningen, gjennem en fornyet, fortrolig Tilknytning
til Oldtiden, i en \emph{Overgangstid} havde vundet Styrke og
Indhold, og udviklet en Rigdom af Former, fremtidsbebudende,
tankerige, frihedsstræbende; men uklare, ubevidst afhængige
af Fortiden, og endnu ikke hvilende paa et selvstændig ud\-%
præget eiendommeligt Princip eller gjennemarbeidede til indre
Consequents.

Søge vi i de philosophiske Forsøg i den Periode, hvilken
% GREEK, p. 24, read on the image at 1200 dpi.  The acute over η is plain.
% The breathing over ε is the SAME sort as the elision apostrophe after
% κατ on the same line -- a bar at the top with the stroke descending on
% the right -- and the elision mark is by definition the smooth sort, so
% the breathing is smooth.  The elision mark is typed U+1FBD.
vi betegne som \emph{Overgangsperioden} \gk{κατ᾽ ἐξοχήν}, \emph{Tænk\-%
ningens Grundpræg}, saa bliver det væsentlig \emph{det samme}
som i den \emph{græske} Philosophie. I hele Overgangsperioden
finde vi den samme \emph{umiddelbare} Behersken af det Reale
ved en med Phantasien forbunden Tanke, den samme \emph{meta\-%
physiske} Retning ogsaa indenfor det Physiske, idet ogsaa de
physisk \emph{benævnede} Principer blive behandlede og benyttede
som metaphysiske Principer og kun blive disses Symboler (jfr.
% --- p. 25 ---
Telesio og Bruno); --- den samme Grundopfattelse af \emph{Tilvæ\-%
relsens Totalitet}. Og gjennem hele Overgangsperioden
finde vi ogsaa Bevidstheden om Slægtskabet med Oldtidens
Aand, og vi see Tænkerne væsentlig bundne til det fra Oldti\-%
dens forskjellige Systemer optagne Tankeindhold. Vel stræbes
der, og navnlig ved Overgangsperiodens Slutning, ud over Old\-%
tiden, baade med Hensyn til \emph{Methoden}, idet der lægges
Vægt paa Erfaringen og Experimentet, og med Hensyn til
\emph{Tankeindholdet}, idet Problemerne antage nye Former. Men
\emph{Erfaringen} og \emph{Experimentet}, der fordres, maa dog be\-%
standig vige Pladsen for den \emph{metaphysiske} Tænkning, og
selv hvor Problemerne komme frem i nuancerede Skikkelser,
bliver den \emph{umiddelbare} Sammenhæng med Oldtidens Tanke\-%
stof dog bestandig kjendelig.

Den gjennem Overgangsperiodens Tankeudvikling forberedte
\emph{nye} Hovedform for Philosophien: Philosophiens \emph{anden}, nu
\emph{fuldt afsluttede} Hovedepoche, betegne vi som \emph{den nyere
Tids Philosophie} eller som Philosophien i Middelalderens
Opløsningsperiode.

Betragte vi denne \emph{anden} Hovedepoche fra \emph{Bevidst\-%
hedsudviklingens} og de \emph{historiske Tidsbetingelsers}
Synspunkt, saa viser dens betegnende Eiendommelighed sig
at være den ved Selvbevidsthedens Enhed bestemte \emph{Stræben
efter en Tankens Enhed, ud fra Forudsætningen af
den absolute Forskjel mellem Tilværelsens Grund\-%
former: Naturen og Aanden}, Tanken og Tankens Andet.
Denne Eiendommelighed betingedes ved det historisk Forud\-%
gaaende. Ved Middelalderens Livsanskuelse var der sat den
skarpe \emph{dualistiske} Adskillelse mellem Aanden, det ene Sande,
og Naturen, det Væsenløse og Usande. Overgangsperiodens
Systemer havde kæmpet mod denne Dualisme, men havde,
bundne til den græske Tænknings Grundform, og ubevidst paa\-%
virkede af Middelalderen, ikke mægtet at overvinde den. De
efterlode denne Opgave \emph{uløst} til den nyere Philosophie, og
% p. 25 l. 35: "med den den ubevidste" -- "den den" as printed, confirmed
% on the image; both witnesses agree.
med den den ubevidste Eftervirkning af Middelalderens Forud\-%
% --- p. 26 ---
% PENCIL, p. 26: a "+|" and a cursive word in the left margin beside
% ll. 5--6, and a grey "?" written between "uvilkaarligt" and "droges" in
% l. 32 (which ABBYY renders "uvilkaarligi#flroges" and tesseract
% "uvilkaarligttdroges").  None of it is printed type; not transcribed.
sætninger. Forholdet mellem Grundmomenterne kunde altsaa
for den nyere Philosophie ikke mere, som hos Grækerne, blive
den umiddelbare Harmonies, der skjuler over den i Grunden
værende Splid; Forholdet er fra Begyndelsen af \emph{Modsætnin\-%
gens} mellem den i sin Væsentlighed anerkjendte Aand og den
i sin Væsentlighed restituerede Natur. Denne Modsætning skal
bringes til Enhed, men Enheden kan, netop paa Grund af den
i Udgangspunktet liggende Forudsætning, kun blive en \emph{ensi\-%
dig} Enhed, en saadan, i hvilken den ene af Modsætningens
Sider er \emph{forsvunden}, og den tilbageblivende, der synes at
omfatte begge, i Virkeligheden kun faaer den Bestemmelse,
den havde \emph{indenfor Modsætningen}. \emph{I selve Enhedens
Art} maa saaledes den \emph{uforsonede Modsætning} vise sin
Indvirkning. Udviklingens Slutningspunkt maa i denne philo\-%
sophiske Hovedepoche aabenbare den i Udgangspunktet skjulte
Mangel.

I den ved Begyndelsen værende Modsætning laae Nødven\-%
digheden af, at denne Hovedform af Philosophien maatte ud\-%
vikle sig i en \emph{dobbelt Række af Systemer}, at den, fra et
\emph{dobbelt Udgangspunkt}: fra Naturen og fra Aanden,
maatte stræbe hen efter Enheden, og ved det dobbelte Ud\-%
gangspunkt maatte \emph{Erkjendelsens Charakteer} og \emph{væsent\-%
lige Gjenstand}, og \emph{Rækkens Slutningspunkt og Re\-%
sultat} blive forskjellig bestemt. Den fra \emph{Naturen} udgaaende
Retning maatte faae Naturen til sin væsentlige Gjenstand, Er\-%
faringen til Erkjendelsens Form, den maatte, gjennem Empi\-%
risme og Sensualisme, udvikle sig til, som Materialisme at lade
Aanden forsvinde i Materien som en enkelt Existentsform af
denne, at sætte den aandløse Materie som Tilværelsens Enhed.
Den fra \emph{Aanden}, fra Selvbevidstheden, udgaaende Retning
maatte paa sin Side, om den end, ifølge hele den nye Tids
Retning, uvilkaarligt droges hen imod Naturen, sætte Aanden
som sin høieste og væsentligste Gjenstand; den maatte til Er\-%
kjendelsens Form faae den aprioriske Tænkning; dens sidste
Maal maatte blive en Panlogisme, der lader al virkelig Væren
% --- p. 27 ---
forvandles til Tanke-Væren, der lader Naturen og det Mate\-%
rielle blive en forsvindende Væreform indenfor Aandens ideelle
Proces, den evigt virkelige eneste Sandhed, og altsaa sætter
den naturløse Tanke som Tilværelsens Enhed. Den vil ved
sit Slutningspunkt faae en Lighed med Middelalderens natur\-%
fornægtende Spiritualisme, men den vil adskille sig fra denne
derved, at den ikke \emph{principielt} og \emph{bevidst} fornægter Na\-%
turen, at den netop har en Stræben efter, en Villie til, at an\-%
erkjende den, men ved sin dualistiske, ubevidst indvirkende
Forudsætning forhindres i at anerkjende den i dens Eiendom\-%
melighed \emph{som} Natur.

Gjennem \emph{begge} den nyere Philosophies Udviklingsrækker
% p. 27 l. 13: a thin, wavy raised mark follows "fælleds" (ABBYY and
% tesseract both read it as an apostrophe).  It is far lighter and finer
% than the type -- a hair or scratch, not a sort -- and is not transcribed.
gaaer saaledes som det \emph{fælleds} Kjendemærke \emph{en af Dua\-%
lismen paavirket og hæmmet Stræben efter Mod\-%
sætningernes Enhed}. Den stræber hen mod det Maal: at
forene Tilværelsens to positive Grundmodsætninger til en sand,
indholdsrig Enhed; men Enheden, den høieste, altomfattende
Væsentligheds Tanke, bestemmes kun under den skjulte Ind\-%
virkning af den dualistiske Forudsætning, og Opfattelsen af
% p. 27 l. 20: "Ud-" is NOT spaced (gaps 2, 3), nor is "viklingsrækker"
% on the next line; only "Begge" is.
den bliver derfor \emph{ensidig} og \emph{uconsequent}. \emph{Begge} Ud\-%
viklingsrækker strande paa denne Ensidighed, og idet de ved
den røbe deres Mangel paa Evne til at løse deres Opgave:
Dualismens Overvindelse, anvise de \emph{Fremtidens Philoso\-%
phie} dens Opgave, om end, efter Sagens Natur, i abstract og
negativ Form.

I denne \emph{anden} philosophiske Hovedepoche reflecterer sig
det Grundforhold i Videns Idee, i hvilket Andetheden, der
først saaes som umiddelbart idealiseret, umiddelbart behersket
af Subjectiviteten, sættes \emph{som} Andethed, altsaa gjør sig gjæl\-%
dende som Modsætning, som det med Subjectiviteten \emph{lige Op\-%
rindelige}, som hvilket den bliver at sætte trods det Subjec\-%
tives ideelle Principat, idet Andetheden, Objectivitetsmomentet,
i sin selvgyldige Væsentlighed bliver Betingelse for Viden som
den absolute Subjectivitets sande, concrete Viden af sig som
idealt-realt Princip.

% --- p. 28 ---
Den \emph{tredie} philosophiske Hovedepoche, til hvilken der
vistes hen gjennem den andens mangelfulde Resultat, bliver,
ifølge Begrebsforholdene, at see som den \emph{sidste mulige},
men tillige som den, der i sit Princip indeslutter Betingelsen
for en \emph{uendelig} Tidsudvikling.

De tre Epochers indbyrdes Forhold bliver nemlig saaledes
at bestemme. Den \emph{første} Epoche: den \emph{græske} Philosophie,
havde havt den Opgave, at finde \emph{Ideens} Tanke, og at søge
Enheden ved \emph{at hævde Ideen som det i Tilværelsens
Harmonie Væsentlige og Sande}. Den \emph{anden} Epoche:
den \emph{nyere} Tids Philosophie, der fremstod efter Philosophiens
store Interregnum i den ved sit Princip fra Frembringelsen af
en eiendommelig Tænkningens Grundform udelukkede Middel\-%
alder, fik den Opgave: \emph{ud fra den dualistiske Forudsæt\-%
ning at søge Modsætningernes Enhed}. Idet denne Op\-%
gave mislykkedes; idet det viste sig, at saavel den i Materia\-%
lismen, som den i Idealismen naaede \emph{ensidige Enhed} kun
var en \emph{skjult Form} for den \emph{Dualisme}, hvis Overvindelse
fordredes som en Tankens Nødvendighed, saa blev \emph{Dualis\-%
mens Overvindelse} atter Opgaven for \emph{Fremtidens} Phi\-%
losophie, og Formen for Løsningen af Opgaven bestemmes ved
det ene mulige Forhold mellem Modsætningerne, der er tilbage
efterat \emph{den umiddelbare Enhed} og den \emph{abstracte Mod\-%
sætning} have viist sig uigjennemførlige: ved \emph{Modsæt\-%
ningsenheden}. Forat løse sin Opgave, maa Fremtidens
Philosophie anerkjende Aandens og Naturens, det Ideelles og
det Reelles, \emph{væsentlige Eiendommelighed}, og see deres
\emph{Væsensdifferents} som en \emph{Modsætning}, men \emph{som en
Modsætning, der begrundes ved og sammensluttes
til Enhed i det som absolut Subjectivitet bestemte
ideelt-reale Princip}.

I den derved betegnede Grundbestemmelse: \emph{Enhed i
væsentlig Modsætning af Aand og Natur}, af Tanken
og Tankens Andet, af Selvbevidstheden og Bevidsthedsgjenstan\-%
den, reflecterer sig det \emph{sidste} Grundforhold i Videns Idee.

% --- p. 29 ---
% SEAM: p. 29 line 1 is set FLUSH LEFT (measured against the page's own left
% margin), so it continues the paragraph running from p. 28.  The splice must
% not introduce a paragraph break here.
Idet Andethedens Fremtræden \emph{som Andethed} var Betingelsen
for Viden som Principets \emph{concrete} Viden af sig selv, var en
Sætten af \emph{Modsætningen} fordret som \emph{nødvendig}, men
idet den sættes som nødvendig for \emph{Principets Viden} af
\emph{sig selv}, er herved allerede viist tilbage til, at Andetheden
hører til \emph{selve Principets Væsen}, at den er \emph{det ved det
overgribende ideelle Princip Begrundede}, Modsætnin\-%
gen en Modsætning \emph{indenfor} den overgribende Subjectivitet,
at altsaa Modsætningen umiddelbart peger hen paa den ind\-%
holdsfyldige Enhed, paa Videns Idee i dens \emph{absolute Con\-%
cretion} som Udtrykket for det absolut concrete ideelt-reelle
% p.29 l.12: a faint speck stands between `Forsigværen' and `gjennem'
% (ABBYY reads it as a full stop, ".gjennem").  It is not type -- much
% lighter than the ink and with no letter shape -- so it is not transcribed.
% RECON.md §9 lists p. 29 as a pencil-pollution site.
Princips Forsigværen gjennem den som Andethed satte An\-%
dethed.

To Epocher af Philosophien, i deres Princip bestemte ved
væsentlige Ideeforhold og i deres Følge ved en Udviklings Nød\-%
vendighed, foreligge saaledes historisk som \emph{fuldkomment
afsluttede}, og ved det \emph{Grundproblem}, de have stillet,
men ikke løst, vise de hen paa en \emph{tredie} Hovedform, hvis
Udvikling, der ved selve den abstracte Formel antydes som
betinget af de til det Reale umiddelbart sig forholdende Viden\-%
skabers Tilegnelse og Assimilation ved den philosophiske Tanke,
tilhører \emph{Fremtiden}, den \emph{hele} Fremtid.

\emph{Schematisk} kan Grundforholdet mellem de tre Hoved\-%
epocher, i hvilket \emph{Inddelingens Fuldstændighed} ligger,
angives i følgende Form, der viser Grundopfattelsen af det Ab\-%
solute i de forskjellige Epocher, og det Grundforhold, de sætte
mellem Tilværelsens Hovedformer.

% =====================================================================
% SCHEMA-TODO p.29 -- see RECON.md §3.  Words captured; the construction
% (the swash ornament, the two columns, the column positions) is set in a
% later pass.  THIS MUST NOT BE LEFT AS RUNNING TEXT.
%
% What is printed, exactly, and where (read at 600 dpi on the image; the
% whole displayed block is set to a measure WIDER than the body text, its
% left edge standing about 1.5 cm outside the body's left margin):
%
%   line 1  BOLD, centred, with about two ems between the figure and the
%           words:            "1.   Den græske Philosophie."
%   line 2  LETTERSPACED roman, centred, the Greek word in the book's
%           italic Greek fount and NOT letterspaced:
%             "Ideen aabenbaret i Verdenskunstværket (Κόσμος)."
%   line 3  LETTERSPACED roman, centred, in parentheses:
%             "(Umiddelbar Enhed af Natur og Aand)."
%   line 4  BOLD, centred, same two-em gap:
%             "2.   Den nyere Tids Philosophie."
%   line 5  A WIDE SWASH ORNAMENT (a swelled "wave" rule), centred, about
%           three quarters of the measure wide.  It is hairline-thin at
%           both ends and swells into two heavy lobes, one under each of
%           the two columns below, which meet at a small UPWARD CUSP at
%           the exact centre, directly under the middle of head 2.  It
%           therefore works as a wide two-column brace whose apex points
%           up to the head.  It is NOT a divisional rule and NOT a
%           footnote rule.
%   lines 6--7  TWO COLUMNS of two lines each, side by side, separated by
%           a wide gap.  ONLY THE FIRST WORD OF EACH COLUMN IS
%           LETTERSPACED ("Natur", "Aand"); everything else in the two
%           columns is ordinary roman -- checked at 600 dpi, "Materie",
%           "Ikke-Aand.", "(Jeg --- absolut Aand)" and "Ikke-Natur = det
%           Ikke-Materielle." are all set close.  The "=" is the text
%           sort, not math; the dash in the right column is an em-dash.
%           Each column's second line is centred under its first.
%
%             left column                right column
%             "Natur = Materie ="        "Aand (Jeg --- absolut Aand) ="
%             "Ikke-Aand."               "Ikke-Natur = det Ikke-Materielle."
%
%           i.e. the two columns read "Natur = Materie = Ikke-Aand." and
%           "Aand (Jeg --- absolut Aand) = Ikke-Natur = det
%           Ikke-Materielle."
%   line 8  LETTERSPACED roman, centred under BOTH columns:
%             "(abstract Modsætning af Natur og Aand)."
%
% The schema runs on across the page break: head 3 and its matter stand at
% the top of p. 30, and the SCHEMA-TODO block there continues this one.
% =====================================================================

% --- p. 30 ---
% =====================================================================
% SCHEMA-TODO p.30 -- the continuation of the p. 29 schema; see the block
% above and RECON.md §3.  Words captured; construction set in a later pass.
%
%   line 1  BOLD, centred, ending in a COLON, not a full stop:
%             "3.   Fremtidens Philosophie:"
%   lines 2--3  LETTERSPACED roman, set to the same wide measure as the
%           p. 29 block (left edge about 1.5 cm outside the body margin),
%           the first line filling the measure and broken with a hyphen,
%           the turnover indented about one and a half ems:
%             "Den absolute Subjectivitet som ideelt-reelt Tilværel-"
%             "     sesprincip."
%   line 4  LETTERSPACED roman, centred, in parentheses:
%             "(Modsætningsenhed af Natur og Aand)."
%   line 5  a SHORT CENTRED DIVISIONAL RULE closing the schema, measured
%           at 1.55 cm and centred on the schema's (wider) measure.  That
%           rule IS text and is set below as \divrule.
% =====================================================================
\divrule
% The rule above closes the displayed schema.  It is the book's ordinary
% divisional rule -- short and CENTRED -- not the left-flush footnote
% separator (RECON.md §3).

Hvad der her er udviklet med Hensyn til Grundforholdet
mellem de enkelte Hovedformer, maa faae sin \emph{concretere}
Bestemthed gjennem en Charakteristik af dem fra de for en
Tankeforms Eiendommelighed afgjørende Hovedsynspunkter. Men
en saadan Charakteristik, ligesom ogsaa den bestemtere Paa\-%
viisning af Forholdet mellem det i Udviklingen Tidligere og
Sildigere, mellem det Afsluttede og det Fremtidige, vil hen\-%
sigtsmæssigst kunne gives, efterat de i historiske Exempler paa\-%
viste teleologiske Grundforhold have givet Stof til en rigere og
fyldigere Anskuelse.

Derimod vil allerede her den \emph{Lovmæssighed}, der be\-%
herskede Udviklingsgangen \emph{i de store Hovedepocher}, blive
at paavise \emph{concretere} indenfor \emph{en enkelt} af dem: den
\emph{græske} Philosophie, og det paa den Maade, at \emph{Epochens
Princip} sees som \emph{Epochens Lov}, som dens indre tilskyn\-%
dende og bevægende Udviklingsenergie, som det, der i sig, som
sine Begrebsmomenter, skjuler \emph{Principerne for Epochens
enkelte Systemer}. I de enkelte Systemers Principer ud\-%
folder Heelhedens Princip sig, vorder, culminerer og opløses
% p.30: PRINTED AS SHOWN -- a hyphen stands between `ubevidst' and `for',
% where the sense wants a space.  The mark is much thinner and lighter than
% this page's ordinary hyphen (compare `be-' two lines below), but it is
% INK, not pencil: measured at 600 dpi its darkest pixels reach grey 38
% against 30 for a certain hyphen and 48 for ordinary text on the same page,
% and it carries the same warm brown as the type.  Both OCR witnesses record
% it.  Read as an under-inked or worn hyphen sort and transcribed as printed;
% the expected reading is `ubevidst for Tænkerne'.
gjennem dem, og aabenbarer, idet det, ubevidst-for Tænkerne,
behersker dem, det mere Omfattende, ved hvilket det selv be\-%
stemmes, og som behersker dets egen Udfoldelse og Opløsning.

% p.30: a light pencil stroke underlines `søge' in the line below, and a
% further faint tick stands under `ciper' four lines on.  Modern marginalia;
% not transcribed.
Idet vi søge den begrebsbestemte Udviklingsgang i den
enkelte Epoche, søge vi altsaa dens \emph{Nødvendighed i den
centrale} Tanke, der behersker Epochen \emph{som Heelhed}, og
hvis Begrebsanalyse maa frembyde de Momenter, der som Prin\-%
ciper beherske de enkelte Systemer, Epochen omslutter.

% --- p. 31 ---
% L1 head (transcription.tex §3): large BOLD centred head with a SHORT
% CENTRED RULE.  Note that here the rule stands BELOW the head, not above
% it: this head opens p. 31 and there is no rule above it.
\begin{center}
\textbf{\Large Den græske Philosophies Udvikling.}
\end{center}
\divrule

Naar vi skulle søge den Tanke, der er at bestemme som
den græske Aands Grundtanke, som \emph{Grækenlands Tanke},
ledes vi til at finde den gjennem en Reflexion, der ligesom af
sig selv frembyder sig. Vi finde hos det græske Folk, som
Form for Anskuelsen af det \emph{Guddommelige}, Menneskevæ\-%
senets Ideal, som udpræget \emph{i alle} Livets Former, som det
hele dets Tilværelse gjennemtrængende Grundpræg, \emph{Skjøn\-%
heden}, Harmonien. Men det, der charakteriserer et heelt
% p.31: two small pencil/dirt marks in the line below -- a dot before
% `aandelige' and an apostrophe-like tick after `betegner' (both OCR
% witnesses pick the second one up as `betegner''; ABBYY also gives
% `Folks -aandelige').  Neither is type; neither is transcribed.
% RECON.md §9 lists p. 31 as a pollution site.
Folks aandelige Selv og betegner dets formende Princip, det
maa, naar Folket uforstyrret har udviklet en \emph{selvstændig},
dette Folk væsentlig tilhørende, Tænkning, og fremfor Alt, naar
Folket har udviklet sig saa naturmæssigt og harmonisk som
Grækerfolket, Philosophien som Folkets \emph{ideelle Selverkjen\-%
delse} paa sit høieste Udviklingstrin blive sig bevidst, og søge
at udtrykke. Og det, der er den \emph{menneskelige} Tilværelses
af Tanken opfattede formende, beherskende Princip, det maa
Philosophien, idet Menneskeaanden maa sætte sig selv som
% p.31: \gk{τέλος} -- tau, epsilon with a clear ACUTE (a stroke rising to
% the right), lambda, omicron, final sigma.  Read at 1200 dpi.  `Tilværelsens'
% before it is NOT letterspaced (spacing.py flags it; the image at 600 dpi
% shows it set close).
Tilværelsens \gk{τέλος}, stræbe at opfatte som al Tilværelses be\-%
stemmende Princip. Derfor maatte den græske Tænkning paa sit
Høidepunkt, idet den vilde opfatte \emph{Tilværelsens Heelhed},
opfatte den fra \emph{Skjønhedens} Synspunkt, som det \emph{altom\-%
fattende Kunstværk}, som \gk{Κόσμος}, og i denne den græske
\emph{Aands høieste Anskuelse} maatte \emph{Ideens Tanke} komme
tilsyne som Tanken om \emph{Skjønhedsformens Grund}, om
% p.31 l.26: ERRATUM SITE -- the Trykfeil leaf asks for a comma after
% `Høieste'.  The print has NO comma; it stands as printed.  See
% transcription.tex §5.
\emph{det absolut Høieste det i Sandhed Værende}.

Ved at analysere den Anskuelse, der for Grækerne over\-%
førtes paa Tilværelsens Heelhed i dens harmoniske Væren, føres
vi til Sondringen af de Grundmomenter, hvis indre Forhold
det var den græske Tænknings Opgave at bestemme, og dette
Forhold indeslutter det Normerende for den græske Philosophies
Udvikling og Spiren til dens Opløsning.

% --- p. 32 ---
\setcounter{footnote}{0}
Tage vi, idet vi ville analysere \emph{Kunstværket}, vort Ud\-%
gangspunkt fra Opfattelsen af det som det \emph{enkelte, be\-%
grændsede} Kunstværk, og gjøre vi Anskuelsen endnu be\-%
stemtere ved at tænke os det for Grækenlands Kunst meest
charakteristiske Kunstværk: et \emph{plastisk} Kunstværk, en Bil\-%
ledstøtte f. Ex. af Marmor eller Erts, saa viser det sig, at
dette Kunstværk \emph{er}, hvad det \emph{er}, \emph{ved sin Form}, at Stoffet
er gaaet op i denne, og i Kunstværket kun er som det \emph{for\-%
mede Stof}, saa at det der ingen selvstændig Gyldighed og
Væren har ved Siden af Formen. Det viser sig, at Formen,
der udtrykker den ideelle Betydning, ikke blot er det Fuld\-%
endende, det Besjælende og Beaandende, men at den er det,
der giver Kunstværket \emph{Væren} ved at give det \emph{Skjønhedens
Væren} ɔ: \emph{den sande} Væren; thi Skjønhed er her Sandhed
og Virkelighed. Men saalænge vi blot tænke paa det enkelte
\emph{begrændsede} Kunstværk, faaer dog bestandig Stoffet en
Væren, der er \emph{relativ uafhængig af Formen}, fordi den er
\emph{relativ uafhængig af Kunstværket}. Marmoret og Ertset
have ogsaa en Væren som uformet Materiale før det dannes
og formes af Kunstneren, og den \emph{samme} Form kan finde sit
\emph{Udtryk i begge} Stoffer: Sammenhængen er ikke uopløselig.
Udvide vi derimod hiin Anskuelse og ombytte Forestillingen
om det \emph{enkelte, begrændsede} Kunstværk med Forestillingen
om det \emph{altomfattende}, om \emph{Verdenskunstværket}, om
% p.32: PRINTED AS SHOWN -- a hyphen stands between `Enhed' and `omfatter',
% where the sense wants a space.  Same phenomenon as `ubevidst-for' on
% p. 30: thinner and lighter than the ordinary hyphen, but ink (darkest grey
% 16 against 48 for ordinary text on this page, same warm brown), and
% recorded by both OCR witnesses.  Transcribed as printed.
Universet som \gk{Κόσμος}, saa faaer Stoffet ikke længer denne
ligegyldige Væren; thi Universet i dets Enhed-omfatter \emph{alt}
Stof og Formen \emph{som Totalitet}. Stoffet overhovedet bliver
saa, idet det ikke længer kan finde nogen Plads ved Siden af
den altomfattende ɔ: alt Stof omfattende, Form, \emph{kun et Væ\-%
rende ved Formen}, ved Verdenskunstværkets besjælende Form\footnote{Den Mulighed, at Stoffet kunde forestilles som havende \emph{Væren
som uformet} forud for Formningen, falder bort naar \emph{Verdens\-%
kunstværket} sees som Verdenstotalitetens \emph{væsentlige} Existents\-%
form: Formen sees da \emph{umiddelbart} som det Formende.}.
% The note above is marked *) in the print, and the mark stands after
% `Form' and before the full stop.  The short LEFT-FLUSH rule above the note
% area is the footnote separator and is not transcribed.
% --- p. 33 ---
Og dette vil med andre Ord sige: saa bliver \emph{Formprinci\-%
pet Værens og Virkelighedens Princip}, \emph{Stoffet} bliver
det, der, løsrevet fra Formen, er et \emph{Ikke-Værende}, og kun
faaer Væren og Virkelighed i Enhed med denne. Men \emph{For\-%
men} er det Legemløse, det \emph{Ideelle}; den er ligesom af \emph{Aan\-%
dens Slægt}, og ved Formen forstod ogsaa den \emph{anskuende}
græske Tænkning, for hvilken Formens og det Formendes Be\-%
greb gik i Enhed, hvad den moderne Tænkning forstaaer ved
\emph{Aanden}, medens \emph{Stoffet} bliver tilsvarende Udtryk til \emph{Na\-%
turen}, tænkt som Materie. I Bestemmelsen af Formen som
det Høieste og Herskende er saaledes \emph{Aandens Principat}
udtalt, og Formens Sammenhæng med Aanden udsiges ud\-%
% p.33: \gk{εῖδος} -- a broad circumflex-shaped mark stands over the iota
% and is plainly visible at 1200 dpi (it matches the tilde over the upsilon
% of Νοῦς four lines below).  Whether a SMOOTH BREATHING is combined with it
% in the single sort cannot be separated from the blob; ἶ/ῖ undecided.  The
% circumflex is transcribed because its shape is visible; nothing is supplied
% for the breathing.  Cf. RECON.md §4 on the compound mark.
trykkeligt af selve den græske Philosophie. Hvor \emph{Aristoteles}
naaer til Begrebet om den høieste \gk{εῖδος}, om Formen i dens
Reenhed, til det Afslutningspunkt i hans System, paa hvilket
den græske Tænkning ligesom for et Øieblik hæver sig over
% p.33: \gk{Νοῦς} (twice) -- the circumflex over the upsilon is an
% unmistakable tilde at 1200 dpi.  \gk{νόησις τῆς νοήσεως} -- acutes over
% the omicron of νόησις and the eta of νοήσεως, circumflex over the eta of
% τῆς; all resolved.
sig selv, der \emph{bytter Formen Navn} og hedder \gk{Νοῦς}, og
denne \gk{Νοῦς} bliver bestemt som den sig selv tænkende gud\-%
dommelige Tankes Energie, som den evige \gk{νόησις τῆς νοήσεως}.

Naar altsaa \emph{Verdenskunstværket} var den græske Aands
\emph{høieste Anskuelse}, saa bliver \emph{det} af Kunstværkets Mo\-%
menter, der er \emph{Værens Princip}, dens \emph{høieste Indhold},
og \emph{Tankeudtrykket} for dette Moment dens \emph{høieste Tanke}.
Men Tankeudtrykket for dette Moment, der opfattedes, ikke
af en abstract, men af en anskuende Tænkning, var \emph{Ideen},
den sandselige Tilværelses af Rum og Tid og Vorden uafhæn\-%
gige intelligible Urform og Virkelighedsgrund.

\emph{Ideens} Tanke bliver saaledes den \emph{centrale} Tanke, der
maa beherske den græske Philosophies Udvikling, idet den,
gjennem sit Forhold til sin i Verdenskunstværket med den til
harmonisk Enhed sammensluttede Modsætning, \emph{Stoffet}, be\-%
stemmer og charakteriserer Udviklingens Perioder, og gjennem
det samme Forhold betinger den græske Tænknings uoverstige\-%
lige Grændse, gjør dens Opløsning uundgaaelig.

% The sheet signature `3' stands at the foot of p. 33.  Not transcribed --
% and not in RECON.md §3's signature list, which begins at `3*' (p. 35).
I Anskuelsen af Kunstværket som Formens og Stoffets
% --- p. 34 ---
Enhed er der, som allerede ovenfor antydet, givet en \emph{umid\-%
delbar} Forening af to Momenter, af hvilke intet er \emph{sat ved
det andet}. \emph{Stoffet} kunde ikke bestemmes som sat, frem\-%
bragt ved Formen; thi en saadan Sætten, der, ifølge Bestem\-%
melsen af det Ideelle som kunstnerisk \emph{Form}, ikke kunde blive
seet som oprindeligt Væsensforhold i Ideen, maatte af Grækerne
have været opfattet som en Skaben; men denne Forestilling
var ikke blot bestandig fremmed for Grækerne, selv for Ny\-%
platonikerne, der forsøgte at udlede Materien af det Ideelle ved
en dynamisk Emanation, men den \emph{maatte} være det efter
Grundanskuelsens, Kunstværkets, Natur; thi ved Kunstværket er
Stoffet ɔ: det, der \emph{skal formes, Forudsætning} for at For\-%
men kan komme til Virksomhed. Formen som Kunstværkets
Form, i sin \emph{umiddelbare} Væren for Anskuelsen, har \emph{umid\-%
delbart} sit Stof, som den former. Og naar Enheden løsnes
af Tanken, og Formen som Idee sondrer sig fra Stoffet, saa
\emph{forudsætter} den --- ikke i Tidens, men i Begrebets Betyd\-%
ning --- i sin hvilende Væren udenfor dette netop Stoffet
som det, den, som formende, \emph{skal forme}, selv om det
i denne Sondring \emph{ikke} kan holdes fast som en \emph{Virkelighed}
af Tanken. Men naar saaledes \emph{Stoffet} ikke kunde tænkes
som sat ved Formen, saa viser det sig endnu lettere, at \emph{For\-%
men} ikke kunde tænkes som sat ved Stoffet, thi Formen er
ikke blot det Høiere og Herskende, men efter sit Begreb, \emph{umid\-%
delbart}, det af Tilbliven og Vorden Uberørte, det Evige. Mo\-%
menternes Enhed er altsaa ikke fremgaaet af det enes Begrun\-%
dethed i det andet; intet af dem er det Frembringende, der
udvikler det andet af sig. Enheden er den \emph{umiddelbart
for Anskuelsen værende}, men den \emph{umiddelbare Enhed} af
hvad der ikke er seet som fremgaaende af hinanden, er en
umiddelbar Enhed af et \emph{oprindelig Forskjelligt}, der
atter for Tanken maa aabenbare denne Forskjellighed, adskilles
i uforsonet Modsætning. Vel er efter Begrebet Formen al sand
Værens Grund; der kan ikke tænkes Væren og Virkelighed
udenfor den; men idet Skjønhedsformen, forat være som for\-%
% --- p. 35 ---
\setcounter{footnote}{0}%
% p.35: \gk{μὴ όν} -- the mark over the eta is a GRAVE (a stroke falling to
% the right), positively resolved at 1200 dpi.  Over the omicron there are
% two elements: the right one is plainly an ACUTE, the left one is the
% breathing, and its shape (comma versus "C") CANNOT be resolved -- this is
% the known-unresolvable compound over a low vowel, RECON.md §4.  The omicron
% is therefore set with the acute alone and the breathing is NOT supplied:
% ὄ/ὅ undecided.  Cf. the same word set legibly on p. 36, where the breathing
% reads as smooth.
mende Form, forat give sig Virkelighed som saadan, behøver
et forudsat Stof, saa faaer dette, der i sin Adskilthed fra For\-%
men maatte begribes som et \gk{μὴ όν}, dog ved denne for Tanken
uundgaaelige Adskillelse bestandig en Art af \emph{egen Væren},
af forudsat Selvstændighed, det bliver det modsigelsesfulde \emph{væ\-%
rende Ikke-Værende}.

Ved det angivne Forhold mellem Formen og Stoffet er nu
\emph{den græske Philosophies Udviklingsgang} bestemt og
tillige det Problem antydet, paa hvilket denne Tænkningsform
maatte brydes.

Idet den græske Tænkning søger at finde sin eiendomme\-%
lige Verdensanskuelse og at erkjende \emph{Tilværelsens høieste
Princip}, og under denne Stræben ledes af det, der er dens
eiendommeligste Væsen, maa den begynde med \emph{Stoffet} som
det af de to Momenter, der i sin umiddelbare Væren maa være
\emph{forudsat} for at Formen kan virke, og gjennem hvilket alene
% p.35: a small dark speck sits over `sig' in the line below (tesseract
% reads it as `stg').  Not type; not transcribed.
det Formende kan aabenbare sig; den maa begynde med at
sætte Stoffet, det \emph{Naturlige}, som \emph{Principet}, men saaledes,
at fra Begyndelsen af Stoffet ligesom \emph{forvandles} og røber
et \emph{Ideelt}, og at Tanken, ved selve de Problemer, der reise
sig idet Tilværelsen i dens Concretion skal forklares ved Stof\-%
principet, maa drives fremad til at opdage det andet Moment:
\emph{Formen, Ideen, Aanden}. Idet Stoffet skal opfattes \emph{som
Princip}, faaer det Charakteren af et \emph{Alment}, og heri
skjuler sig en \emph{Idealitets}bestemmelse, idet den Bestemthed,
den Particularitet, der betinger Stoffets \emph{virkelige} Væren \emph{som
dette Stof}, maa forsvinde under hiin Opfattelse, og idet Fore\-%
stillingen om et Chaos som en værende \emph{almeen Materie} ---
en Forestilling, der er utilstrækkelig til at forklare det Mang\-%
% p.35: \gk{πρώτη ὕλη} -- the ROUGH breathing over the upsilon is a clear
% "C" opening to the right, with the acute beside it; positively resolved at
% 1200 dpi and plainly a different shape from the smooth breathings
% elsewhere.  This is the calibration case of RECON.md §4.
foldiges Vorden\footnote{Hverken den chaotisk ubestemte Materie, eller den som \gk{πρώτη ὕλη}
tænkte abstracte Materie, kan tænkes at bestemme sig selv til det
Mangfoldige. Ved en saadan Bestemmen vilde jo Materien ved sig
selv vinde Form.}, --- ikke kan fastholdes. Bestemmelsen af
% The sheet signature `3*' stands at the foot of p. 35; not transcribed.
% --- p. 36 ---
% p.36 is NOT in RECON.md §4's list of Greek-bearing pages, but it carries
% \gk{τὸ ὄν}.  The grave over the omicron of τὸ falls to the right and is
% unmistakable.  Over the omicron of ὄν the two elements separate cleanly at
% 1200 dpi: the left one is a comma-shaped SMOOTH breathing (visibly not the
% "C" of ὕλη on p. 35) and the right one an acute.  Both read off this image;
% nothing supplied.
Principet som Stof vil derfor uundgaaeligt skride frem mod
ideellere Bestemmelser: Principer som Vand, Luft, ubegrændset
og ubestemt Stof, ville afløses af Principer som Tal, Væren
(\gk{τὸ ὄν}), begge opfattede som det Stof, hvoraf Tingene bestaae,
som den legemlige Tilværelses Substants, og idet Problemet
om \emph{Vorden}, om det Mangfoldiges Tilbliven af Principet,
kommer frem som \emph{det væsentlige}, bliver Vordens Begreb
det, der, medens det, uforklaret, tilegner sig \emph{det Sandse\-%
liges} Omraade, luttrer Begrebet \emph{Væren} til \emph{Ideen}. Og saa\-%
snart \emph{Ideens} Begreb er fundet: det Begreb, til hvilket selve
Stoffet viste hen som til sin Fuldendelse, bliver Ideen sat som
\emph{den sande, væsentlige Virkelighed}, og Tankens Opgave
bliver: \emph{at søge et Udtryk for Enheden af de to Mo\-%
menter, i hvilket denne Forudsætning er fastholdt}.
Men idet denne Opgave kommer frem, saa forudsætter den, at
\emph{Tanken bevidst} har \emph{sondret} de to Momenter, der \emph{for
Anskuelsen} ere forenede i Kunstværkets Enhed, og denne
Enheds Art umuliggjør, som det ovenfor blev paaviist, en \emph{ny}
Forening ved den græske Tanke, for hvilken \emph{Harmoniens
Umiddelbarhed} var Enhedens \emph{eneste} Form. Skjønhedens
Umiddelbarhed taaler ikke Tankens sondrende Magt, derfor
bringer Tanken Splid ind i det, hvorpaa den hviler, og gjen\-%
nem denne Splid maa denne Tankeform gaae tilgrunde. Mo\-%
menterne, der i deres Sondring \emph{begge} tiltvinge sig en Be\-%
stemmelse af \emph{Væren}, staae \emph{dualistisk} imod hinanden, og
den græske Philosophie maa ende i den frugtesløse Bestræbelse
for, efterat Bevidstheden om Momenternes Væsensforskjel er
vakt, atter at finde en Forsoning.

I det Fremstillede ligger, at den græske Philosophies Ud\-%
vikling maa gjennemløbe \emph{tre} Hovedstadier. Det \emph{første} bliver
det Stadium, paa hvilket Principet for Alt søgtes i \emph{Naturen,
i Stoffet}; det \emph{andet} det Stadium, paa hvilket \emph{Ideens} Be\-%
greb var fundet, og Enheden i det Værende søgtes ud fra
Forudsætningen af, at Ideen var den sande Virkelighed; det
\emph{tredie} endelig bliver det Stadium, hvor Enheden søgtes ud
% --- p. 37 ---
\setcounter{footnote}{0}
fra den klarere eller dunklere Bevidsthed om Momenternes
Splidagtighed. Søge vi Navne for disse tre Perioder, der
kunne betegne deres Forhold til det Centrale, der er Udvik\-%
% p.37 l.5: ERRATUM SITE -- the Trykfeil leaf asks for `Vorden' where the
% print has `Vordens'.  It stands as printed.  See transcription.tex §5.
lingens bevægende Magt, kunne vi benævne dem som \emph{Idee\-%
lærens Vordens, Ideelærens Udfoldelse og Ideelærens
Opløsning}.

Disse Udviklingens Hovedtrin faae nu deres bestemtere
Charakteer ved den Bevægelse, der foregaaer indenfor hvert
enkelt af dem.

Den \emph{ældste} Periode begynder med den begyndende Phi\-%
losophies \emph{første} Spørgsmaal, der er begrundet i Aandens
Stræben efter at drage Bevidsthedsindholdet ind under Selv\-%
bevidsthedens Enhed, med Spørgsmaalet: hvad \emph{er det Værende}?
den spørger om det blivende Ene i den vexlende urolige
Mangfoldighed. Svarene paa dette Spørgsmaal skride hurtigt
fremad: naar \emph{Jonierne} nævne et sandseligt Stof, et bestemt
Element eller det ubegrændsede, ubestemte Stoflige, nævne
\emph{Pythagoræerne} Tallet, det Stofliges Maal, \emph{Eleaterne} det
Værende i dets Enhed, og i disse Bestemmelser, der allerede
% p.37: the note below is another Greek site not listed in RECON.md §4.
% \gk{Κόσμος} at NOTE size; the acute over the omicron is clear at 1200 dpi
% and the word ends in a final sigma.
føre en Kreds af væsentlige Begreber\footnote{Allerede hos Pythagoræerne finde vi saaledes Begrebet om \gk{Κόσμος}.}, af afgjørende Proble\-%
mer, med sig, antyde de ubevidst de tre Sphærer, i hvilke
for den udviklede græske Tænkning Tilværelsen sondrer sig.

I Spørgsmaalet om, \emph{hvad Alt er}, ligger det \emph{andet}
Spørgsmaal: \emph{hvorledes er dette, som er det Værende i
Alt, blevet til alt det Mangfoldige?} --- med andre Ord:
hvorledes er Alt blevet til, hvorledes forklares Mangfoldigheden
i Tilværelsen? Dette Spørgsmaal kommer frem som et \emph{secun\-%
dært} Spørgsmaal i hine ældste philosophiske Forsøg, indtil
\emph{Eleaterne} bestemme Væren saaledes, at de ved deres Be\-%
stemmelse afskjære Spørgsmaalet om Vorden, uden dog at
kunne fortrænge Vorden fra den sandselige Forestilling, gjen\-%
nem hvilken den ligesom tillister sig en Væren. Men mod
% --- p. 38 ---
denne Væren, der ikke kan tilintetgjøre sin Modstander, stiller
saa \emph{Heraklit} selve Vorden, Tilblivelsens og Forandringens
Omskrivning, som Tilværelsesmangfoldighedens Forklarings\-%
grund, og denne \emph{umiddelbart} som \emph{sandseligt} Stof, som
Ild, anskuede Vorden, bliver ikke blot det Princip, \emph{hvorved},
men det, \emph{hvoraf} Alt \emph{er}; den bliver \emph{Kraft} og \emph{Stof}, og
% p.38: both guillemet pairs on the line below point INWARD, «…», read at
% 1200 dpi.  ABBYY gives the first pair as «Vei opad« -- an OCR error, not
% the compositor's.
bærer i sig, under sin «Vei opad» og sin «Vei nedad», som
det eneste Blivende og Faste, hvad Grækeren maatte fordre:
\emph{Forandringens lovmæssige, harmoniske Orden}. Men
\emph{Vordens Grund} forklaredes ikke, idet Vorden opfattedes \emph{ikke}
som det af \emph{en Væren Omfattede}, i en \emph{Væren} Begrundede;
men som sin \emph{egen Forudsætning}, som Tilblivelsens tau\-%
tologiske Omskrivning, og ligesaalidt forklaredes dens \emph{Regel
og Lov}, fordi Vorden, der kun sees som det \emph{Sandseliges
Vorden}, \emph{ikke} forklarer det \emph{iboende Maal}. Der stilledes
derfor en \emph{ny} Opgave for Tanken: \emph{at paavise Vordens
Grund og Nødvendighed}, og idet denne Opgave gribes af
Philosopher, der kjende \emph{Eleaternes} Lære, faaer den den be\-%
stemtere Form: \emph{i selve det Tilgrundliggende at søge
Vordens Grund uden at Værens Begreb opgives}. Og
Forsøget paa at løse denne Opgave bliver da med det Samme
et Forsøg paa at \emph{mægle mellem Eleaternes og Hera\-%
klits Grundbegreber}, under hvilket disse omformes. \emph{Em\-%
pedokles, Atomisterne og Anaxagoras} udgjøre den
Gruppe af Philosopher, der griber denne Opgave, og den sidste
af dem \emph{berører}, under Forsøget paa at løse den, Begrebet
% p.38: another Greek site not listed in RECON.md §4.  \gk{Νοῦς} -- the
% circumflex over the upsilon is a clear tilde at 1200 dpi.
om det, der er høiere end det Materielle: om \gk{Νοῦς}, som han
dog endnu ikke med Klarhed kan sondre qvalitativt fra det
Stoflige, eller, medens han tilskriver den Ordenen i Universet,
hævde en anden Virksomheds\emph{form} end Frembringelsen af den
mechaniske Bevægelse. Mæglingsforsøgene mislykkedes, idet
de hverken kunde hævde Væren eller forklare Vorden, og de
maatte mislykkes fordi de modsatte Principer skulde mødes
paa det \emph{samme} Tilværelses-Gebet, indenfor den \emph{sandse\-%
lige} Tilværelse. Og idet Mæglingsforsøgene mislykkedes,
% --- p. 39 ---
maatte de to Principer atter sondre sig abstract fra hinanden,
og i Gjennemførelsen af dem i denne skarpe Sondring, under
hvilken Eleaternes Enhedsprincip, ved at gjøre Væren be\-%
stemmelsesløs, altsaa lig Intet, umuliggjør en Erkjendelse af
den, medens Heraklits Vorden fører til samme Resultat: Er\-%
kjendelsens Tilintetgjørelse, ved at lade al blivende Bestemthed
forsvinde, opløses den græske Philosophies ældste Form og en
høiere Udviklingsform forberedes. Denne opløsende, og, indi\-%
rect, gjennem det Negative, forberedende Virksomhed tilfalder
\emph{Sophistiken}, der, som \emph{philosophisk} Phænomen, repræ\-%
senterer den ældre Philosophies Opløsning gjennem den en\-%
sidige exclusive Fastholden og dialektiske Gjennemføren af
dens modsatte Grundbegreber.

Den ældste græske Philosophie, i hvilken de mange for\-%
skjellige philosophiske Forsøg samle sig i \emph{fire} Grupper: \emph{den},
hvis Grundtanke er Væren; \emph{den}, hvis Grundtanke er Vorden;
\emph{den}, der søger at mægle mellem disse Modsætninger og at
bringe dem i Overensstemmelse med hinanden, og \emph{den}, der
ved disse Grundtanker som Modsætninger lader dette Udvik\-%
lingstrin af Philosophien opløse sig selv, --- betegne vi som
\emph{Ideelærens Vorden}, og vi søge Berettigelsen dertil i hvad
der allerede er antydet og i det Efterfølgende nærmere vil
blive paaviist: at det, der er \emph{Udviklingens Maal} indenfor
den græske Philosophie, gjør sig gjældende, \emph{ubevidst}, i den
endnu uudviklede Tænknings første Forsøg, og at der vises
hen til det gjennem deres \emph{Grundbegreber} og deres \emph{Pro\-%
blemer}, og gjennem \emph{Tænkningens Art} og \emph{eiendomme\-%
lige Charakteer}.

Trods det altsaa, at vi see den \emph{første} Periode søge det
Værendes Princip i \emph{Stoffet} og af den Grund blive \emph{Natur\-%
philosophie}, uden skarp Sondring af det \emph{Legemlige} og
det \emph{Aandelige}, og derfor uden Tankens Reflexion over \emph{Vi\-%
dens Betingelser}, saa er det dog den \emph{samme ideelle}
Stræben, der behersker den og den efterfølgende Periode, den
græske Philosophies Culminationstid, i hvilken, efterat \emph{Aan\-%
% --- p. 40 ---
\setcounter{footnote}{0}%
dens Tanke} var brudt frem, det \emph{Ideelle} sættes som \emph{det
Høieste}, og, gjennem Spørgsmaalet om \emph{Videns Betin\-%
gelser}, \emph{Ideen} bestemmes som \emph{Videns egentlige Gjen\-%
stand}, saa at de nye, ved Erkjendelsen af Ideen betingede
Videnskaber: \emph{Dialektiken og Ethiken}, faae Forrangen for
\emph{Physiken}.

Til denne \emph{anden} Periode regne vi ikke \emph{Sokrates},
skjøndt det var ham, der \emph{positivt} forberedte den ved i selve
Tanken, der i Sophistiken kun havde været nedbrydende, at
paavise det dybere Grundlag for Videnskab og Sædelighed.
Hos Sokrates var Dialektiken \emph{i det Ethikes Tjeneste}: idet
Tænkningen, med Inductionen til Redskab, fra det vexlende
Mangfoldige søger hen til det almene Væsen, der er udtrykt
i Begrebet, naaer den kun \emph{det Negative}, og ender ironisk
i Negativitetens Uendelighed uden atter at kunne gaae ud fra
denne, saa at Individet, hvis \emph{Væsens}bestemmelse er \emph{Viden},
gjennem Erkjendelsen af sin \emph{Uvidenhed} tvinges over i det
\emph{Ethisk-Religiøse}, hvor Tanken, med \emph{Uvidenhedens
Mellembestemmelse} som Forudsætning, finder \emph{et ved det
Guddommelige betrygget Indhold}\footnote{Jfr. Udførelsen af disse Bestemmelser i Afsnittet om Sokrates.}. Sokrates vælger
saaledes selv det Ethisk-Religiøse til sit Gebet, og stiller sig,
som uvidende, udenfor Philosophiens Gebet; men ved det, der
for ham kun var det Ethiskes Redskab: \emph{Methoden, Dia\-%
lektiken}, giver han Incitamentet til den græske Philosophies
hele senere Udvikling. Og idet den sokratiske \emph{væsentlige
Viden} afkaster \emph{Uvidenhedens} Form og bliver til \emph{den exi\-%
sterende Tænkers Viden af Begrebet}, og Begrebet,
der hos Sokrates, medens det sattes som \emph{Videns uopnaaede
Maal}, nærmest gjordes gjældende som det \emph{subjective Livs}
Regulator, bliver opfattet som \emph{metaphysisk Væsenhed},
som \emph{iogforsigværende objectiv Virkelighed}, som Na\-%
turens og Menneskelivets Væsen og Sandhed, fremgaaer, hos
% --- p. 41 ---
% p.41 is another Greek site not listed in RECON.md §4: \gk{Κόσμος}, acute
% over the omicron, read at 600 dpi.
\emph{Plato, Ideelæren} af den sokratiske \emph{forsøgte Begrebs\-%
dannelse}, og \emph{Speculationen}, som den græske Aand, der
fandt sit \emph{Væsen i Viden}, ikke kunde opgive, hævder atter
sin Gyldighed.

I \emph{Ideebegrebets} Genesis ligger Muligheden til en
\emph{dobbelt} Bestemmelse af Ideens Væsen og af dens Forhold
til Stoffet. Idet Tanken, der finder Ideen, i Ideen gjenfinder
\emph{sig selv}, og nu, opgivende Stoffet som Princip, bestemmer
\emph{Ideen som det sande Værende}, som det, der giver Til\-%
værelsen Væren, som det Evige, der ikke berøres af Vorden,
% p.41: the letterspacing stops INSIDE the broken word: `Aa-' at the end of
% the line below is spaced, `benbaring,' on the next line is not.  Hence the
% \emph{} closes on the hyphen.
saa bliver der paa dette Trin, ved denne Ideens \emph{første Aa\-%
}benbaring, ingen anden Bestemmelse for det fra Ideen for\-%
skjellige Moment, end det \emph{Ikke-Værende}. Ideen sondrer
sig skarpt fra dette i Væsen og en gabende Kløft adskiller
det Aandelige, Erkjendelsens usandselige Gjenstand, fra det
Materielle og som saadant Ikke-Erkjendelige. Men dette \emph{Stof\-%
agtige}, der er \emph{Sandseverdenens Grundlag}, kan Tanken,
om den end, ligesom gjennem en Subtraction af den Væren,
Sandseverdenen skylder Ideen, fra Sandseverdenens Væren, be\-%
stemmer det som et \emph{Ikke-Værende}, dog ikke fastholde i
denne Bestemmelse, ikke trænge udenfor Væren. Idet den
endelige Tilværelse skal begribes, hævder det sin Plads som
\emph{medbestemmende} Factor i denne, og nu kommer den
\emph{Modsigelse} frem, \emph{at det Ikke-Værende dog faaer en
Art af Væren og Virksomhed}, og en \emph{uforklaret} Phæ\-%
nomenernes Verden med en \emph{ufuldkommen} Væren, der dog
\emph{ikke} kan frakjendes Charakteren af \emph{Væren}, kommer til, seet
\emph{i sin Sandselighedsbestemthed}, at sondre sig fra Idee\-%
verdenen med dens fuldkomne og fuldstændige Væren; medens
i Anskuelsen af \gk{Κόσμος} den reale Tilværelses Heelhed sees i
harmonisk Enhed med Ideeverdenen. Disse Vanskeligheder
ere de, ved hvilke \emph{Plato} standser, og de føre \emph{Aristoteles}
til at opstille en anden Formel for Forholdet, der lader \emph{baade}
Ideen og dens Modsætning, \emph{baade} Formen og Stoffet, som
han benævner dem, opfattes under \emph{Værens} Bestemmelse,
% --- p. 42 ---
\setcounter{footnote}{0}
% p.42: \gk{ἐνεργείᾳ} -- the breathing over the epsilon is a comma-shaped
% SMOOTH breathing (plainly not the "C" of ὕλη on p. 35); acute over the
% iota; the final alpha carries an IOTA SUBSCRIPT, clearly visible at
% 1200 dpi.  \gk{δυνάμει} -- acute over the alpha, no breathing.  Both read
% at 1200 dpi.  (This occurrence of ἐνεργείᾳ is correctly set; the γ-for-ν
% misprint of the Trykfeil leaf is at p. 99, not here.)
men indenfor dette Begreb gjør en Adskillelse mellem en
Væren \emph{efter Virkelighed} (\gk{ἐνεργείᾳ}) og en Væren \emph{efter
Mulighed} (\gk{δυνάμει}). Derved betegnes den \emph{oprindelige
Væsenssammenhæng}. Stoffet er Muligheden til det, der
i Formen er blevet til udfoldet, sand Virkelighed, og Stoffet
attraaer Formen, stræber hen imod den som mod sit Væsen,
medens omvendt Formen er den \emph{iboende} Form, der som
% p.42: PRINTED AS SHOWN -- the `i' of `virker i Stoffet' below is set from
% a HEAVIER (bold) fount than the rest of the line: its stem and its dot are
% markedly thicker and wider than the `i' of `virker' immediately before it.
% A WRONG SORT, not emphasis (bold does not otherwise occur inside a
% paragraph in this book, RECON.md §2).  Transcribed as an ordinary `i'.
Entelechie virker i Stoffet, som Udviklingsprincip fører det til
Virkelighed og i Verdenstilværelsens Heelhed virker som det
harmoniske Verdensalts guddommelige Enhedsprincip. Men
ogsaa hos \emph{Aristoteles} kan Systemet ikke afsluttes, Ener\-%
% p.42: \gk{Νοῦς} twice, circumflex over the upsilon, clear at 600 dpi.
giens sidste Grund ikke bestemmes, uden at \emph{Formen i sin
Reenhed}, som guddommelig \gk{Νοῦς} og som activ \gk{Νοῦς} \emph{i
Mennesket, sondrer sig fra alt} Stof og fra al Bevægelse,
og ved denne Sondring, paa det Punkt hvor Væsenet \emph{renest}
fremtræder, faaer \emph{Dualismen} atter \emph{Herredømmet}.

Gjennem disse to Opfattelser af Ideen i dens Forhold til
det Stoflige udfolder \emph{Ideelæren} sig, og \emph{kun} disse to Op\-%
fattelser blive mulige hvor de \emph{lige oprindelige} Tilværelses\-%
% =====================================================================
% RUN-OVER FOOTNOTE -- REPORT THIS.  The note attached below BEGINS on
% p. 42 and CONTINUES ON p. 43, where the note area opens WITHOUT a mark
% and runs:
%   "faaer Bestemmelsen τὸ μὴ ὄν, der bliver F o r u d s æ t n i n g, ikke
%    Resultatet af en Analyse af et Væsensgyldigt, en anden Belysning, og
%    Materiens Stilling til det Ideelle bliver, trods Oprindeligheden, en
%    anden end hos Plato."
% Only the p. 42 portion is transcribed here, because this fragment covers
% pp. 29--42 only; the pp. 43--56 batch must append that continuation
% INSIDE this \footnote{} (and must NOT open a new note on p. 43).
%
% THIS CONTRADICTS RECON.md §6 AND BATCH-AGENT.md §5, both of which state
% that no run-over note exists in this book.  p. 42 was not among the nine
% candidates they checked.  Batch seams at note-bearing pages are therefore
% NOT free.
% =====================================================================
momenters Forhold\footnote{Trods \emph{Platos} Bestemmelse af Materien som det Ikke-Værende, ad\-%
skiller hans Opfattelse af den sig dog væsentlig fra de spiritualistiske
og dualistiske Systemers i den efteraristoteliske Tid. Platos Be\-%
% p.42 note: \gk{μὴ όν} twice, at NOTE size.  The grave over the eta is
% resolvable; the breathing over the omicron is not -- RECON.md §4 warns
% that the smooth/rough test is unproven at note size, and here the left
% element of the compound is a shapeless hook even at 1200 dpi.  The omicron
% is set with the acute alone and the breathing is not supplied; ὄ/ὅ
% undecided at both sites.  Cf. p. 36, where the word is set legibly at body
% size and reads smooth.
stemmelse af Materien som et \gk{μὴ όν} er \emph{Resultatet} af hans Ana\-%
lyse af Sandseverdenens Væsen; men netop gjennem Anskuelsen af
\emph{Verdensheelheden}, fra hvilken Plato kommer til Bestemmelsen
af Materiens Begreb, faaer Materien, om den end sættes som et
\gk{μὴ όν}, fra først af \emph{Gyldighed}, og \emph{lige Oprindelighed} med
Modsætningen. I de \emph{spiritualistiske} Systemer bliver Materien
deduceret af den ideelle Enhed som dens svageste Udtrædelsesform,
og idet den som saadan bestemmes som et Ikke-Værende, bliver den
altsaa ikke sat i lige Oprindelighed med hiin. I de \emph{dualistiske}
Systemer, i hvilke Opfattelsen af Verdenstilværelsen er en anden,
idet Anskuelsen af \gk{Κόσμος} i dens uforstyrrede Herlighed viger for
Anskuelsen af Verdenstilværelsen som en ufuldkommen, splidagtig
Væren; i hvilke Forestillingen om en \emph{dobbelt Verdenssjæl} faaer
Magt, og den \emph{transscendente Gud} bliver Materiens Modsætning,
% --- the note continues on printed p. 43, where the note area opens with
% no mark.  Joined here, at the call, per transcription.tex §4.  The p. 43
% portion was read on the image at 500 and 1200 dpi by the pp. 43--56
% batch, which confirmed the pp. 29--42 batch's quotation word for word
% and added that `Forudsætning' is letterspaced.
% \gk{τὸ μὴ όν}: the graves over the omicron of τὸ and the eta of μὴ are
% both plainly resolvable (they match the certain grave of τὸ πλέον on
% p. 52 and are visibly not the up-right acute of πλέον's epsilon).  Over
% the omicron of ὄν the mark is the compound blob RECON.md §4 warns of:
% the acute is there, the left element is a shapeless hook, and ὄ/ὅ cannot
% be separated even at 1200 dpi.  The acute alone is set and the breathing
% is NOT supplied -- as at this same note's two earlier occurrences.
faaer Bestemmelsen \gk{τὸ μὴ όν}, der bliver \emph{Forudsætning}, ikke
Resultatet af en Analyse af et Væsensgyldigt, en anden Belysning, og
Materiens Stilling til det Ideelle bliver, trods Oprindeligheden, en
anden end hos Plato.}
% =====================================================================
% SEAM at the foot of p. 42.  The letterspaced run and
% the broken word BOTH continue onto p. 43.  The print reads
%   "... skal bestemmes u d  f r a  F o r u d s æ t -
%    n i n g e n  a f  d e t  I d e e l l e s  V æ s e n s p r i n c i p a t ."
% two \emph{} groups either side of the page break have been MERGED into
% the single run the print shows, and the broken word joined.  Verified
% letter-gap by letter-gap at 500 dpi by the pp. 43--56 batch.
% =====================================================================
skal bestemmes \emph{ud fra Forudsæt\-%
% =====================================================================
% (The p. 43 continuation of this note has been merged into the p. 42
% \footnote{} at its call site; see the comment there.)
% =====================================================================
% SEAM at the head of p. 43 -- the other half of the p. 42 seam note.
% p. 42 ends "skal bestemmes \emph{ud fra Forudsæt\-}%" and p. 43 opens
% with the rest of the broken word, still letterspaced.  The calling
% conversation must DROP the closing brace the p. 42 fragment put after
% "Forudsæt\-" and DROP the \emph{ that opens the first line below, so
% that the joined text reads
%   skal bestemmes \emph{ud fra Forudsæt\-%
%   ningen af det Ideelles Væsensprincipat}.
% The group is closed and reopened here only so that this fragment's
% braces balance on their own; the two \emph{} groups are ONE run in the
% print, verified letter-gap by letter-gap at 500 dpi.
% =====================================================================
% --- p. 43 ---
% p. 43 bears NO note of its own.  Its note area opens WITHOUT a mark and
% carries only the tail of the p. 42 note (set above); hence NO
% \setcounter{footnote}{0} here and no \footnote{} on this page.  The
% left-flush hairline above that note area is the footnote separator and
% is not transcribed.
ningen af det Ideelles Væsensprincipat}. \emph{Enten} bliver
det Ideelle det \emph{Eneste}, der faaer Værens Bestemmelse, det
\emph{Eneste}, der giver Tilværelsen Virkelighed og Sandhed, \emph{eller}
det bliver det, der ene er et Værende i \emph{høieste} Betydning,
men ved sin Natur staaer i nødvendigt Forhold til Stoffet som
til det ved Existentsformen Forskjellige, ved Væsensbeslægtet\-%
heden Lige. Men naar den \emph{første} Opfattelse ved sin \emph{abso\-%
lute} Modsætning mellem det Værende: Ideen, og det Ikke-%
Værende: Stoffet, maatte føre til, at Sandseverdenens Existents
blev uforklarlig og Materiens Væren modsigende, saa viser den
\emph{anden}, der, ved at opfatte begge Modsætningerne under
Værens Bestemmelse, undgaaer hins indre Modsigelse, og ved
at søge det Ideelle i selve den reale Virkelighed, kan aner\-%
% p.43, l. 15: "at den" is NOT letterspaced.  The a--t gap measures with
% the ordinary body gaps, well below the run of "a n d e n" two lines
% above, and spacing.py flags no run here.  Left plain.
kjende denne i dens Væsentlighed, en \emph{dobbelt}, i selve den
græske Tænknings Væsen begrundet Mangel: at den Formens
Virksomhed i Stoffet, ved hvilken Enheden skulde aabenbare
sig, og Tilværelsens Concretion forklares i sin Tilbliven, dog
bliver ubegreben, og at med det Samme den energisk tilstræbte
og ligesom berørte principielle Enhed atter opløser sig i Dua\-%
lisme. Men idet den Aristoteliske Opfattelse af Principernes
Forhold, der vil omfatte de fra hinanden sondrede Modsæt\-%
ninger i Værens Enhed og derved betegner den græske Phi\-%
losophies høieste Stræben, er den sidste alsidig græske fordi
det er den eneste ægte græske ved Siden af Platos, bliver
\emph{Philosophien efter Aristoteles}, der med klarere eller
dunklere Bevidsthed om Bruddet stræber efter en, den ene af
Modsætningerne opoffrende, og derfor Modsætningerne ikke i
Sandhed forsonende Enhed, \emph{den græske Philosophies Op\-%
løsningsperiode}.

I sin Stræben efter en \emph{principiel} Enhed udvikler den
% --- p. 44 ---
% p.44: the printed folio is garbled in the ABBYY layer (it comes through
% as an apostrophe); the image reads "44" plainly, and pagemap.py gives
% PDF 53.  No note, no rule, no Greek on this page.
\emph{efteraristoteliske Philosophie} sig fra en \emph{Materia\-%
lisme}, gjennem en \emph{principielt udtalt Dualisme} som
kortvarigt Overgangsled, til en \emph{spiritualistisk Enhedslære},
der magtesløst stræber ud over den græske Aands Grændser,
og ender i phantastisk Scholastik.

I det almindelige Forhold mellem den græske Tænknings
Grundbegreber, og i den eiendommelige Retning, som Philo\-%
sophien havde taget hos \emph{Aristoteles}, laae Grunden til, at
hiin Stræben efter en principiel Enhed maatte begynde \emph{fra
Materialismen}. Denne Materialisme er, i væsentlig for\-%
skjellige Former, udtalt i \emph{Stoicismen} og \emph{Epikuræismen},
den skinner frem under \emph{Skepticismen}. Den \emph{principielle
Dualisme}, der, gjennem Modsætningen mellem den \emph{trans\-%
scendente aandelige Gud} og \emph{Materien}, forbereder den
% p.44: throughout this page the compositor leaves the conjunction "og"
% UNSPACED inside an otherwise letterspaced run -- at "Gud og Materien"
% (l. 14), "ræerne og Philo" (l. 23) and "alt det Værende og selve
% Materien" (l. 32).  Measured at 500 dpi the o--g gap is the ordinary
% body gap in each case, against 2--3 times that inside the neighbouring
% words.  Set as printed, i.e. as separate \emph{} groups; not
% regularised into one run.  (Contrast p. 44 l. 29 and p. 51 l. 6, where
% "og" IS spaced and the run is unbroken.)
\emph{spiritualistiske Monisme}, afløser ved den christelige Tids Be\-%
gyndelse hiin Materialisme, hos hvis dogmatiske Repræsentanter
Dualismen havde viist sig uovervunden: hos Stoikerne allerede
ved selve Principets Bestemmelse, hos Epikuræerne i Anthro\-%
pologien, hvor den Modsætning, der kunde holdes skjult i den
overfladisk behandlede almindelige Deel af Systemet, brød frem
gjennem de concretere Bestemmelser. Den principielt udtalte
Dualisme, hvis væsentligste Repræsentanter vare \emph{Nypythago\-%
ræerne} og \emph{Philo}, blev, netop fordi den var Dualisme, af
ringe Varighed. Den maatte atter vige Pladsen for en \emph{mo\-%
nistisk} Philosophie, der nu, efterat de andre Muligheder vare
forsøgte, maatte blive af \emph{spiritualistisk} Art, og som fuld\-%
endtes i \emph{Nyplatonismen}, i hvilken den græske, af Orien\-%
tens Aand paavirkede Tanke, af den \emph{rent aandelige, over
Tænken og Væren transscendente Guddom}, til hvis
\emph{forskjelsløse Enhed} Mennesket søgte tilbage for i \emph{eksta\-%
tisk Hengivelse} at guddommeliggjøres, stræbte at udlede
\emph{alt det Værende} og selve \emph{Materien}. Under Kampen mod
Christendommen, der consequentere gjennemfører det af Nyplato\-%
nismen Tilstræbte, afslutter denne den \emph{efter sit hele Anlæg}
fuldstændig udfoldede græske Philosophie.
% --- p. 45 ---
% p.45 is one of the pencil-pollution sites of BATCH-AGENT §7.  Three
% marks were found and NONE is transcribed: a small dot between "Lys,"
% and "seer" (l. 7 -- ABBYY renders it "Lys ^"), a dot inside "Enheden"
% (l. 13 -- ABBYY "Enhed'en"), and a faint tick after "Jorden." at the
% foot of the second paragraph.  None is type; the printed readings are
% "Lys, seer", "Enheden" and "Jorden.".

I hele denne Udvikling, i hvilken intet Udviklingstrin er
oversprunget, bliver eet Grundforhold det Ledende og Be\-%
stemmende, Udfoldelsens og Opløsningens Kilde. --- Be\-%
gyndende i Naturen stræber Tanken hen imod at finde sig
selv; den anskuer sig i Skikkelsen af den naturlige Tilværelses
harmonisk formende Idee-Princip, den seer Tilværelsens Skjøn\-%
% p.45: \gk{τέλος} -- acute over the epsilon, no breathing, final sigma
% ς; read at 2500 dpi (500 dpi render enlarged 5x).  Not letterspaced.
hedshele i Aandens Lys, seer Naturen som idealiseret i dens
Stræben hen mod Menneskeaanden, dens \gk{τέλος}; den sondrer
de Momenter, som Anskuelsen af hint Hele umiddelbart for\-%
binder; den søger, under den dobbelte Form, der paa dette
Standpunkt er mulig, at bestemme deres Forhold, og idet de,
engang sondrede ved Tanken, ikke atter lade sig føre tilbage
af den til Enhed, søger den Enheden gjennem en Underordnen
af det ene under det andet. Den gjør Forsøget med en ma\-%
terialistisk Enhed, men idet denne brister, kaster den sig,
inden Forsøget gjøres fra det modsatte Udgangspunkt, for et
Øieblik over i Dualismen, der netop paa den Tid træder
mægtig op udenfor dens Gebet; men den finder ikke Hvile i
dens Modsigelse; den samler da sine Kræfter til et sidste
Forsøg, og vover sig ikke blot dristigt ud over den græske
Aands Grændser, men den stræber ud over Tanken selv, søger
den Forsoning, som Tanken og den virkelige Tilværelse med
dens uforsonede Splid ikke giver den, over Tanken og over
Virkeligheden; den drømmer om en Aandens Magt, der ikke
er Tanke, men mere end Tanke; idet den stræber mod hvad
der er Gjenstand for dette Høiere, forvandles den til \emph{Phan\-%
tasie}; med Phantasiens Ikarusflugt bæres den op over den
endelige Tilværelse, over Gudeverdenens Vrimmel, over selve
% p.45: \gk{Νοῦς} -- capital nu, circumflex (a clean tilde) over the
% upsilon, final sigma; unmistakable at 2500 dpi and identical with the
% setting at p. 42.  Not letterspaced.
\gk{Νοῦς}, den styrter sig blændet mod det uendelige formløse
Enes favnende Dyb, og synker med mattede, forbrændte Vinger
udtømt og døende tilbage til Jorden.

% DIVISIONAL RULE, p. 45, between "... tilbage til Jorden." and "Hvad
% der gjælder ...": short, CENTRED (spans x738--1120 at 500 dpi against a
% measure centre of x1105 -- i.e. centred on the measure, not flush left),
% with white space above and below.  No display head follows it.  It is
% the book's divisional rule, not the footnote separator; this page bears
% no note.
\divrule

Hvad der gjælder om den \emph{græske} Philosophie, hvis Ud\-%
foldelses klare Simpelhed gjør Udviklingsgangen overskueligere,
% --- p. 46 ---
vil ogsaa stadfæstes ved den \emph{nyere} Philosophie. Men da
Fremstillingen af dennes Udvikling staaer i nøie Forbindelse
med Synspunkter, der først senere kunne fremhæves, vil Paa\-%
% p.46, l. 4: ERRATUM OF THE TRYKFEIL LEAF, NOT APPLIED.  The print
% reads "lovbestemte Følge blive given"; the author's leaf (PDF 306)
% asks for "først blive given".  The printed reading is kept, per
% transcription.tex §5 and BATCH-AGENT §8.  Confirmed on the image.
visningen af dens Systemers lovbestemte Følge blive given
efterat disse Synspunkter ere gjorte gjældende.

% DIVISIONAL RULE, p. 46, between "... ere gjorte gjældende." and "Ved
% det Foregaaende ...": short, CENTRED, white space above and below, no
% display head under it.  Not the footnote separator; this page bears no
% note.
\divrule

Ved det Foregaaende er der saaledes med Hensyn til Epo\-%
chernes Rækkefølge og Eiendommelighed, og med Hensyn til
Periodernes og Systemernes Orden indenfor den første Epoche
paaviist en Lovmæssighed og Nødvendighed, begrundet i Tan\-%
kens Princip, der bestemmer Tankens Tidsudvikling, og derved
er den \emph{første} af de ovenfor stillede Opgaver løst indenfor de
der angivne Grændser. I denne Udvikling har Principet alle\-%
rede viist sig som en \emph{teleologisk} bestemmende Magt, der
gjennem Menneskeaandens Tidsudvikling giver sig Virkelighed
i Menneskeaanden. Men denne teleologiske Virksomhed bliver
nu --- og dette var den \emph{anden} Opgave, der stilledes --- til\-%
lige at see fra et andet, allerede ovenfor berørt, Synspunkt:
fra Synspunktet af \emph{Forholdet mellem det Bevidste og
det Ubevidste}, mellem den \emph{individuelle} Tankes \emph{bevidste
Stræben} og det \emph{overgribende Almenes, Individet par\-%
tielt eller totalt ubevidste, Magt}, der, gjennemvirkende
den individuelle Tænken, virker bagved og gjennem hiin be\-%
vidste Stræben, fremskyndende den eller hæmmende den. I
al teleologisk Udvikling er, paa Udviklingens fremskridende
% p.46: \gk{τέλος} twice on this page (here and eight lines below), both
% set exactly as at p. 45 -- acute over the epsilon, no breathing.  The
% Greek is not letterspaced; spacing.py's flag on it is the artefact
% BATCH-AGENT §4 rule 5 describes.
Stadier, Udviklingens \gk{τέλος} \emph{ideelt tilstede i sin Heelhed},
som det Udviklingsgangen Bestemmende og Ledende, og denne
Maalets ideelle Tilstedeværen glimter ligesom frem, medens det
i det enkelte Stadium kun er \emph{realt} i den til Stadiet svarende
% p.46: PRINTED AS SHOWN -- in "Bevidsthedsudviklingen" the compositor
% letterspaces only the first element of the compound: "B e v i d s t h e d s"
% is spaced (55 px per letter at 500 dpi) and "udviklingen" is set solid
% (29 px per letter, the same as the surrounding body).  The break is
% mid-word, not at a line end.  Transcribed as printed.  The same habit
% recurs at p. 51 ("T a n k e omsætningen").
\emph{begrændsede} Form. Anvendt paa \emph{Bevidstheds}udviklingen
% p.46: "gjennem" alone is letterspaced in the line below; "det for
% Bevidstheden Værende," is set solid, checked letter by letter at
% 500 dpi.  The emphasis on a bare preposition is the author's, and
% recurs at p. 44 ("fra Materialismen").
vil dette da sige, at det \gk{τέλος}, der i sin Heelhed ideelt er tilstede
som styrende Udviklingen, paa de enkelte Bevidsthedsstadier
aabenbarer sig \emph{gjennem} det for Bevidstheden Værende,
medens det kun er realiseret for Bevidstheden i denne \emph{be\-%
% --- p. 47 ---
stemte} Form, saaledes at altsaa det, der i sidste Instants be\-%
stemmer Bevidstheden, ligger bagved Bevidstheden, som skjult
for Bevidstheden, hiin Aabenbaring altsaa er \emph{ubevidst}. Føies
nu den Bestemmelse til, at dette Styrende er et \emph{universelt},
ikke et \emph{individuelt} \gk{τέλος}, saa bestemmes Forholdet nærmere
saaledes, at den individuelle Bevidsthedsudvikling, selv hvor
dens \emph{individuelle} \gk{τέλος} er realiseret, dog staaer i Forhold
til det, der, \emph{ubevidst} for Individet, gjennemvirker Individets
Bevidsthed. Saadant bliver altsaa Forholdet hvor det Bestemmende
\emph{er det universelle Tankens Princip}, der styrer den in\-%
dividuelle Tænkning. Tankens Princip er \emph{for} Tanken, men
\emph{kun partielt}, efter en \emph{enkelt} Bestemthed, et \emph{enkelt}
Grundforhold, eller paa \emph{abstract} Maade. Men virkende \emph{som
concret Totalitet} efter sin \emph{metaphysiske} Nødvendighed
er det tilstede, \emph{ubevidst for Individet}, som det tænkende
Individs \emph{Skjæbne}, der forklares, men \emph{kun relativt} forklares,
til \emph{bevidst} Aandslov, saa at ogsaa selve \emph{Tankeudviklin\-%
gens Lovmæssighed}, der nærmest forholder sig til \emph{Be\-%
vidsthedens} Bestemmelse, faaer det \emph{Ubevidste} som væ\-%
sentlig Bestemmelse i sig.

Det \emph{almene} Forhold mellem det Ubevidste og Bevidste
i Tankeudviklingen, der udtrykker Forholdet mellem det tæn\-%
kende Individ og det Udviklingen teleologisk bestemmende
Princip, bliver nu at see som trædende ud i concretere For\-%
mer, idet de særegne Forhold, det indeslutter, sondre sig fra
hverandre, og som begrundende en \emph{Flerhed af Udviklings\-%
love}, hvis Gyldighed bliver at paavise i historiske Exempler,
der, som det tidligere blev antydet, ville blive saaledes valgte,
at gjennem dem væsentlige, hidtil utilstrækkeligt oplyste Punk\-%
ter i Philosophiens Historie kunne stilles i et nyt og klarere
Lys.

Idet Vexelforholdet mellem det Bevidste og Ubevidste er
seet i dets Sammenhæng med den \emph{teleologiske} Udvikling,
der har en \emph{Tidssuccession}, vise dets concretere Former,
der paa den ene Side ere betingede ved Forholdet mellem det
% --- p. 48 ---
Universelle, det relativt Almene og det Individuelle, sig paa
den anden Side at være det ved \emph{Tidsforholdene}. Der
bliver et \emph{firedobbelt} Forhold at fremhæve, og med det be\-%
grundes en \emph{firedobbelt} Lov, der udtrykker det i Forholdet
\emph{Nødvendige}.

Vi fremhæve \emph{først} det i selve Begrebet \emph{Udvikling} lig\-%
gende almindelige \emph{Forhold mellem det Sildigere og
Tidligere}, ifølge hvilket hint sees som bestemmende dette
og som aabenbarende sig gjennem det. I det i Udviklingens
Begreb indesluttede almene Forhold mellem Anlæg, Mulighed
og det ideelt Præexisterende, efter Tidsforholdet \emph{Fremtidige},
ligger det, at det, der endnu er et Fremtidigt, ideelt maa
virke, og som virkende \emph{aabenbare} sig som nærværende i det
i Tiden Forudgaaende, der teleologisk er anlagt paa det. In\-%
denfor den \emph{bevidste Tankeudvikling} maa dette Forhold,
efter det allerede ovenfor Berørte, gjøre sig gjældende i den
Form, at det, der er Udviklingens \gk{τέλος}, og til hvilket Be\-%
vidstheden, ifølge sit Væsen som Bevidsthed, som værende i
Almindelighedens Form, altid implicite forholder sig om end
% p.48, ll. 20 f.: ERRATUM OF THE TRYKFEIL LEAF, NOT APPLIED.  The print
% reads "gjennem sit partielt realiserede τέλος"; the author's leaf asks
% for "gjennem sit (individuelle) τέλος, der partielt realiserer det".
% The printed reading is kept, per transcription.tex §5.  "partielt" and
% "realiserede" are NOT letterspaced here, checked at 500 dpi.
indenfor en \emph{Abstractions} Grændser gjennem sit partielt
realiserede \gk{τέλος}, \emph{ubevidst} maa finde sit Udtryk gjennem det
\emph{bevidst} Tænkte og Udtalte; saa at altsaa en Udviklings\-%
rækkes Grundtanke under sin Vorden, som det endnu for det
tænkende Individ Ubevidste, præformerer og anticiperer sin
udfoldede og fuldendte Bestemthed.

% p.48: a small dark speck stands between "Sildigere" and "fra" in the
% second line of the paragraph below (ABBYY and tesseract both ignore
% it).  Reddish-brown, the colour of the paper's foxing and not of the
% type; not transcribed.
Vi fremhæve dernæst, idet vi see Forholdet \emph{mellem det
Tidligere og det Sildigere fra et nyt}, ved det \emph{relativt
Almenes} Forhold til det \emph{Universelle} og det \emph{Individuelle}
bestemt Synspunkt, det Forhold, ifølge hvilket et endnu ikke
naaet mere fremskredet Udviklingstrins Princip, der antydes
paa et tidligere Stadium, atter \emph{maa} trænges tilbage og for\-%
dunkles. I dette Forhold virker det realiserede \emph{bestemte}
Stadiums \emph{eiendommelige, for Bevidstheden frem\-%
traadte} Princip, --- samtidigt med at det selv, ifølge sin Sam\-%
menhæng med den \emph{totale} Udviklings Princip, ved sin \emph{indre
% --- p. 49 ---
Universalitet inciterer Individets Stræben ud over
Stadiets og dets Princips Grændse}, --- som \emph{hævdende
sig i sin Eiendommelighed}, og bliver den individuelle
bevidste Stræbens \emph{ubevidst} hæmmende Magt, men bliver det
saaledes, at det, gjennem denne Inciteren og Hæmmen, bevir\-%
ker sin egen \emph{fuldstændigere} Virkeliggjørelse og sikkrer Ud\-%
viklingen af den \emph{høiere} Form. Ogsaa dette Forhold er et
ved Udviklingens Væsen \emph{nødvendigt}, og vil derfor finde sit
Udtryk i en tilsvarende Lov.

% p.49: "Tanke-Udviklingen" below carries a PRINTED hyphen, not a
% line-break hyphen: the break falls at the hyphen, but the second
% element is capitalised ("Udviklingen"), so the compositor set the
% compound and not a broken "Tankeudviklingen".  (The solid form
% "Tankeudviklingen" does occur, at p. 47.)  Both readings are the
% author's; neither is normalised.  Set with a real hyphen and a % to
% eat the newline, so that no space is introduced.
Vi fremhæve dernæst det Forhold, i hvilket \emph{det i Tanke-%
Udviklingen Tidligere, efterat Udviklingen er skre\-%
det videre frem, men endnu ikke har naaet sit afslut\-%
tende Stadium, ud fra sin oprindelige Eiendomme\-%
lighed paavirker det Sildigere.} Idet nemlig det bestemte
sildigere Udviklingsstadium staaer i Relation til det for\-%
udgaaende, af hvilket det er fremgaaet, og hvis Bevidstheds
Indhold det har optaget i sin, men saaledes, som det \emph{kan op\-%
tages paa dette bestemte} nye Udviklingsstadium, saa \emph{maa},
idet det tænkende Individs \emph{bevidste} Stræben gaaer ud paa
at hævde det nye Stadiums Princip i Kamp mod det gamles,
% p.49: the letterspaced run below stops at "for"; "Individet," on the
% next line is set solid, checked at 500 dpi.  The sense wants "ubevidst
% for Individet" as one unit, but the print does not give it, and it is
% not regularised.
dette, \emph{som endnu ikke fuldt behersket, ubevidst for}
Individet, virke hæmmende ind paa dets bevidste Stræben.
Men denne Hæmmen bliver, ifølge Sammenhængen mellem
Fortidsprincipet og den omfattende Totalitet, i hvilken det er
Moment, naar den sees fra det Universelles Synspunkt, at op\-%
% p.49: only "fremskridende" is letterspaced in the line below --
% "Udviklings Moment," is set solid.  A small reddish speck stands after
% the comma of "Moment,"; ABBYY reads it as a dash ("Moment,- som").  It
% is foxing, not type, and is not transcribed.
fatte som en \emph{fremskridende} Udviklings Moment, som di\-%
recte eller indirecte fremkaldende \emph{en dybere og alsidigere
Tankeudvikling}, en Tilnærmelse til Tankens \gk{τέλος}. Dette
% p.49: PRINTED AS SHOWN -- "Udviklingstadiers" with a single s, where
% the sense and the author's usage elsewhere want "Udviklingsstadiers"
% (cf. "Udviklingsstadium" nine lines above on this same page).  Both
% OCR witnesses read it so and the image confirms it at 500 dpi.  Not on
% the author's Trykfeil leaf.  Transcribed as printed.
Forhold er med Nødvendighed begrundet i de eiendommelige
Udviklingstadiers Forhold til hinanden og til Totaliteten.

Vi fremhæve endelig Forholdet mellem det i den teleolo\-%
giske Udvikling Samtidige og Sideordnede. Idet Philoso\-%
phiens teleologiske Udvikling fører til en Articulation i \emph{coor\-%
dinerede og samtidige} Udviklingsformer, maa \emph{Vexelfor\-%
holdet} mellem disse Former nødvendigen blive et \emph{bevidst-%
% The sheet signature `4' stands at the foot of p. 49; not transcribed.
% p.49/50: "bevidst-ubevidst" is a compound with a PRINTED hyphen that
% falls at the page break; both elements are letterspaced and the run is
% one \emph{} group.  The hyphen is followed by a % so that no space is
% introduced across the page marker.
% --- p. 50 ---
ubevidst}, idet hvert af de modsatte Principer, medens de
% p.50, l. 2: ERRATUM OF THE TRYKFEIL LEAF, NOT APPLIED.  The print
% reads "bevidst hævde sig imod hinanden"; the author's leaf asks for
% "hævdes".  The printed reading is kept, per transcription.tex §5.
\emph{bevidst} hævde sig imod hinanden, udtrykker \emph{den begge
omfattende Heelhed}, og derfor i Udviklingen \emph{ubevidst}
bestemmes af Modsætningen, saaledes, at under den gjensidige
\emph{Bestemmen et mere Omfattende, skjult for Indivi\-%
derne, forbereder Fremtidsudviklingen.}

Ved \emph{Tidsforholdene} i den teleologiske Udvikling betinges
\emph{altsaa en Række af afledede Udviklingslove}, hvilende
\emph{paa eet og samme almene} Grundforhold: \emph{det Univer\-%
selles teleologiske, ved Nødvendighed bestemte,
Selvvirkeliggjørelse gjennem den individuelle
Tankeactivitet}, og disse afledede Love blive nu at oplyse
og at paavise i deres Gyldighed ved Exempler.

% DIVISIONAL RULE, p. 50, between "... ved Exempler." and the L2 head
% below: short, CENTRED, white space above and below.  Not the footnote
% separator; this page bears no note.
\divrule

% L2 HEAD, p. 50 (transcription.tex §3): the FIRST of the book's four
% theses, I.--IV. (pp. 50, 79, 137, 170).  BOLD antiqua, preceded by the
% divisional rule above.  It is set as a bold PARAGRAPH on the full
% measure -- lines 1 and 2 justified, line 3 centred -- broken in the
% print as
%   "I. En Udviklingsrækkes Princip præformerer og an-"
%   "ticiperer under sin Vorden, de tænkende Individer"
%   "ubevidst, sin udfoldede Virkelighedsform."
% The line division is recorded here rather than forced with \\, so that
% no line-break hyphen is frozen into the text.  Bold, so it takes NO
% \emph{} (BATCH-AGENT §3: a bold head is flagged by spacing.py exactly
% like Sperrsatz and must not be marked from that flag).
\begin{center}
\textbf{I. En Udviklingsrækkes Princip præformerer og
anticiperer under sin Vorden, de tænkende Individer
ubevidst, sin udfoldede Virkelighedsform.}
\end{center}

Som Exempler, der oplyse og bekræfte denne Lov, vælge vi:

1. Anticipationen af Ideelærens Grundbestemmelser i de
ældste Former af den græske Philosophie.

2. Anticipationen af de Tanker, i hvilke Overgangsperio\-%
den culminerer, under denne Periodes Udvikling.

3. Den kritiske Idealismes Grundbegrebs Gjennembrud
under den Kantiske Philosophies Udvikling.

% DIVISIONAL RULE, p. 50, between item 3 and the L3 head below: short,
% CENTRED, white space above and below.
\divrule

% L3 HEAD, p. 50 (transcription.tex §3 assigns the 1./2./3. heads under
% each thesis to L3).  REPORT THIS: §3 describes L3 as a "BOLD centred
% head", and this head is neither bold nor centred.  It is LETTERSPACED
% ROMAN at body size, set as a paragraph on the full measure with the
% first line indented and the last line running short and flush LEFT --
% i.e. typographically it is §3's L4 setting, at §3's L3 level.  The
% companion head "2." on p. 54 is set identically.  Transcribed at the
% level §3 gives, with the letterspacing rendered \emph{} and the
% paragraph setting preserved.
\emph{1. Anticipationen af Ideelærens Grundbestem\-%
melser i den græske Philosophies ældste Former.}

Ideelærens Forberedelse i den græske Philosophies ældste
Former er allerede berørt i det Foregaaende. Den græske
Philosophies \emph{første} Periode kan ikke blot forsaavidt betegnes
som Ideelærens Vorden, som denne Lære factisk har udviklet
sig af den, men ogsaa forsaavidt som den frembyder bestemte
Antydninger af den udfoldede Ideelæres Grundbestemmelser og
Grundproblemer og af den fuldt udviklede græske Tænknings
% --- p. 51 ---
Grundpræg. Saaledes svare de allerede tidligere angivne Be\-%
stemmelser af \emph{Værensprinciperne} hos de gamle \emph{Joniere},
hos \emph{Pythagoræerne} og hos \emph{Eleaterne} til den \emph{platoniske
Sondring} mellem det \emph{Materielles}, det \emph{Mathematiskes} og
\emph{Ideernes Sphære}; saaledes ere de to Grundbegreber: \emph{Væren
og Vorden}, der beherske den gamle Philosophie, og om hvilke
hele dens Tankebevægelse dreier sig, i en dybere Opfattelse
netop den udviklede Ideelæres Grundbegreber; saaledes er \emph{Be\-%
stræbelsen for at bringe dem til Enhed} den samme,
der i en høiere Form gjentager sig i Ideelæren, og gjentager
sig med den samme Mangel paa et Resultat; saaledes er det
\emph{Subjectivitetsprincip}, der i Sophistiken virkede som op\-%
løsende de ældre Tankeformer, i sin ufuldkomne Form Forbe\-%
redelsen for det mere udviklede, der er Betingelsen for, at
Ideens Tanke kan udklares, Ideen sættes som Tilværelsens
Væsen og en ved denne Grundtanke bestemt Philosophie
udvikle sig; saaledes er endelig \emph{Tænkningens Art} og
\emph{Grundpræg} og \emph{Dialektikens Form} overensstemmende med
den, der betinges ved Anskuelsen af Værensprincipet som hy\-%
postaseret Idee, som Form uden Negationens Bestemmelse.

Paa dette Punkt, hvor Forholdet mellem den \emph{bevidste}
Tankeudvikling og den gjennem denne, \emph{ubevidst} for Indivi\-%
det, teleologisk virkende Tanke-Nødvendighed afgiver Syns\-%
punktet for Opfattelsen, bliver den ældste Philosophies \emph{forbe\-%
redende} Forhold til den udviklede Ideelære bestemtere at see
paa saadanne Punkter, hvor Antydninger af denne med ube\-%
vidst Nødvendighed vikler sig ud af det med Bevidsthed Til\-%
stræbte, og det Fremtidige anticiperes.

En saadan ubevidst Anticipation kunde allerede synes at
komme frem hos en af de ældre Joniere, hos \emph{Anaximander}.
Idet han forlader \emph{Thales}'s Bestemmelse af Principet som et
% p.51: "det almene Stof" is NOT letterspaced -- "almene" measures with
% the body fount at 500 dpi -- but "T a n k e omsætningen" in the next
% line IS, and only in its first element: "Tanke" runs 55 px per letter
% against 29 px for "omsætningen" and 31 px for the neighbouring
% "mythiske".  The same partial letterspacing of a compound occurs at
% p. 46 ("B e v i d s t h e d s udviklingen").  Transcribed as printed.
enkelt, bestemt Stof, og som Princip sætter det almene
Stof, \emph{Tanke}omsætningen af det mythiske \emph{Chaos}, giver
% p.51: \gk{τὸ απειρον} -- THE KNOWN-UNRESOLVABLE CASE (transcription.tex
% §GREEK, RECON.md §4; the recon list names p. 52 for this word but not
% p. 51, where it stands first).  The grave over the omicron of τὸ is
% plainly resolvable.  Over the alpha of ἄπειρον the breathing and the
% accent are printed as ONE compound sort and at 1200 dpi it is a single
% blob: neither element can be given its own shape, and ἄ/ἅ cannot be
% separated.  Per BATCH-AGENT §4 rule 2 the alpha is set UNACCENTED and
% nothing is supplied.  The colon after it is an ordinary colon, not the
% ɔ: id-est mark -- checked at 1200 dpi, both dots are round and there is
% no reversed c.
han, forat sikkre Tilblivelsen en uudtømmelig Kilde, Stof\-%
fet Bestemmelsen af \gk{τὸ απειρον}: det \emph{Grændseløse}. Men
% The sheet signature `4*' stands at the foot of p. 51; not transcribed.
% --- p. 52 ---
i samme Øieblik som denne Betegnelse er given, fordobbler
dens Betydning \emph{sig}; \gk{τὸ απειρον} faaer Betydning af det \emph{Be\-%
stemmelsesløse}, \emph{Formløse}, og i denne Betydning falder
det sammen med det Begreb, der hos \emph{Plato} bliver Udtrykket
for Ideens Modsætning. Dog er vel hos Anaximander den
dobbelte Betydning af \gk{τὸ απειρον} at udlede af den \emph{bevidste}
Reflexion: at det, der skal være i Alt og vorde til alt det
Bestemte, selv ikke kan være noget bestemt Stof.

Sikkrere viser sig derimod den \emph{ubevidste} Anticipation
hos de Tænkere, der repræsentere \emph{Hovedmodsætningerne}
i den første Periode: hos \emph{Eleaterne} og \emph{Heraklit}, og hos
den Tænker, der først udtaler det \emph{høiere} Princip, men som
et endnu Ubegrebet: hos \emph{Anaxagoras}.

\emph{Eleaterne} søge, ligesom Jonierne og Pythagoræerne,
\emph{Tilværelsens Stofprincip}: de spørge; \emph{hvad er det, der
er i Tingene}? og deres Svar er det Omspurgte selv i \emph{tau\-%
tologisk} Omskrivning: det, der er i al Tilværelse, er \emph{Væren}
som Enhed, som udelukkende al Ikke-Væren, og derfor som
utilbleven, uforgængelig, uberørt af Tidsforskjel, udelelig, ufor\-%
anderlig, ubevæget. Og denne Væren, der opfattes som det
% p.52: \gk{τὸ εῖναι} -- over the iota stands a compound sort (breathing
% plus circumflex) of exactly the kind recorded at p. 33 for \gk{εῖδος}.
% The circumflex element is broad and visible and is transcribed; the
% breathing cannot be separated from it at 1200 dpi and is NOT supplied;
% ἶ/ῖ undecided.  Same treatment as p. 33.
% p.52: \gk{τὸ ἐὸν} -- the smooth breathing over the epsilon is a
% comma-shaped hook (visibly not the "C" of the rough breathing over the
% rho of ῥεῖ on p. 53).  PRINTED AS SHOWN: the omicron carries a GRAVE,
% not the acute the position before a comma calls for.  The mark slopes
% down to the right and is identical with the certain grave of the τὸ
% beside it, and visibly opposite to the up-right acute over the epsilon
% of πλέον on the next line -- the two are directly comparable in one
% render.  Expected reading ἐόν; not corrected.
% p.52: the sign between "er" and \gk{τὸ πλέον} is an ordinary EQUALS,
% two short parallel bars, set as text (BATCH-AGENT §9); ABBYY doubles
% it ("==").
legemligt Værendes Væsen, bestemmer sig til et \emph{anskuet}
Værende: \gk{τὸ εῖναι} bliver til \gk{τὸ ἐὸν}, der, som Tilværelsens
materielle Substants, er = \gk{τὸ πλέον}, det Rumopfyldende. Men
under denne Bestemmelse trænger \emph{Ideens} Begreb sig ligesom
uimodstaaeligt frem. Det Værendes Substants, i sin \emph{evige,
uforanderlige Enhed, befriet fra Tidens og Rummets
Udstrækning}, kun anskuet som begrændset fordi det Ube\-%
grændsede som et Uformet vilde være et Ufuldkomment, bliver
\emph{Eet med Tanken} og aabenbarer i denne Enhed sin \emph{høiere
ideelle Natur}. Værens sande Begreb: \emph{Ideen}, er ligesom
ifærd med at bryde frem i sin egen Skikkelse, men Ideen
bliver det \emph{ubevidst} Virkende, skjult under det anskuede, ved
Rummet bestemte Værende, til hvilket Tanken bevidst føres
tilbage.

Hos \emph{Heraklit} fremtræder Anticipationen paa følgende
% --- p. 53 ---
% p.53: \gk{Παντα ῥεῖ} -- PRINTED AS SHOWN, with NO accent over either
% alpha of Παντα.  Nothing stands above the word at 1200 dpi, where the
% neighbouring accents on this page are unmistakable; the expected
% reading is Πάντα.  Not on the author's Trykfeil leaf; not corrected.
% Over the rho stands a clear "C" opening to the right -- the ROUGH
% breathing, positively of the other kind from the comma-shaped smooth
% breathing over the epsilon of ἐόν on p. 52 -- so ῥ is read, not
% supplied.  The mark over the iota of ῥεῖ is a broad circumflex; no
% breathing is legible with it and none is supplied.
Maade. \gk{Παντα ῥεῖ} er hans Grundtanke: Tilværelsen er som
en Strøm, i hvis rastløse Bølgen intet Punkt er i Hvile, men
Alt i ustandselig Vexlen, gjennemløbende «Veien opad» og
«Veien nedad», Forvandlingens dobbelte Retning. Og denne
Vorden er ikke et Værendes Vorden, men er selv sit eget
Subject: den sandselige Tilværelses Uro er ligesom hævet til
en høiere Potents ved denne Bestemmelse, der ingen Væren
lader blive tilbage. Og dog trænger et \emph{Blivende} sig ind,
\emph{uforklaret} og \emph{uforklarligt} af det \emph{Sandseliges} Væsen;
% p.53: PRINTED AS SHOWN -- a COMMA stands after "Nødvendighed" below
% where the sense wants a full stop ("... er en uafviselig Nødvendighed.
% I den rastløse Vorden ...").  The mark carries a descending tail and is
% plainly a comma at 500 dpi; both OCR witnesses read it so.  Not on the
% author's Trykfeil leaf; not corrected.
men som en Bestemmelse, der for den græske Tænkers Be\-%
vidsthed er en uafviselig Nødvendighed, I den rastløse Vor\-%
den postuleres der et \emph{Maal}, en beherskende \emph{Harmonie}, en
\emph{Forandringens Lov}, der kun mangler \emph{Navnet} forat aaben\-%
bare sin Enhed med den \emph{Ideens Lov}, hvilken \emph{Plato} lader
bringe Grændse og Maal ind i det «Større og Mindre», i det
urolige Stofagtige.

Hos \emph{Anaxagoras} endelig kommer Anticipationen frem
idet han søger et Princip, der kan forklare Vorden, hvilken
han opfatter som den Stedbevægelse, af hvilken den sondrede,
ordnede Tilværelse resulterer. I selve det Stofagtige med dets
Ligegyldighed for Forbindelse og Adskillelse, med dets oprin\-%
delige chaotiske Blanding, kan han ikke finde et saadant be\-%
vægende og ordnende Princip: han tvinges da, skjøndt han
endnu kun ved en qvantitativ Forskjel kan sondre det Søgte
fra Stoffet, skjøndt han endnu kun kan bestemme det som
det fineste og reneste \emph{Stof}, og kun kan give det en Virk\-%
somhed liig den, det ene mechanisk bevægede Stofagtige udøver
paa det andet, ud over Stoffet, idet han nævner det \emph{Navn}, der
betegner \emph{Ideen}, hvis Indhold endnu ikke er opgaaet for hans
Tanke, ikke er forstaaet af ham. Navnet \gk{Νοῦς} bliver \emph{Var\-%
selet om Ideen}, og paa Ideelærens \emph{høieste} Udviklingstrin
kommer det igjen som Betegnelse for \emph{Formen}, der i sin
Reenhed aabenbarer sig som \emph{Aand}, som \emph{guddommelig
Tanke}.

% DIVISIONAL RULE, printed at the FOOT of p. 53 (short centred hairline,
% white space above and below), standing over the L3 head that opens
% p. 54.  It is at the foot of 53, not the head of 54 -- p. 54 carries no
% rule.  Not the footnote separator; neither page bears a note.
\divrule
% --- p. 54 ---
% L3 HEAD, p. 54 (transcription.tex §3).  Set exactly as its companion
% "1." on p. 50: LETTERSPACED ROMAN, NOT bold and NOT centred, a
% paragraph on the full measure with the first line indented and the
% third line ("Udvikling.") running short and flush LEFT.  See the note
% at the p. 50 head: §3 calls this level "BOLD centred", and the print
% does not bear that out at either site.  REPORT.
\emph{2. Anticipationen af de Tanker, i hvilke Over\-%
gangsperioden culminerer, under denne Periodes
Udvikling.}

Ved \emph{Overgangsperioden} forstaae vi den halvandet\-%
hundrede Aar omfattende Periode i Philosophiens Historie, der
ligger mellem Scholasticismens Opløsning og den nyere Phi\-%
losophies selvstændige Udvikling. Betegnelsen: Overgangspe\-%
riode, viser deels hen til dens Stilling som Mellemled mellem
Middelalderen og den nyere Tid, deels hen til, at den danner
det \emph{philosophiske} Mellemled mellem de to store Epocher i
% p.54: \gk{κατ᾽ ἐξοχήν} -- the elision mark after κατ is the raised
% comma-shaped koronis (U+1FBD), set as at the earlier occurrence of this
% phrase in the file; the smooth breathing over the epsilon of ἐξοχήν is
% a comma opening to the right and the acute over the eta rises to the
% right.  All three read at 2000 dpi.
Philosophiens Historie. Med Hensyn til dette sidste Forhold
bliver den Overgangsperiode \gk{κατ᾽ ἐξοχήν}.

En saadan Overgangsperiode var \emph{nødvendig} ved Mid\-%
delalderens Slutning før Philosophiens nye Hovedform kunde
fremtræde i sin udviklede Bestemthed. Tænkningen havde i
Scholasticismen, idet dens Gjenstand var det for Tanken Ube\-%
gribelige, kun vundet en \emph{formel} Udvikling uden at kunne
udfolde et selvstændigt Tankeindhold. Idet den frigjorde sig
fra Kirkens Herredømme, og, forkyndende en ny Verdens\-%
anskuelse, vendte sig mod Scholasticismen, maatte den
derfor søge det for denne Kamp nødvendige Indhold, der
maatte være et saadant, i hvilket den befriede menneskelige
Tanke kunde gjenfinde sig og gjenkjende Virkeligheden, hvor
det forelaae som givet, og den fandt det da hos \emph{Oldtidens}
Philosophie, med hvilken den følte sig i Slægtskab. Den om\-%
byttede da Scholasticismens \emph{tvungne} Underkastelse under en
\emph{fremmed} Autoritet med en \emph{selvvalgt} Underkastelse under
en \emph{beslægtet}, i hvilken den nye Tids Tanke, der nu kunde
have Virkelighed og Nærværelse i sit Indhold, kunde vinde
Fylde og Selvstændighed gjennem Tilegnelsen og Udformningen
af dette Virkelighedsindhold, i hvilket den først skulde leve sig
ind. I alle de mangfoldige og forskjelligartede philosophiske
Forsøg, der falde i dette Tidsrum, er det fælleds Særkjende,
ved Siden af Oppositionen mod den tidligere Philosophie, denne
% --- p. 55 ---
\emph{Afhængighed af Oldtiden}, trods Bestræbelsen for og den
efterhaanden fremtrædende Fordring paa Uafhængighed.

Søge vi at opfatte Periodens \emph{væsentlige Eiendomme\-%
lighed} gjennem de \emph{Tanker}, der i den ere de \emph{centrale}, saa
ligger der en stor Vanskelighed i Mangfoldigheden og den til\-%
syneladende totale Uensartethed af de philosophiske Forsøg, der
udfylde Perioden. \emph{Med Bevidsthed} stemme de alle kun
overeens i \emph{det Negative}: i Oppositionen mod den tidligere
Philosophies Art og Retning; men denne Opposition antager
talrige Skikkelser, --- saameget talrigere som den ved Mangelen
paa Selvstændighed betingede Optagelse af Fortidens Tanke\-%
former fører til forskjellige Combinationer, paa hvilke Tænkernes
individuelle Retning faaer en overveiende Indflydelse ---; og det
\emph{Positive} i den aandelige Retning, der gaaer igjennem de
forskjellige Nuancer, kommer først successivt til Klarhed. Kun
idet man fæster Øiet paa de Tiders Eiendommelighed, mellem
hvilke denne Periode danner Overgangen, og paa det Maal,
til hvilket Udviklingen historisk har ført, og idet man fordyber
sig i Systemerne, ligesom gribende deres skjulte Tanker, bliver
det muligt at finde det ledende Princip. Det gjælder i denne
Tid, ligesom i enhver stærkere bevæget Tid, om, gjennem den
aandelige Strømning, der viser sig paa Overfladen og paavirkes
af hver Luftning, at see den, der bevæger sig i Dybet og som
let unddrager sig Iagttagelsen. Gjennem hiin, der er Billedet
paa Individernes bevidste Stræben, maa man speide efter de
Magter, af hvilke de ubevidst ledes; man maa see Fortidens og
Fremtidens skjulte Magters hemmelighedsfulde Indvirkning
som hæmmende og tilskyndende; kun paa denne Maade for\-%
staaer man en saadan Gjæringstid, rig paa mægtige Kræfter,
paa energiske Bestræbelser, men fattig paa Besindighed, paa
Begrændsning og paa Bevidsthedens Klarhed.

See vi altsaa hen til de Tiders Eiendommelighed, som
dette Tidsrum forbinder, saa har \emph{Scholasticismens}, den
«christelige Philosophies», Tid viist sig som ufri, splidagtig og
usammenhængende i sit Tankeliv. Scholasticismen optog sine
% --- p. 56 ---
Tankeformer fra en \emph{fremmed} Livsanskuelse, optog dem der\-%
for \emph{kun som Former}, i hvilke den bevægede sig ved en
Forstandsreflexion, der mere og mere \emph{fjernede sig fra Vir\-%
keligheden}, og idet den kun kunde anvende dem paa det
\emph{over} Tanken som urokkelig Autoritet staaende christelige
Dogme som \emph{ydre} Former, i hvilke det ikke forvandledes til
et Tankeindhold, eller tilegnedes som begrebet af Bevidstheden,
saa blev dens Charakteer \emph{Ufrihed}, \emph{Virkelighedsløshed} og
\emph{Formalisme}. Den opløstes ved Modsigelsen mellem dens
Elementer, ved den tænkende og religiøse Aands Utilfredsstil\-%
lethed i den, ved Aandens i mangfoldige Phænomener frem\-%
trædende, gjennem Oldtidens Gjenoplivelse styrkede Stræben
henimod Naturen og mod den frie Selvbevidstheds ideale Virke\-%
lighed. Trangen til \emph{Frihed} og Trangen til \emph{Virkelighed},
Tankelivets dobbelte Betingelse, førte Bevidstheden bort fra
Kirkens Autoritet og fra det naturløse transscendente Aandelige,
og idet den førte den bort derfra, førte den den hen til \emph{Selv\-%
bevidstheden} som til Sandhedens Kilde, og til \emph{Naturen} som
Tankens Indhold, som det i den objective Virkeligheds Form
existerende Fornuftige og Sande, som Grundlaget for en concret
Udfoldelse af det menneskelige Liv. Ved denne forandrede
Retning vandt Aanden, som det allerede tidligere blev viist,
Betingelsen tilbage for Udviklingen af en virkelig Philosophie,
og denne Retnings \emph{methodiske, bevidste Gjennemførelse}
% p.56: the single letter "i" between "Philosophie" and "strængere"
% cannot itself show letterspacing and is left plain; the two spaced
% groups are set separately, as at p. 2 l. 15.
constituerer Begyndelsespunktet for den \emph{nyere Philosophie}
i \emph{strængere} Forstand, for den Tænkningens Hovedform, der
forberedes gjennem Overgangsperioden.

Den nyere Philosophies Eiendommelighed bliver den, at
den, gaaende ud fra \emph{Tvivlen}, i hvilken Overgangsperioden
ender, idet den ikke magter de Modsætninger, den vil forene,
% =====================================================================
% SEAM at the foot of p. 56 -- REPORT THIS.  The letterspaced run AND
% the broken word BOTH continue onto p. 57.  The print reads
%   "... og det uomstødelige Udgangspunkt betrygget, e n s-
%    a r t e t  T a n k e i n d h o l d ."
% The \emph{} group is closed here only so that this fragment's braces
% balance.  Whoever splices pp. 57--70 must join "ens\-" to "artet" and
% merge the two \emph{} groups so that the run reads
%   \emph{... og det uomstødelige Udgangspunkt betrygget, ensartet
%   Tankeindhold}.
% p. 56 bears NO footnote at all -- its note area is empty white paper --
% so nothing runs over from 56 to 57 (BATCH-AGENT §5 check 2).
% =====================================================================
sætter sig den Opgave: \emph{uafhængig af Kirkens og af Old\-%
tidens Autoritet, ud af Selvbevidsthedens Indre,
eller gjennem den mod Naturen rettede Erfaring,
at udvikle et selvstændigt, ved den exacte Methode
og det uomstødelige Udgangspunkt betrygget, ens\-}%
% --- p. 57 ---
% =====================================================================
% SEAM at the head of p. 57 -- REPORT THIS to the pp. 43--56 batch.
% (1) A WORD IS BROKEN ACROSS THE BATCH BOUNDARY.  p. 56 ends
%     "... ved den exacte Methode og det uomstødelige Udgangspunkt
%      betrygget, e n s-"  and p. 57 opens "a r t e t  Tankeindhold."
%     The joined reading is "ensartet Tankeindhold."; whoever splices
%     pp. 43--56 must leave `ens\-%` there so the halves join.
% (2) A LETTERSPACED RUN CROSSES THE SAME SEAM.  The programmatic
%     sentence on p. 56 ("uafhængig af Kirkens og af Oldtidens
%     Autoritet ... ensartet Tankeindhold.") is set in Sperrsatz, and
%     the run does not stop at the page break: "artet Tankeindhold."
%     is letterspaced on p. 57.  The \emph{} group is REOPENED here so
%     that this fragment's braces balance on their own; the two groups
%     should be merged when the batches are joined.
% (3) NO RUN-OVER FOOTNOTE.  Checked per BATCH-AGENT §5: p. 56 bears no
%     note at all, and p. 57's type block is filled by body text down to
%     the last line, with no note area and no unmarked matter.  The
%     p. 56/57 seam is clean in that respect.
% =====================================================================
\emph{artet Tankeindhold.} \emph{Naturens Sandhed og Selvbe\-%
vidsthedens Autonomie} bliver dens Løsen, og ved dette
Løsen kommer den i principiel Modsætning til Middelalderens
Retning og Livsanskuelse.

Ved disse to Yderpunkter bestemmes indirecte Charakteren
af det, der forbinder dem. Det maa vise hen i \emph{samme Ret\-}%
% p.57, l. 6--7: the compositor letterspaces "s a m m e  R e t-" at the
% end of l. 6 but sets the carried syllable "ning" at the head of l. 7 in
% ORDINARY spacing (confirmed at 2x on the 300 ppi render: n-i-n-g are
% tight, against the plainly spaced "R e t-" above).  This is a defect,
% not the rule: on l. 1--2 of this same page the run "S e l v b e-
% v i d s t h e d e n s" IS carried across the break.  Transcribed as
% printed, so the \emph{} closes before the hyphen.  The same defect
% recurs at the p. 62/63 page break ("V æ s e n t-" / "lighed").
ning som det, der fremgaaer af det som Udviklingens Maal,
det maa i sine Grundtanker aabenbare Slægtskabet, men i
Grundtankernes Form, i deres eiendommelige Farvning, vise
dem som fremgaaede af, som endnu ikke løste fra Sammen\-%
hængen med Fortidens Tanker.

Dette stadfæster sig ogsaa, naar vi, ledede ved hiin An\-%
visning, som de Grundtanker, der, klarere eller dunklere, be\-%
herske Overgangsperiodens Systemer, og mod hvilke disses
Retning viser hen, selv hvor Afstanden tilsyneladende er størst,
% p.57, l. 16--17: the conjunction "og" between "U n i v e r s u m" and
% "d e t" is NOT letterspaced (o and g touch), although both flanking
% phrases are.  Confirmed at 2x.  Contrast l. 1, where the "o g" between
% "S a n d h e d" and "S e l v b e-" IS spaced.  The compositor is
% inconsistent about short conjunctions inside a run; both readings are
% set here as printed, so the emphasis is broken into two groups.
finde de nøie forbundne Tanker om det \emph{guddommelige
Universum} og \emph{det guddommelige mikrokosmiske
Menneskevæsen}. Disse Tanker, der \emph{klart} udtales i Perio\-%
dens \emph{meest udviklede Systemer}, fremgaae \emph{med Nød\-%
vendighed} af den Retning, der kalder den \emph{nye} Tænkning
tillive. Idet \emph{Naturen} atter faaer Betydning og ikke længere
bliver det blot Væsenløse, maa den sættes i et \emph{indre, væ\-%
sentligt Forhold til Guddommen}; Forestillingen om en
% p.57, l. 24: a small speck of dirt stands between "end" and "med"
% ("om end med Inconsequentser"); ABBYY reads it as a hyphen
% ("om end- med") and tesseract as a full stop ("om end. med").  It is
% not type -- there is no hyphen or point in the printed line -- and it
% is not transcribed.
vilkaarlig Skabelse i Tiden maa, om end med Inconsequentser,
% p.57, l. 25: the ɔ: ID-EST MARK, confirmed on the image at 4x
% (U+0254, a reversed c, followed by a colon).  Neither witness reads it
% ("o:" in both).  This site is not the one named in the recon list as
% confirmed (p. 69); it is a second, independent occurrence.
gaae over til Tanken om en \emph{evig} Skabelse ɔ: en evig, nød\-%
vendig Udfoldelse af Guddommens Liv og Væsen, saa at Gud\-%
dommen og Universet bliver den samme Tanke, kun seet fra
to Sider. Og i dette guddommelige Universum, i denne
væsensgyldige besjælede Natur, bliver \emph{Mennesket}, hvis
Aand er det opfattende Speil for \emph{alle} Universets Tilvæ\-%
relsesformer, \emph{Centrum som mikrokosmisk} Væsen. Den
Guddommelighedens Nimbus, der omstraaler Universet som
Heelhed, omstraaler Menneskenaturen, og med sværmerisk Be\-%
geistring, med den unge Friheds Overvætteshed, troer Tidsalde\-%
ren paa \emph{Menneskevæsenets Uendelighed}, paa dets
% --- p. 58 ---
\setcounter{footnote}{0}
Uendelighed i \emph{Magt}, der magisk griber ud over alle Tilvæ\-%
relsens Sphærer; paa dets Uendelighed i \emph{Væren}, der ikke be\-%
grændses af nogen Tid; paa dets Uendelighed i \emph{Erkjendelse},
der ikke hæmmes af nogen Skranke.

I disse Tanker komme den \emph{nyere} Philosophies centrale
Tanker tilsyne, men paa den ene Side --- hvad der paa et sildi\-%
gere Sted vil blive paaviist --- som paavirkede af \emph{Middelal\-%
derens} Forestillinger; paa den anden Side som bestemte gjen\-%
nem en Reflex af \emph{Oldtidens} Tanker. \emph{Naturen}, hvis Enhed
% p.58: \gk{Κόσμος} (three times on this page) and \gk{τέλος}.  Read off
% the image at 1200 dpi; neither witness contains a Greek character
% (ABBYY gives "K 6 o> o c", "Kocsfxos", "Korr^oc", "réloc").
%   Κόσμος -- capital kappa (no breathing, the word begins with a
%   consonant); the mark over the first omicron is an unmistakable ACUTE,
%   a short stroke rising to the right, resolved at 1200 dpi.
%   τέλος -- tau, then an ACUTE over the epsilon, equally clear.
% Neither word is letterspaced (the letters sit tight inside the cursive
% Greek fount), so neither takes \emph{} -- cf. BATCH-AGENT §4.5, which
% warns that spacing.py flags all Greek.
den Erfaringen anticiperende Tanke fordrer, anskues af \emph{Phan\-%
tasien som en uendeliggjort} \gk{Κόσμος}, og bliver gjennem
hele denne Periode væsentlig den \emph{metaphysiske} Tankes og
\emph{Phantasiens}\footnote{Man erindre \emph{Magiens} Betydning i denne Tid.} Gjenstand. Og \emph{Mennesket}, der danner
Centrum i denne \gk{Κόσμος}, bestemmes paa \emph{antik} Maade ad
den \emph{metaphysiske} Tænknings Vei som Naturtilværelsens
harmoniske \gk{τέλος}, som det den hele Tilværelse i sig Concen\-%
trerende, dog saaledes, at den moderne Tids Accentueren af
\emph{Inderlighedens Uendelighed} nuancerer den antike Tanke.
Thi det er, hvormeget end \emph{Natursiden} accentueres, dog
Menneskets \emph{Aand}, der bliver den uendelige Tilværelses ideelle
Centrum, ligesom for Oldtidens Tanke Jorden var det materi\-%
elle i den begrændsede \gk{Κόσμος}. Overhovedet maa det Antike,
idet det gjenoplives i den moderne Tid, bestandig modificeres
ligesom ved en Gjennemgang gjennem brydende Medier\footnote{Et interessant Exempel paa en saadan Brydning finde vi i \emph{Dantes}
divina commedia, hvor Oldtidens Tankeindhold viser sig først som
brudt gjennem Scholasticismens, den christelige philosopherende Be\-%
vidstheds, og saa gjennem den digteriske Bevidstheds Medium.}.
% p.58, note **): "divina commedia" is set in the SAME upright roman as
% the surrounding Danish and is NOT letterspaced, so per BATCH-AGENT §3
% it takes neither \emph{} nor \textit{}.  Checked on the image.

Medens Overgangsperiodens to Hovedtanker saaledes an\-%
tyde den \emph{nyere} Philosophies Retning, danne de med det Samme
paa een Gang Modsætningen og Parallelen til \emph{Middelalde\-%
rens} centrale Forestillinger: den naturløse Gud og det histo\-%
riske Gudmenneske. Den Verdens- og Livsanskuelse, der ud\-%
gaaer fra dem, maa være \emph{principiel} forskjellig fra den fore\-%
% --- p. 59 ---
gaaende Tids, men det er først \emph{successivt}, at Modsætningen
% p.59, l. 2: a small speck stands between "og" and "i"; both witnesses
% read it as punctuation ("og.i").  It is not type and is not transcribed.
ret bliver klar, og i den Form, i hvilken hine Grundtanker
fremtræde, viser der sig fra Begyndelsen af en Indvirkning af
et \emph{theologisk Element}, under hvilket \emph{Naturelementet}
ligesom maa arbeide sig frem. Det er netop denne Proces,
denne Udvikling af det Tilstræbte af det Svøb, i hvilket det
fra først af ligger indviklet, der afgiver Synspunktet for en
Sondring mellem \emph{to Hovedafsnit i Overgangsperioden},
mellem hvilke den \emph{kirkelige Reformations Begyndelse}
bliver Grændseskjellet. I det første Afsnit havde fra Middel\-%
alderen af det theologiske Element bevaret sig og tilhyllet Na\-%
turelementet, der, hvor det hævdede sig som saadant, gjorde
det under Former, optagne fra den gamle Philosophie. Idet
Reformationen giver det theologiske Element et Afløb og ud\-%
skiller det ved at anvise det et særskilt Gebet, fremtræder Ret\-%
ningen henimod det Naturlige indenfor Philosophien med kla\-%
rere Bevidsthed og med afgjørende Overvægt. Naturvidenska\-%
bernes hurtige Opblomstren i denne Tidsalder, der forbereder
Uafhængigheden af Oldtiden ved en udvidet og berigtiget Na\-%
turanskuelse og ved Uddannelsen af en Methode, der kan op\-%
tages af Philosophien, befrier tillige fuldstændigen fra Afhæn\-%
gigheden af Theologien, og de tilbageblevne theologiske Ele\-%
menter udsondres bestemtere, endog paa fjendtlig Maade. Op\-%
positionen, der i første Afsnit af Perioden kun havde været
rettet mod Middelalderens \emph{Philosophie}, bliver i dette ogsaa
rettet mod Middelalderens \emph{Religion}. Ved Siden af Systemer,
der forholde sig indifferente til Christendommen eller formelt
underordne sig under dens Autoritet, opstaae saadanne, i hvilke
Modsætningen til den fremtræder med klar Bevidsthed, endog
som et lidenskabeligt Had, og der opstaaer en Kamp, i hvilken
Philosophien faaer sine Martyrer (Bruno, Vanini).

Søge vi de ovenfor angivne Grundtanker i deres \emph{reneste}
Form, saa ville vi finde dem i denne hos den Mand, der ikke
blot var Overgangsperiodens største og eiendommeligste Tænker,
men den fuldstændigste Type paa Tidens aandelige Charakteer:
% --- p. 60 ---
hos \emph{Giordano Bruno}. Han er et Billede paa sin Tid, idet
han, ligesom den, paa een Gang har sit Liv i tvende Verde\-%
ner: den grublende Tankes og den digteriske Phantasies; han
er dens Billede ved det Gjærende og Fremtidsbebudende hos
ham, ved sin stærke Lidenskab, ved sit Livs rastløse Uro, ved
sin holdningsløse Vexlen mellem det Ophøiede og det Burleske,
mellem begeistret Selvhengivelse og forfængelig Selvforgudelse,
mellem ideel Spiritualisme og grov Sandselighed, mellem man\-%
% p.60, l. 9: a modern pencil tick stands after "Underdanighed;"; both
% witnesses pick it up as an apostrophe or quote ("Underdanighed;'",
% "Underdanighed;\"").  Not type; not transcribed.
digt Mod og Stolthed og uværdig Underdanighed; han afspei\-%
ler, med eet Ord, i sig hele Tidens indre Splid, dens Elemen\-%
ters uklare Gjæring.

Idet \emph{Giordano Bruno} tilegner sig den Opfattelse af
\emph{Gud som Modsætningernes Enhed}, der ved Periodens
Begyndelse var gjort gjældende af Nicolaus Cusanus, og idet
han nærmere bestemmer Gud som Enheden af de fire aristo\-%
% p.60 carries FOUR quotations and ALL FOUR OPEN WITH » -- the
% compositor's anomalous direction (BATCH-AGENT §9).  Each was read at
% 1200 dpi: »med alle Former svangre», »absolut uendelige»,
% »heroiske Kjærlighed», »guddommeligt Raseri».  The closing marks are
% all » and therefore normal; only the openings are reversed.  Set as
% printed and NOT regularised, so this page contributes four unbalanced
% guillemets to check.py's count.
teliske Aarsager, sætter han ogsaa den »med alle Former
svangre» \emph{Materie i Gud}, og tilnærmer derved Gudsbegrebet
til Begrebet om den selvvirksomme, altomfattende Natur, det
\emph{besjælede Verdensalt}. Af \emph{Guds} »absolut uendelige» Væ\-%
sens Nødvendighed udfolder sig evigt \emph{Universets Uendelig\-%
hed}, i hvilket Modsætningernes Strid, under Delenes uafbrudte
Bevægelse, Forandring, Tilbliven og Forgaaen, frembringer det
uendelige Heles Harmonie. Universets Begreb som dette uende\-%
lige, levende, altomfattende Hele dækker Gudsbegrebet saa
nøie, at kun den skjulte Indvirkning af Fortidens Forestillinger,
som vi i det Efterfølgende skulle omtale, forhindrer deres fuld\-%
stændige Identificeren, og at et adskillende Kjendemærke kun
tvungent kan bringes tilveie. I dette uendelige guddommelige
Universum, der omslutter Totaliteten af \emph{Monaderne}, de af\-%
ledede ideelle Enheder, er \emph{Mennesket} Tilværelsens Baand;
dets uendelige og udødelige Væsen omfatter \emph{alle} Tilværelsens
Sphærer, og finder sin høieste Livsform i Erkjendelsen af Gud,
der forberedes ved den trinvise Opstigen fra Ideernes sandse\-%
lige Billeder, og naaes i Contemplationen, gjennem den »he\-%
roiske Kjærlighed», der i »guddommeligt Raseri», i Selvhen\-%
% --- p. 61 ---
givelsens Lidenskab, glemmer den endelige Existentses altid
varende Begrændsethed.

Disse Tanker, der henimod Periodens Slutning saa klart
udtales af den Tænker, som i begeistret Anskuen af det
overalt besjælede Universums Uendelighed, der lader Tilværel\-%
sescentret være overalt, føler Uendeligheden i sig, og idet
denne Følelse mødes med den spirende Tvivl om Endeligheds\-%
skrankernes Overvindelse, søger ud over Endeligheden i Tan\-%
% p.61: \gk{μανία}.  A GREEK SITE NOT IN THE RECON LIST (BATCH-AGENT §4.6
% warns that the list is a lower bound); ABBYY and tesseract both give
% "pavia".  Read at 1200 dpi: mu, alpha, nu, then an unmistakable ACUTE
% over the iota -- a short stroke rising to the right, well separated
% from the dot -- then alpha.  No breathing (the word begins with mu).
% Not letterspaced.
kens Lidenskab, i den modige Selvforglemmelse, hvilken Plato
betegnede som en \gk{μανία}, som Guddommens Væren i Individet,
ville vi nu see \emph{som antydede ligefra Periodens Be\-%
gyndelse, og som det Samlende og Forenende i den.}

Allerede hos \emph{Nicolaus Cusanus} arbeide Tidsalderens
Grundtanker sig frem i hans Lære om en \emph{høiere Materie i
Gud, i hans Opfattelse af Erkjendelsen af Universet}
% p.61, l. 16--17: the run breaks and resumes.  At 2x, "B e t i n g e l s e"
% and "G u d," are letterspaced but "som", "for", "Erkjendelsen" and "af"
% between them are not -- the compositor spaces only the two substantives.
% Set as printed.
som \emph{Betingelse} for Erkjendelsen af \emph{Gud}, og i hans Begreb
om \emph{Mennesket} som \emph{Mikrokosmos}. Hos de betydeligste
Repræsentanter for det \emph{platoniske Akademie i Florents}:
hos \emph{Marsiglio Ficino} og \emph{Pico af Mirandola}, ere Tan\-%
kerne om \emph{Menneskets Uendelighed} i Erkjendelse og Exi\-%
stents, om det menneskelige Væsen som \emph{Tilværelsens
Baand}, som det, der omfatter Alt i Erkjendelse og Kjær\-%
lighed, det Midtpunkt, om hvilket hele deres Philosophie
samler sig. Og mod Tanken om \emph{Universets Guddomme\-%
lighed} sigtes der hen, idet Gud, medens han opfattes som
den absolut frie, af den evige Urmaterie uafhængige Fornuft,
dog bestemmes som Tingenes evige, \emph{indre} Urform, der som
saadan har ordnet Alt. Hos \emph{Theosopherne} i dette Afsnit
antydes Grundtankerne gjennem den \emph{Væsentlighed}, der
tillægges det \emph{Materielle}, Naturlige, og i Læren om den
\emph{tredobbelte Magie}, ved hvilken \emph{Mennesket} staaer i
Herskerforhold til den tredelte Tilværelse. Hos \emph{Aristoteli\-%
kerne} træde de frem gjennem den Tilbagevenden til den ægte
% p.61/62: A LETTERSPACED RUN AND A BROKEN WORD BOTH CROSS THE PAGE
% BREAK.  The print reads "... Anerkjendelsen af N a-" / "t u r e n s
% V æ s e n t l i g h e d,  o g  a f  d e n  s o m  a l m e e n
% u f o r-" (p. 61, last two lines) and "g æ n g e l i g e
% M e n n e s k e a a n d s  S a m m e n h æ n g  m e d  d e n s
% H e e l h e d." (p. 62, first two lines).  Both pages are inside this
% batch, so the \emph{} group simply spans the page marker below and the
% word "ufor-gængelige" is joined with \-%.
aristoteliske Lære, der atter fører til Anerkjendelsen af \emph{Na\-%
turens Væsentlighed, og af den som almeen ufor\-%
% --- p. 62 ---
\setcounter{footnote}{0}%
gængelige Menneskeaands Sammenhæng med dens
Heelhed.} Netop hos dem træder klarest Bevidstheden frem
om Modsætningen til det Christelige, om end denne Bevidsthed
ledsages af en, idetmindste formel, Underkastelse under Troes\-%
lærdommenes Autoritet (Pomponazzo). Og selv hos den sidste
Gruppe i dette Afsnit, hos de \emph{Philologer}, der ogsaa havde
en, skjøndt underordnet, Betydning som Philosopher (\emph{Valla,
Agricola}), er der, i Bestræbelsen efter at vende tilbage til
\emph{Naturlighed i Tænkningen}, en Henvisning til Naturens
Væsentlighed og til Tankens Sandhed i dens Harmonie med
Naturen, der ligesom fjernt antyder Grundretningens Maal:
Grundtankerne.

% p.62, l. 13: only "e f t e r" is letterspaced -- "falder" before it and
% "Reformationen," after it are both set tight (checked at 2x).  Set as
% printed; the emphasis is on the preposition alone.
I det \emph{andet} Afsnit, det, der falder \emph{efter} Reformationen,
kunne vi gruppere de philosophiske Phænomener saaledes, at
vi skjelne imellem dem, der kun optræde som Fortsættelser
af det tidligere Begyndte, og kun adskille sig \emph{quantitativt}
fra dette ved den \emph{større} Vægt, der lægges paa \emph{Erfaringen},
og ved den \emph{større}, om end ikke consequent fastholdte Be\-%
tydning, der tillægges det \emph{Materielle} (\emph{Philologerne, de
platoniske og aristoteliske Skoler}), --- og dem, der have en
for dette Tidsrum charakteristisk \emph{Eiendommelighed}, der
træder frem gjennem Bestræbelsen for at finde et \emph{nyt Er\-%
kjendelsesprincip}, eller gjennem en \emph{ny Bestemmelse
af Tilværelsens Grunde}, eller paa begge disse Maader
(\emph{Paracelsus}\footnote{\emph{Paracelsus} lader Philosophien blive \emph{Naturphilosophie}. Hans
Opfattelse af Tilværelsen viser sig gjennem hans Opfattelse af den
\emph{yderste Dom} som en \emph{chemisk Adskillelsesproces}.}, \emph{Cardano, Telesio, Vanini, Bruno, Cam\-%
panella, Jacob Bøhm}). Og medens om hiin første Gruppe
det Samme gjælder som om de beslægtede Phænomener i første
Afsnit, saa vise i den anden Gruppe de, der ere fjendtlige
eller ialfald indifferente mod Christendommen, hen paa An\-%
skuelsen af det \emph{uendelige Universum} som \emph{Tilværelsens
Heelhed}, de fremhæve \emph{Materiens Betydning og Væsent\-}%
% --- p. 63 ---
% p.62/63: the same defect as p. 57 l. 6--7 -- the run "M a t e r i e n s
% B e t y d n i n g  o g  V æ s e n t-" is letterspaced to the foot of
% p. 62, but the carried syllable "lighed," at the head of p. 63 is set
% tight (checked at 2x).  The \emph{} therefore closes before the hyphen.
lighed, de søge \emph{physiske} Principer for de endelige Existents\-%
former, de give \emph{Sandser} og \emph{Erfaring} en afgjørende Betyd\-%
ning for Erkjendelsen, og lade \emph{Menneskeaanden} i sig have
Sandhedskilden og Evnen til at omfatte Tilværelsen. Og hos
de Tænkere, hvis aandelige Retning er en anden, \emph{deels} en
mæglende eller restaurerende, der fremkaldtes ved Katholicis\-%
mens Reform (\emph{Campanella}, der vil philosophere «efter \emph{Guds
og Sandsernes Vidnesbyrd}»), \emph{deels} en tilsyneladende med
de christelige Læresætninger overensstemmende og kun fra en
christelig-religiøs Stemning udgaaende, men dog ubevidst af
en fra Christendommen forskjellig Interesse paavirket (\emph{Jacob
Bøhm}), træder dog det \emph{Naturliges} Betydning, selv om det
bliver en «\emph{Natur i Gud}», og \emph{Erkjendelsens Uendelighed},
selv om det er som en «Christi Viden i os», tydelig frem.
Ja endog i de \emph{skeptiske} Former (\emph{Montaigne, Charron,
Sanchez}), i hvilke denne Overgangsperiode opløses, fordi den
i sig bærer \emph{uensartede} Elementer og en uforsonet \emph{duali\-%
stisk} Adskillelse af det Aandelige og det Legemlige, bevares
Tidsalderens Retning: der vises hen til \emph{Selvbevidstheden}
og til \emph{Naturloven} som det \emph{relativt} Sikkre og som Tankens
eneste paalidelige Norm og Udgangspunkt.

Men medens vi saaledes ligefra Begyndelsen af see de
angivne Grundtanker antydede, see vi tillige, at de først
komme frem, ligesom under Skjulet af Fortidens Tanker, og
at deres Udtalen er \emph{ubevidst} idet de umiddelbart føre til og
ligesom udtrykkelig vise hen paa Consequentser, der \emph{bevidst}
kaldes tilbage. Vi søge Beviser herfor hos de Mænd i Perio\-%
dens \emph{første} Halvdeel, der, idet de brøde med Fortidens Phi\-%
losophie, meest energisk og selvstændigt arbeidede hen mod et
Nyt: hos \emph{Nicolaus Cusanus} og \emph{Florentinerne}. De andre
Philosopher i dette Afsnit af Perioden faae en ringere Betyd\-%
ning, idet de enten, som \emph{Aristotelikerne}, med en langt
større Uselvstændighed sluttede sig til en bestemt, enkelt
Form af Oldtidens Tænkning, og ligesom kun stirrede \emph{tilbage}
paa Oldtiden, eller, som \emph{Theosopherne}, kun gave et Gjæ\-%
% --- p. 64 ---
\setcounter{footnote}{0}%
rende og Uklart, eller som Philologerne, kun gave et Frag\-%
mentarisk, intet virkeligt Sammenhængende i philosophisk
Retning.

Hos \emph{Nicolaus Cusanus} fandt vi Tanken om en \emph{høiere
Materie i Gud}, --- en Tanke, der indeslutter det Naturliges
Væsentlighed, --- om Naturerkjendelsens Nødvendighed for
Erkjendelsen af Gud, og om Mennesket som Mikrokosmus.
Men betragte vi nærmere, hvorledes disse Tanker komme
frem, saa vise de sig i en eiendommelig Belysning. \emph{Grund\-%
interessen} er hos Nicolaus Cusanus \emph{Erkjendelsen af
Gud}. Gjennem den positive, negative og mystiske Theologie
vil han naae dertil som til Tankens høieste Maal, og idet
dette Maal opfylder Tanken, bliver \emph{Universet}, der var nærmet
til Enhed med Gud, kun en \emph{secundær} Tanke, en \emph{Conse\-%
quents} af den skabende Guds Tanke. Vel skal det give et
Bidrag til Erkjendelsen af Gud, men Naturerkjendelsens Be\-%
tydning svinder dog atter ind til et \emph{Intet} ligeoverfor den
\emph{overnaturlige Aabenbaring og Troen}. Og naar Cusa\-%
neren anskuer \emph{Menneskenaturens Uendelighed}, saa anskuer
han den i \emph{Gudmennesket}.

Hos \emph{Platonikerne i Florents}, navnlig hos \emph{Ficino},
fremhævede vi Opfattelsen af Menneskenaturens Uendelighed
og det Naturliges indre Forhold til Guddommen. Der lægges
en saadan Accent paa det Naturlige, at det, at Erkjendelsens
Uendelighed, som Aanden attraaer, kun bliver fuldstændig ved
Erkjendelsen af det Enkelte, Naturlige, for Ficinus bliver For\-%
% p.64, note *): the ɔ: ID-EST MARK at note size, confirmed at 2x on the
% 300 ppi render and again at 1200 dpi (U+0254 + colon).  Both witnesses
% give "o:".  Note that "oprindelig" immediately before it is NOT
% letterspaced, though the eye at half-page scale suggests it is.
klaringen af Menneskets Legemlighed\footnote{Netop i denne Lære træder den \emph{Dobbelthed} frem, der er eien\-%
dommelig for Perioden. Aandens Stræben drager den, som med
magisk Magt, hen imod \emph{Naturen}; men den Aand, der saaledes
drages, er en oprindelig ɔ: efter sit Væsen, legemløs, immateriel
Aand, og det sidste Maal for den, idet den drages hen mod Naturen,
er en Erkjendelse, for hvilken Naturen kun bliver \emph{Midlet}.}. Physiken faaer (jfr.
\emph{Pico af Mirandola}) sin Betydning ved Siden af Religionen.
Men idet Menneskeaandens \emph{Uendelighed} væsentlig søges i
% --- p. 65 ---
Retning af \emph{Udødelighed}, faae \emph{Reminiscentserne fra
Dogmet} Overvægten, og idet Physiken uundgaaelig skal lede
til at see Endelighedens Splid, søges der tilbage til en «Skuen
af Alt i \emph{Gud}», der bliver Indholdet for den Udødelighed, i
hvis Opfattelse Eftervirkningen af Nyplatonismens Lære blandes
med Dogmets.

% DIVISIONAL RULE, p. 65, between "... med Dogmets." and the head "3.
% Den kritiske Idealismes ...": short, CENTRED (measured at about a third
% of the measure), with generous white space above and below.  It is the
% book's ordinary divisional rule standing over a display head, not the
% left-flush footnote separator; this page bears no note.
\divrule

% p.65: the L3 head under thesis I (transcription.tex §3 lists it as the
% "3." of the sequence 1./2./3. at pp. 50, 54, 65).  BUT WHAT IS PRINTED
% DOES NOT MATCH §3's DESCRIPTION OF L3, AND THE DISCREPANCY SHOULD GO
% INTO RESUME-NOTES: §3 describes L3 as a "BOLD centred head", whereas
% this head is
%   (a) NOT BOLD -- the strokes are exactly the weight of the body fount,
%       checked at 2x against the bold L1 head on p. 31;
%   (b) NOT CENTRED -- it is set as a full-measure paragraph, three lines,
%       with the first line indented by the ordinary paragraph indent and
%       the last line ("ling.") short at the LEFT;
%   (c) LETTERSPACED throughout, at body size.
% It is therefore set here as a letterspaced paragraph, broken exactly as
% printed, and not with \begin{center}.  Whoever writes the p. 50 and
% p. 54 heads should check whether they are printed the same way.
\emph{3. Den kritiske Idealismes Grundbegrebs Gjen\-%
nembrud under den Kantiske Philosophies Udvik\-%
ling.}

Den nyere idealistiske Philosophies Maal er betinget ved
dens Udgangspunkt og ved dettes Forhold til den modsatte
Retnings. Under Bestræbelsen for, at føre den Tilværelse,
hvis dualistiske Delthed var Middelalderens Forudsætning,
tilbage til Enheden, gaaer Idealismen ud fra \emph{Selvbevidst\-%
heden}, der kun finder sig selv \emph{som det fra Naturen
absolut Sondrede}. Ud fra dette Udgangspunkt, i denne
Opfattelse, kan der ikke naaes til et andet Maal end et saa\-%
dant, der, om det end synes at forene Modsætningerne, dog
kun gjør det paa den Maade, at den enes Eiendommelighed
forsvinder, at Enheden bliver en \emph{ensidig} Enhed, tilveiebragt
ved at opoffre den ene af Modsætningerne. Dette Maal naaes
i den \emph{Hegelske} Philosophie. Der er den rene logiske Tanke,
der som Idee evigt er Aand, sat som det, der i sin evige
Proces omslutter al Virkelighed, som det, der er Enhed af
Tanke og Væren, som det, der som Aand har Naturen som
et evigt ophævet Moment i sig. Men denne Enhed, i hvilken
Forsoningen mellem Væren og Tanke allerede er tilveiebragt
indenfor det Logiskes Gebet, inden Væren træder frem som
Naturens Væren, i Realitetens Bestemmelser, er netop kun
Enheden af Tanken og af Værens abstracte Tanke, en Tankens
Enhed med sig selv, i hvilken Væren som virkelig Væren ikke
finder Plads, og derfor, hvor den gjør sig gjældende i sin reale
Concretion, paa Naturens Gebet, bliver sat som et ved sin
% p.65: the sheet signature "5" stands at the foot of the page (the "5*"
% of the same sheet is at the foot of p. 67).  Not transcribed.
% --- p. 66 ---
Væreform Usandt, medens Naturen, der kun er ved denne
Værens Form, bliver opfattet som det, der er afmægtigt til
at fastholde Begrebet, og i hvilket der derfor bliver Plads for
et existerende Ufornuftigt, --- et Modsigende, ved hvilket Dua\-%
lismens Uovervundethed ufrivilligt røbes.

Men i det Hegelske System er \emph{Enheden}, Begrebets
og Virkelighedens sande Enhed i den absolute Idee, det
\emph{bevidst} satte Maal; mod dette Maal stræber Udviklingen hen,
og det er interessant og lærerigt at iagttage, hvorledes Tanken
derom ved den kritiske Idealismes Begyndelse arbeider sig frem
ubevidst og ligesom kæmpende mod sin Modsætning.

% p.66, l. 26: the print reads "udviklede R e a l i s m e" -- ABBYY reads
% "Idealisme", tesseract "Realisme"; the image at 2x is unambiguous, an R.
\emph{Kant} begrunder den \emph{kritiske} Idealisme efterat den
nyere Idealismes \emph{første} Form: den \emph{dogmatiske} Idealisme,
var udtømt og magtesløs trængtes tilbage af den i det 18de
Aarhundrede til sine sidste Consequentser udviklede Realisme.
Han finder en ny Mulighed for Tanken ved anden Gang at
føre den tilbage til Jeget, til Selvbevidstheden, og denne Gang
med Klarhed over, at Begyndelsen maatte gjøres med at be\-%
stemme Forholdet til Objectet, at Philosophiens ved \emph{Hume}
tvivlsomme Mulighed maatte reddes ved skarpt at sondre
mellem hvad der i Erkjendelsen som et udenfra Givet tilhører
Objectet, og hvad der er Subjectets oprindelige Eiendom,
begrundet i dets Tænknings Natur, udtrykkende dens Love.
Denne Sondring gjennemføres ved i al Erkjendelse at skjelne
mellem \emph{Stoffet} som det \emph{empirisk Givne} og \emph{Formerne}
som det \emph{Aprioriske}, der tilhører Aanden, og under deres
Enhed sammenfatte det Givnes Mangfoldighed. I Alt, hvad
der oprindelig tilhører Aanden, er der \emph{en Yttring af Spon\-%
taneiteten, en Sammenfatten til Enhed. Anskuel\-%
sens} Former: Rummet og Tiden, samle det sandselige Ind\-%
tryks empiriske Mangfoldighed; \emph{Forstandens} Former: Kate\-%
gorierne, sammenknytte Anskuelserne til en Tanke-Enhed;
\emph{Fornuftens} Former: Ideerne, sammenknytte Forstandens Er\-%
kjendelser til en høieste, omfattende Enhed. Men \emph{alle disse
Enhedsvirksomheder og alle disse Spontaneitetens
% --- p. 67 ---
Yttringer} vise tilbage til et \emph{fælleds Udspring: til et
sidste Enhedspunkt, en sidste Activitetens Kilde,}
% p.67 carries MODERN PENCIL (a known site, BATCH-AGENT §7).  Visible:
% a dot between "vise" and "tilbage" in l. 1 (ABBYY reads "vise _tilbage");
% a dot in the left margin beside l. 3; a dot before "Kant" in l. 25; and
% a dot in the right margin beside l. 27.  None is type; none is
% transcribed.
og denne finder Kant i \emph{Selvbevidsthedens oprindelige
Enhed}, i det, han kalder den \emph{transscendentale} --- ikke\-%
empiriske --- \emph{Apperception}, Selvbevidsthedens oprindelige
Synthese. Denne transscendentale Apperception, der ikke har
noget empirisk Indhold, men kun den \emph{Selvbevidsthedens
Grundform}: jeg tænker, jeg er, der er Forudsætningen for
al empirisk Bevidsthed om mig selv, ligesom de afledede For\-%
mer vare det for al Erkjendelse af Gjenstanden, er en Selv\-%
bevidsthed, der gaaer \emph{forud for} al Erfaring. Den empiriske
Apperception, Bevidstheden om Selvet som udfyldt med et
empirisk Indhold, var kun en passiv Perciperen af vore indre
Tilstande, af os selv som Phænomen. Den rene Apperception
derimod er en reen Spontaneitetens Act, og har derfor slet
Intet at gjøre med Sandseligheden, altsaa ikke heller med
dennes mangfoldige Indhold, med vort phænomenale Jegs
mangfoldige vexlende Tilstande, men indeholder kun dette: \emph{at
jeg er}, ikke en Bestemmelse af denne Værens Hvad; thi det,
at: \emph{jeg er}, er ligt med: \emph{jeg tænker}, skal ikke give Væren
et Erkjendelses-Indhold, da dette: jeg tænker, kun betegner
\emph{Tankens abstracte Form}, der i sig bærer \emph{al Erkjen\-%
delses Mulighed, men ingen} Erkjendelsens \emph{Virkelighed}.
Det er den samme Adskillelse, der hos \emph{Kant} gjør sig gjæl\-%
dende i den almindelige Sondring mellem Form og Indhold,
mellem Tænken og Erkjenden; den samme Adskillelse, der
overhovedet gjør det \emph{Intelligibles} Verden, til hvilken hint
ikke-phænomenale Jeg hører, utilgængeligt for al Erkjenden.

\emph{Denne Selvbevidsthedens oprindelige Enhed er
det egentlige Midtpunkt i Kants Kritik der reinen
Vernunft.} Som den \emph{dogmatiske} Idealisme søgte tilbage
til \emph{Jeget} for i det som tænkende Jeg at finde Philosophiens
sikkre Fodfæste, dens betryggede Udgangspunkt, saaledes
gaaer den kritiske Idealisme tilbage til det samme Punkt for
der at aabne en ny Tankens Kilde, og den gjør det ved i
% p.67: the sheet signature "5*" stands at the foot of the page.  Not
% transcribed (BATCH-AGENT §6).
% --- p. 68 ---
Jeget at see det active almene Enhedspunkt, fra hvilket alle
Erkjendelsens Former skulle fremgaae. Betydningen af dette
Punkt som det afgjørende for hele Problemet om Erkjendelsen
er Kant sig klart bevidst, og han udtaler den med en saadan
Bestemthed, at man maa forbauses over at finde Fremstillinger
af hans Lære --- f.\ Ex. det udførlige Værk af Rosenkranz ---
i hvilke den aldeles er overseet. Til dette Punkt er det ogsaa,
at \emph{Fichte} knytter sig i sin Omformen af Kriticismen: hans
Læres hele Charakteer forstaaes ved dette Begreb, i hvilket
\emph{Spontaneiteten er Væsensbestemmelsen}; den eien\-%
dommelige \emph{Tvetydighed} i hans Grundbegreb: \emph{Jeget}, der
bestandig har vanskeliggjort Opfattelsen af hans Philosophie
og endog gjort ham selv usikker, forklares ved det.

Idet nemlig \emph{Kant}, under Udviklingen af sin Philosophie,
tager sit Udgangspunkt fra \emph{Erkjendelsesproblemet}, be\-%
gynder han sin Bevægelse \emph{indenfor Jeget}; han skjelner
mellem \emph{vor} Erkjendelses Form og Stof, mellem det, der er
\emph{Jegets} oprindelige Eiendom, og det, der er givet det udenfra;
han finder de \emph{aprioriske} Former som \emph{Jegets Former}, og
idet deres Enhed og Kilde søges, søges den i \emph{Selvbevidst\-%
heden}. Denne Selvbevidsthed har \emph{Individualitetens Navn}
og \emph{Individualitetens Skikkelse}; men ligesom i den dog\-%
matiske Idealisme \emph{Descartes's}: jeg tænker, ved en uund\-%
gaaelig Consequents hos \emph{Spinoza} blev til: min Tankes Grund,
den uendelige tænkende Substants, tænker, saaledes gaaer hos
\emph{Kant} hiin Individualitetens Form med ubevidst Nødvendighed
over i et \emph{høiere} Begreb: den faaer \emph{Almindelighedens
Omfang og Objectivitetens Indhold.}

Derved, at Selvbevidstheden bestemmes som \emph{Selvbe\-%
vidsthedens rene Form}, bliver den opfattet som \emph{det i
alle Selvbevidstheder Identiske}, den bliver «\emph{Bevidst\-%
hed overhovedet}». Oversætte vi denne Bestemmelse fra
Fornuftkritikens Sprog i Udtryk, der ligge Kants sildigere
Værker og de efterfølgende Former af den kritiske Idealisme
nærmere, saa vilde vi sige: den betegner \emph{Grundlaget,}
% --- p. 69 ---
\setcounter{footnote}{0}%
«\emph{Substratet}» for den \emph{individuelle Bevidsthed}; den er
den \emph{almene Tanke}, der er \emph{Væsenet}, det \emph{substantielle
Ene og Forbindende i det Individuelle.} Og denne
\emph{almene} Bevidsthed faaer i Kritik der reinen Vernunft endnu
en ny Bestemmelse, der antyder et fyldigere Begreb. Idet
\emph{Kant} gjør den transscendentale Apperception til \emph{Erfaringens
Kilde}, og ved Erfaring forstaaer en Erkjendelse med \emph{objec\-%
% p.69, l. 8: the ɔ: ID-EST MARK, the site named in the recon list, here
% in its FIRST of two occurrences on this page.  Read at 1200 dpi:
% U+0254 (a reversed c, its opening to the right) followed by a colon.
% Both witnesses give "o:".
tiv Gyldighed} ɔ: med \emph{nødvendig Almeengyldighed}, saa
mødes hiin \emph{det Subjectives abstracteste Form}, Syn\-%
thesernes Kilde, med \emph{det Objectives abstracteste Be\-%
stemmelse: Tingen i og for sig}, i at være \emph{det med
Nødvendighed Bestemmende}, der gjør vor Erkjenden
uafhængig af Subjectets tilfældige Tilstande, og paatrykker
den Lovmæssighedens Præg. Og naar det udtales (i \emph{første}
Udgave af Fornuftkritiken), at hiin Ting i og for sig, hint
ubekjendte X, turde være \emph{den samme} i alle Erkjendelser, saa
nærme de to, tilsyneladende modsatte Sider, \emph{som Enheder},
sig endnu mere afgjort til hinanden og deres Identitet udtales
hypothetisk. Og idet de sees som gaaende sammen, viser \emph{det Sub\-%
jective} sig som det i Enheden \emph{Primitive}, og bestemtere: som
\emph{det Overgribende}, der \emph{omfatter} det Objective. Ud fra
den Kantiske Forudsætning, der lader al Enhed være betinget
ved det Subjective, kan Objectet kun sættes som værende ved
hiin oprindelige Enhed, der ikke blot bliver Kilden til Enheden
% p.69 carries MODERN PENCIL (a known site, BATCH-AGENT §7): a tick after
% "Gjenstande" in the line below (ABBYY splices it in as "Gjenstande^"),
% and further specks in the lower margin.  Not type; not transcribed.
af de \emph{flere} Gjenstande, som Forstanden forbinder, men ogsaa
til den Enhed, ved hvilken Phænomenets Mangfoldige \emph{først}
bliver til \emph{Gjenstandens Enhed}. Gjenstanden \emph{som Gjen\-%
stand, som} Object er først til ved den, ligesom Gjenstandene
% p.69: the SECOND ɔ: on this page, in the line below -- the reading the
% recon list quotes ("... som Erfaringsobjecter ɔ: som bestemte ved
% Forbindelsens Almeengyldighed og Nødvendighed*)").  ABBYY misreads the
% word as "Erfaringsobjeeter"; the image gives a clear c.
som Erfaringsobjecter ɔ: som bestemte ved Forbindelsens
Almeengyldighed og Nødvendighed\footnote{Naturen, som «Natur», i \emph{kantisk} Betydning, er først værende
ved den transscendentale Apperception.}. Og naar Kants Forud\-%
% p.69, note *): the German «Natur» is set in the SAME upright roman as
% the Danish and is NOT letterspaced, so it takes neither \emph{} nor
% \textit{} (BATCH-AGENT §3).  The guillemets were read at 1200 dpi and
% are the NORMAL inward pair, « … » -- not the reversed opening found
% four times on p. 60.  "N a t u r e n" at the head of the note and
% "k a n t i s k" are letterspaced.
sætninger gjøre Anvendelsen af Anskuelsens Former og For\-%
standens Kategorier, altsaa ogsaa af \emph{Causalitetsbegrebet},
% --- p. 70 ---
% p.70, l. 1: the guillemets around «Ting i og for sig» were read at
% 1200 dpi and are the NORMAL inward pair; ABBYY's closing « is an OCR
% error.  The phrase is NOT letterspaced here, although the same phrase
% IS letterspaced on p. 69, l. 11--12.
paa en objectiv «Ting i og for sig» umulig, saa svinder det
som den ydre Paavirknings Kilde supponerede X, hvad \emph{Fichte}
paaviste, som et ved sin Modsigelse Opløst, ind i Jeget, og
Begrebet om dette X som det udenfor os Værende, vor Er\-%
kjendelse Paavirkende og Bestemmende, sees som det, der
uvilkaarlig danner sig idet Tanken bestemmes ud fra hiin
transcendentale Enhed, Nødvendighedens Grund. Dette For\-%
hold er allerede, paa en mere indirecte Maade, ligesom fra en
større Afstand, antydet, idet den rene Selvbevidsthed er sat
ogsaa som \emph{Anskuelsesformernes} (Rummets og Tidens)
Kilde, og idet den, som Væsensbestemmelse, er opfattet som
begrundende et \emph{Phænomenalt}.

Af hint centrale Punkt, der oprindelig benævnes med
\emph{Selvbevidsthedens} Navn, er altsaa den \emph{høiere} Tanke
ifærd med at udfolde sig, der først sammenfatter det Indivi\-%
duelle i dets begrundende, \emph{almene} Væsen, og saa lader dette
det \emph{Subjectives} Væsen vise sig som det, i hvilket det \emph{Ob\-%
jective} har sin Rod.

% p.70: "homo noumenon" is LETTERSPACED ROMAN, set in the same upright
% fount as the Danish, and therefore takes \emph{} and NOT \textit{}
% (BATCH-AGENT §3; confirmed at 1200 dpi).  So are the German titles
% "r e i n e n" and "p r a k t i s c h e n" -- but note that in both
% cases the following word "Vernunft" is set TIGHT, so the emphasis
% covers the adjective alone.  All four guillemet pairs on this page are
% the normal inward « … ».
Gaae vi fra Kants \emph{første} Hovedværk: Kritik der \emph{reinen}
Vernunft, over til det andet: Kritik der \emph{praktischen} Ver\-%
nunft, for i dette, ligesom i hint, at søge det centrale Punkt,
saa finde vi dette i Begrebet om «\emph{homo noumenon}», hvis
Væsen er \emph{Frihed}. Gjennem Morallovens kategoriske Impera\-%
tiv, der viser hen til denne Frihed som sin forklarende For\-%
udsætning, føres vi ud over de Skranker, indenfor hvilke Kri\-%
tik der reinen Vernunft, endnu ved sin Slutning, holdt os:
det \emph{Intelligibles} Verden, Væsenets Verden i Modsætning til
Phænomenets, aabner sig for os, ikke som en Mulighed, men
som en \emph{Virkelighed}, og denne Virkelighed møder os som
% p.70: the single letter "i" before "o s  s e l v" cannot itself show
% letterspacing and is left plain, on the convention already used at
% p. 2, l. 15--16 of this edition.
Virkelighed i \emph{os selv}. Moralloven, Fornuftloven, er vort
fornuftige, intelligible Væsens Lov, «vor egen Fornufts Lov\-%
givning», «en \emph{indre}, nødvendig Fornuftens Lovgivning», den
er fremgaaet af Frihed, men i Lovens Form fordi den for\-%
pligter os som dem, der høre til Phænomenets Verden, fordi
% =====================================================================
% SEAM at the foot of p. 70 -- REPORT THIS to the pp. 71--84 batch.
% The sentence runs on into p. 71 ("... modtage vort eget" / "Bud som et
% fremmedt Bud, paa een Gang altsaa blive Lovgivere og Loven
% Undergivne."), but NO word is broken at the break and NO \emph{} group
% is open across it -- the last three words are plain.  So the only thing
% the next batch must not do is start a new paragraph: joints.py will
% flag the blank line splice.py leaves above the marker.
% p. 70 BEARS NO FOOTNOTE, and p. 71's note area is empty, so there is no
% run-over note at this seam (BATCH-AGENT §5, checked).
% =====================================================================
vi, paa een Gang hjemme i to Verdener, modtage vort eget
% --- p. 71 ---
% SEAM (batch boundary pp. 57--70 | 71--84): p. 71 line 1 is set FLUSH LEFT
% against the page's own left margin, so it continues the paragraph running
% from p. 70.  The splice must not introduce a paragraph break here.
% p. 70 was read at the foot on the image: it carries NO footnote and its last
% line ends "modtage vort eget" with no word-break hyphen, so nothing runs over
% from p. 70 onto p. 71 (BATCH-AGENT §5 check, done).
% p. 71 bears no note and no Greek.
Bud som et fremmedt Bud, paa een Gang altsaa blive Lov\-%
givere og Loven Undergivne.

Her er det atter hint centrale Begreb fra Kr. d. r. Ver\-%
nunft, der træder frem, men med et forandret Udtryk og lige\-%
som voxende hen mod sin Fuldendelse. Den \emph{rene Selvbe\-%
vidsthed,} Udtrykket for vort \emph{intelligible} Væsen, det, der
bestemte sig til \emph{det Theoretiskes særegne Former,} er
% p. 71: "frie" is NOT letterspaced -- measured 5 px between glyphs at 600 dpi
% against 14--16 px in the spaced runs on the same two lines.  spacing.py's
% single-word flag (1.20) is noise.
bleven til det fornuftige, frie Menneskevæsen, der ved sin Lov
bestemmer \emph{Handlingens Form,} og idet dette fornuftige
% PRINTER'S ERROR, p. 71: the printed reading is "homo n o n m e n o n" -- an
% `n' where `u' belongs.  Both OCR witnesses read "nonmenon"; confirmed on the
% image at 1200 dpi (the letter is a plain n, not a broken or inked-in u).
% p. 70 sets the same phrase correctly as «homo n o u m e n o n» (checked on
% the image), so this is an isolated wrong sort, not the author's spelling.
% Transcribed as printed per the repo's diplomatic stance.
Væsen med sin \emph{almeengyldige} Lov betegnes som \emph{homo
nonmenon,} vises der ud \emph{over det Individuelle.} Det
Lovgivende, der er i mig, er ikke Jeget, men Jegets Slægt\-%
begreb, det er \emph{«Menneskeheden i det enkelte Menne\-%
ske».} Og dette Intelligible, der aabenbarer sig i Fornuft\-%
lovens Form, og aabenbarer sig som Lov for det ogsaa til
\emph{Phænomenets} Verden hørende Væsen, forkynder Loven
som den, der skal virkeliggjøres i \emph{Naturen,} og virke\-%
liggjøres idet Dyd og Lyksalighed harmonere. Atter her skju\-%
ler sig altsaa, ligesom i Kr. d. r. Vernunft, \emph{Forholdet til
det Naturlige} i det Intelligible, og det Intelligible, der er
i os, bliver atter, i dette Forhold, seet som det \emph{Beherskende}
og \emph{Bestemmende.} Men dette udtales endnu ikke i begrebs\-%
mæssig Form: en \emph{Gud} træder til forat tilveiebringe Harmo\-%
nien mellem de to \emph{adskilte} Sphærer, og denne Gud viser
sig for Forestillingen som et \emph{Tredie} ved Siden af begge, og
betegnes som det \emph{Høiere,} medens han dog --- hvad der be\-%
tegner det egentlig Centrale i Kr. d. pr. Vernunft --- sættes
som det, der skylder \emph{Frihedsbevidstheden,} vor \emph{Væsens\-%
bevidsthed,} sin Vished, og kan blive usikkert, medens Mo\-%
ralloven er det for Bevidstheden Umistelige, der, som den
aldrig nedgaaende Stjerne, aldrig forsvinder fra Menneskeaan\-%
dens Horizont.

Gjennem eiendommelige Betegnelser, der vise hen paa et
fyldigere Indhold, gjennem nye Problemer, til hvilke det træ\-%
der i Forhold, er saaledes Centralpunktet i Kritik der pr. Ver\-%
% --- p. 72 ---
\setcounter{footnote}{0}%
nunft hævet høiere end det Begreb, der var det centrale i det
ældre Hovedværk. Men hvad der i Kr. d. pr. Vernunft kun
er antydet, saaledes, at Consequentserne ligesom standses i
samme Øieblik, som de ere ifærd med at bryde frem, det faaer
en høiere Klarhed i \emph{Kritik der Urtheilskraft,} og i dette
beaandede og tankedybe Værk see vi \emph{Kant,} som baaren af en
Tankens Strømning, han ikke selv gjennemskuer, og mod hvil\-%
ken han kæmper, blive ført videre til det Punkt, hvor den
kritiske Idealismes dybeste Tanke kommer tilsyne, og til hvil\-%
% p. 72: the "og" between the two spaced names is NOT itself letterspaced
% (measured 3 px against 13--14 px in the names).  The compositor spaces the
% names only; set as printed.
ket \emph{Schelling} og \emph{Hegel} kunne knytte deres Begreb om
\emph{det Absolute.}

\emph{Kritik der Urtheilskraft} er det Værk, i hvilket den
Kantiske Philosophies \emph{fulde} Eiendommelighed fremtræder.
Denne er kun ensidigt og ufuldstændigt betegnet, naar den
% p. 72 note *) -- printed mark *), superscript.  Letterspacing inside the
% note measured separately from the body (BATCH-AGENT §3): note words run
% 4--7 px unspaced, 12--14 px spaced, so the two modes still separate.
betegnes som \emph{kritisk}\footnote{Som Betegnelse for den Periode i den nyere
Idealismes Udvikling, der begynder med Kant og afsluttes med Hegel, benytte vi
Udtryk\-%
ket: \emph{kritisk} Idealisme, idet vi stille den \emph{som Heelhed} mod den
forkantiske \emph{dogmatiske} Idealisme. Hvorledes Idealismens sidste Periode
nærmere articulerer sig, vil blive paaviist i de efterfølgende Afsnit.}; den
udtales først fuldstændigt, naar
% p. 72, l. 5--6 of this paragraph: the spaced run stops at "Speculative."; then
% "Den" is unspaced, "kritiske" spaced, "Sondring er Kants" unspaced, "første"
% spaced.  Measured glyph by glyph at 600 dpi and checked at 1200:
%   Speculative. 13 | Den 5 | kritiske 15 | Sondring 5 | er 5 | Kants 2 |
%   første 14 | Opgave: 5.   The discontinuity is the compositor's.
den charakteriseres som \emph{kritisk med Retning mod det
Speculative.} Den \emph{kritiske} Sondring er Kants \emph{første}
Opgave: det er kun ved den i dens Strænghed, at han hævder
sin Plads som Voldgiftsmand mellem de to Retninger i det
18de Aarhundredes Tænkning, og ved sin Voldgift skaffer Phi\-%
losophien idetmindste Noget af dens Eiendom tilbage. Derfor
sondres der paa alle Gebeter og Sondringen faaer Modsætnin\-%
gens Form. Mellem det Subjective og det Objective, mellem
% p. 72: this X is set in UPRIGHT roman (checked at 1200 dpi), unlike the X on
% p. 77 l. 26, which is italic.  Each occurrence follows its own page; set as
% printed.  spacing.py's flag on "X," measures the gap to the comma.
Jeget og det ubekjendte X, der er vore Fornemmelsers Kilde,
mellem Form og Stof, mellem Apriorisk og Empirisk, mellem
Ting i og for sig og Phænomen, mellem Anskuelsesformer og
Tankeformer, mellem Spontaneitet og Receptivitet, mellem
Uendeligt og Endeligt, mellem Aand og Natur, gjøres der en
Adskillelse, der haardnakket fastholdes, endog der, hvor Fast\-%
% --- p. 73 ---
\setcounter{footnote}{0}%
holdelsen af den fører til de haardeste Consequentser, som i
% p. 73 note *) -- printed mark *), superscript.
Selvets Opfattelse af sig selv\footnote{Saaledes i det Kantiske Spørgsmaal:
existerer \emph{jeg} selv, som Phæ\-%
nomen for den indre Sands, udenfor \emph{min} Forestillingskraft i Tiden?
og i den Slutning, at da Jeget som rent Jeg kun bestemmes ved
Kategorier, bestemmes det \emph{altsaa} ikke efter dets Væren i og for
sig, --- og det uagtet Kategorierne, som Tankeformer, der opfattes
som fremgaaede af den rene Selvbevidsthed, ifølge dette Forhold
maae være det rene Jegs væsentlige Udtryksformer, Væsensbestem\-%
melser for det.}. Men medens \emph{Sondringen}
overalt træder frem i Forgrunden, aabner Kant paa alle af\-%
gjørende Punkter en Udsigt til det Fjerne, hvor \emph{Enheden} i
\emph{Modsætningerne} viser sig, om end kun viser sig for atter
% p. 73: a small speck stands after "tilhylles." (tesseract reads it as a
% second stop, ABBYY ignores it).  Much lighter than the ink and with no
% letter shape -- not type, and not transcribed.
at tilhylles.

Netop dette er Eiendommeligheden i Kr. d. Urtheilskraft.
Mellem det Theoretiskes og det Praktiskes Sphære, mellem
Natur og Aand, mellem en Lovgivning efter Causalitetens Nød\-%
vendighed og en Lovgivning efter Frihed var der sondret ved
de to ældre kritiske Hovedværker; der var sat ligesom en
uoverstigelig Skranke mellem Naturgebetet og Frihedsgebetet.
% p. 73: BOTH witnesses read "Kr:" here (a colon after Kr).  At 1200 dpi there
% is a stop on the baseline plus a blob at x-height standing just clear of the
% r's arm -- i.e. the r's bulbed terminal, not a second point.  Per
% TRANSCRIPTION-PLAYBOOK §2 rule 3 the burden of proof is on the claim of a
% defect, and "Kr." is the form the book uses everywhere else, so "Kr." is set.
Denne Skranke gjennembryder Kr. d. Urtheilskraft. I \emph{Hen\-%
sigtsmæssighedens} Begreb, der indeslutter Tanken om det
Mangfoldiges Bestemmethed ved Øiemedets Enhed, om det
Reales Beherskethed ved et Ideelt, søges det Mæglende mellem
Nødvendighed og Frihed; i Skjønhedens og i det Organiskes
Sphære søges et Gebet, i hvilket en Lovgivning, der paa een
Gang gjælder Natur og Aand, er raadende. Og idet der gaaes
ud fra en Sondring af en dobbelt Hensigtsmæssighed: den
æsthetiske, subjective, og den (teleo)logiske, objective, saa vises
der ogsaa hen til denne dobbelte Forms Enhed: fra det æsthe\-%
tiske Gebet gjennem Udviklingen af æsthetiske Grundbegreber
som Geniet, de æsthetiske Ideer, den kunstneriske Production;
og fra det Teleologiskes Gebet idet der fra Opfattelsen af den
i det Organiske fremtrædende, til Naturens Heelhed udvidede,
Naturhensigtsmæssighed vises tilbage til den æsthetiske Hen\-%
% --- p. 74 ---
% p. 74 bears no note and no Greek.
sigtsmæssighed som gaaende ind under denne Naturhensigts\-%
mæssigheds Synspunkt.

I dette Begreb om en Naturhensigtsmæssighed som en
\emph{indre} Hensigtsmæssighed, som et indre Naturøiemed, ligger
det Begreb skjult, ved hvilket Kr. d. Urtheilskraft hæver sig
over de tidligere Værker. Der naaes først til den indre Hen\-%
sigtsmæssigheds Begreb indenfor den \emph{enkelte} Organisme: i
den findes det Teleologiskes, for Mechanismen uforklarlige,
Grundbestemmelser: Heelheden som Forudsætning for Delene,
Enhed af Aarsag og Virkning, af Middel og Øiemed. Ved
disse i den individuelle Organismes Kjendsgjerning givne Be\-%
stemmelser kommer Problemet frem om det Teleologiskes
Forhold til det Mechaniske. Det Teleologiske viser sig som
et \emph{nyt} Princip, under hvilket det Mechaniske, der paa Na\-%
turgebetet ikke kan tænkes som forsvundet, men kun som util\-%
gængeligt for vor Opfattelse, bliver at sætte som underordnet.
Og idet dette nye Princip paa eet Punkt er anerkjendt, idet
en anden Lovgivnings Mulighed paa Naturens Gebet er sta\-%
tueret, ifølge den i Phænomenet givne Nødvendighed, om end
ei forklaret, saa udvides dette Princip uundgaaeligt til \emph{Naturens
Heelhed;} denne bliver seet som bestemt ved det \emph{altomfat\-%
tende Øiemed,} under hvilket \emph{alt det Mechaniske} gaaer
ind \emph{som tjenende Middel.}

Men naar denne Bestemmelse er given, saa reiser der sig
strax et nyt Spørgsmaal: \emph{hvorledes er dette Forhold
muligt?} --- og gjennem Besvarelsen af det naaer \emph{Kants}
% p. 74: "Kants" is spaced but "Tænkning" is not (14 px against 3 px), and on
% p. 72 above it is "Kant," alone that is spaced.  The book letterspaces the
% proper name itself (BATCH-AGENT §3); set as printed.
Tænkning til sit Høidepunkt. For \emph{vor} Erkjenden skilles de
to Principer uundgaaelig ad; den begriber ikke deres Enhed;
men ligesom \emph{Tanken} i Kr. d. r. Vernunft og i Kr. d. pr.
Vernunft ilede ud over \emph{Erkjendelsens} Grændse til det over
Phænomenet værende Intelligible, saaledes søger ogsaa her
Tanken hen til det Enhedspunkt, hvor Modsætningerne mødes.
I den for vor Erkjendelse utilgængelige \emph{Naturens Grund}
(den «oversandselige Realgrund for Naturen»), i det \emph{intelli\-%
gible Substrat,} der skjuler sig bag Naturphænomenet, an\-%
% --- p. 75 ---
\setcounter{footnote}{0}%
tydes \emph{Enheden} af Mechanismens og Teleologiens Princip, og
dette «intelligible Naturens Substrat», der allerede ved Udvik\-%
lingen af den æsthetiske Hensigtsmæssigheds Begreb var traadt
i en nøie Forbindelse med det, der er den til Skjønhedsphæ\-%
nomenet i harmonisk Forhold staaende Menneskeaands og alle
dens Evners Substrat: «det oversandselige Substrat for alle
% p. 75, the «fælleds ... Menneskeaanden» quotation: the letterspacing is
% DISCONTINUOUS in the print and is set here exactly as measured at 1200 dpi:
%   «fælleds 13 | intelligible 5--6 | Substrat 4 | for 15 | Naturen 14 |
%   og 14 | Menneskeaanden», 12--14
% i.e. the two words "intelligible Substrat" stand unspaced inside an otherwise
% spaced phrase.  Verified twice by direct glyph-gap measurement and once by
% eye at 1200 dpi; it is not a measurement artefact.  Not regularised.
Subjectets Evner», --- det bliver nu her sat som det \emph{«fælleds}
intelligible Substrat \emph{for Naturen og Menneskeaanden»,}
som det Intelligible, Oversandselige, der er begges fælleds dy\-%
% p. 75 note *) -- printed mark *), superscript.  The reading of the figures
% was taken off the image at 1200 dpi, NOT from ABBYY, which gives
% "Kr. a. U. S. 246, 268 f. smign." -- wrong on four counts.  The page prints
% "Kr. d. U. S. 245, 258 f. smlgn. S. 352 (3die Udgave, fra 1799)."
beste Grund.\footnote{Jfr. Kr. d. U. S. 245, 258 f. smlgn. S. 352 (3die
Udgave, fra 1799).}

Sammenligne vi nu dette Begreb, der bliver det \emph{høieste}
% p. 75: a light brown foxing stain sits over the "i" of "i Kr. d."; it is not
% type.  Both witnesses garble the phrase ("Pkr.", "PkKr.").  Read on the image:
% "Begreb i K r.  d.  U r t h e i l s k r a f t".  The single letter "i"
% cannot itself show letterspacing and is left plain (cf. the note at p. 2).
Begreb i \emph{Kr. d. Urtheilskraft,} med de \emph{centrale} Begreber
% p. 75: both witnesses read an accent on "dét".  At 1200 dpi the mark over the
% e is a faint, roughly horizontal smudge far lighter than the ink and with no
% stroke shape -- an offset or speck, not an acute.  "det" is set.
i de \emph{andre} Kritiker, saa viser det sig først, at den \emph{Indivi\-%
dualitetens} Form, der endnu klæbede ved Grundbegrebet i
Kr. d. reinen Vernunft, her er forsvunden, og at Fornuftkri\-%
tikens Antydning af en Gaaen sammen til Enhed af det Sub\-%
jective og det Objective her har faaet den \emph{bestemtere} Skik\-%
kelse af en Enhed i et fælleds Substrat, der muliggjør den
Møden af Natur og Aand, af Nødvendighed og Frihed, som vi
finde i Kunstens og det Organiskes Sphære. Og sammenligne
vi saa Begrebet med Centralbegrebet i Kr. der praktischen
Vernunft, saa finde vi atter et lignende Fremskridt. I Kr. d.
pr. Vernunft var Udvidelsen af Grundbegrebet til det \emph{Almene}
klarere, og Problemet om det \emph{Aandeliges Forhold til
Naturen} bestemtere berørt gjennem Fordringen af det Sæde\-%
liges Virkeliggjørelse i denne. Men den \emph{kritiske} Sondring,
der abstract adskilte Pligtens \emph{Form} fra den pligtmæssige
Handlings \emph{Indhold,} Menneskets \emph{Fornuftside} fra dets \emph{Na\-%
turside:} Driften, gjorde det uundgaaeligt, at denne det Sæde\-%
liges Virkeliggjørelse \emph{kun} kunde sættes som mulig gjennem
\emph{et postuleret Tredie:} en almægtig \emph{Gud.} Denne Guds\-%
forestilling --- thi et Gudsbegreb kan den ikke kaldes, --- viger
nu i Kr. d. Urtheilskraft Pladsen for Begrebet om det, der
% --- p. 76 ---
\setcounter{footnote}{0}%
\emph{ikke,} som staaende udenfor Modsætningerne, kun paa \emph{vil\-%
kaarlig} Maade kan forbinde dem, men som er deres \emph{En\-%
hedsgrund,} om det end nærmest bestemmes \emph{negativt,} i en
% p. 76 note *) -- printed mark *), superscript.
% PRINTER'S DEFECT, p. 76 note *): the figures "362." are set from a MARKEDLY
% HEAVIER fount than "259." on the same line -- the strokes are about half
% again as thick, checked at 1200 dpi.  This is a wrong sort, not emphasis:
% bold does not otherwise occur inside a paragraph or a note in this book
% (cf. the wrong-sort "i" logged at p. 42).  Transcribed as printed, i.e. plain.
Form, der anticiperer det \emph{Schellingske} Indifferentsbegreb.\footnote{Kr.
d. U. S. 259. 362.}
Idet \emph{Aandens Lov} og \emph{Naturens Lov} have deres Udspring
i det \emph{samme} Væsen, bliver \emph{det Sædeliges Virkelig\-%
gjørelse} i \emph{Naturen} mere end blot en Fordring\footnote{Sammest.
S. 245.}, og naar
Kr. d. Urtheilskraft i \emph{begge} sine Hovedafdelinger, i Læren
om den æsthetiske og i Læren om den teleologiske Dømme\-%
kraft, ligesom munder ud i det \emph{Ethiske,} gjennem Opfattel\-%
sen af det Skjønne som det Sædeliges Symbol, og gjennem
% GUILLEMET DEFECT, p. 76: »Endzweck» in the BODY opens AND closes with a
% right-pointing guillemet -- both marks read at 2400 dpi.  Cf. »afspeilet»
% (p. 4) and »christelige Philosophie» (p. 23).  In the ***) note below, the
% same German word is set with the normal inward pair, «Endzweck».  Neither is
% regularised.  p. 76 therefore contributes +2 to the guillemet balance.
% The German is set in the same upright roman as the Danish and is NOT
% letterspaced (measured 5 px), so it takes no \emph{} and never \textit{}.
% p. 76 is one of only two pages in the book carrying a ***) note (with p. 115).
Begrebet »Endzweck»\footnote{Sammest. S. 299. Naturens «Endzweck» skal søges
ude over Na\-%
turen selv.}, saa fører dette os ikke tilbage til
Standpunktet i Kr. d. pr. Vernunft; men hint for det Sædeli\-%
ges Virkelighed afgjørende Grundforhold mellem de to Love
bliver opfattet speculativt, begrundet ved det dybere Tilgrund\-%
liggende, det Omfattende, hvis Væsens Enhed ikke blot er Kil\-%
den til Modsætningerne, men ogsaa til deres Forsoning.

Ved Begrebet om et for Naturen og Aanden fælleds in\-%
telligibelt Substrat, der i sit Væsen bærer Kilden til Natur\-%
phænomenets Nødvendighedslov og til Aandslivets Frihedslov,
er \emph{den kritiske Idealismes Høidepunkt} berørt, dens
\emph{Grundtanke} nævnet. Men denne \emph{høiere} Tanke, der hos
\emph{Kant} bryder frem gjennem hans Grundproblems Consequent\-%
ser, kommer tilsyne ligesom \emph{imod hans Villie,} og idet han
forholder sig \emph{modstridende} til den og kun modstræbende
udtaler den. Og i samme Øieblik som den ligesom undslipper
ham, advarer han imod ethvert anmassende Forsøg af Erkjen\-%
delsen paa at drage den ind i \emph{sin} Kreds, og han søger lige\-%
% p. 76: the ɔ: id-est mark, read on the image at 1200 dpi (neither witness
% gives it -- both read "o:").  U+0254 followed by a colon, per the preamble.
som at kalde den tilbage ved at forvandle den til en \emph{blot
subjectiv} gyldig ɔ: kun for \emph{vor} Tænkning efter dens særegne
% --- p. 77 ---
\setcounter{footnote}{0}%
Indretning gyldig, Consequents, ligesom han gjør det ved hele
Begrebet om \emph{Hensigtsmæssighed,} fra hvilket der naaedes
til hiin afsluttende Tanke. Det \emph{høiere} Standpunkt, som en
skjult Tankemagts Nødvendighed har drevet ham til at antyde,
vil han føre tilbage \emph{til Fornuftkritikens subjective
Standpunkt} i Behandlingen af Erkjendelsesproblemet; men
den Tanke, der mod hans Villie behersker ham, er allerede
\emph{indirecte anerkjendt} i den forudgaaende Tankebevægelse,
og idet den engang er udtalt, er den udtalt med en saadan
Bestemthed, at de \emph{Consequentser,} som \emph{Kant} vil skjule for
sig selv og tilhylle for Andre, dog maae trænge igjennem.

% p. 77: a small raised speck stands between "subjective" and "Vending"
% (tesseract reads it as a quote, ABBYY as an apostrophe).  It sits at cap
% height, is much lighter than the ink and has no letter shape -- not type,
% and not transcribed.  "subjective" is spaced, "Vending" is not.
Hvad først den hele \emph{subjective} Vending angaaer, ved
hvilken \emph{Kant,} i samme Øieblik, som han i Kr. d. Urtheils\-%
kraft bringer de nye og afgjørende Begreber frem, ligesom vil
tage dem tilbage og afskjære deres Consequentser, saa hænger
den sammen med, og er i Formen identisk med, den \emph{Grund\-%
sondring mellem Subjectivt og Objectivt,} fra hvilken
% p. 77: "ud," is NOT letterspaced -- spacing.py's 1.10 flag is the two-letter
% artefact; measured 10 px on a single gap against 15--17 px in the certain
% runs on the same page, and the word reads close at 1200 dpi.
han efter hele sin historiske Stilling gaaer ud, og som gjør
sig gjældende overalt i hans Lære. Men denne \emph{Sondring,}
der kun bliver til en \emph{uforenelig Modsætning} ved den
\emph{Feilslutning,} at det, der viser sig som \emph{oprindelig} tilhø\-%
rende Subjectet, derfor maa sættes som \emph{kun} tilhørende dette
% p. 77 note *) -- printed mark *), superscript, standing after "Væsen" and
% before the comma.  There is NO dash after that comma: tesseract's "-" is a
% speck, read out at 1200 dpi.
og \emph{som udelukket fra det Objectives Væsen}\footnote{Jfr. Kants Deduction
af de aprioriske Anskuelsesformers subjective
Natur, i Kr. d. r. Vernunft.}, falder
% p. 77: a faint speck stands between "Plads" and "for" (tesseract reads a
% colon).  Not type; not transcribed.
med denne Feilslutning og giver Plads for den \emph{Enhed,} til
hvilken Kant selv i sine dybsindige Antydninger om Forholdet
% p. 77: THIS X IS ITALIC on the page (checked at 1200 dpi -- the slant and the
% italic X form are unmistakable), where the X of "det ubekjendte X" on p. 72
% is upright roman.  It is the same compositor's practice as the algebraic
% variables a, b, c, x on p. 5, which the edition sets \textit{} as printed
% (see transcription.tex §2).  Each occurrence follows its own page; do NOT
% harmonise p. 72 and p. 77.
mellem Tingen i og for sig, det ubekjendte \textit{X}, og Selvbevidst\-%
hedens rene Form, viser hen, og som allerede gjør sin Nød\-%
vendighed gjældende, saasnart der reflecteres paa \emph{Umulig\-%
heden} af, at det af ham som virkeligt og selvstændigt be\-%
handlede Objective kunde --- hvad han statuerer --- optages i
\emph{Subjectivitetens Former,} være for \emph{Subjectets Be\-%
vidsthed} naar hiin Modsætning havde Gyldighed.
% --- p. 78 ---
% p. 78 bears no note and no Greek.
Hvad det andet Punkt angaaer: de \emph{Consequentser,} til
hvilke det antydede Centralbegreb viser hen, saa komme de
tilsyne i de efterfølgende Systemer, og den afgjørende Betyd\-%
ning, dette Begreb har havt for disse, samt dets Forhold til
% GUILLEMET DEFECT, p. 78 (two of them): »fælleds intelligible Substrat for
% Naturen og Aanden» and »Substrat» each OPEN and CLOSE with a right-pointing
% guillemet; all four marks read at 1200 dpi.  This is the p. 78 site the recon
% pass had already confirmed.  Not regularised: p. 78 contributes +4 to the
% guillemet balance.
deres Grundbegreber, er let at paavise. I det »fælleds intelli\-%
gible Substrat for Naturen og Aanden», der hverken maa ud\-%
trykkes ved den ene eller den anden Modsætnings Bestemthed,
er den \emph{Schellingske Identitetsphilosophies} Grundbe\-%
greb \emph{umiddelbart} givet. Det Absolute, der er Naturens og
Aandens, det subjectivt og det objectivt Fornuftiges, Indiffe\-%
rents, Subject-Objectet, der ikke berøres af de quantitative
Differentser, men er som det Absolute i disse, er hint for Er\-%
kjendelsen utilgængelige, altsaa bestemmelsesløse, Intelligible,
der aabenbarer sig i de to forskjellige Former: Naturens og
Aandens, uden at disses Fremgaaen deraf eller deres Enheds
Mulighed begribes af Tanken. Og den \emph{intellectuelle An\-%
skuelse,} den \emph{intuitive Forstand,} som Kant hypothetisk
tillægger hint »Substrat», bliver for Schelling den postulerede
adæquate Form for Opfattelsen af det Absolute. Den \emph{Hegel\-%
ske} Philosophies Grundbegreb: den \emph{absolute Idee,} der i
sin evige Proces omfatter \emph{Naturen,} er hiin Enhed, seet fra det
Synspunkt, til hvilket Kant i sin Bestemmelse af Naturbegrebet
% p. 78: "(Endzweck)" is set in the same upright roman as the Danish and in
% ROUND BRACKETS here, not in guillemets as at p. 76.  Not letterspaced.
og i sin Opfattelse af Naturen som den, hvis sidste Maal (End\-%
zweck) maa falde udenfor Naturen, saa stærkt hælder over,
at det Overveiende deri, Grundbestemmelsen, er \emph{det Ideelle,
Subjectiviteten,} der saa opfattes som det \emph{absolut Spon\-%
tane,} sig selv Bestemmende, saaledes som den maatte tænkes
i Analogie med Opfattelsen af den rene Selvbevidsthed som
det Aprioriskes Kilde, og af Friheden som den bestemte ethiske
Handlens Udspring. --- I \emph{Kants} Philosophie have altsaa \emph{alle}
de Systemer, der betegne \emph{Hovedtrin} i den kritiske Idealis\-%
mes Udvikling, fundet deres Grundtanker. \emph{Fichtes} hele Phi\-%
losophie ligger som i en Spire i Begrebet om den transscen\-%
dentale Apperception; \emph{Schellings} i Begrebet om Naturens
og Aandens fælleds intelligible Substrat; \emph{Hegels} i det samme
% --- p. 79 ---
% p. 79 bears no note.  Greek: Κόσμος (one word).
Begreb, levendegjort ved at sees som Ur-Synthesen, i hvis
Energie al Virkelighed er begrundet. Den Tanke, Kant an\-%
tydede ved at sætte Selvbevidstheden som de aprioriske For\-%
mers Kilde, er gjennemført i den Hegelske Logik, og i denne
Gjennemførelse, i hvilken de aprioriske Grundformer faae Virke\-%
lighedsindholdets Betydning, er det Kantiske \emph{Jegbegreb,} der
allerede hos ham voxede ud over sig selv, udfoldet til det Be\-%
greb, der er den nyere idealistiske Philosophies Maal: til Be\-%
grebet om \emph{den absolute Idee,} der, som Idee, evigt er \emph{den
absolute, Naturen som sit Moment omfattende, Aand.}

% DIVISIONAL RULE, p. 79: short, CENTRED, white space above and below,
% standing between the paragraph above and the II. head below.  p. 79 bears no
% footnote, so this cannot be the (left-flush) footnote separator.
\divrule

% HEAD, p. 79 (L2, transcription.tex §3): the SECOND of the book's four theses
% (I. 50, II. 79, III. 137, IV. 170).  Set BOLD, centred, over seven lines, at
% a size clearly larger than the body but smaller than the p. 31 L1 head; set
% \large here.  It is NOT letterspaced -- bold letters are wider than
% spacing.py's per-glyph fit predicts, so a bold head is flagged exactly like
% Sperrsatz (BATCH-AGENT §3); measured glyph gaps in the head are 2--5 px,
% i.e. the body's unspaced value.
\begin{center}
\textbf{\large II. En Epoches historisk fremtraadte Grundtanke virker,
hævdende sig i sin eiendommelige Bestemthed, som Skjæbne, som ubevidst
hæmmende og begrændsende Magt, ligeoverfor den individuelle Tankes Stræben ud
over Udviklingsstadiets Grændse, henimod Realisationen af et mere fremskredet
Udviklingstrins Tankeform.}
\end{center}

Vi oplyse dette Forhold ved historiske Exempler fra den
\emph{græske} Philosophies \emph{Culminationstid} og dens \emph{Slutning,}
idet vi, efter først at have paaviist, hvorledes den græske Phi\-%
losophies ved dens Princip betingede Grundretning og Grund\-%
% p. 79 Greek: Κόσμος -- capital Κ, ACUTE over the first ο (the stroke rises
% left-to-right, read at 1200 dpi), cursive Greek fount.  Not letterspaced.
% Neither OCR witness contains a Greek character; both give "KoOfiog"/"Køouoc".
charakteer udpræger sig i den \emph{platoniske} Philosophies Grund\-%
begreber (Ideen --- Materien --- \gk{Κόσμος}) og i dens Bestem\-%
melse af disses indbyrdes Forhold, see den indvirke som til\-%
bageholdende og ubevidst hæmmende Magt paa den \emph{Aristo\-%
teliske} Philosophie under Grundbegrebernes videre Udvikling
og under Antydningen af en moderne Tids Tanker, og paa
\emph{Nyplatonismen} under dens Stræben hen mod en spiritua\-%
listisk Verdensanskuelse, der ligger udenfor den græske Aands
Grændse.

% DIVISIONAL RULE at the FOOT of p. 79, below the last line of text and BEFORE
% the p. 80 head -- the same placement as the rule at the foot of p. 15
% (RESUME-NOTES "things the recon pass got wrong", item 3), not the usual
% rule-above-head.  Centred, short.  There is consequently NO rule above the
% p. 80 head.  (The BIBLIOTHECA REGIA HAFNIENSIS stamp below it is a library
% stamp, not type, and is not transcribed.)
\divrule
% --- p. 80 ---
% HEAD, p. 80 (L3, transcription.tex §3): the first subdivision under thesis II.
%
% *** DISCREPANCY WITH transcription.tex §3, REPORTED: §3 describes the L3
% level ("1./2./3. under each thesis: ... 80, 94, 127 ...") as a BOLD centred
% head.  THIS HEAD IS NOT BOLD.  At 1200 dpi it is ordinary roman at body
% weight and body size, LETTERSPACED, centred over three lines -- i.e. the
% style §3 assigns to L4 (the α./β. heads at pp. 240, 250).  Compare the
% unmistakably bold II. head on p. 79 immediately above.  Measured glyph gaps
% here are 10--16 px against 2--5 px in the p. 79 bold head and 3--6 px in the
% body.  It is set \emph{} accordingly.  Whoever writes pp. 50, 94 and 127
% should check those heads the same way, and §3 should be corrected once the
% pattern is known. ***
\begin{center}
\emph{1. Den græske Philosophies Grundeiendommelighed, som udprægende sig i den
platoniske Philosophie.}
\end{center}
% p. 80 bears no note.  Greek: τόπος νοητός.
Den platoniske Philosophies Charakteer og Betydning er
ovenfor i Almindelighed angivet. Det var \emph{Platos} philosophi\-%
ske Gjerning, at han gjennem Dialektiken opdagede \emph{Ideen,}
Grækenlands Grundtanke; at han, ved i den at give Tanken
et sandt og evigt, ideelt Værende til Gjenstand, gjenfødte den
opløste Philosophie som en ny, Fortidens væsentlige Indhold
ligesom gjennem en Erindring i sig bevarende \emph{universel
Videnskab,} der ved Principets Idealitet helt igjennem faaer
Idealitetens Præg; og at han lader den for Erkjendelsen aaben\-%
bare Idee beherske \emph{hele} den menneskelige Tilværelse og forme
den til kunstnerisk Enhed. Idet han finder den \emph{Grundtanke,}
der er den græske Tænknings blivende Midtpunkt, og seer
hele Tilværelsen i dens Lys, staaer hans Philosophie som den
\emph{egentlige Type} for den \emph{græske} Tankes Udfoldelse; i hans
Grundbegreber udpræger sig dens \emph{Retning} og \emph{Charakteer,}
dens \emph{Stræbens Maal} og dens \emph{indre Begrændsning.}
Dette vil vise sig gjennem en kort Angivelse af hans Philoso\-%
phies Grundbegreber i deres indre Sammenhæng.

Vi tage Udgangspunktet fra den græske Tænknings Cen\-%
tralbegreb, fra \emph{Ideens} Begreb, til hvis Bestemmelse de al\-%
mindelige Grundtræk allerede ere givne i det Foregaaende.
Idet Ideen for første Gang aabenbarer sig for den \emph{Tingenes
Væsen} søgende Tanke, aabenbarer den sig som den \emph{høiere
Væren,} der hæver sig op over den \emph{sandselige} Tilværelses
rastløse Vorden; som frigjort fra Vordens Uro viser den sig
som den Væren, der \emph{ene} i Sandhed fortjener dette Navn, som
den \emph{evige} og evigt med sig selv lige Væren, som det for sig
% p. 80: a small raised speck follows "Usandselige," (both witnesses give it as
% an apostrophe).  Not type; not transcribed.
værende \emph{Substantielle,} der \emph{negativt} bestemmes som det
% p. 80 Greek: τόπος νοητός -- acute over the first ο of τόπος and over the ό
% of νοητός, both read at 1200 dpi (the strokes rise left-to-right).  Cursive
% Greek fount, not letterspaced.  ABBYY gives "zonog vor/iog", tesseract
% "zomog voride"; neither carries a Greek character.
\emph{Usandselige,} som det i \gk{τόπος νοητός} værende \emph{Ulegem\-%
lige.} Ideens Begreb udfolder sig til Begrebet om en \emph{Ideer\-%
nes Verden} idet Eleaternes Ene for den platoniske Dialektik
bestemmer sig til Kredsen af de concrete Ideer, i hvilke en
% --- p. 81 ---
\setcounter{footnote}{0}%
\emph{Flerhed} faaer Plads, men som \emph{behersket af Enheden,}
en \emph{Ikke-Væren,} men kun som \emph{Forskjel i Væren,} som
Andet-Væren. Ideeverdenen bliver \emph{Enheden af de orga\-%
nisk forbundne Ideer,} og Platos \emph{Stræben} gaaer ud paa
at opfatte denne Ideeverden, hvis Grændser omfatte det \emph{hele}
% p. 81: ABBYY reads a comma and tesseract a colon after "Almindeligt".  At
% 1200 dpi the mark is a very pale baseline speck, far lighter than the ink,
% and the sense requires no stop.  Per TRANSCRIPTION-PLAYBOOK §2 rule 3 it is
% treated as a speck and not transcribed.
Gebet, paa hvilket et \emph{Almindeligt} kan forfølges, og først
sættes hvor det Almenes Spor taber sig i Vordens begrebløse
Mangfoldighed og Uro, som en \emph{Trinrække,} der udfolder sig
af \emph{det Godes} høieste Idee, Værens og Videns Kilde, og naaer
ned til de laveste Artbegreber.

Idet den høieste Idee sættes som Værens \emph{Aarsag,} inde\-%
slutter dette, at overhovedet Ideen, hvis primitive Bestemmelse
% p. 81 Greek: δύναμις -- acute over the υ, read at 1200 dpi.
er Væren, sees fra et nyt Synspunkt, som \emph{virksom Aar\-%
% p. 81 Greek: αἰτία -- SMOOTH breathing over the ι (a comma opening to the
% right, not the rough breathing's "C"), and an acute over the second ι.  Both
% read positively at 1200 dpi.  ABBYY gives "ahia", tesseract "afrta".
sag,} som Kraft (\gk{δύναμις}), som det, der \emph{begrunder} Tilvæ\-%
relse (som \gk{αἰτία}), og der træder et Forhold frem mellem Ideen
og det, der i sin Tilværen kun har Væren ved den og som deel\-%
agtigt i den. Ideen bliver, i Forholdet til dette, seet som \emph{Ur\-%
billedet for de sandselige Afbilleder,} eller som det
% p. 81: "i det" between the two spaced phrases is NOT spaced -- measured, on
% the same line, blivende 17 | Ene 16 | i (single letter) | det 6 |
% vexlende 16 | Mangfoldige. 14 | Der 6.  The compositor leaves short function
% words unspaced inside a spaced run (cf. "og" at p. 72, "af" at p. 83); set as
% printed rather than smoothed into one \emph{} group.
\emph{blivende Ene} i det \emph{vexlende Mangfoldige.} Der søges
saa en Begrebsbestemmelse for det, der, skjøndt i Væsen og
Væren forskjelligt fra Ideen, har \emph{Deelagtighed} i Ideen: for
den \emph{sandselige} Tilværelse, og Kjendemærkerne for den som
seet fra \emph{dens Sandselighedsside} blive: \emph{Vorden, Deelt\-%
hed} og \emph{relativ Væren.} Men disse det Sandseliges Bestem\-%
melser, der, idet de betegne det som en, om end ufuldkommen,
\emph{Væren,} paa den ene Side vise hen til \emph{Ideen,} det \emph{Værende,}
som en Factor, vise paa den anden Side, gjennem den ved
dem betegnede \emph{Ufuldkommenhed,} hen til en \emph{anden} Fac\-%
tor, der maa være \emph{det Ideen Modsatte.} Forat absolut Væ\-%
ren kan blive til relativ, Lighed med sig selv til Vorden og
Forandring, Enhed til Deelthed, fordres et Princip, der til Be\-%
stemmelser har Ikke-Væren, absolut Foranderlighed og enheds\-%
løs Mangfoldighed, og disse Bestemmelser charakterisere det
% p. 81 note *) -- printed mark *), superscript.  RUN-OVER NOTE: this note is
% cut off at the foot of p. 81 after "naar man," and CONTINUES at the head of
% p. 82, where the note area opens with matter carrying NO mark ("gaaende ud
% fra det Vordende, ... et Virkeligt vorder.").  The whole note is set here in
% one \footnote{} at its call, per transcription.tex §4; a % comment at the
% p. 82 page marker below records that it continues there in the print.
%
% p. 81 Greek: τὸ ἄπειρον -- the GRAVE over the ο of τὸ is positive (the stroke
% falls left-to-right; calibrated against the acutes of Κόσμος on pp. 79 and 83,
% which rise).  THE MARK OVER THE α IS NOT RESOLVABLE.  At 1200 dpi (4x
% interpolation over the 300 ppi original, no 1-bit mask) it is a blob about
% 5x5 source pixels with a hook at the left and a falling stroke at the right;
% the breathing and the acute are ~2 px apart and cannot be separated, so the
% smooth ἄ cannot be told from the rough ἅ.  Per TRANSCRIPTION-PLAYBOOK §2
% rule 2 nothing is supplied: the α is set UNACCENTED.
% This is the known-unresolvable case named in transcription.tex §2 and met
% before at p. 52.  IT COULD NOT BE RESOLVED HERE EITHER, so p. 52 stands as
% it is; the two sites should be treated alike.
formløse, altid vordende \emph{Stof}\footnote{\emph{Plato} bestemmer
\emph{Materien} ved det, der bliver tilbage, naar man,
gaaende ud fra det Vordende, tager Alt bort fra dette, der tilhører
Formen; han kommer altsaa til Materiens Begreb \emph{ved en Slut\-%
ning tilbage fra Vorden;} \emph{Aristoteles} bestemmer Materien
ved det, der er \emph{Betingelsen for, at en Vorden kan blive
mulig} som den, hvorigjennem \emph{et Virkeligt} vorder.} (\gk{τὸ απειρον}),
der som saa\-%
% --- p. 82 ---
% RUN-OVER FOOTNOTE: the note area at the foot of p. 82 opens with unmarked
% matter -- "gaaende ud fra det Vordende, ... hvorigjennem et Virkeligt
% vorder." -- which is the continuation of p. 81's *) note.  It is set in full
% inside that \footnote{} above and is NOT repeated here.  p. 82 opens no note
% of its own, so no \setcounter is needed.
% Greek on p. 82: Κόσμος (three times, none letterspaced) and αὐτοζῶον.
dant, i sin fuldstændige Mangel paa bestemt Form og Væren,
hverken kan være Gjenstand for Sandsning eller Tænkning.

I \emph{Sandseverdenen,} der ved sit Væsen viser hen til et
\emph{dobbelt} Princip, \emph{mødes} saaledes Modsætningerne, og, \emph{som
Modsætninger,} maae de \emph{begge} gjøre sig gjældende; ikke
blot Ideen maa være virksom, men Materien maa blive virk\-%
som i Modsætning til Ideens Virken; det vil sige: som \emph{be\-%
grændsende} og \emph{hæmmende} Ideen. Men medens denne
hæmmende Virksomhed fremhæves hvor den sandselige Tilvæ\-%
relse sees fra \emph{sin Sandselighedsside,} d. v. s. i det Sand\-%
seliges Udstykning, saa træder den tilbage naar Sandseverdenen
sees \emph{som Totalitet;} Modsætningen bliver saa \emph{af Anskuel\-%
sen} ført tilbage til en harmonisk Enhed, Materiens hæmmende
Virken bliver ligesom skjult under Ideens Afglands, og Sandse\-%
verdenen bliver seet i en anden Belysning, som \emph{den af Ideen
formede og besjælede levende Heelhed,} som det \emph{gud\-%
dommelige Verdenskunstværk,} som \gk{Κόσμος}. Seet som
\gk{Κόσμος} er Sandseverdenen opfattet som det, i hvilket det sande
Værende, \emph{Ideen,} der i sig som Form bærer Maal og Grændse, er
det \emph{Herskende,} og ved Ideen, der gjennem Materien aaben\-%
barer \emph{sig selv,} er den et \emph{Guddommeligt,} Uforgængeligt,
i sig harmonisk Afsluttet, Levende, Besjælet, Fornuftbegavet.
% p. 82 Greek: αὐτοζῶον.  Read at 1200 dpi: a SMOOTH breathing over the υ (a
% comma-shaped mark opening right, not the rough breathing's "C"), and a
% CIRCUMFLEX (tilde) over the ω.  NO IOTA SUBSCRIPT is visible under the ω at
% 1200 dpi, so none is supplied and the word is set ζῶον, not the ζῷον of
% modern orthography; if a subscript is there it is below what this scan can
% show.  ABBYY gives "avzo£wov", tesseract "adroløo»".
\gk{Κόσμος} er, efter Plato, det fuldkomne levende Væsen, dannet
til Lighed med det Levendes Idee (\gk{αὐτοζῶον}), efter Guddom\-%
mens Billede, den indbefatter i sit Legeme Heelheden af det
% p. 82: the ɔ: id-est mark, read on the image (both witnesses give "o:").
Legemlige, er ved sin Sjæl i Besiddelse af et eget ɔ: den selv
tilhørende, Liv og af guddommelig Fornuft; den berøres aldrig
af Alder eller Undergang; den er ikke blot det fuldkomne Af\-%
billede af den evige og usynlige Guddom, men \emph{selv en salig}
% --- p. 83 ---
\setcounter{footnote}{0}%
% p. 83 line 1 is set FLUSH LEFT: it continues the sentence and the spaced run
% ("men selv en salig / Gud,") from p. 82.  No paragraph break.
% Greek on p. 83: μονογενὴς, and Κόσμος four times (three in the body, all
% unspaced; one in the note, LETTERSPACED).
% p. 83 Greek: μονογενὴς -- the mark over the η is a GRAVE (the stroke falls
% left-to-right), calibrated on the same page against the acute of Κόσμος two
% lines below, which rises.  Standard Greek would give the acute μονογενής
% before the closing parenthesis, so this is the compositor's, and it is set as
% printed.  ABBYY gives "jxovoysvrji;", tesseract "uovoyevyc".
\emph{Gud,} eneste i sin Art (\gk{μονογενὴς}), \emph{sig selv nok, og ikke
trængende til noget Andet.}

I denne Anskuelse af \gk{Κόσμος}, den af Ideen beherskede
Sandseverden, som den salige i sig hvilende «tilblevne Gud\-%
dom», Ideens og den evige Guddoms fuldkomne Afbillede, frem\-%
træder tydeligt, naar vi fjerne det ved \emph{Forestillingsad\-%
skillelsen} mellem Ideeverdenen og den ideebestemte Sandse\-%
verden hæftende, paa Ideen uanvendelige Rumsforhold, og
den af Plato selv ophævede mythiske Forestilling om en Ver\-%
densdannelsens \emph{Historie,} det Maal, efter hvilket den græske
% p. 83: within this long spaced run "Opfattelsen", "af" and "som" are NOT
% spaced (measured 5--6 px) while "al", "Tilværelse", "et" and everything from
% "altomfattende" to "Idee." are (13--19 px).  Set as printed.
\emph{Tanke} stræbte hen som det høieste: \emph{Grundmodsætnin\-%
gernes Enhed,} Opfattelsen af \emph{al Tilværelse} som \emph{et alt\-%
omfattende, levende, plastisk Kunstværk, værende
% p. 83: "har al ... omfatter alt Stof" is NOT letterspaced (checked at
% 1200 dpi); spacing.py's flags on the two-letter "al"/"alt" here are the
% single-gap artefact.
ved og besjælet af den iboende Idee.} \gk{Κόσμος} har al
Ideens Fylde i sig, omfatter alt Stof, og det saaledes, at Ideen
behersker det og gjør det deelagtigt i Væren. Fastholdtes
Ideens Begreb strængt i \emph{Tankens} Form, saa vilde \emph{Ideever\-%
denen} og \emph{Sandseverdenen} ikke kunne sættes som to virke\-%
lige Verdener, der stode ved Siden af hinanden. Al Virkelig\-%
hed tilkommer \emph{Ideen;} det Sandselige vilde altsaa, forsaavidt
Væren tillægges det, kun kunne sættes som værende i \emph{En\-%
heden med Ideen.} Fra dette Synspunkt vilde Rumsforhold
\emph{ikke} blive adskillende (jfr. Parmenides), Ideen og Phænome\-%
net ikke blive sidestillede, Urbilledets Forsigværen ikke blive
forskjellig fra dets Væren som iboende Afbilledet. Og idet
Ideen som Form har Materien, Stoffet, til sit \emph{nødvendige}
Supplement, maatte deres Forhold bestemmes som et \emph{evigt,}
og det ved Forholdet Satte: Verden, maatte faae den samme
Bestemmelse af Evighed, saa at den mythiske Forestilling om
en Verdensdannelse ved Verdensformeren (Demiurgen) vilde
% p. 83 note *) -- printed mark *), superscript.  RUN-OVER NOTE: it breaks off
% at the foot of p. 83 after "til Problem af" and CONTINUES at the head of
% p. 84, where the note area opens with unmarked matter ("Plato, og dette
% Problem vilde, ... ved Platos Forudsætninger.").  The whole note is set here
% in one \footnote{} at its call; see the % comment at the p. 84 marker.
% The Κόσμος inside this note IS letterspaced (its glyph gaps run 17 px against
% 5--7 px for the unspaced words of the same note, and against the unspaced
% Κόσμος of the body above), so it takes \emph{\gk{...}} per BATCH-AGENT §3.
forsvinde i Ideens nødvendige Begrebsforhold, Tidsbestemtheden
i et Evighedsforhold\footnote{Forholdet mellem den \emph{evige Væren} og
\emph{Tidssuccessionen in\-%
denfor den værende} \emph{\gk{Κόσμος}} er neppe gjort til Problem af
Plato, og dette Problem vilde, fra den græske Tænknings Stand\-%
punkt, \emph{ikke} kunne løses. Med Eleaterne at lade Tiden være et
\emph{Skin,} var udelukket ved Platos Forudsætninger.}. At \gk{Κόσμος}
var den \emph{tilblevne} Gud,
% p. 83: COMPOUND LETTERSPACING -- only the first element is spaced,
% "B e g r e b s forhold", exactly as at pp. 4, 10, 14 (RESUME-NOTES,
% "Conventions settled during the run").  Measured at 1200 dpi.
vilde da kun udtrykke dens \emph{Begrebs}forhold til dens Factorer.
% --- p. 84 ---
% RUN-OVER FOOTNOTE: the note area at the foot of p. 84 opens with unmarked
% matter -- "Plato, og dette Problem vilde, ... ved Platos Forudsætninger." --
% the continuation of p. 83's *) note, which is set in full inside that
% \footnote{} above and is NOT repeated here.  p. 84 opens no note of its own.
% The continuation ends with a FULL STOP after "Forudsætninger" (read at
% 1200 dpi; both witnesses waver between a comma and a stop).
% p. 84 carries no Greek but the one word Κόσμος, and no ɔ:.
Men netop paa det Punkt, hvor Tænkningens Stræben
efter Enhed ligesom berører sit Maal, kommer den græske
Tankes, ved Grundanskuelsens Art betingede, Eiendommelighed
tilsyne som det, der atter løsner Enheden, og, tiltrods for Tan\-%
kens Stræben, fører til en blivende Sondring af de Grundmo\-%
menter, der som ved et Anskuelsens Postulat vare forbundne
i Verdenskunstværket.

I den græske Aands Grundretning: at opfatte Tilvæ\-%
relsen i \emph{Skjønhedens} Form, ligger en Impuls, der be\-%
standig maa drage Modsætningerne henimod hinanden forat
forbinde dem i \emph{Anskuelsen} af \gk{Κόσμος}, og som maa
drive \emph{Tanken} til at give den anskuede Enhed sin Form,
til at udtrykke den gjennem Begrebet. Men for at op\-%
fatte Enheden som \emph{Begrebsenhed,} maa Tanken først give
Enhedens \emph{Momenter Begrebets} Form, hvilket den kun kan
gjøre ved at løse dem fra deres \emph{umiddelbare,} for Anskuel\-%
sen givne Forbindelse. Men ifølge det \emph{Anskuelsens} Mo\-%
ment, som den i Skjønheden, i Harmoniens Umiddelbarhed,
hvilende Aands Tanke bestandig maa bære i sig og ved
hvilket dens Form bestemmes, maae de sondrede Momen\-%
ters Begreber \emph{fæstne} sig i Sondringen som \emph{hvilende}
Modsætninger i Umiddelbarhedens Form, hvilke Tanken ikke
mægter at bringe i Bevægelse og at føre tilbage til Enhed,
hvormeget den end forsøger det. Forholdet maa blive dobbelt
og tvetydigt, Tankens Virksomhed deelt mellem Forening og
Sondring, men saaledes, \emph{at Sondringen bliver herskende,}
at Tanken, hvis Bevægelse bindes i Anskuelsens hvilende
Umiddelbarhed, ikke mægter, i \emph{Begrebets Form} at virke\-%
liggjøre sin Enhedsstræben.

% =====================================================================
% BATCH SEAM, p. 84 | p. 85.  The last line of p. 84 breaks the word
% "Opfat-telsen" across the page: p. 85 opens "telsen af Ideen, af Materien,
% ..." (read on the image).  The discretionary hyphen and the newline-eating %
% are written here; whoever splices pp. 85--98 must join `Opfat\-' to `telsen'
% and must NOT let splice.py's blank line above the marker introduce a
% paragraph break -- p. 85 line 1 is flush left and continues this sentence.
% Run joints.py --fix after the splice.
%
% p. 85 was also checked for a run-over note: it opens its own *) note
% ("Dialektiken gaaer altid ud fra givne bestemte Begreber som
% Forudsætninger."), so nothing runs over from p. 84 (BATCH-AGENT §5 check).
% =====================================================================
Vi ville see Indvirkningen af denne Tankens eiendomme\-%
lige Bestemthed ved Anskuelsen gjøre sig gjældende i Opfat\-%
% --- p. 85 ---
% p. 84 ends mid-word "i Opfat-"; this page opens with "telsen".  p. 84's
% note area was read on the image: it closes with a full stop
% ("...var udelukket ved Platos Forudsætninger.") and does NOT run over.
% p. 85's note area opens with its own *) mark.  BATCH-AGENT §5 satisfied.
\setcounter{footnote}{0}
telsen af Ideen, af Materien, af deres indbyrdes Forhold, og
af det Begreb, der udtrykker deres Forbindelse.

Idet \emph{Ideens} Begreb først søgtes, blev det fundet gjen\-%
nem den Dialektik, der knyttede sig til den sokratiske, som
\emph{Tingenes Begreb}, deres \emph{almene Væsen}, og derfor faldt
Ideeverdenens Grændser netop sammen med det Almenes. Og
% p. 85, l. 7: a short horizontal mark stands in the space between "som"
% and "det".  It is not type -- no sort of that shape occurs in this
% book -- and is not transcribed (§7).
idet Ideen opfattes som det almene, begrebsmæssige Væsen,
saa viser denne Bestemmelse hen paa \emph{nødvendige Begrebs\-%
relationer}, den antyder \emph{Væsenets indre Bevægelse},
% p. 85, ll. 10--11: "dets" is set UNSPACED both times inside the spaced
% run, as printed; the same happens with "deres" in l. 5 above.
dets \emph{Virksomhed}, dets \emph{intellectuelle Charakteer}; men
idet Tanken, der giver den denne Bestemmelse, er den Tanke,
\emph{der er i Enhed med Anskuelsen}, saa bliver \emph{Ideen}, som
ogsaa \emph{Navnet} antyder, til \emph{Urbilledet, det ideelt ansku\-%
ede hypostaserede Væsen}; den faaer ligesom plastiske
Omrids, den \emph{hvilende Værens Form}, medens \emph{Virksom\-%
hedens} Bestemmelse, der søges hævdet for Ideen, træder til\-%
bage og kun viser sig som et Postulat; og Ideen faaer en Plads
% p. 85: the accent on the omicron of νοητὸς is positively a GRAVE -- the
% stroke runs DOWN TO THE RIGHT at 1200 dpi -- where an acute would be
% expected before the following comma.  Calibrated against the acutes of
% μόγις πιστόν (p. 86) and Κόσμος (p. 89), whose strokes run down to the
% LEFT.  Transcribed as printed.
i det Rum, der, om end det end betegnes som \gk{νοητὸς}, dog,
som den ideelle \emph{Anskuelses Schema}, beholder Rummets
\emph{Navn og Væsen}. Og idet Ideerne skulle sees \emph{som Enhed},
saa bliver, netop paa Grund af det \emph{Anskuelsens Moment},
der gjør sig gjældende i Opfattelsen af dem som en \emph{adskilt}
Mangfoldighed, denne Enhed \emph{ikke} forklaret; de danne en ord\-%
% p. 85, l. 24: a small tick stands before "hvilende" (ABBYY reads it as
% an apostrophe).  Not type; not transcribed.
net \emph{Heelhed}, en \emph{Stat}, en \emph{Verden af hvilende Ideer}; men
trods Platos Tankes Fordring: at see dem fremgaae af den
høieste, fyldigste Idee, bliver Enheden aldrig udledt, Ordningen
kun en \emph{udvortes} Ordning af det fra Anskuelsen af den ende\-%
lige Tilværelse Tagne, gjennem Inductionen Bestemte. De hvi\-%
lende Urbilleder gruppere sig, ordne sig under hverandre, faae
% p. 85, l. 30: the run ends with a single "i" at the line end, which
% cannot itself show letterspacing; "hverandre," on l. 31 is unspaced, so
% the "i" is left outside the \emph{}.
idethøieste \emph{gjennem Sammenligningen en Deelagtighed} i
hverandre, men fremgaae ikke med indre Nødvendighed, gjen\-%
nem en virkelig dialektisk Udvikling\footnote{Dialektiken gaaer altid ud
fra givne bestemte Begreber som Forudsætninger.}, en Begrebsudfoldelse,
% --- p. 86 ---
af en høieste Idee. De staae ved Siden af hverandre som de
olympiske Guder, og den herskende blandt dem, Ideernes Zeus,
det Godes Idee, bliver dog kun en Idee ved Siden af de
andre. Ideen faaer ikke Tankens indre Bevægelse; den bliver,
skjøndt intelligibel og af Tankens Art, Tankens \emph{objective},
ved hvilende Væren bestemte \emph{Gjenstand}.

Naar saa \emph{Ideens Modsætning} kommer frem som et
\gk{μόγις πιστόν}, gjennem en Slutning fra den \emph{endnu ikke
forklarede} Sandseverdens Bestemmelser, faaer den, som det
\emph{Værendes} absolute Modsætning, \emph{Ikke-Værens} Bestemmelse.
Men idet den sættes som Modsætning til den for \emph{en intel\-%
lectuel Anskuelse} hypostaserede Idee, \emph{hypostaseres
Ikke-Væren selv}, og bliver, som det anskuede Værendes
\emph{Modsætning}, et anskuet \emph{Ikke-Værende}, der, netop som
\emph{absolut} Modsætning til hint, staaer som en \emph{uafledt}, altid
tilværende Existents, og ligesom med Magt hævder sig for den
anskuende Tanke som \emph{et objectivt existerende Ikke-%
Værende}, som hint \gk{ἄπειρον}, der i sin uophørlige Vorden
% p. 86: the mark over the alpha of ἄπειρον IS resolvable here.  At
% 1200 dpi it separates into two components with a clear gap: the right
% one is a stroke running down to the LEFT (acute), the left one a curl
% whose opening faces LEFT.  That left shape is the SMOOTH breathing, on
% the calibration built in this batch -- it is the same glyph as on
% Ἀνάγκη (p. 87), οὐσία (p. 97), ταὐτὸν and θἄτερον (p. 88) and ἀρίστου /
% ἰδέαν (p. 92) -- and it is NOT the ROUGH breathing of ὕλη (p. 98),
% whose curl opens to the RIGHT.  The file header records this compound
% as unresolvable at pp. 52 and 81; those two sites should be re-read.
unddrager sig Sandsningens og Tankens bestemte Former,
men i denne Mangel paa Maal og Begrændsning, i denne
Bestemmethed som usynlig og formløs Væsenhed, dog bliver
som et Efterbillede for Aandens blændede Øie, som et uroligt
Bølgende, og idet den forvirrede Bevægelse bringes til Hvile,
idet de vexlende Skygger af Væsener svinde, ligesom klarer
sig til det anskuede Rum, Vordens formløse Moderskjød.

\emph{Ideens absolute} Modsætning til Materien ophæves ikke
derved, at \emph{Plato}, i sin Polemik mod \emph{Eleaternes} abstracte
Enhedslære, sætter en Ikke-Væren ogsaa i Ideerne; thi Ideernes
Ikke-Væren er kun Forskjellens relative Ikke-Væren, bestem\-%
tere: den kun ved en sammenlignende Reflexion satte Ikke-%
Væren, der ikke faaer sin Plads som Væsensbestemmelse i det
blot positivt Værende, som udelukker al indre Negativitet, og
bliver det mod denne Andet-Væren ligegyldige Forskjellige.
Denne Ikke-Væren er altsaa selv Væren, den er netop Ikke-%
Væren af et Andet som en Væren af dette Bestemte. Mate\-%
% --- p. 87 ---
riens Ikke-Væren derimod er den for Ideen fremmede absolute
Ikke-Væren, en iogforsigværende Ikke-Væren. Derfor bliver
Materien altid, fra Ideens Synspunkt, \emph{en ubegreben For\-%
udsætning}, der vel kan omskrives ved andre Udtryk, f.\ Ex.
\gk{Ἀνάγκη}, men ikke forklares, ikke deduceres ud fra den blot
positivt bestemte Idee.

Naar nu saaledes Materiens Begreb kommer frem, uafledt
af Ideens, som dens absolute Modsætning, saa faaer det ikke
blot den \emph{Modsigelse i sig} at betegne \emph{et værende Ikke-%
Værende}, der \emph{fæstner sig ligeoverfor Ideens abso\-%
lute Væren}, \emph{og i Sandseverdenen}, i hvilken denne Ideens
absolute Væren bestemmes til relativ, bliver \emph{virksom Fac\-%
tor}; men Materiens Forhold til Ideen bliver, ifølge hiins
% p. 87, l. 13: a small superscript tick stands after "til" (ABBYY reads
% it as an apostrophe, tesseract as a closing quote).  Not type; not
% transcribed.
Grundbestemmelse, et saadant, \emph{at en Enhed ikke kan til\-%
veiebringes}, og at derfor \emph{Tilværelsen af den Sandse\-%
verden}, fra hvis \emph{uafviselige Givethed} der naaedes til
Materiens Begreb, og gjennem hvilken dette bestandig hævder
sig en Gyldighed, bliver \emph{uforklarlig}, idet Sandseverdenens
Tilværelse netop forudsætter Enheden. Her bliver den prin\-%
cipielle Dualisme aabenbar.

\emph{Sandseverdenens Væren}, der \emph{ikke} kan forsøges dedu\-%
ceret ud fra Ideen, fordi den indeslutter den fra den rent
positive Idee absolut udelukkede Materiens Ikke-Væren, kan
\emph{Plato} kun stræbe at forklare paa den Maade, at han søger
\emph{et Forbindende, et Sammenknyttende} mellem de to
som \emph{færdige} forudsatte Modsætninger; men et Forsøg af
den Art viser sig let som efter Sagens Natur uigjennem\-%
førligt.

\emph{Plato} søger det Sammenknyttende deels som \emph{det virk\-%
somt Forbindende}, der finder sit Udtryk \emph{i Verdens\-%
kunstneren} (Demiurgen), deels som \emph{det begrebsmæssigt
Mæglende}, der repræsenteres ved \emph{Verdenssjælen}.

Forestillingen om \emph{Demiurgen} er imidlertid \emph{kun det
mythiske} Udtryk for det af Plato i Ideen \emph{fordrede}, men
\emph{ikke realiserede} Bevægelses- og Virksomhedsprincip, som
% --- p. 88 ---
% p. 88 is the worst pencil page in the book (§7).  A two-line cursive
% pencil note stands in the LEFT MARGIN level with ll. 6--7, and a faint
% pencil stroke runs under "af Ideens" in l. 6.  The ABBYY layer splices
% the marginal note straight INTO the text, producing "'pv^hj £", a stray
% "U" on a line of its own, and "'/fåfl't $(/*•*"; tesseract renders the
% same marks as "fa", "(", "i /".  tesseract also reads a stray "70 Å"
% and "ø !" at the left edge.  NONE of it is type.  It is all deleted.
% The true reading of ll. 6--7, verified on the image at 1200 dpi, is
% "Forklaringen af Ideens og Materiens Forbindelse maatte / da søges i
% Verdenssjælen.", with "Verdenssjælen" letterspaced.
\setcounter{footnote}{0}
derfor bliver anskuet \emph{udenfor} Ideen og personificeret,
men uden at dog den verdensformende Guddoms Væsen og
Væren kan blive nogen anden end netop \emph{Ideens}, og hiin
Personification kan altsaa ikke forklare hvad der ikke var
forklarligt ved Ideens Begreb.

Forklaringen af Ideens og Materiens Forbindelse maatte
da søges i \emph{Verdenssjælen}. Hvorledes Plato opfatter dennes
Væsen, fremgaaer af den Maade, paa hvilken dens Dannelse
motiveres, og dens oprindelige Forhold til Stoffet beskrives.
Universet, der skulde dannes saa fuldkomment som muligt,
skulde være fornuftigt, men da Fornuften ikke kunde meddeles
det uden en Sjæl, satte Demiurgen, da han dannede Verdens\-%
altet, Fornuften i en Sjæl og Sjælen i dens Legeme. Forat
Sjælen kan forbinde Fornuften med Legemet, det Ideelle med
det Materielle, bliver den dannet af en sammentvungen Blan\-%
ding af Ideens og det Ikke-Værendes Substants, til hvilken \gk{ταὐτὸν}
% p. 88: ταὐτὸν carries the CORONIS (the smooth-breathing sort) over the
% upsilon and a GRAVE over the omicron -- the stroke runs down to the
% right, as on νοητὸς (p. 85).  θἄτερον likewise carries the coronis over
% the alpha, followed by an acute: the two components separate cleanly at
% 1200 dpi.  Modern editions print ταὐτόν / θάτερον.
og det dermed svært forenelige \gk{θἄτερον} sættes; denne Blanding
deles, formes og udbredes efter de harmoniske Talforhold, og i dens
Ramme indføies det chaotisk bølgende Materielle efterat være
fordelt i de physiske Elementers mathematisk bestemte Rums\-%
former, hvorved det Jorden som hvilende Midtpunkt omslut\-%
tende System af stjernebærende Himmelsphærer opstaaer.
\emph{Verdenssjælen}, der i sig indeslutter alle \emph{de harmoniske
Forhold}, bestemmes tillige som \emph{det ved sig selv Be\-%
vægede}, der, som bevæget, har \emph{Liv og Bevidsthed}, og
ifølge denne sin Natur bliver den det, der er \emph{Aarsag til
den i Universet herskende Orden}, det, der \emph{bevæger
Verden}, og det, der er \emph{Bevidsthedens Princip}\footnote{De forskjellige
Momenter i Verdenssjælens Væsen: Tal, Bevidsthed
og Bevægelse, har den græske Philosophie allerede tidligere bragt
i Forbindelse. I \emph{Tallet saae Pythagoræerne} tillige Betingelsen
for Bevidstheden, og med Bevægelsens Begreb bringes Bevidstheden
sammen af Heraklit og Anaxagoras. Paa lignende Maade lader
Plato selv, i Sophisten, en \gk{κίνησις} være Betingelse for Erkjendelse,
Liv, \gk{Νοῦς} og \gk{φρόνησις}, og i Timæos kalder han ligefrem Tænk\-%
% RUN-OVER FOOTNOTE.  This note begins at the foot of p. 88 and its
% continuation fills the note area of p. 89, where it stands UNMARKED.
% Per the file header §4 the whole note is set here, inside this single
% \footnote{}; a % comment at the p. 89 page marker records the break.
% The print divides it at "Tænk- | ningen".
ningen en Bevægelse, og bestemtere: en Kredsbevægelse
(\gk{περιφορά}) i Sjælen. Verdenssjælens Erkjenden udleder Plato deels af dens
Sammensætning; deels deraf, at den, som inddelt efter harmoniske
Talforhold, er sammensluttet med sig selv; deels af dens i sig selv
tilbagevendende Kredsbevægelse.}. Alle
% --- p. 89 ---
% p. 89 carries NO footnote call of its own.  Its note area holds only the
% continuation of p. 88's note ("ningen en Bevægelse ... tilbagevendende
% Kredsbevægelse."), which is set in full inside the \footnote{} at its
% call on p. 88 above.  Hence no \setcounter here.
disse Bestemmelser vise hen paa \emph{Ideen}; de harmoniske Tal\-%
forhold constitueres netop ved at Ideen, det, der har Maal og
Form i sig, gaaer ind i Blandingen; Bevægelse (Selvbevægelse)
og Tænkning ere Bestemmelser, der begge ligge i Ideens
Natur, saaledes som \emph{Plato har søgt at opfatte den}. Gjennem
Bestemmelsen af det de mathematiske Forhold Indesluttende
bliver Verdenssjælen Udtrykket for Ideen, \emph{seet fra den Side,
efter hvilken den staaer i Forhold til Phænomenet},
og, som bestemmende dette efter de mathematiske Talforhold,
fremtræder som det blivende Maal, som det Formende og
Ordnende i den legemlige Tilværelse. I \emph{det Mathematiske}
seer Plato en \emph{Mellemsphære} mellem den \emph{forsigværende}
\emph{Idee og det Legemlige}: Tallene udtrykke et Blivende og
% p. 89, l. 14: the "i" of "men i det Sandselige" is set in BOLD -- a
% heavier stem and a larger dot than the surrounding roman, measured at
% 1200 dpi.  It is the compositor's way of emphasising a single letter,
% which cannot be letterspaced; the same device recurs twice on p. 95.
% The file header states that bold occurs only in display heads; that is
% wrong, and should be corrected.
\emph{Uforanderligt}, men \textbf{i} det Sandselige, hæftende ved dette og
dets Flerhed og dets Væren i Rummet.

Betragte vi nu den Maade, paa hvilken Plato skildrer
Dannelsen af det \emph{Mellemled}, der skulde forklare Muligheden
af en Forbindelse mellem Ideen og dens Modsætning, saa er
det strax paafaldende, \emph{at det selv kun bliver til ved den
Forbindelse, det skulde forklare}, og, ligesom Mod\-%
sætningerne ere hypostaserede hver for sig, saaledes bliver det
det ogsaa og kommer til at staae som \emph{et Tredie} ved Siden
af Ideen og Materien. Hvad der, hvis ikke Ideen og Materien
vare sondrede som oprindelige, hypostaserede Modsætninger,
maatte blive til det Hele og falde sammen med den i sin
Enhed \emph{begrebne} \gk{Κόσμος}, nemlig: den harmoniske Enhed, i
hvilken Materien var fuldstændig optagen i Ideen, det bliver
nu paa den ene Side til \emph{en Deel af Tilværelsen}, i den
Forstand, at det i sig ligesom kun omfatter Noget af Ideens,
% --- p. 90 ---
\setcounter{footnote}{0}
Noget af Materiens Heelhed, og derfor kun bliver Sandse\-%
verdenens \emph{Ramme}, og paa den anden Side bliver det til et
\emph{eiendommeligt Led}, der i sit Væsen indeholder Elementer
af begge Modsætningers Væsen, men falder \emph{udenfor} dem
begge. \emph{Verdenssjælen} skal være \emph{usynlig og uforander\-%
lig} medens det Vordende er synligt og foranderligt; den skal
\emph{være udbredt gjennem Rummet} medens Ideen er uden
\emph{Forhold til Rummet}. Men i denne Mellemstilling staaer
Verdenssjælen \emph{i Væsen Ideen} nærmere end den staaer Mate\-%
rien, ja, den staaer Ideen saa nær, at det bliver ugjørligt at
fastholde dens Adskillelse fra denne. Det udføres af Plato
ikke videre, hvorledes Fornuften kom ind i den: det betragtes
som givet ifølge den første Blanding, ved hvilken Ideernes
Natur meddeltes den; derimod skildres det særskilt, hvorledes
det Legemlige forbandtes med den, og ved denne Forbindelse
blev harmonisk formet i sin Heelhed\footnote{Elementernes Grundformer
udledes ikke af Verdenssjælen.}. \emph{Ideen er det Be\-%
stemmende i den}, derfor er den «usynlig og deelagtig i For\-%
% p. 90: the ɔ: id-est mark, confirmed on the image at 1200 dpi (a
% reversed c followed by a colon).  Both witnesses read it as "9:".
nuft og i intelligible og udødelige Væseners», ɔ: Ideernes,
«Harmonie». Ligesom der i Ideen er Mangfoldighed, men
som behersket af Harmonien, Ikke-Væren, men som Forskjel
i Væren, som Form for Væren, saaledes er i Verdenssjælen
Flerheden bestemt ved Talforholdene, Andetheden behersket
af Enheden. Det Adskillende fra Ideen vilde derfor kun blive
Forholdet til Rummet. Men Verdenssjælens Udbredelse i
Rummet bliver egentlig kun \emph{et Billede}. Idet den «usynlig»,
\emph{som en Enhed}, gjennemtrænger Verden, er væsentlig Af\-%
hængighedsforholdet til Rummet ophævet, og det saameget
desto udtrykkeligere, som Verdenssjælen bestemmes som \emph{ule\-%
gemlig og tænkende}. Verdenssjælen er \emph{det Existerendes
Harmonie som hypostaseret og sondret fra det har\-%
monisk Existerende}, men som \emph{Harmonie er den uaf\-%
hængig af Rummet}.
% --- p. 91 ---
% p. 90's note ("Elementernes Grundformer udledes ikke af
% Verdenssjælen.") is a single line closing with a full stop; it does NOT
% run over.  RECON's run-over candidacy for it is cleared a second time.
% p. 91 bears no note.

De \emph{Bindeled} altsaa, der hos Plato skulde forene Idee
og Materie, vise sig som \emph{utilstrækkelige}. Verdensbyg\-%
mesterens mythiske Skikkelse kan kun for Phantasien faae
Magt og Betydning. Den Virksomhed, der blev ubegreben i
den til Urbillede hypostaserede Idee, bliver ubegreben i ham,
og ubegreben bliver, selv om Virksomheden forudsattes, dens
% p. 91, l. 7: a short mark stands between "det" and "ved" (tesseract
% reads "det-ved").  Not type; not transcribed.
Magt til at forbinde det ved en abstract Modsætning Adskilte.
Her skulde Begrebsmellemleddet: Verdenssjælen, udfylde Mang\-%
len. Men dette Mellemled forklarede ikke Forbindelsen fordi
det selv kun blev forklaret ved Forbindelsen, og Modsætnin\-%
gernes Uforenelighed viste sig deri, at dets postulerede Be\-%
stemmelser dog ikke kunde holde det ude fra den ene af
Modsætningerne: Verdenssjælen blev Universets almene \emph{Form}.
Det i Væsen \emph{absolut} Forskjellige \emph{kan intet sandt Mel\-%
lemled} have: ethvert Forsøg paa at fastholde et saadant maa
føre over til en af Siderne.

Problemet om Sandseverdenens Mulighed bliver altsaa
uløst. Forat det skulde have kunnet løses, maatte \emph{Virk\-%
somhedens} Bestemmelse ikke blot gjennem et Postulat være
% p. 91: in "Væsensbestemmelse" only the FIRST element is letterspaced;
% "bestemmelse" is set close, as printed.  A second such split within one
% word occurs at p. 98 ("indeslut- | ter").
tillagt den hvilende Idee, men afgjørende være hævdet som
Ideens \emph{Væsens}bestemmelse, og dette kunde kun være skeet
derved, at \emph{Negationens} Bestemmelse, Bevægelsens, Dialek\-%
tikens, Incitament, var sat i Væsenet, i Ideen. Kun da vilde
Muligheden være given til at forklare Sandseverdenens Væren
og concrete Indhold. Nu er dens relative Væren ved Siden af
Ideens absolute Væren, --- et Forhold, der sættes idet Ideen,
det almene Væsen, hypostaseres som Urbillede, --- en Modsigelse,
og dens Existents ikke blot uforklaret af Plato, men, ifølge hans
Tænknings Eiendommelighed, uforklarlig for ham.

I Platos Philosophies Grundbestemmelser see vi saaledes
den \emph{græske} Tænknings ovenfor paaviste \emph{almene} Eiendom\-%
melighed at gjøre sig gjældende. I Platos Tanke, som i den
græske Tanke overhovedet, kæmper Enheden med Dualismen
og den sidste bliver den mægtigste. I den Erkjendelsens
Umiddelbarhedsform, i hvilken det Ideelle umiddelbart beher\-%
% --- p. 92 ---
sker sin Modsætning, er Andetheden tilstede som skjult, men
som det, der maa sættes som Andethed, der uundgaaelig maa
gjøre sig gjældende. Og idet det sættes, er den \emph{umiddel\-%
bare} Enhed brudt og kun en \emph{høiere} Tænkningens Form kan
atter bringe Enhed tilveie. Plato søgte i Ideen, i det sande
Værende, Bestemmelser, der ligesom droge det hen imod Mod\-%
sætningen, han søgte forbindende Mellemled, han sammen\-%
knyttede Modsætningerne i Anskuelsen af \gk{Κόσμος}. Men i
samme Øieblik afskar han sin Stræben fra dens Maal: han
ligesom tilbagekaldte hine Bestemmelser i Ideen; han lod Mod\-%
sætningerne blive absolute, abstract sondrede Modsætninger;
han umuliggjorde ved denne Modsætningernes Bestemthed
Mellemleddenes forbindende Virksomhed; hvor han satte En\-%
heden, gav han dobbelte, modsigende Bestemmelser af Grund\-%
momenternes Forhold, og selv paa det Punkt, hvor Enheden
fandt sit fyldigste, ligesom definitive, uigjenkaldelige Udtryk,
hvor Tanken søgte Hvile i Anskuelsen af Tilværelsens harmo\-%
niske Enhed, hvor den berørte Maalet og i den digteriske Be\-%
geistrings Sprog udtalte Verdenskunstværkets Guddommelighed,
lod den dog i samme Øieblik Anskuelsen, der fæstner Begreberne,
holde Ideen sondret fra dette, fra dens harmoniske Aabenbarelse,
og hensætte den i en høiere Tilværelsens Sphære, i et af
usandselige Urbilleder udfyldt intelligibelt Rum, hvorved Ver\-%
densvirkeligheden atter blev ufattelig for Tanken, da al virke\-%
lig Væren indesluttedes i den Ideernes Verden, der sattes som
et \gk{χωριστόν}.

De her fremhævede Punkter ere i det Væsentlige opkla\-%
rede i det Foregaaende; kun den Dobbelthed, der træder frem
hvor Modsætningerne berøre hinanden, maa endnu gjøres ty\-%
deligere ved at de derhenhørende Træk samles. \emph{Modsætnin\-%
gerne enes i} \gk{Κόσμος}: \emph{Ideen hersker} gjennem \emph{Verdens\-%
sjælen}, der ligesom er dens Frænde og Statholder. \emph{Aar\-%
sagerne i} Verdensvirkeligheden stræbe, forsaavidt de ere
\emph{egentlige} Aarsager i dette Begrebs strængere Betydning, efter
at virkeliggjøre \gk{τὴν τοῦ ἀρίστου ἰδέαν}; de gjøre det fordi de,
% --- p. 93 ---
\setcounter{footnote}{0}
som ideelle, Øiemedsaarsager, henhøre til \gk{Νοῦς}, til Ideen, det,
efter Begrebet, ene Værende og ene Virksomme, hvis Maal er
den selv. \emph{Men ved Siden} af disse virkelige Aarsager staae
saa dog \emph{Medaarsager} (\gk{ξυναίτια}, Betingelser, det, uden
hvilket Aarsagen ikke vilde være Aarsag)\footnote{Jfr.\ Stoffets Forhold
til Formen i Kunstværket.}: \emph{de mechaniske
Aarsager}, der have deres Udspring fra hiin \gk{Ἀνάγκη}, som under
sit mythiske Navn skjuler Materiens Bestemmelse, og de \emph{be\-%
grændse og hæmme} hines fuldstændige Virken, saa at gjen\-%
nem dem det Ikke-Værende, det Materielle, den uundgaaelig
fremtrædende Andethed, faaer en eiendommelig Virksomhed,
bliver Principet for det Ondes Disharmonie, --- en Disharmonie,
der dog atter ligesom bliver uhørlig i Heelhedens Harmonie.

Den samme Dobbelthed i Opfattelsen af Forholdet mellem
\emph{det Materielle og Ideelle} træder frem, naar \emph{Sandsever\-%
denen} snart efter Begrebsbestemmelserne maa omfattes som
\emph{værende omsluttet af og i Enhed med Ideeverdenen},
der ikke adsplittes af det beherskede Sandseliges Mangfoldig\-%
hed, snart betegnes som værende som Sandseverden \emph{uden\-%
for den adskillelige Idee}; --- den træder frem i de af
Principerne afledede Bestemmelser, \emph{hvor Sjæl og Legeme}
snart sees som nødvendigt, altsaa oprindeligt, sammenhørende,
snart som forbundne i Tiden, efter den legemløse Præexistentses
Ophør; og i en nuanceret Form, \emph{hvor den ethiske Hand\-%
lens høieste Maal} snart bestemmes, i Analogie med Op\-%
% p. 93: "ene" is set UNSPACED inside the letterspaced run, as printed
% (the same device as "dets" on p. 85 and "det, der" on p. 88).
fattelsen af Ideen \emph{som det} ene \emph{Virkelige}, som den fuld\-%
komne Befrielse fra det Sandselige; snart, i Analogie med Op\-%
fattelsen af Ideen \emph{som Princip for en harmonisk Væren
i det Sandselige}, som den harmoniske Udvikling af Menne\-%
% p. 93: PRINTER'S ERROR -- "sandelig-aandelige" as printed, where
% "sandselig-aandelige" is expected; confirmed on the image at 1200 dpi
% and read the same way by both OCR witnesses.  Not on the Trykfeil leaf.
skets hele sandelig-aandelige Tilværelse, --- altsaa snart be\-%
stemmes fra Dualismens, snart fra Enhedens Synspunkt. Men
paa den meest frapperende og naive Maade fremtræder Dob\-%
beltheden ved Opfattelsen af \emph{Verdenstilblivelsen}. Opfat\-%
telsen af Verdens Begyndelse fører Plato ind i en uundgaaelig
% --- p. 94 ---
\setcounter{footnote}{0}
Modsigelse. Naar nemlig Materien sees som det, der efter
Begrebet hører sammen med Ideen, som den formende Idees
nødvendige Supplement, maa Verden \emph{tænkes som evig};\footnote{Man
jevnføre ogsaa Platos Opfattelse af Ideen som efter dens Be\-%
greb værende \emph{Værens Grund}, hans Opfattelse af Vordens Forhold til
Væren og hans Udsagn om Tidens Begyndelse.}
men idet Materien \emph{anskues} som det oprindelige Ikke-Værende,
\emph{der er ved Siden af det Værende}, maa der danne sig en
\emph{Forestilling} om \emph{et begyndende} Forhold og en \emph{succes\-%
siv} Verdensdannelse. Denne sidste Forestilling er det, der
kommer frem i Timæos's mythiske Skildring. Men \emph{indenfor}
denne Forestilling kommer nu atter det, at de to Momenter
oprindelig vise hen paa hinanden, frem i den naive Form, at
der i \emph{Materien}, allerede \emph{før} Verdensdannelsen begyndte, var
en uordnet, forvirret Blanding af \emph{Formerne}, i en uophørlig
% p. 94: the first καί of the quotation carries an ACUTE and the second
% and third καὶ a GRAVE, all read at 1200 dpi; the sense (καὶ before
% χώραν) would want a grave in the first place too.  Transcribed as
% printed.  εἶναι carries the smooth-breathing-plus-circumflex compound
% sort (a curl followed by a tilde), distinct from the plain circumflex
% of τοῦ; τριχῇ has the iota subscript.
Bevægelse (\gk{ὄν τε καί χώραν καὶ γένεσιν εἶναι, τρία τριχῇ,
καὶ πρὶν οὐρανὸν γενέσθαι}).

% DIVISIONAL RULE, p. 94: short, CENTRED hairline, measured on the image
% at x 674--919 against a text block of 240--1349 (rule centre 796.5,
% block centre 794.5) and 2.1 cm long, with white space above and below.
% It stands over the L3 head that follows.  NOT the footnote separator,
% which on this page is left-flush at x 230--479.
\divrule

% p. 94: L3 head under thesis II (transcription.tex §3).  IT IS NOT BOLD.
% Measured against the body of the same page its ink density is LOWER
% (0.153 against 0.234), because it is LETTERSPACED roman, and it is not
% centred either: it is set full measure and justified over three lines,
% x 240--1348, with a first-line indent and a nearly full last line.  §3
% describes the 1./2./3. heads as bold centred heads; that description
% does not fit this one, and the site should be checked against pp. 80
% and 127.  Set here as a letterspaced paragraph, which is what is
% printed; the line breaks of the print are 2. ... ud / over ... hæmmet /
% ved ... Væsenseiendommelighed.
\emph{2. Den Aristoteliske Philosophies Stræben ud
over den græske Tænknings Skranker som hæmmet
ved denne Tænknings Væsenseiendommelighed.}

\emph{Naar vi i den platoniske} Philosophie have fundet det
\emph{ligesom typiske} Udtryk for den græske Tænknings eiendom\-%
melige Art og eiendommelige Stræben, og seet disses Conse\-%
quentser i Bestemmelsen af dens Grundbegreber, i den Form,
i hvilken Grundproblemerne komme frem, og i det Mangelfulde
i deres Løsning; --- saa ville vi hos \emph{Aristoteles}, hos
hvem den individuelle, bevidste Tanke, hvilende paa det af
Plato givne Grundlag, stræber hen mod et Maal, der allerede
ligger ude over den græske Aands Grændser, see den græske
Tænknings i dens Grundtanke aabenbare Grundeiendommelig\-%
hed gjøre sig gjældende som den bevidste Tankes Skjæbne,
som hæmmende og begrændsende Magt.
% --- p. 95 ---
% p. 94's note closes with a full stop and does not run over.
% p. 95 bears no note.
% p. 95, l. 1: PRINTER'S ERROR -- a full stop after "tydeligt" where the
% sense wants a comma, with lowercase "bestemme" following.  The glyph is
% a round dot on the baseline with none of the comma's tail (1200 dpi),
% and both OCR witnesses read a full stop.  Transcribed as printed.
Forat gjøre dette tydeligt. bestemme vi først gjennem en
Sammenligning med den platoniske Philosophie den aristote\-%
liskes \emph{almindelige Charakteer og Retning}. \emph{Aristo\-%
teles} deler med Plato Opfattelsen af Formen, \emph{det Ideelle,
som det sande Værende og som Kilden til al Virkelig\-%
hed}; men idet han i Platos Lære om Ideerne som Tingenes
hypostaserede, adskillelige Urbilleder kun seer en unødvendig,
paa modsigende Consequentser rig, Fordobbling af Virkelig\-%
heden, og i den hverken finder en Forklaring af Tingenes
Væren eller af deres Tilblivelse, saa opfatter han Formen,
medens han ogsaa bringer Stoffet ind under dens Grundbe\-%
stemmelse: Væren, altsaa under den Bestemmelse, der hos Plato
\emph{uvilkaarlig} omfattede ogsaa det Ikke-Værende, som den
\emph{iboende Form, som det i Stoffet virksomme Formende}.
I selve den concrete reale Tilværelse, ikke gjennem en Flugt fra
den eller en Tilintetgjørelse af den, vil han derfor finde det Ideelle,
og han søger \emph{en Enhed af Modsætningerne} idet han seer
dette som realiserende sig selv i det væsensbeslægtede, kun ved
Existentsens Art forskjellige Stof, der har sit Maal i Formen.

Under Modsætningen til Plato, under Polemiken mod den
\emph{særegne Form}, i hvilken det ideelle Princip opfattedes af
denne, er derfor den aristoteliske Philosophies Grundpræg \emph{det
samme} som den platoniskes. Forskjellen er kun den, at naar
det ideale Princips Magt over Tanken hos \emph{Plato} gjør sig
gjældende i Bestræbelsen for at lade det Endeliges Verden
opløses gjennem Dialektiken for at Ideens Verden kunde aaben\-%
bare sig i sin Reenhed og Uafhængighed, saa gjør den sig
% p. 95, l. 29: the two "i" of "i det Sandselige, i den concrete" are set
% in BOLD, as on p. 89 l. 14 -- the compositor's device for emphasising a
% single letter, which cannot be letterspaced.  Measured at 1200 dpi
% against the ordinary "i" of "i hans Stræben" on the line above.
hos \emph{Aristoteles} gjældende i hans Stræben efter at see Ideen,
Formen, \textbf{i} det Sandselige, \textbf{i} den concrete naturlige Tilværelse,
som iboende Aarsag til dens Væren, som Kilde til dens Liv,
som virksomt realiserende sig selv i det Stofagtige. Der er
hos de to, Grækenlands største, Tænkere som en dobbelt Tanke\-%
strømning: \emph{tilbage til Ideen og ud fra Ideen}, og i dette
Kredsløb pulserer den ideebesjælede græske Tanke. Aristoteles
er Platos Supplement: i Forening omfatte de den græske Phi\-%
% --- p. 96 ---
\setcounter{footnote}{0}
losophies \emph{hele} væsentlige Indhold; de staae i Midtpunktet af
Rækken af dens store Repræsentanter, hvor deres ideale
Skikkelser ligesom smelte sammen til en Dobbeltherme med et
Ansigt, der erindrende seer tilbage og bag det Nærværende
skuer Ideen, og med et, der seer fremad og øiner de Tanker,
som, virkeliggjorte, skulle blive en anden Tidsalders Grund\-%
tanker. Det er den største Misforstaaelse, naar man har stil\-%
let Aristoteles som Empiriker imod Plato som Idealist: han
er ikke blot Idealist som hiin, men han er det forsaavidt
endog i høiere Forstand, som han har Forvisningen om den
det Ideelles Magt, der bevarer dets Væren ogsaa midt i Phæ\-%
nomenernes Verden, i Sammenknytningen med Andetheden.

I det angivne Forhold til Plato er den \emph{Stræben} hos
Aristoteles begrundet, der, naar den var naaet til sit Maal,
vilde have ført \emph{ud over den græske} Tænknings Skranker.

Idet for \emph{Aristoteles} Formen er \emph{iboende Form}, og
Stoffet ikke, som hos \emph{Plato}, det modsigelsesfulde værende
Ikke-Værende, men \emph{et Værende}, nemlig \emph{det efter Mulig\-%
heden Værende}, ere Principerne satte i et saadant Grundfor\-%
hold til hinanden, at den græske Tankes høieste Stræben: i
\emph{Begrebets Form} at udtrykke den Tilværelsens Enhed, der
umiddelbart er værende for Anskuelsen i Verdenskunstværket,
har et virkeligt Grundlag. Idet Formen sees som \emph{Virkelig\-%
hedsprincip, Udviklingsprincip}, er Løsningen af \emph{Vor\-%
dens Problem}, den græske Tænknings Anstødssteen, ligesom
berørt, og Opfattelsen af \emph{Verdensaltet} under den \emph{hvi\-%
lende Skjønheds Form} ifærd med for Tanken at ombyttes
med Opfattelsen af det som et uendeligt \emph{levende Hele}, i
hvilket det Ideelle er den reale Tilværelses \emph{Livsprincip}\footnote{Hos
\emph{Plato} berøres denne Bestemmelse idet \gk{Κόσμος} opfattes som
dannet efter det levende Væsens Idee og selv som et levende Hele,
men det Levendes Begreb træder dog for hans kunstneriske Anskuelse
strax i Skygge for Skjønhedens. Først ved den \emph{Aristoteliske Phy\-%
siks} Accentueren af det \emph{Organiskes og Udviklingens Bestem\-%
% RUN-OVER FOOTNOTE.  This note begins at the foot of p. 96 and its
% continuation fills the note area of p. 97, where it stands UNMARKED.
% The print divides it inside a word, at "Bestem- | melser"; the
% letterspacing stops at the page break, "melser" being set close.
melser} faaer Begrebet om Universet som \emph{det levende Hele} sin
egentlige Betydning.}.
% --- p. 97 ---
% p. 97 carries NO footnote call of its own; its note area holds only the
% continuation of p. 96's note, which is set in full inside the
% \footnote{} at its call on p. 96 above.  Hence no \setcounter here.
Og idet endelig Naturen, der har Formen i sig, rykker Ideen
nærmere end hos Plato, saa bliver \emph{det Værende i dets En\-%
kelthed, det enkelte} formbestemte Naturlige, \emph{der som
saadant er} \gk{οὐσία}, i sin Enkelthed, efter sit \emph{hele Omfang}, af
% p. 97: the breathing over the upsilon of οὐσία is SMOOTH -- a small
% curl whose opening faces LEFT.  It is the calibration word for the
% distinction in this batch; the ROUGH breathing (ὕλη, p. 98) is a "C"
% whose opening faces RIGHT.
\emph{væsentlig} Betydning for Erkjendelsen; det bliver, som det
Ideelles første Fremtrædelsesform for os, sat som det «for os
\emph{vissere}» Erkjendelsens Udgangspunkt, som Betingelsen for
den Virksomhed, ved hvilken Tanken \emph{ved sig selv} hæver sig
til det «\emph{i og for sig vissere}» Almene, og som det, til hvilket
Erkjendelsen altid stræber tilbage fra dette Almene. Saaledes
er \emph{Erfaringen}, den Erkjendelsens Form, der forholder sig
\emph{til det enkelte} Reale, sat som en \emph{væsentlig} Erkjendelsens
Form; der tillægges et Værd, hvorved en Erfaringsphilosophie
i \emph{moderne} Forstand antydes.

Men dog ville vi see, at \emph{Aristoteles}, trods denne saa
klart udprægede Stræben hos ham og trods hans Tankes vid\-%
underlige Energie, dog bindes af hvad der var den græske
Aands indre Begrændsning, at hans bevidste Stræben hæmmes
paa de afgjørende Punkter, og at han ikke mægter at løse de
Problemer, der indenfor disse Skranker kunne øines, men først
løses naar Skrankerne gjennembrydes.

Dette ville vi gjennemføre ved en Betragtning af saadanne
Hovedpunkter i den Aristoteliske Lære, der vise sig som af\-%
gjørende for dens historiske Stilling.

% p. 97: a RUN-IN head not in transcription.tex §3 -- "a)" in ordinary
% roman followed by "Opfattelsen af Vorden." in BOLD, set inside the
% opening line of the paragraph, then an em-dash and the running text.
% This is the second site on this batch's pages where bold occurs inside
% a paragraph (see p. 89 l. 14 and p. 95 l. 29); §3 and the file header's
% statement that bold is confined to display heads need correcting.
a) \textbf{Opfattelsen af Vorden.} --- De to Hovedbegreber, der i
den græske Philosophies ældste Periode bestemte Grupperingen
af de philosophiske Theorier: \emph{Væren og Vorden}, reflectere
sig ogsaa, i den græske Philosophies Culminationstid, i For\-%
holdet mellem de to store Tænkere. I \emph{Platos} hvilende, evige
Urbillede har Væren faaet en høiere Idealitetens Form; gjen\-%
nem \emph{Aristoteles's} virksomme, bevægende Formprincip, Ud\-%
viklingens Kilde og Maal, Tilværelsens Livsprincip, er Vorden
% --- p. 98 ---
% p. 98 bears NO footnote: the text block runs to the foot of the page
% and there is no footnote separator on it, so nothing can run over onto
% p. 99.  (The sheet signature 7* stands at the foot of p. 99, not here,
% and is not transcribed in any case.)
hævdet og knyttet til Virkelighedens ideelle Grund. \emph{Aristo\-%
teles} accentuerer \emph{fra Begyndelsen} af Vordens Begreb, der
er af afgjørende og almeen Betydning for Philosophien, fordi
det, som det første abstracte Udtryk for en Forbindelse af
Ikke-Værens, Negationens, Bestemmelse med Væren, overho\-%
% p. 98: the letterspacing begins in the MIDDLE of "indeslutter", at the
% start of l. 7 -- "og indeslut-" is set close and "ter" onward spaced.
% The same split within one word occurs at p. 91 ("Væsensbestemmelse").
vedet repræsenterer \emph{Modsætningernes Enhed} og indeslut\-%
\emph{ter Betingelsen for Begrebets Dialektik, for den con\-%
crete Tilværelses Udfoldelse af Principet.} Gjennem
Løsningen af Vordens Problem vilde Aristoteles have kunnet
sikkre Begrebernes Enhed, og finde en Forklaring for den con\-%
crete Tilværelse. I Følelsen af den Vægt, der ligger paa
dette Begreb, hævder han det med Energie mod dem, der
havde nægtet det, og med det Udviklingens og Mulighedens.
Han seer Vordens Mulighed, der ikke kan forklares af det
Værende alene eller det Ikke-Værende alene, som betinget ved
Værens og Ikke-Værens Forbindelse, nemlig ved en Væren,
\emph{der endnu er ikke-værende}, og en saadan Væren: en
\emph{Væren efter Mulighed}, er netop \emph{Materien}. Enhver Vor\-%
den forudsætter et \emph{Substrat}, i hvilket der skeer en Overgang
fra en Tilstand til en anden, fra en Egenskab til en anden,
og dette Grundlag, \emph{der kan} modtage \emph{de forskjellige Be\-%
stemmelser, er som Grundlag} netop kun \emph{Muligheden til}
begge, har ingen af dem virkeliggjort i sig. Og hvad der
gjælder med Hensyn til \emph{Vorden i det Enkelte}, gjælder med
\emph{Hensyn til Vorden overhovedet}: den kræver \emph{et Sub\-%
strat, der Intet er, men kan blive til Alt; der er
Alt efter Mulighed, Intet efter Virkelighed}. Dette
er \emph{den oprindelige Materie}, \gk{πρώτη ὕλη}, der som Betin\-%
% p. 98: the breathing over the upsilon of ὕλη is a compound blob at
% 1200 dpi, but its left component is a "C" whose opening faces RIGHT --
% the ROUGH breathing, and positively not the LEFT-facing curl of the
% smooth on οὐσία (p. 97).  It agrees with the rough on ὕλη at p. 35
% recorded in the file header.
\emph{gelse for al Vorden ikke selv kan være vordet}. Denne
Materie, der i sin \emph{absolute Bestemmelsesløshed} ikke er
Gjenstand for Erkjendelsen, men kun opfattes af Tanken ved
en Analogieslutning fra det formbestemte Materielle, kan ikke
% p. 98: PRINTER'S ERROR -- "Tidspuukt" as printed, a u where an n
% belongs; read the same way by both OCR witnesses and confirmed on the
% image.  Not on the Trykfeil leaf.
i noget Tidspuukt existere for sig alene; thi i denne Løsrevet\-%
hed fra Formen vilde den blive det absolut Uvirkelige, det
absolut Ikke-Værende; og naar den bestemmes som \emph{Mulig\-}%
% p. 98 ends MID-WORD: p. 99 opens with "heden", and the letterspacing
% carries across the page break.  The \emph{} is closed here so that the
% fragment's braces balance; the pp. 99--112 batch must open its own
% \emph{} on "heden" if the spacing continues there, and must not insert
% a paragraph break at this seam.
% --- p. 99 ---
% SEAM p. 98/99.  p. 98 bears NO footnote (confirmed again here: p. 99's note
% area is empty and there is no footnote separator on either page), so nothing
% runs over.  But p. 98 ends MID-WORD inside a letterspaced run: "Mulig\-" /
% "hed".  The \emph{} was closed at the foot of the pp. 85--98 fragment for
% brace balance and is reopened here on "hed".  No paragraph break at the seam.
% The sheet signature "7*" stands at the foot of p. 99 and is not transcribed.
\emph{hed}, saa har den kun denne Bestemmelse ved det oprindelige,
væsentlige Forhold til Formen. Derfor bliver Materien vel en
\emph{Betingelse for Vorden}, men denne forklares først idet
Materien sees i sit \emph{reale Forhold til Formen}: det efter
\emph{Virkeligheden Værende}, det \emph{Actuelle} (Actualiserede).
Stof og Form, der altsaa \emph{begge} bestemmes ved Væren, hint
% p. 99 l. 7: THE AUTHOR'S ERRATUM no. 7 (Trykfeil og Rettelser, PDF 306):
% the print reads \gk{ἐγεργείᾳ}, a GAMMA where a NU belongs; the leaf asks for
% ἐνεργείᾳ.  TRANSCRIBED AS PRINTED.  Verified on the image at 1200 dpi: the
% second letter carries a full descender below the baseline and is identical
% in stroke and shape with the γ of the following -γεί-, whereas the ν of
% \gk{δυνάμει} earlier on the same line sits wholly on the baseline with no
% descender.  The iota subscript under the alpha is visible.  Compare the
% correctly set ἐνεργείᾳ at p. 42 and \gk{ἐνέργεια} at l. 25 below, both with ν.
ved en Væren \gk{δυνάμει}, denne ved en Væren \gk{ἐγεργείᾳ}, staae
ved denne Bestemmelse i \emph{nødvendigt} Forhold til hinanden.
% p. 99: the ɔ: id-est mark occurs TWICE on this page, at l. 9 and l. 14;
% both read as a reversed c followed by a colon at 800 dpi.
Stoffet som det Mulige \emph{stræber hen} mod Formen ɔ: mod
det, til hvilket det er Mulighed, og som selv er det Muliges
Virkelighed; Stoffet «attraaer» Formen; men denne dets Stræ\-%
ben hen efter Formen er kun vakt \emph{ved Formen} og faaer
kun Virkeliggjørelse ved denne, der i dette Forhold bestemmer
sig som \emph{Entelechie} ɔ: som det iboende virksomme For\-%
% p. 99 ll. 15--16: \gk{τό ἐν ἑαυτῷ τὸ τέλος ἔχον}.  Breathings read at
% 1200 dpi and calibrated against each other in one render: over the ε of ἐν
% a thin oblique comma (SMOOTH); over the ε of ἑαυτῷ a broad closed "C" bowl
% opening right (ROUGH).  The two are unmistakably different sorts.  The ω of
% ἑαυτῷ carries a circumflex and an iota subscript; τό carries an acute and
% τὸ a grave, read directly against each other.
% \gk{ἔχον}: the mark over the epsilon is a COMPOUND, but -- unlike the single
% blob of ἄπειρον at pp. 52 and 81 -- it resolves into two elements separated
% by clear paper at 1200 dpi.  The left element is a narrow oblique tick with
% no bowl, matching the certain smooth over ἐν above and NOT the closed C of
% the rough over ἑ beside it; the right element rises to the right (acute).
% Read ἔ.  Nothing is supplied from knowledge of Greek.
mende, der i sig har Virksomhedens Maal (som \gk{τό ἐν ἑαυτῷ}
\gk{τὸ τέλος ἔχον}). Som \emph{Entelechie} viser Formen sig fra den
\emph{Side}, at den er det, der \emph{fører Stoffet fra Mulighed til
Virkelighed}. Denne Overgang fra Mulighed til Virkelighed
% p. 99 l. 20: the compositor spaces "som" but sets "muligt." close, so the
% emphasis stops in the middle of the phrase.  Set as printed.  The same
% half-spaced phrase occurs at p. 100 ("Ved Formen") and p. 110 note
% ("Fore-stillingsbillede"); cf. the mid-word starts at pp. 91, 98, 104.
er \emph{Bevægelsen}: Virksomheden (Entelechien) i det Mulige
\emph{som} muligt.

\emph{Bevægelsen}, der i sit \emph{almene} Begreb omfatter al
\emph{Forandring og Vorden}, betegner altsaa hiin Overgang,
Mellemstadiet mellem Muligheden som relativ eller absolut og
Virkeligheden; den er Forbindelsen af Væren og Ikke-Væren,
% p. 99 l. 25: \gk{ἀτελής} -- THE BREATHING OVER THE ALPHA IS RESOLVABLE HERE.
% At 1200 dpi it is a single oblique tick c. 55 px long (an em is 167 px),
% with no bowl, identical in form with the smooth over the ε of ἐνέργεια two
% words earlier in the same render and plainly unlike the closed C of the
% rough over ἑαυτῷ at l. 15.  Read SMOOTH.  What is unresolvable at pp. 52
% and 81 is the breathing PLUS ACUTE compound of ἄπειρον, not a breathing
% standing alone over α; here the mark is single and unambiguous.
der stræber hen imod Væren, en \gk{ἐνέργεια ἀτελής}, hvis Grændse
er den fuldstændige Virkeliggjørelse. Den forudsætter et Be\-%
vægende og et Bevægeligt, en actuel og en potentiel Væren;
den kan ikke forklares blot ved det Mulige eller blot ved det
Virkelige; men kun ved en Indvirkning af dette paa hint. Det
Virkende maa, som bevægende Aarsag, have en Prioritet i
Forhold til Bevægelsen, ogsaa i det Tilfælde hvor Bevægelsen
synes at udgaae fra det Mulige, f. Ex. hvor et organisk Indi\-%
vid bliver til af Sæden. Overalt hvor en actuel Væren støder
sammen med en potentiel, og ingen ydre Hindringer virke
hæmmende, opstaaer Bevægelsen med Nødvendighed i det Po\-%
% --- p. 100 ---
tentielle som en Fælledsvirkning af det Bevægendes \gk{ποίησις} og
% p. 100 l. 2: \gk{πάθησις} is set with the cursive fount's ϑ variant of theta;
% plain θ is typed per the file header.
det Bevægedes \gk{πάθησις}. I Bevægelsen er Formens og Stoffets
\emph{oprindelige} Forhold udtrykt; Bevægelsen maa altsaa være \emph{evig}
% p. 100 ll. 4--6: these three lines are set with markedly tightened word
% spaces so as to hold more words to the measure.  The TYPE SIZE is unchanged
% (x-heights measured against ll. 3 and 7 at 400 dpi); this is not a smaller
% setting and not a block quotation.
ligesom Forholdsleddene. Og det, at enhver Bevægelse maa have en
Bevægelse til Betingelse for sin Tilbliven og sit Ophør, viser ogsaa,
at Bevægelsen i sin Heelhed maa være uden Begyndelse og Ende.
Men forat Bevægelsens \emph{Virkelighed} overhovedet skal kunne
tænkes, maa der statueres \emph{et første actuelt Udgangs\-%
punkt}, der selv maa være \emph{ubevæget}, ikke blot i den Betyd\-%
ning, at det ikke bevæges af et Andet, men ogsaa i den, at
det ikke bevæger sig selv. Dette \emph{første Bevægende} er
\emph{Guddommen}, den \emph{rene Form}, den \emph{fuldkomne Energie},
og Guddommens Begreb er det \emph{afsluttende} i den aristote\-%
liske Metaphysik, det, i hvilket al Virksomheds, al Vordens og
al Bevægelses sidste Grund søges.

Saaledes søger Aristoteles at forklare Vorden, dette for
den græske Tænkning vanskeligste Begreb, af Formens og
Stoffets Væsen og Forhold. Men naar vi nærmere betragte de
Bestemmelser, han giver med Hensyn til de to Grundmodsæt\-%
ningers Forhold til Vorden, vise de sig utilstrækkelige til at
forklare dennes Virkelighed. I \emph{Stoffets} Natur ligger der
% p. 100: the single capital "I" before the spaced "Stoffets" cannot be
% measured for letterspacing (one letter, and the word space inside a spaced
% run is itself enlarged); it is left plain.  Same decision at "afsluttende i
% den" above, where "den" is measurably close.
\emph{intet} Forklarende; thi Stoffet er opfattet som det Bestem\-%
melsesløse, Passive, i hvilket altsaa ingensomhelst Virksomhed,
end ikke en \emph{saadan} som den, der ligger i en \emph{Stræben}, en
«Attraa», kan tænkes. Derfor maa \emph{Formen} være det, ved
hvilket Attraaen, Bevægelsesdriften, skal vækkes i Stoffet;
Formen, det Virkelige, sees som Bevægelsesprincip og \gk{τέλος};
den sees som det, der fremkalder Vorden i det efter Mulighed
% p. 100: "Ved" is letterspaced and "Formen" immediately after it is set
% close, at 400 dpi and again at 800 dpi.  Set as printed; the emphasis is
% not extended over "Formen".
Værende og fører den til Maalet. \emph{Ved} Formen er altsaa, efter
% p. 100: the "i" of "Vorden, i det Stofagtige" is from a MARKEDLY HEAVIER
% FOUNT -- stem and dot visibly thicker than the neighbouring i of "i sig"
% two lines below, at 400 dpi.  This is the fourth occurrence of the wrong
% sort recorded in the file header (pp. 42, 89, 95); it is a defect of the
% compositor's i-box, not an emphasis, and is set as an ordinary i.
Aristoteles's Opfattelse, Bevægelse, Vorden, i det Stofagtige,
og det, der vorder, er det af Form og Stof Sammensatte, der
% p. 100: immediately after the "i" that opens this line stands a short,
% thick, ink-black horizontal stroke at mid x-height, c. 1/6 em -- far too
% short for a dash and quite unlike the em-dashes on pp. 103 and 105.  It is
% a risen space or quad (a work-up), not a sort; both OCR witnesses pick it
% up ("i-").  NOT TRANSCRIBED.  The identical artefact stands after "derfor."
% at p. 101 l. 10.
i Bevægelsen har Formen i sig som Entelechie. Men idet
Vorden ved disse Bestemmelser skal være ikke blot hævdet
som Kjendsgjerning, men begreben, undslipper den dog paa
een Gang for Tanken, naar de nærmere Bestemmelser, Ari\-%
% --- p. 101 ---
stoteles føier til, tages med i Betragtning. Idet Formen frem\-%
kalder Bevægelse, skal Bevægelsen \emph{kun være i det Passive,
ikke i det Active}, og idet \emph{Udviklingen er naaet til sit
% p. 101 ll. 5--6: BOLD ANTIQUA INSIDE A PARAGRAPH.  "er" and "er vordet"
% are set in a heavy fount -- thick strokes, letters CLOSE -- standing inside
% an otherwise letterspaced run.  Measured at 600 dpi: stem widths about
% double those of the surrounding roman, and inter-letter gaps of 2--4 px
% against 10--16 px in the spaced words on the same line.  This is bold, not
% Sperrsatz, and it contradicts §2 of this file's header, which records that
% bold does not occur inside a paragraph except for the stray heavy "i".
% THE HEADER SHOULD BE CORRECTED.  Set \textbf{}.
Maal, og Formen er realiseret, saa skal der, efter
Aristoteles, om den kun kunne siges, at den} \textbf{er}\emph{, ikke,
at den} \textbf{er vordet}. Det, som er \emph{Vordens} Princip, falder
saaledes \emph{udenfor Vorden}, som \emph{uberørt af den}; men med
\emph{det Samme bliver netop Vorden uforklaret}; thi ufor\-%
klaret er den, naar det, der skal bevirke den, ikke i sig har
% p. 101 l. 10: a short thick ink-black stroke at mid x-height stands after
% "derfor.", c. 1/6 em long -- a risen space or quad, the same artefact as at
% p. 100 l. 32.  ABBYY reads it "-", tesseract "«".  NOT TRANSCRIBED.
en Plads derfor. \emph{Aristoteles} har ikke forklaret og kan ikke
forklare, hvorledes \emph{Formens ufuldstændige Realisation
under Udviklingen}, der ligesom er en Deelthed af Formen,
en partiel Negation, skal tænkes, naar Formen \emph{selv} holdes
udenfor Vorden, og hvorledes Formen saa kan blive Entelechie.
Hans Forsøg paa, at tilnærme Formens Begreb til Mulighedens
ved at bruge Materiens Begreb i relativ Betydning, saaledes
at den lavere Form opfattes som Stof i Forhold til den høiere,
en Energie altsaa gaaer over i Bestemmelsen af \gk{δύναμις},
hjælper ikke til at hæve Vanskeligheden, og ligesaalidt hæves
denne ved de Bestemmelser, han giver med Hensyn til Bevæ\-%
gelsens Realisation i det ene Existerende ved et andet actuelt
Existerende; thi denne Bevægelsens Realisation bliver kun et
Postulat, saalænge ikke netop det indre Bevægelsesprincips
Virksomhed er forklaret. Og vilde man forsøge paa at finde
en Forklaring ved at gaae tilbage til Metaphysikens afsluttende
Begreb: den \emph{første Bevæger}, \gk{Νοῦς}, Bevægelsens \emph{definitive}
Udspring, saa vilde ikke blot Vanskeligheden ikke heller her\-%
ved blive hævet, men man vilde kun støde paa nye Vanske\-%
ligheder og Modsigelser. Vi skulle tænke os \gk{Νοῦς} som Bevæ\-%
% p. 101 l. 31: \gk{χωριστόν} is itself LETTERSPACED -- gaps of 12--16 px at
% 800 dpi between every letter, against 2--4 px in \gk{δύναμις} thirteen lines
% above and in the three \gk{Νοῦς} on this page.  This is a real emphasis, not
% the spurious flag spacing.py raises on all Greek, and it is set
% \emph{\gk{...}} accordingly.  The rho is the cursive ϱ variant; plain ρ typed.
gelsens og Vordens \emph{oprindelige} Kilde skjøndt den er et
\emph{\gk{χωριστόν}}; vi skulle tænke os den som bevirkende Bevægel\-%
sen ved at fremkalde \emph{det absolut Bestemmelsesløses
Attraa}, og tillige skulle vi tænke os den som \emph{berørende
Verden}, skjøndt den er \emph{absolut immateriel}, som fremkal\-%
dende Bevægelsen fra \emph{Verdens Omkreds, skjøndt den intet}
% --- p. 102 ---
% p. 102 l. 1: the letterspaced run that ends p. 101 STOPS at the page break.
% "Forhold har til" is set close and only "Rummet," is spaced (checked at
% 600 dpi).  No paragraph break at this seam.
Forhold har til \emph{Rummet}, som Bevægelsens første og som
saadant \emph{almene} Princip, skjøndt \emph{flere} særegne \emph{evige Be\-%
vægelsesprinciper} staae \emph{uafledte} ved Siden af den.

Paa intet Punkt finde vi saaledes hos Aristoteles en vir\-%
kelig Forklaring af Vorden, og en saadan Forklaring \emph{afskjæres
ligefrem} ved de Udsagn, der udelukke Bevægelsen \emph{fra den
virksomme Form} og lade Formen i dens Fuldendelse ikke
være vordet, men kun værende. Saaledes faae vi nemlig atter
i Formen et \emph{Værende}, et \emph{Virkeligt, udenfor hvilket
Negationen falder, og som derfor bliver uden indre
Bevægelighed}, og med det Samme forvandles Muligheden
% p. 102: "Ikke-Væren" is broken at its OWN hyphen at the line end; the
% hyphen is printed and stays, so the newline is eaten with % rather than
% resolved with \-.
til en \emph{tom Mulighed}, en \emph{Mulighed til Intet}, en \emph{Ikke-%
Væren}. Den platoniske \emph{bevidste} absolute Modsætning
mellem Væren og Ikke-Væren kommer \emph{uvilkaarlig} tilbage
hos Aristoteles. Materien bliver, som denne Mulighed, hvis
Udviklingsprincip er tilintetgjort, ikke den Væren, der endnu
ikke er, men det Ikke-Værende, hvis Ikke-Væren er \emph{bli\-%
vende}.

Idet den Opgave, der for Aristoteles var den nærmest
liggende: at forklare Tilblivelse, Vorden, mislykkedes for ham,
og mislykkedes, uagtet han allerede angav de Begrebsbestem\-%
melser, ved hvilke den skulde løses, indeslutter Manglen paa
dens Løsning, at ogsaa den Opgave: at forklare \emph{den reale
Tilværelses Concretion}, maa mislykkes for ham, og den
viser tilbage til, at \emph{Dualismen}, som han søgte at fjerne ved
sin Bestemmelse af Grundmomenterne, ikke i Virkeligheden
er fjernet, og atter maa komme tilsyne under Systemets Ud\-%
vikling.

Hvad Forklaringen af \emph{den reale Tilværelses Con\-%
cretion} angaaer, saa var den fordret ved Aristoteles's hele
Stræben efter, med Tanken at omfatte \emph{al} Tilværelsens Rigdom
og at optage den i Erkjendelsen, og i Begrebet \emph{Udvikling}
og \emph{Udviklingsprincip} syntes han at besidde Betingelsen
derfor. Men at hans Stræben ogsaa paa dette Punkt \emph{maa}
% --- p. 103 ---
% p. 103 l. 1: the letterspacing stops at the page break -- "hæmmes," is set
% close (checked at 400 dpi).  No paragraph break at this seam.
hæmmes, ville vi see ved en Betragtning af Grundmodsæt\-%
ningerne.

Vilde man først søge en Begrundelse af Tilværelsens
concrete Mangfoldighed ved en \emph{directe, positiv} Virksomhed
af \emph{Materien}, saa vilde dette ikke blot stride mod den al\-%
% p. 103: the comma after "Materien" is lightly inked and both witnesses
% double it ("Materien,,som" / "Materien..som"); on the image it is a single
% comma with a small speck beside it.  Read as one comma.
mindelige Bestemmelse af Materien, som det Passive, men
ogsaa mod Aristoteles's Udsagn, at den Vorden, ved hvilken
\emph{det Bestemte} vorder, --- og al \emph{virkelig} Vorden er \emph{be\-%
stemt} Vorden --- bestandig forudsætter en \emph{bestemt Mod\-%
sætning som allerede værende} Egenskab eller Tilstand
i Stoffet, altsaa en \emph{forudgaaende Formning} af dette; thi
saaledes kan Materien som Materie i dens Ubestemthed og
Uformethed ikke forklare det concrete Virkeliges Tilblivelse,
men vi vilde vises hen til en \emph{allerede} formbestemt Materie.
Vilde vi søge en Modsætning til \emph{Materien som saadan},
saa vilde vi ikke faae nogen anden end \emph{Formen som saa\-%
dan}, og derigjennem altsaa kun kunne naae til, at Materien
gik op i \emph{den rene Forms Enhed}. --- Vilde man da forsøge,
istedenfor en \emph{positiv} Virkning af Materien, at sætte en
\emph{negativ}, en \emph{passiv Modstand}, og at forklare Tilværelsens
Mangfoldighed ved denne, saaledes, at man, knyttende sig til
de aristoteliske Bestemmelser om Materiens \emph{hæmmende Vir\-%
ken}, lod Formen i dens Enhed, idet dens Realisation i Til\-%
værelsen hæmmedes ved Stoffet, under den \emph{successive Over\-%
vinden} af denne Hæmmen sondre sig i en \emph{Trinrække} af
Former, saa maatte ogsaa dette Forsøg mislykkes. Det vilde
strande derpaa, at \emph{Modstanden} kun kan begribes, \emph{naar
der allerede forudsættes en Bestemthed}, men denne
kunde, som det ovenfor er fremhævet, kun have fundet Sted
ved en bestemt \emph{Form}, og vi vilde saaledes atter vises hen
til en anden Factor. Som den \emph{absolut bestemmelsesløse}
Mulighed maatte Materien modstandsløs gaae op i den rene
Form uden at kunne fremkalde en Mangfoldighed af Tilvæ\-%
relsesformer. Dette har Aristoteles ogsaa en Følelse af, og
% p. 103/104: the letterspaced run "fra første Færd af" CROSSES the page
% break; the \emph{} is left open across the page marker.  No paragraph
% break at this seam.
derved bringes han til at lade Materien, ligesom \emph{fra første
% --- p. 104 ---
\setcounter{footnote}{0}
Færd af}, træde op med en Bestemmethed; men denne Be\-%
stemmethed: Bestemmetheden med de \emph{første elementære
Modsætninger}, er allerede en \emph{Formbestemthed}, der
staaer som uforklaret.

Naar saaledes hverken en positiv eller en negativ Virken
af Materien kunde forklare Tilværelsens Concretion, vilde denne
\emph{kun kunne} forklares af \emph{Formen}. Som en \emph{Realitet}, et
virkeligt Værende, maatte den jo ogsaa have sin Grund i det,
der i \emph{egentlig} Forstand er \emph{Væren}, og til Formen vise
ogsaa Aristoteles's Udsagn om Betingelserne for det qvalitativt
Bestemtes Tilværen tilbage. Forklaringen \emph{ved Formen} vilde
da nærmere være en Forklaring ved \emph{den høieste Form
som Enhed}; thi \emph{Tilværelsens Mangfoldighed} vilde,
naar ikke \emph{Formernes Mangfoldighed}, der betinger den,
førtes tilbage til Enheden, staae som et \emph{ubegrebet} Givet.
Men hvad der ogsaa her umuliggjør Forklaringen, er det, at
den \emph{høieste Form}, idet, som ovenfor viist, \emph{Negationen}
falder \emph{udenfor} Formen overhovedet, \emph{ikke kan blive dia\-%
lektisk}, ikke faae den \emph{indre Udfoldelses} Mulighed i sig.
Vel gaaer Aristoteles's \emph{Stræben} ud paa, at opfatte den
\emph{som Virksomhed}; men denne Stræben hæmmes ikke blot ved
hiin Opfattelse af Formens Forhold til Vorden, men allerede
ved det \emph{Grundeiendommelige i den græske} Tænkning,
for hvilken Formen af den Grund ikke kunde blive bevægelig,
at den oprindelig er \emph{Kunstværkets Form} og som saadan
% p. 104: \gk{εῖδος} (twice on this page) -- over the iota stands the same
% compound sort recorded at p. 33 and p. 52: a broad, plainly visible
% circumflex whose left end hooks, in which a smooth breathing cannot be
% separated from the circumflex at 1200 dpi.  The circumflex is transcribed
% because its shape is visible; the breathing is NOT supplied; ἶ/ῖ undecided.
% Same treatment as pp. 33 and 52.  \gk{μορφή} carries a clear acute over the
% eta; \gk{Νοῦς} a clear circumflex over the upsilon.
\emph{hvilende} Form: et Forhold, der ogsaa skinner igjennem hos
Aristoteles, naar \gk{εῖδος}, der, som Guddommens Bestemmelse,
blev til \gk{Νοῦς}, som \emph{Naturens} \gk{εῖδος} til Synonym faaer Be\-%
stemmelsen \gk{μορφή}\footnote{Man erindre ogsaa Opfattelsen af den
  høieste Forms \emph{Rumsforhold} til \gk{Κόσμος}.}. Selv naar Aristoteles opfatter den rene
\gk{Νοῦς} som \emph{reen, uafbrudt Energie}, saa er den dog forsaa\-%
vidt \emph{hvilende} som den evigt, uforanderligt, kun tænker sig
% p. 104: the note is self-contained and ends with a full stop; nothing runs
% over onto p. 105.  Printed mark *), superscripted, after μορφή.
% p. 104/105: the letterspaced run "sin egen forskjelsløse Afspeiling."
% CROSSES the page break, and the word "Af-speiling" is broken across it.
% "kun frembringer denne" between "selv," and "sin" is set CLOSE (checked at
% 600 dpi) and is left plain.
\emph{selv}, kun frembringer denne \emph{sin egen forskjelsløse Af\-%
% --- p. 105 ---
\setcounter{footnote}{0}%
speiling}. En Mulighed til at løse Problemet om \emph{den ene
Forms Udfoldelse til de mange Former}, kunde synes
at vise sig, naar vi forudsatte --- hvad vi dog efter det histo\-%
risk Foreliggende ikke ere berettigede til at \emph{paastaae}, selv
om enkelte Antydninger og Analogier\footnote{Jfr.\ nedenfor S.\ 106.} pege hen derpaa ---
at den Tanke har foresvævet Aristoteles, at den guddommelige
Tankes, den rene Forms, \emph{Tænken af sig selv} skulde være
\emph{en Tænken, der omsluttede alt det Intelligible som
en Enhed}, i hvilken hele Verdensvirkelighedens teleologisk
sammenhængende Indhold var sammentrængt, saa at altsaa
den guddommelige Tankes Tænken af sig selv blev en \emph{fyldig
Enhed}, i hvilken alt det, der i Verden fremtræder adskilt i
% p. 105: a small pencil tick stands over the final "t" of "sammensluttet"
% (ABBYY reads it as an apostrophe).  Modern pencil; not transcribed.
Mangfoldighed, var sammensluttet uadskilleligt, og altsaa Ver\-%
den, idet de enkelte i den realiserede Former bleve Bestem\-%
melser af den ene guddommelige Form, ikke blev et fra Gud\-%
dommen Løsnet, Guds Tænken af den ikke en Tænken af et
Andet. Men denne Mulighed vilde dog atter forsvinde, idet
Aristoteles ikke vilde kunne fjerne de Vanskeligheder, der
ifølge hans Opfattelse af Guddommens Sondring fra Verden
og af Formens Forhold til Negationen, og ifølge hans Paral\-%
leliseren af de almene Begrebers\footnote{Man sammenligne ogsaa
  Aristoteles's Udsagn om de for hver Videnskab eiendommelige
  Principer og om Kategoriernes indbyrdes Forhold.} Forhold til de særegne med
Forholdet mellem \gk{δύναμις} og \gk{ἐνέργεια}, maatte stille sig
iveien for den Udvikling og Gjennemførelse af hiin Tanke,
der vilde være nødvendig til en Løsning af Problemet.

% p. 105: TWO notes, *) and **), both self-contained; neither carries
% letterspacing and neither runs over onto p. 106.  spacing.py's fit was
% split at the footnote rule before it was trusted, as this file's §3
% requires for this page.
I Aristoteles's Opfattelse af det Begreb, der udtrykker den
\emph{absolute Virkelighed}, viser sig saaledes paa det Klareste
den aristoteliske Metaphysiks og overhovedet den græske Tænk\-%
nings \emph{Grundmangel}: Manglen paa Evne til at gjøre det
Værende levende og bevæget ved Negationen. Ved at lade
\emph{Formens} Begreb blive til den guddommelige \emph{Tankes}, ved
% --- p. 106 ---
at sætte denne som \emph{reen Energie} og som den af ydre Be\-%
% p. 106: \gk{νόησις τῆς νοήσεως} occurs TWICE on this page and the two are
% set differently: the first (l. 2) is LETTERSPACED, the second (l. 13) is
% set close.  Measured at 800 dpi in one render: 12--17 px between letters in
% the first against 2--4 px in the second.  Both readings are as printed.
% Acutes over the omicron of νόησις and over the eta of νοήσεως; circumflex
% over the eta of τῆς.
tingelser uafhængige evige \emph{\gk{νόησις τῆς νοήσεως}}, ved at
bestemme den med den guddommelige identiske \emph{menneske\-%
lige} \gk{Νοῦς} som \emph{\gk{εῖδος εἰδῶν}}, som i sin \emph{Enhed} omfattende
% p. 106: \gk{εἰδῶν} -- the iota carries a breathing WITHOUT an accent, and it
% is a plain comma-shaped hook at 800 dpi: SMOOTH, positively not the C-bowl
% of the rough.  The omega carries a circumflex.  (\gk{εῖδος} beside it is the
% compound sort; treated as at p. 104.)
\emph{de lavere} Sjæleformer og som \emph{tænkende sig selv i sin
Gjenstand}, var Aristoteles naaet saa langt henimod en
Opfattelse, der kunde give Formprincipet \emph{indre Bevægelse}
og en deraf fremgaaende \emph{Fylde}, og føre til at see \emph{Tilvæ\-%
relsens Concretion} udvikle sig af Principet, som det for
den græske Tænkning var muligt; men paa dette afgjørende
Punkt hæmmes Tankens Stræben som ved en usynlig Magt:
den \emph{absolute Væren}, der \emph{udskiller Negationen} af sig,
bliver \emph{blot positiv Væren}, den evige \gk{νόησις τῆς νοήσεως}
bliver kun som en \emph{Tankens Stirren paa sig selv i sin
Ensformighed}, Reflexionen synker sammen i \emph{Hvilens
Umiddelbarhed}, og Verdensvirkelighedens Rigdom staaer
uforklaret ved Siden af Principet.

% p. 106: A DIVISIONAL RULE stands here, mid-page -- short, centred, with
% white space above and below, c. 1.7 cm at 400 dpi, and set off from the
% left margin (so not the footnote separator, which is left-flush; there is
% no note on this page).  IT CARRIES NO DISPLAY HEAD: the matter that follows
% it is ordinary running text, not one of the heads listed in §3 of this
% file's header.  This is a section division marked by the rule alone, and it
% is NOT in the heading list; reported to the calling conversation.
\divrule

Hvad der gjennem denne Undersøgelse af Aristoteles's
Opfattelse af Vorden har viist sig som hans \emph{Metaphysiks
Begrændsning}, ville vi nu, inden vi gaae over til de andre
for hans Philosophies historiske Stilling afgjørende Hoved\-%
punkter, forat stille den indre Sammenhæng mellem det ari\-%
stoteliske Systems forskjellige Dele i et klarere \emph{Lys}, see som
reflecterende sig i Opfattelsen af Grundbegreber fra det \emph{Ethi\-%
skes Omraade} og fra det Omraade, der danner det Ethiskes
Grundlag og Forudsætning: det \emph{Anthropologiskes}.

Det er en Betragtning, der frembyder sig ved det første
Blik paa den \emph{græske Ethik}, at den har en \emph{æsthetisk}
Charakteer, nærmere: en \emph{æsthetisk-metaphysisk}. Lige\-%
som i den græske almindelige Bevidsthed det \emph{Godes} og det
\emph{Skjønnes} Begreb smeltede sammen i det for Grækerne
eiendommelige, i ethvert andet Sprog uoversættelige, Begreb:
% --- p. 107 ---
% p. 107 l. 1: \gk{καλοκἀγαθία} -- the CORONIS over the alpha after the second
% kappa is a clear comma-shaped mark at 800 dpi, and the iota carries an
% acute.  The theta is the cursive ϑ variant; plain θ typed.
\gk{καλοκἀγαθία}, og ligesom det dannede græske Sprogs Udtryk
for en \emph{sædelig} Fuldendthed overalt minder om \emph{æsthetiske}
Bestemmelser, saaledes opfatter ogsaa den \emph{videnskabelige}
græske Ethik det Ethiske \emph{under Skjønhedens Synspunkt}:
seer Individets \emph{ethiske Opgave} som den: at udforme sig
selv til et \emph{individuelt Kunstværk}, og den opfatter den
\emph{ethiske Udvikling} som den af \emph{Skjønheden} beherskede,
i en Række af \emph{discrete} Momenter harmonisk fremtrædende
successive Aabenbarelse af den uforstyrrede Menneskenatur.
Den ethiske \emph{Handlen} bliver for den en \emph{Foregaaen}, en \emph{Be\-%
givenhed}, og \emph{det Godes} Begreb faaer, hvad der svarer
hertil, en \emph{metaphysisk} Opfattelse.

Denne \emph{Grundcharakteer} i den \emph{græske Ethik}, der
forhindrer den \emph{udtømmende} Opfattelse af det Ethiskes
Væsen og Grundbestemmelser, gjenfinde vi ogsaa hos \emph{Ari\-%
stoteles}, trods hans, med hans Tænknings hele Art sammen\-%
hængende, Stræben efter at hævde det Ethiske i dets Eien\-%
dommelighed, og at skærpe Opfattelsen af de ethiske Hoved\-%
begreber.

Søge vi, inden vi nærmere udvikle dette, Grunden til den
videnskabelige græske Ethiks gjennemgaaende Eiendommelighed,
og saaledes ogsaa Grunden til hvad der \emph{definitivt} bestem\-%
mer den aristoteliske Ethiks Charakteer, saa maae vi \emph{nærmest}
og \emph{umiddelbarest} finde den i \emph{Anthropologien}, hvis
Sammenhæng med Ethiken hos Grækerne, netop fordi de i
det Sædelige saae \emph{den i sit Væsen uforstyrrede Men\-%
neskenaturs} Skjønhedsudfoldelse, var inderligere end i den
moderne Tid. Men Anthropologien er i den græske Philoso\-%
phie \emph{Physikens} Afslutning, og gjennem de physiske Grund\-%
bestemmelser maa, i de fuldt udfoldede græske Systemer, An\-%
thropologiens Charakteer bestemmes ved \emph{Metaphysikens}.

Den Eiendommelighed ved \emph{Anthropologien}, der gav
den græske Ethik overhovedet dens eiendommelige Præg, var
væsentlig \emph{Manglen paa en dybere Opfattelse af Per\-%
sonligheden i dens Enhed og indre Uendelighed}, og,
% --- p. 108 ---
hvad der paa det \emph{Nøieste} hænger sammen dermed, af \emph{Villien
i dens Frihed og Selvbestemmen}.

\emph{I Aristoteles's Anthropologie} see vi nu først, lige\-%
som i enhver Deel af hans System, en energisk Stræben ud
over den tidligere Philosophies Skranker, men idet vi følge
denne Stræben, viser det sig dog, at den paa de \emph{afgjørende}
Punkter bliver hæmmet, saa at den \emph{almene Grundmangel}
ogsaa gjør sig gjældende hos ham. Dette vil først vise sig
ved Opfattelsen af \emph{Personlighedens Enhed}.

Om et \emph{Begreb} om Personlighedens Enhed kan der først
være Tale hvor en \emph{videnskabelig} Anthropologie udfolder sig.
Men hos \emph{Plato}, den Første, hos hvem en saadan begynder,
% p. 108: PRINTER'S ERROR -- "kun han blive" as printed, where the sense
% requires "kan".  Both OCR witnesses read "han"; confirmed on the image at
% 400 dpi.  NOT on the author's Trykfeil leaf.  A modern reader has pencilled
% "Kan" in the left margin and underlined the printed "han"; the pencil is
% NOT transcribed and p. 108 is not among the pollution sites listed in this
% file's header -- a NEW pencil site.
er der udtrykkelig sat \emph{Dele} af Sjælen, saa at Enheden, i Ana\-%
logie med Enheden i Ideeverdenen, kun han blive \emph{det Ad\-%
skiltes Harmonie}, ikke den \emph{indre levende Væsensen\-%
hed}. Hos \emph{Aristoteles} stræbes der henimod en Hævden af
Enheden. Han taler ikke om Dele af Sjælen, men om \emph{For\-%
mer} af Sjælelivet, og disses Forhold opfatter han saaledes, at
den høiere \emph{forudsætter} den lavere, og, idet den gjennem en
\emph{teleologisk} Udvikling kommer frem, \emph{optager} den lavere i
\emph{sin Enhed} og \emph{bestemmer den} i denne Optagen. Derfor
bliver hos Mennesket, hvis \emph{eiendommelige} Sjæleform er den
\emph{tænkende Sjæl}, \gk{Νοῦς}, dennes Bestemmelse Betegnelse for
\emph{alle} lavere sjælelige Former som \emph{optagne i den fornuftige
% p. 108: \gk{Νοῦς παθητικὸς} is letterspaced, AND the omicron of παθητικὸς
% carries a GRAVE, not the acute.  The mark slopes DOWN to the right and is
% directly opposite in slope to the certain acute over the alpha of
% \gk{διάνοια} nine lines below on the same page, the two compared in renders
% at the same scale.  The same grave-for-acute stands on ποιητικὸς twice at
% p. 109 and on ψυχὴ and θεὸς in p. 109's note, while παθητικός and
% ποιητικός carry a true acute elsewhere on p. 109 -- the compositor is
% inconsistent, and each occurrence is set as printed.  The theta is the
% cursive ϑ; plain θ typed.
Sjæls Enhed}. \emph{\gk{Νοῦς παθητικὸς}} --- dette saa hyppigt mis\-%
forstaaede Begreb --- betegner hos Aristoteles netop Indbegre\-%
bet af de endnu ikke til det sidste Maal, til den fuldendte
Virkeliggjørelse, naaede Udviklingstrin af den menneskelige
Sjæl, og omfatter ikke blot alle de lavere Bevidsthedsformer,
der blive Momenter i den fornuftige Sjæls Enhed, men ogsaa
den reflecterende, middelbare, paa Sandsning og Erfaring støt\-%
tede \emph{Tænkning} (\gk{διάνοια}, til Forskjel fra \gk{νοῦς}). Gjennem
alle de lavere Udviklingstrin gaaer \emph{eet Væsens Enhed}:
den \emph{Form}, der ved Udviklingens \emph{Afslutning} aabenbarer sig
% p. 108: "som" before Νοῦς is letterspaced but the Greek beside it is set
% close (the same width as the unspaced Νοῦς twelve lines above); as printed.
\emph{som} \gk{Νοῦς}, som \emph{tænkende Aand}; de ere derfor alle \emph{deel\-%
% --- p. 109 ---
% NB the % on the \setcounter line: p. 108 ends mid-word ("deel\-agtige") and
% without it TeX would turn the line end into a space.
\setcounter{footnote}{0}%
agtige i dens Væsen og gaae ind under dens Bestem\-%
melse}; men idet de, som \emph{Trin i den endnu ikke fuldt}
realiserede Entelechies Udvikling, endnu have \emph{Mulighed} og
\emph{Vorden} i sig og derved staae i Berøring med det \emph{Materi\-%
elle}, ere de kun \gk{Νοῦς} som \emph{den i Udvikling og Vorden
værende}, som \emph{\gk{παθητικός}}. Fra \gk{Νοῦς} i denne Bestemthed
sondrer sig saa \gk{Νοῦς} som \emph{\gk{ποιητικός}}, som \emph{reen Energie}.
\gk{Νοῦς ποιητικός} er, for strax her at vise tilbage til \emph{Metaphy\-%
sikens} Grundbegreber, \emph{Menneskets} \gk{εῖδος}, den er Maalet,
det \emph{\gk{τέλος}}, hvorefter hele Udviklingen stræber hen, og som
Maalet det \emph{virkende Princip}. Den forholder sig til hvad
der sammenfattes under Bestemmelsen \gk{νοῦς παθητικός} som
\emph{Energien} til den endnu ikke fuldt realiserede \emph{Entelechie}\footnote{Følgende
  Formel kan oplyse Sammenhængen mellem de Aristote\-%
liske Bestemmelser paa \emph{Metaphysikens, Psychologiens} og
  \emph{Physikens} Gebeter: \gk{ἐνέργεια}: \gk{ἐντελέχεια} = \gk{νοῦς}: \gk{ψυχὴ} =
  \gk{θεὸς}: \gk{φύσις}.}.
% p. 109 note: the "=" are ordinary text equals signs, not math.  ψυχὴ and
% θεὸς both carry GRAVES where the position calls for acutes (each stands
% before a colon); read at 800 dpi and directly comparable with the acute of
% φύσις in the same render.  Set as printed.  The note is self-contained.
Ved disse Bestemmelser er Aristoteles kommen \emph{Erkjendel\-%
sen af Enheden} nær. De sjælelige Functioners Udviklings\-%
række stræber, behersket af den iboende Form, hen imod det
Maal, der omfatter Trinnene i sin Enhed. Men netop paa
dette Punkt brister dog Enheden. Ligesom vi i \emph{Metaphy\-%
siken} have seet \emph{Formen} som \emph{Entelechie} Skridt for Skridt
bestemme Udviklingen og dog ved Udviklingens Slutning staae
\emph{uberørt af Vorden}, i sin Færdighed unddrage sig denne,
saaledes at den kun maatte bestemmes som \emph{et Værende},
\emph{ikke som et Vordet}, saaledes sondrer sig i \emph{Psychologien}
\gk{νοῦς ποιητικὸς} som Menneskets \gk{εῖδος} fra Udviklingen og dens
Trin. \gk{Νοῦς ποιητικὸς} er den \emph{evige, af Vorden uafhæn\-%
gige Form, uafbrudt Energie}, efter sit \emph{Væsen Eet med
den guddommelige} \gk{Νοῦς}, som den udtrykker i den men\-%
neskelige Existentses Form, den er \emph{enkelt og ublandet,
uafhængig af Legemet, et eiendommeligt, adskille\-%
ligt} Princip, \emph{udødelig}, men \emph{ikke i Personlighedens
% --- p. 110 ---
\setcounter{footnote}{0}
Form}. Skjøndt Aristoteles's Stræben bestandig gaaer ud paa
at fastholde \emph{Enheden i det Menneskelige}: Enheden af det
Sjælelige og det Legemlige, i hvilket Sjælen er Entelechie, og
Enheden mellem de teleologisk sammenhængende sjælelige
Udviklingstrin og Functioner\footnote{Derfor lader Aristoteles selve
  \emph{Tanken} være ledsaget af et \emph{Fore\-}%
stillingsbillede, et \gk{φάντασμα}, der forholder sig til Tanken ligesom
  den mathematiske Tegning til det, der bevises, og som ogsaa bliver
  \emph{Mellemleddet}, gjennem hvilket Fornuften indvirker paa \emph{At\-%
traaen}.}, og skjøndt han ved at be\-%
% p. 110 note: "Aristoteles" and "Tanken" are letterspaced, and so is "Fore-"
% at the line end -- but "stillingsbillede," on the next line is set CLOSE,
% so the emphasis stops inside the word.  Set as printed; cf. p. 99 l. 20 and
% p. 100 ("Ved Formen").  The note ends with a full stop; nothing runs over.
stemme \gk{Νοῦς} som det, vi, som Mennesker, egentlig ere (Eth.
Nic.\ IX.\ 4. X.\ 7) viser hen til et \emph{Centralt, Forenende},
saa brydes dog, paa det \emph{afgjørende} Punkt, Enheden atter
derved, at \gk{Νοῦς}, der skulde være det \emph{fyldige Enhedsud\-%
tryk for det menneskelige Væsen, dualistisk} skiller
sig fra \emph{det Legemlige} og fra \emph{alle de ved det Legemlige}
betingede \emph{sjælelige} Former, den skulde sammenfatte i sig,
og staaer som \emph{det af Udviklingen Uberørte, Uperson\-%
lige, Abstracte}. Ved denne Sondring bliver \emph{Udviklingen
selv} og \emph{Trinnenes Sammenhæng} uforklaret, og Begrebet
om Selvets, om \emph{Personlighedens Enhed} undslipper Ari\-%
stoteles. Enheden bliver saa hos ham idethøieste hvad den
var hos \emph{Plato: Harmoniens Enhed}, den umiddelbare, har\-%
moniske Forbindelse af det \emph{Forskjellige}; den bliver liig med
den Idealitetens Enhed, der i det plastiske Kunstværk ligesom
er udgydt over den \emph{hele} besjælede Skikkelse, ikke, som i Ma\-%
leriet, concentreret i Blikkets Inderlighed.

Idet Begrebet om \emph{Selvets indre Enhed} mangler, mang\-%
ler tillige det derved betingede Begreb om en \emph{indre Uende\-%
lighed}, der kun ved en for Aristoteles \emph{uigjennemførlig}
Antydning er berørt ved Bestemmelsen af \emph{\gk{νοῦς}} som \emph{\gk{εῖδος
εἰδῶν}}. Den Harmoniens Enhed, der træder i Stedet for en
indre Væsensenhed, har en uopløst \emph{Umiddelbarhedens} og
% p. 110 ll. 29 f.: THE AUTHOR'S ERRATUM no. 8 (Trykfeil og Rettelser):
% the print reads "For\-brydelse"; the leaf asks for "Fordybelse".
% TRANSCRIBED AS PRINTED.  In this copy a modern hand has struck "brydelse"
% through in pencil and written "dybelse" below it in the margin -- a reader
% applying the printed errata.  THE PENCIL IS IGNORED.  Neither half of the
% word is letterspaced (checked at 400 dpi), although "Ethiske i strængere"
% after it is.
\emph{Naturlighedens} Rest, der forhindrer Selvets uendelige For\-%
brydelse i sig, Forudsætningen for det \emph{Ethiske i strængere}
% --- p. 111 ---
Forstand. Og denne Mangel bliver atter \emph{medbetingende}
for den mangelfulde Opfattelse af \emph{Villiens Begreb}, et Begreb,
der hos Grækerne ogsaa af \emph{den} Grund maatte blive mangel\-%
fuldt opfattet, at de, for hvem \emph{Erkjendelsen} var Menneskets
\emph{egentlige Væsen}, Dyd Viden, Synd Uvidenhed, ikke kunde
hævde Villien \emph{i dens primitive Eiendommelighed} ved
Siden af Erkjendelsen. \emph{Aristoteles} stræber vel ogsaa med
Hensyn til Villien at naae videre end sine Forgængere, og sø\-%
ger saaledes, hvad Forholdet til Erkjendelsen angaaer, at hævde
Villien en særskilt Plads, en relativ Uafhængighed, ved at ac\-%
centuere \emph{Attraaens} Element og det \emph{Praktiskes} Eiendom\-%
melighed. Men dog bliver Villiens Begreb kun svævende og
usikkert. Naar Villien opfattes som en \emph{fornuftbestemt
Attraa}, saa maa det, da, efter aristoteliske Forudsætninger,
al Attraa kun kan blive til idet den fremkaldes af \emph{Formen},
og en \emph{bestemt Grændse} derfor ikke kan drages mellem den
\emph{blotte Attraa og Villien}, blive \emph{uklart}, hvad der i Vil\-%
liesvirksomheden er \emph{det egentlig Afgjørende: Attraaen},
Villiens Grundlag, eller den Indvirkning, som \emph{Fornuften}
igjennem Forestillingsbilledet udøver paa hiin, og ved hvilken
% p. 111: PENCIL.  A small stray stroke stands after the final r of "eller"
% at the end of l. 22 (ABBYY reads "uklart^" and "deo''"), and a pencil
% flourish is drawn off the "der" that ends the last paragraph's first line.
% Both are modern pencil and are NOT transcribed.  p. 111 is on the list of
% pollution sites in this file's header.  The comma after "uklart" IS printed
% type, though set low and slightly displaced; it is transcribed.
Attraaen først skal blive til Villie. Ligeledes bliver det uklart,
hvorledes Menneskets Bestemmethed til det Fornuftige \emph{eller}
det Ufornuftige, det Gode \emph{eller} det Slette, kan, hvad Aristo\-%
teles vil, blive en \emph{frivillig} Sag, og hvorledes det \emph{frit væl\-%
gende Subject} skal bestemmes. Der er saaledes kun gjort
et \emph{Tilløb} til at bestemme Villiens Begreb i dets Eiendomme\-%
lighed, og medens Spørgsmaalet om \emph{Villiens Frihed} i Val\-%
get mellem det \emph{Gode og Onde} ikke bliver indtrængende be\-%
handlet, bliver Spørgsmaalet om Villiens Frihed i Forholdet til
det \emph{Guddommelige}, fra hvilket Villiens Bestemmelse ude\-%
lukkes, og til \emph{Naturens} Virken forbigaaet.

Disse Mangler ved Aristoteles's \emph{Anthropologie}, der
hænge sammen med hans \emph{Metaphysiks} eiendommelige Grund\-%
bestemmelser, afspeile sig nu atter i hans \emph{Ethiks Begrænds\-%
ning}.
% --- p. 112 ---
Ogsaa paa \emph{Ethikens} Gebet stræber \emph{Aristoteles} ud
over sine Forgængere, navnlig over \emph{Plato}, deels ved en skar\-%
pere Sondring af det Praktiskes Sphære fra det Poietiskes og
% p. 112: a short pencil stroke crosses the y of "Dy\-den" at the end of
% l. 5 (ABBYY reads "Djr-den").  Modern pencil; not transcribed.  p. 112 is
% NOT among the pollution sites in this file's header -- a new site, and the
% third on this batch's pages after 108 and 111.
Theoretiskes, deels ved omhyggeligere og mere udtømmende
Bestemmelser af ethiske Grundbegreber (f.\ Ex.\ det Gode, Dy\-%
den, Frivillighed, Forsæt, Hensigt). Men trods dette viser dog
hans Ethik sig i sin \emph{hele Grundcharakteer} at være be\-%
stemt ved den \emph{almene græske} Grundanskuelse, hvis over\-%
mægtige Indvirkning vi allerede have seet indenfor hans An\-%
thropologie og Metaphysik. Afgjørende bliver her \emph{Anthro\-%
pologiens} Mangel paa Begrebet om den \emph{Personlighedens
indre Uendelighed}, der er Forudsætningen for det \emph{Ethi\-%
ske i stræng} Forstand. At der ikke naaedes til Begrebet
om Personligheden i dens indre Uendelighed, maatte medføre
væsentlige \emph{Consequentser} for det Ethiske. \emph{Deels} kunde
nemlig saadanne ethiske Begreber, der umiddelbart vise tilbage
til hiin Personlighedens og Villiens uendelige Concentration,
som Begreberne om ethisk Forpligtelse, Samvittighed, Anger,
o.\ desl.\ ikke gjøre sig gjældende. \emph{Deels} maatte, idet Natur\-%
elementet i Mennesket bevarede en blivende Umiddelbarhed,
der begrændsede Aanden, Naturbestemthedens Tilfældighed blive
bestemmende for Aandens Frihedsudvikling, og Begrebet om
det Almeen-Menneskelige kunde ikke vinde Klarhed. \emph{Deels}
maatte det upersonlige Almene, Staten, faae en Overvægt over
det sædelige Individ, analog med den Overvægt, de kosmiske
Principer, Sphærernes Sjæle, havde over den enkelte Menne\-%
skesjæl, og det maatte indvirke \emph{umiddelbart} bestemmende
paa Individets Villie. \emph{Deels} maatte endelig den ethiske Hand\-%
ling opfattes, ikke som udgaaende fra en i sig uendelig gjen\-%
nemsigtig, i udtømmende Reflexion klaret, Personligheds ethi\-%
ske Beslutning; men som blivende til, ligesom pludselig, kunst\-%
nerisk-genialt, ved en Concretion af det Almene, ved den
% p. 112: \gk{αἴσθησις} -- over the iota stands the same COMPOUND
% (breathing plus acute) as over the epsilon of ἔχον at p. 99 l. 16, and it
% resolves the same way at 1200 dpi: two elements separated by clear paper,
% the left a narrow oblique tick with no bowl (matching the certain smooth
% breathings on this book's pages) and the right rising to the right.
% Read ἴ.  This is NOT the single undifferentiated blob of ἄπειρον at
% pp. 52 and 81; per TRANSCRIPTION-PLAYBOOK §2 rule 2 that page's failure is
% no evidence about this one, and the positive signal here is the separation
% of the two elements.  The theta is the cursive ϑ; plain θ typed.
\emph{umiddelbare} Virkning af \gk{Νοῦς}, der unddrager sig Reflexio\-%
nen, og er Gjenstand for en \gk{αἴσθησις}. Og idet den \emph{ethiske
% p. 112: SEAM p. 112/113.  p. 112 bears NO footnote -- the text block runs
% to the foot of the page and there is no footnote separator -- so NOTHING
% can run over onto p. 113, and p. 113 opens with body text ("vidst ethisk
% Charakteerudvikling"), not with unmarked note matter.
% p. 112 ends MID-WORD and inside a letterspaced run: the printed word is
% "be\-vidst" and "vidst ethisk Charakteerudvikling" on p. 113 continues the
% spacing.  The two fragments' \emph{} groups have been MERGED into the
% single run the print shows, and the broken word joined.
Handling} opfattedes saaledes, maatte Begrebet om en \emph{be\-%
% --- p. 113 ---
% SEAM p. 112 --> p. 113, checked on the image at 600 dpi.  p. 112 bears NO
% footnote (there is no separator rule at its foot and the text block runs
% to the last line), so nothing runs over onto p. 113, whose first line is
% ordinary body text.  p. 112 ENDS MID-WORD: its last line is "Handling
% opfattedes saaledes, maatte Begrebet om en be-", and the "be-" is already
% LETTERSPACED, so the emphasis carries across the page break into "vidst
% ethisk Charakteerudvikling" below.  The run is set as ONE \emph{} across
% the page break, and there is no paragraph break here.
%
% THE p. 113 ERRATUM IS FOUND.  The Trykfeil leaf asks l. 35 for
% "Bostemmelse, læs: Bestemmelse."  The printed "Bostemmelse" actually
% stands on LINE 33, counted from the first line of type -- "Begrebet
% \gk{μεσότης} som Bostemmelse for den ethiske Dyd." -- the last line of
% that paragraph.  The page carries exactly 35 lines of type; l. 35 is
% "Opfattelse af Dyden som Væsenets uforstyrrede Virksomhed og", which
% contains no such word.  The author's line reference is therefore two
% lines out.  The sheet signature "8" stands below l. 35 and is not
% transcribed.  The printed reading is kept; see the % note at the site.
vidst ethisk Charakteerudvikling} mangle og erstattes
ved Forestillingen om det sædelige Liv som en \emph{Række} af
\emph{saadanne Umiddelbarhedens Øieblikke}, i hvert af
hvilke det ethiske Individ, i sin umiddelbare Bestemmethed ved
det Almene, staaer kunstnerisk afsluttet i sig, men saaledes,
at Continuiteten for Bevidstheden mangler. Dette maatte alle\-%
rede blive en Følge af den Grundanskuelse af Menneskets og
Tilværelsens Væsen, der gav det \emph{blot positivt Værende} og
den \emph{hvilende Umiddelbarhed} en saadan Overvægt, at
\emph{Vordens Mulighed} blev ubegribelig og \emph{Udviklingsprin\-%
cipet} forsvandt. Men i \emph{nærværende} Sammenhæng viser
det sig som \emph{uundgaaelig Følge} deraf, at Selvets bevidste
Fordybelse i sin indre Uendelighed, Villiens bevidste Concen\-%
tration i den ethiske Livsbeslutning: de \emph{nødvendige For\-%
udsætninger} for den bevidste ethiske Charakteerudvikling,
som vi have seet, ikke ere begrebne. I Individets Handling
og i dets Udvikling bliver der da, ud fra denne Opfattelse,
Mere af kunstnerisk ubevidst Frembringen end af bevidst Vil\-%
liesenergie, Mere af det Æsthetiske end af det Ethiske. Og
ved den Overvægt, der, \emph{paa Villiens Bekostning}, gives
det \emph{Theoretiske}, drages, principielt, den menneskelige Hand\-%
len hen mod det \emph{Metaphysiskes Nødvendighed}, ligesom
Bestemmelsen af det \emph{ethisk Gode} hos Grækerne altid nær\-%
mede sig til eller endog faldt sammen med det \emph{metaphy\-%
sisk Gode}. Den \emph{aristoteliske} Ethiks Charakteer bliver
% p. 113: "æsthetisk-metaphysisk" is a HARD compound hyphen that happens
% to fall at the line break (the compound recurs unbroken at p. 127,
% "den æsthetisk-metaphysiske Naturerkjendelse"), so the hyphen is kept
% as printed and only the newline is eaten.
derfor, ligesom den \emph{græske} Ethiks overhovedet, \emph{æsthetisk-%
metaphysisk}.

% p. 113 l. 29: both OCR witnesses put a stray apostrophe before "oplyse"
% ("'oplyse"); at 600 dpi there is nothing there but a speck of the
% modern pencil this copy carries.  Not transcribed.
Hvad der er sagt om Charakteren af \emph{den ethiske Be\-%
stemmen af Handlingen}, oplyse vi gjennem Udviklingen
af et vigtigt Begreb i den aristoteliske Ethik, der i høi Grad
har været misforstaaet, og som kun i Sandhed kan forstaaes
% p. 113 l. 33: "Bostemmelse" AS PRINTED -- an o where an e belongs.
% This is the ninth of the author's own errata (Trykfeil leaf, PDF 306:
% "Bostemmelse, læs: Bestemmelse"), and per transcription.tex §5 the
% printed reading is kept.  The leaf cites l. 35; the word stands on
% l. 33.  Confirmed at 600 dpi; a faint modern pencil tick stands above
% the word and is not transcribed.
% \gk{μεσότης}: acute over the omicron, no breathing; read at 1200 dpi.
naar det sees i det \emph{metaphysiske} Grundforholds Lys: af
Begrebet \gk{μεσότης} som Bostemmelse for \emph{den ethiske Dyd}.

Idet \emph{Aristoteles}, i Overensstemmelse med den \emph{græske}
Opfattelse af \emph{Dyden} som Væsenets uforstyrrede Virksomhed og
% --- p. 114 ---
% p. 114: a vertical modern pencil stroke stands in the left margin
% beside ll. 15--17.  Not transcribed.
Hævden af sig gjennem denne Virksomhed, bestemmer den
Dyd, \emph{Ethiken} skal undersøge: den \emph{menneskelige} Dyd, som
den Virksomhed, ved hvilken Mennesket udfører sin eiendom\-%
melige Gjerning vel, og idet det for Mennesket \emph{Eiendomme\-%
lige} af ham er bestemt som det \emph{fornuftige} Sjæleliv; saa
sondrer den menneskelige Dyd sig for ham i den \emph{dianoeti\-%
ske}: den normale Virken af det fornuftige Sjæleliv som saa\-%
dant, og i den \emph{ethiske} Dyd: den \emph{fornuftbestemte Vil\-%
lies Energie}. Den \emph{ethiske} Dyd, den \emph{specifisk menne\-%
skelige} Dyd, i Sammenligning med den \emph{menneskelig\-%
guddommelige} dianoetiske, bliver Ethikens egentlige Midt\-%
% p. 114 Greek, all read at 1200 dpi and nowhere else (neither witness
% carries a Greek character): \gk{φύσις} acute over the upsilon, no
% breathing; \gk{ἔθος} a two-element compound over the epsilon -- a curl
% opening LEFT (smooth) followed by a rising stroke (acute); \gk{λόγος}
% acute.  The theta is printed in the cursive ϑ form and, per
% transcription.tex §GREEK, is typed as plain θ.
punkt. Den forbinder som Momenter i sit Begreb \gk{φύσις},
% p. 114 ll. 13--14: the letterspacing covers only "Natur" in BOTH
% compounds -- "N a t u r bestemthed" and "N a t u r-/driften" -- with
% the second element set close.  Measured at 1200 dpi; the same split
% within one word is recorded at pp. 91 and 98.
\gk{ἔθος} og \gk{λόγος}, idet den oprindelige \emph{Natur}bestemthed, \emph{Natur}\-%
driften, som et Stofagtigt bliver formet med Bevidsthed af
den \emph{fornuftige} Indsigt, og fæstner sig i denne Form gjen\-%
nem \emph{Øvelsen, Vanen}, som gjør det, der oprindelig var den
enkelte frie Beslutnings Sag, til \emph{blivende Beskaffenhed i}
Sindelaget, og endog er Betingelsen for \emph{Erkjendelsen} af det
Sædelige. Forat erkjende dette, maa man nemlig, efter Ari\-%
% p. 114 l. 20: "Aristoteles" is NOT letterspaced here (checked at
% 1200 dpi); only "have" is.
stoteles, \emph{have} handlet sædeligt, fordi Erkjendelsen ogsaa her
% p. 114: \gk{τὸ ὅτι} -- the omicron of τὸ carries a stroke falling
% left-to-right (GRAVE, calibrated against the rising acutes of μεσότης
% on this page and p. 113); over the omicron of ὅτι stands a compound of
% a "C" opening RIGHT (rough) and a rising stroke (acute).
gaaer ud fra \gk{τὸ ὅτι} og bliver real paa samme Maade som paa
det theoretiske Gebet Fornufterkjendelsen af det Almene bliver
real ved Hjælp af Sandsningen af det Enkelte. Den ethiske
% p. 114, THE ARISTOTLE DEFINITION, read letter by letter at 1200 dpi.
% Breathings resolved by direct contrast within the one line: ἡμᾶς, ὡς
% and ὁ carry a "C" opening RIGHT (rough); ἐν and ἂν carry the mirrored
% comma opening LEFT (smooth).  ἕξις: rough + acute, both components
% separable.  οὖσα: the mark over the upsilon is a two-lobed compound
% whose left element curls like the smooth of ἐν; read smooth +
% circumflex.  ἂν: the breathing is positively smooth by the contrast
% just named, but the accent element is a blob; GRAVE is the reading
% set and ἄν is the alternative rejected.  τῇ and λόγῳ both carry a
% visible iota subscript.
% NOTE ON ὡρισμένη: the final eta has NO iota subscript, and this is a
% POSITIVE observation, not an absence -- the subscript under the omega
% of λόγῳ three words later is plainly there at the same magnification.
% The received Aristotelian text has ὡρισμένῃ.  Printed as shown.
Dyds fuldstændige Definition hos Aristoteles bliver: \gk{ἕξις προ\-%
αιρετικὴ, ἐν μεσότητι οὖσα τῇ πρὸς ἡμᾶς, ὡρισμένη λόγῳ
καὶ ὡς ἂν ὁ φρόνιμος ὁρίσειεν}.

I Begrebet \gk{μεσότης}, som Aristoteles anvender i denne
Definition af den \emph{ethiske} Dyd, og som man har villet opfatte
som en rent udvortes og empirisk Bestemmelse, viser der sig
% p. 114 ll. 30--31: spacing.py reports a run "Principer-/nes Forhold";
% on the image "Bestemmelse af Principer-/nes Forhold og af den" is set
% CLOSE and the spacing resumes only at "græske".  False run; ignored.
en Reflex af den \emph{metaphysiske} Bestemmelse af Principer\-%
nes Forhold og af den \emph{græske idealistiske Grundansku\-%
else}. \emph{Aristoteles} tænker sig nemlig Forholdet i Sjælens
\emph{fornuftløse, men for Fornuft modtagelige, Deel} saa\-%
ledes, at den med Lyst- og Ulystfølelsen sammenknyttede Af\-%
fect og Attraa, og den Handlen, der fremkaldes ved dem, ifølge
% --- p. 115 ---
\setcounter{footnote}{0}
% p. 115 CARRIES FIVE NOTES, not three: *) **) ***) †) ††).  Every mark
% was read on the image at 1200 dpi, in the text and again at the head of
% each note.  This CORRECTS transcription.tex §4, which says "Three levels
% are used: ***) occurs on p. 76 and p. 115 and nowhere else" -- the
% dagger and the double dagger are a fourth and a fifth level, and p. 115
% is the only page in pp. 113--126 that uses them.
% The preamble's \thefootnote has been extended to five levels
% accordingly; see it for the note.
deres Paavirkethed ved det \emph{Sandselige} og overhovedet ifølge
deres Sammenhæng med det \emph{Materielles}, det i sig Ubestem\-%
tes, Gebet, maae savne den sikkre \emph{Norm} og den rette Bestemt\-%
hed saalænge indtil de underlægges Fornuftens bevidste formende
% p. 115: the ɔ: id-est mark, confirmed at 1200 dpi -- a reversed c
% followed by a colon.  Neither witness reads it ("o:" in ABBYY).
Indvirkning, og at altsaa først Fornuften bringer \emph{Maal} og \emph{har\-%
monisk Bestemthed} ɔ: en \gk{μεσότης}, ind i Affecten\footnote{Forholdet
til Affecten er det Oprindelige (jfr. Eth. Nic. II, 6), der\-%
for er den \gk{μεσότης}, der tilveiebringes, en \gk{μεσότης πρὸς
ἡμᾶς}. Dog gives der en \gk{ἔλλειψις} og en \gk{ὑπερβολή} i og for sig.}
% p. 115 note *): ἡμᾶς rough + circumflex; ἔλλειψις smooth + acute;
% ὑπερβολή rough over the upsilon and acute on the final eta.  All read
% at note size at 1200 dpi.
og At\-%
traaen og derigjennem i Handlingen. \emph{I dette Forhold virker}
% p. 115 l. 8: the GREEK IS LETTERSPACED here and only here on this page.
% Measured against the same two words set close elsewhere on the page:
% φρόνησις on l. 8 is visibly wider than φρόνησις on l. 11, and νοῦς on
% l. 8 wider than νοῦς on l. 9.  Both therefore take \emph{\gk{...}}.
\emph{\gk{νοῦς}} \emph{som praktisk Indsigt, som} \emph{\gk{φρόνησις}}. Ligesom,
efter den aristoteliske Erkjendelseslære, \gk{νοῦς} \emph{umiddelbart}
griber de høieste Principer, Forudsætningerne for al Viden,
saaledes skal paa det Ethiskes Gebet \gk{φρόνησις} ved en umid\-%
% p. 115: \gk{αἴσθησις} -- over the iota stands a two-element compound.
% The left element is a hook and the right a rising stroke, read as
% smooth + acute; ἴ/ἵ cannot be fully separated at note-adjacent size
% and ἅ is the alternative rejected.  Cf. the same treatment at p. 33.
delbar Takt, en Art \gk{αἴσθησις} (jfr. Eth. Nic. II, 9. VI, 9, 12),
gribe det Gode og Rette i dette aldeles concrete Tilfælde, og
det ved en Virkning af \gk{νοῦς}, der, som den høieste Energie,
% p. 115 l. 15: "Fornuftfunctioner" -- ABBYY reads "Fornuftfunctionor"
% and tesseract "Fornuftfunetioner"; the image is plain at 1200 dpi.
maa være det Virkende i alle Fornuftfunctioner (Eth. Nic. VI,
12). Idet \emph{Handlingen} bestemmes gjennem Affectens og At\-%
traaens Regulerethed ved Indsigten\footnote{Attraaen bliver saaledes
en \gk{ὄρεξις διανοητική, ὄρεξις ὀρθή}.}, og selv gaaer ind under
Bestemmelsen \gk{μεσότης}, saa giver den, som et i det Ydre Frem\-%
trædende, \gk{μεσότης} Skikkelsen af en \emph{Middelvei}, og derfor
kan Begrebet \gk{μεσότης} overføres paa Gebeter, hvor det ikke
hører hjemme efter dets egentlige Betydning. Naar Begrebet
anvendes ved Bestemmelsen af den \emph{ethiske} Dyd, saa er alt\-%
saa derved udtrykt Menneskets sandselige Naturs harmoniske
% p. 115 note ***): \gk{νόμος} alone is LETTERSPACED inside the Greek
% quotation -- its letters are visibly spread where ὁ γὰρ and τὸ μέσον
% on either side are set close.  ὁ rough, γὰρ and τὸ grave, μέσον acute.
Bestemthed i sin Virken ved Formen, det Fornuftige\footnote{Dens
Lovbestemthed: \gk{ὁ γὰρ} \emph{\gk{νόμος}} \gk{τὸ μέσον}.} --- ikke
den umiddelbare Naturbestemmethed, det umiddelbart givne
Naturanlæg\footnote{Derfor regner Aristoteles f. Ex. ikke
\gk{ἐγκράτεια} til Dyderne.}, men den umiddelbare, ligesom æsthetiske\footnote{Jfr.
i Eth. Nic. II, 5 Sammenstillingen af den ethiske \gk{μεσότης} med
Kunstværkets.},
% The sheet signature "8*" stands at the foot of p. 115; not transcribed.
% p. 115 ends in a comma and p. 116 continues the same sentence; no
% paragraph break at this seam.
% --- p. 116 ---
% p. 116 bears no footnote.
Bestemmethed gjennem \emph{Indsigten}. \gk{Μεσότης} udtrykker, med
denne nærmere Bestemmelse af et \emph{Bevidsthedsforhold}, det
Samme, som \emph{Plato} i nogle af sine sildigere Dialoger kaldte
% p. 116: the ɔ: id-est mark, confirmed at 1200 dpi.  τὸ twice with a
% grave; μέτριον and μέσον with acutes.
\gk{τὸ μέτριον} og \gk{τὸ μέσον} ɔ: det, der ligger midt imellem det
Mere og det Mindre, det Stofagtiges Bestemmelser, og hos
\emph{Aristoteles} vexle ogsaa disse Begreber med Begrebet \gk{μεσότης}.
Af denne Betydning af Begrebet følger saa ogsaa, at det hos
Aristoteles kun kan faae Gyldighed for de \emph{ethiske} Dyder,
ikke for de \emph{dianoetiske}, thi ved disse er der ikke noget
Forhold til et Stofagtigt, der skal begrændses; derfor kan der
ved dem ikke blive Tale om et Overmaal eller en Mangel.

I Betegnelsen: \gk{μεσότης}, og i den nærmere Udvikling af
dette Begreb, bliver den \emph{æsthetiske} Opfattelse af den sæde\-%
lige Bestemmen og den sædelige Handling iøinefaldende. Er\-%
kjendelsens Gjenstand: det Almene, Loven, bliver i den enkelte
sædelige Handling et \emph{paa Umiddelbarhedens Maade Vir\-%
kende}; i den sædelige Handling gaaer, paa samme Maade som
i \emph{Instincthandlingen}, den almene Lov umiddelbart i \emph{En\-%
hed} med den individuelle Bestemthed, og Bevidstheden for\-%
holder sig kun \emph{til Resultatet}, ikke til den \emph{Proces}, af hvil\-%
ken det resulterer. Og idet det sædelige Liv bliver Rækken
% p. 116 l. 23: \gk{τέλος} is LETTERSPACED (τ έ λ ο ς, gaps measured at
% 1200 dpi against the close-set Greek elsewhere on the page), and it
% closes a spaced run that begins with "en" at the end of l. 22.
af saadanne Umiddelbarhedens Øieblikke, saa er der vel \emph{en}
\emph{Fremskriden mod et ethisk} \emph{\gk{τέλος}}, men denne Frem\-%
skriden, der foregaaer som ved en harmonisk \emph{Naturproces},
og bliver en \emph{Tildragelse}, ikke en Handling i stræng For\-%
stand, foregaaer paa samme Maade som naar Uhrviseren i eet
Spring flytter sig fra det ene afdelte Punkt til det andet, og
saa staaer stille til næste Spring; den mangler Continuiteten.

\emph{En Udvikling i stræng} Forstand, en \emph{bevidst ethisk}
\emph{Charakteerudvikling} faae vi, som det allerede ovenfor blev
fremhævet, ikke, og vi kunde ikke faae den, ifølge den \emph{ari\-%
stoteliske} og overhovedet den \emph{græske} Tænknings Grund\-%
eiendommelighed, der ikke kunde føre til en Erkjenden af
den ved \emph{Negationens} Optagen i \emph{Væren} betingede \emph{indre}
% --- p. 117 ---
\setcounter{footnote}{0}
\emph{Bevægelse}, af Modsætningernes Udfoldelse af hinanden og
% p. 117 l. 2: spacing.py reports a run "Sammenslutten til"; on the
% image the whole line is set close.  False run; ignored.
af deres Sammenslutten til en levende Enhed\footnote{Gjennem det
Udviklede vil det omtvistede Spørgsmaal kunne afgjø\-%
res, \emph{om det græske Drama kjender en Charakteerudvik\-%
ling i strængere Forstand.} Det maa besvares \emph{benægtende.}
Naturligviis kan den græske Dramatiker ligesaalidt have lukket Øiet
for det Phænomen, at en Charakteer ikke altid er den samme, men
under forandrede Betingelser fremtræder paa en anden Maade, paa
et høiere eller lavere Trin, som den almindelige græske Bevidsthed
kunde lukke Øiet for de Phænomener, der vise tilbage til en For\-%
andring, en \emph{Vorden.} \emph{Den} græske Dramatiker har derfor ikke
fremstillet sine Personer i død Uforanderlighed; dette behøver kun
en Henvisning til en Figur som Neoptolemos (i Sophokles's Philok\-%
tet) forat være uimodsigeligt. Men ligesom det er Et, at \emph{indrømme
Forandringens Phænomen}, et Andet, at \emph{begribe Vorden},
altsaa \emph{for Tanken} at \emph{retfærdiggjøre} hint Phænomen, saaledes
er det Et at \emph{skildre} en Charakteer paa dens \emph{forskjellige} Ud\-%
viklingsstadier, et Andet at \emph{tydeliggjøre} den indre Udfoldelse,
den Udvikling ud fra Personlighedens Inderste, ved hvilken hine
% p. 117 note: \gk{αδυτον} -- the mark over the alpha is a single
% compound blob at note size and 1200 dpi; neither the breathing nor the
% accent can be given its own shape and ἄ/ἅ cannot be separated.  Per
% BATCH-AGENT §4 rule 2 and the ἄπειρον precedent recorded at pp. 51, 52
% and 81, the alpha is set UNACCENTED and nothing is supplied.
Stadier \emph{komme frem.} Dette gjorde den græske Dramatiker ikke;
thi dette Inderste var for den græske Digter i en endnu bestemtere
Betydning end for Digteren overhovedet et \gk{αδυτον}, paa Grund af
det Halvmørke, i hvilket den græske sædelige Bevidsthed i dens
Umiddelbarhedsbestemthed hvilede. Charakteerbestemmelsens Stadier
i det græske Drama forholde sig som de Billeder, der i katholske
Lande fremstille Stadierne i Lidelseshistorien. Der er en \emph{Sammen\-%
hæng} mellem disse «Stationer», men de vise sig hver for sig som
\emph{et hvilende} Hele, og Sammenhængen virkeliggjøres kun i \emph{den
Tilskuers Bevidsthed}, der, idet han, fra den ene Station til den
anden, bringer \emph{Mindet} om det Anskuede med sig og \emph{samler} Bil\-%
lederne i \emph{Erindringen}, ved en \emph{levendegjørende} Phantasies
Magt forvandler de sondrede Billeder til en Heelhed og giver dem
den \emph{Continuitet}, de ikke have ved sig selv. Saaledes faae ogsaa
hine det græske Dramas Billeder kun den sande Continuitet for en
\emph{Bevidsthed}, der i sig har Magten til at bringe dem i Strømning
og derved see dem i indre Enhed som fremgaaende af eet Udvik\-%
lingsprincip.}.
% p. 117: the note is SELF-CONTAINED -- it ends "lingsprincip." at the
% foot of the page and p. 118 opens with its own *) note.  Checked.
%
% p. 117 RUN-IN HEAD.  "b)" is set in the ORDINARY roman and "Enhed og
% Dualisme." in a markedly heavier BOLD; measured at 1200 dpi, the bold
% letters are tight (no letterspacing) and their stems are visibly
% thicker than the "b)" on the same line.  So the head takes \textbf{}
% and the enumerator does not.  The same is true of "c) Det Enkelte og
% Erfaringen." on p. 120, checked the same way.

b) \textbf{Enhed og Dualisme.} --- Naar det, at Løsningen af Vor\-%
dens Problem mislykkedes for Aristoteles, uundgaaelig med\-%
førte, at Løsningen af Problemet om Tilværelsens Concretion
% --- p. 118 ---
\setcounter{footnote}{0}
maatte mislykkes, saa viser det tilbage til, at \emph{Dualismen
ikke er fjernet}; thi netop den uvilkaarlige abstracte Sondring
af Grundmomenterne var det jo, der gjorde Vorden ubegribelig.
Det er allerede antydet, hvorledes Aristoteles, mere energisk
end nogen anden græsk Tænker, stræber at overvinde Dualis\-%
men. Ved Bestemmelsen af Begrebsforholdet mellem Form og
Stof, og fornemmelig ved Opfattelsen af Formen som \emph{Udvik\-%
lingsprincip}, viser han afgjørende hen mod den principi\-%
elle Enhed; men dog bliver det Tilstræbte ikke naaet; hans
bevidste Stræben hæmmes ved den ubevidst virkende Magt,
der ligger i den græske Tænknings Naturbegrændsning, og
Dualismen bryder atter frem og adskiller det tilsyneladende
Forenede. Adskillelsen viser sig fra \emph{Formens} Side navnlig
deri, at den, i sin høieste Fuldendelse, bliver den \emph{fra alt
Stof sondrede} Form, og at den overhovedet, selv hvor den
sees i Forhold til Udviklingen, skiller sig fra det Stofagtiges
% p. 118 l. 20: both witnesses read a period after "røber" ("røber. sig").
% At 1200 dpi the mark is a small ORANGE-BROWN blob sitting above the
% baseline, warmer and lighter than the black of the type -- a foxing
% spot in the paper, not a sort.  Not transcribed; logged here so that a
% later pass does not restore it from the OCR.
\emph{Vorden}; --- fra \emph{Materiens} Side navnlig deri, at der til\-%
lægges den en \emph{Modstandskraft}, der kan foranledige det fra
Formen Afvigende, altsaa en selvstændig Virken, gjennem hvil\-%
ken en \emph{Grundmodsætning} til Formen røber sig. I denne
Grundmodsætning staae Form og Stof som \emph{lige oprindelige};
deres \emph{Forbindelse for Anskuelsen} bliver den \emph{umiddel\-%
bare} Forening af det \emph{Forskjellige}, deres Enhed \emph{for Tan\-%
ken} kun et \emph{Postulat}. Dette ligger ogsaa i Aristoteles's Ud\-%
sagn, at de to Principer: Form og Stof, hverken kunne føres
tilbage til hinanden eller til det samme Slægtbegreb\footnote{\emph{Værens}
Bestemmelse, der bliver fælleds for dem, opfatter nemlig
Aristoteles ikke som \emph{Slægtbegreb}, ligesaalidt som \emph{Enhedens}
Begreb.}.

Paa en slaaende Maade bryder den af Aristoteles bekæm\-%
% p. 118 l. 29: \gk{Νοῦς} is LETTERSPACED, standing inside the spaced run
% "hans Fordobbling af Tilværelsen i Νοῦς og Universet"; the letters are
% visibly spread at 1200 dpi.  The circumflex over the upsilon is an
% unmistakable tilde.
pede Dualisme frem i \emph{hans Fordobbling af Tilværelsen
i} \emph{\gk{Νοῦς}} \emph{og Universet}. Skjøndt hans Stræben gaaer hen
imod en Gjennemtrængen af Stoffet ved Formen, der, conse\-%
quent, vilde føre til at see dem som uadskillelige, saa at den
% --- p. 119 ---
\setcounter{footnote}{0}
høieste Form vilde være den, der fuldstændigst gjennemtrængte
og beherskede Stoffet; at see det Evige som evigt i det ved
Tid og Bevægelse Bestemte, det af Rummet Uafhængige som
paa rumfri Maade værende i Rummet som Enhed; saa hæm\-%
mes dog denne Stræben ved Principernes oprindelige Adskilt\-%
% p. 119: \gk{Κόσμος} (twice) with an acute over the omicron, and
% \gk{Νοῦς} (six times) with a plain tilde-shaped circumflex over the
% upsilon; both read at 1200 dpi, neither letterspaced on this page.
hed. For Aristoteles former Verdensvirkelighedens Heelhed sig
til Anskuelsen af \gk{Κόσμος} som det altomfattende, levende, af
Formen besjælede Kunstværk. Dette formede Universum, hvis
Stof er det Formendes Mulighed, og som evigt bevæges ved
Forholdet til den guddommelige Bevæger, indeslutter i sig en
\emph{fuldstændig} Trinrække af Former, der fra den laveste stiger
op til den med den guddommelige Tanke identiske active
\gk{Νοῦς} i Mennesket, og i denne Fuldstændighed udtrykker
Formens Væsen \emph{heelt}. Aristoteles kommer det her saa nær
som muligt at lade Formen som guddommelig \gk{Νοῦς} faae sin
% p. 119 l. 16: "Heelhed", not the "Ideelhed" ABBYY reports; plain at
% 1200 dpi and confirmed by tesseract.
Plads \emph{indenfor} Verdens Virkelighed. Naturen i \emph{dens Heel\-%
hed} sees som \emph{besjælet}, Formen sættes som Verdens \emph{indre
bevægende Entelechie}, og denne Form, der teleologisk
udfolder sig i de særegne Formers Række, viser, gjennem
Rækkens Afslutning i Menneskets \gk{Νοῦς}, sin \emph{Enhed med den
guddommelige}. \emph{I Mennesket} har altsaa den guddomme\-%
% p. 119 l. 22: WRONG SORT.  The "i" of "en Væren i Verden" is set from
% a markedly heavier fount -- stem and dot both visibly thicker than the
% "i" of "lige" earlier on the same line, measured at 1200 dpi.  This is
% the fourth instance of the defect recorded in the file header at
% pp. 42, 89 and 95 (a stray heavy sort in the compositor's "i" box), not
% an emphasis.  Transcribed as an ordinary "i".
lige \gk{Νοῦς} en Væren i Verden, og da, ifølge dens immaterielle
Væsens Forhold til Rumsbestemtheden, dens rumlige Sondring fra
Universet ikke kan fastholdes af Tanken, saa synes der ingen
Hindring at være for, overhovedet at tænke sig den guddom\-%
melige \gk{Νοῦς} som \emph{immanent i Verden}\footnote{Jfr. Udviklingen
i den peripatetiske Skole.}, og det saameget
mindre som Aristoteles jo ogsaa tænker sig Himmelsphærernes
evige Bevægere som ubevægede i de bevægede Sphærer, Sjælen
som ubevæget i det bevægede Legeme, Formen som ubevæget
og evig i det foranderlige og bevægede Stof, og den active
\gk{Νοῦς} i Mennesket uden Delthed og Vorden.

\emph{Men dog gjennemfører Aristoteles ikke Enheden
til dette sidste Punkt.} \gk{Κόσμος} bliver ikke Tilværelsens
% --- p. 120 ---
\setcounter{footnote}{0}
Heelhed, ikke heller bliver den guddommelige \gk{Νοῦς}, der dog
i sit Væsen indeslutter al Virkelighed, det Ene-Værende; men
Virkeligheden fordobbles i den rene stofløse \gk{Νοῦς} og i Ver\-%
% p. 120 note *): \gk{εῖδος} -- over the iota stands a broad
% circumflex-shaped mark, plainly visible at 1200 dpi, but no separate
% breathing element can be told from it at note size; ἶ/ῖ undecided.
% Only the circumflex is transcribed and nothing is supplied for the
% breathing, exactly as at p. 33 ("den høieste \gk{εῖδος}") and p. 52
% ("\gk{τὸ εῖναι}").  \gk{μορφή} carries an acute; \gk{Νοῦς} a tilde.
% p. 120 note, l. 4: both witnesses hang a stray mark on "sig" ("sig*",
% "sigt").  At 1200 dpi it is a faint tick of the modern pencil, not a
% sort.  Not transcribed.
densvirkeligheden, i hvilken Stoffet, der ikke kunde bortskydes
af Tanken\footnote{Ogsaa af den Grund laae det nær for Aristoteles, at
lade Sondringen trænge ind, at \emph{Naturens} \gk{εῖδος}, opfattet som
\gk{μορφή}, som Universets \emph{physiske Form}, trods \emph{Formbegrebets}
Tvetydighed, uvilkaarlig maatte sondre sig fra Formens \emph{ideelle}
Skikkelse i den menneskelige active \gk{Νοῦς}. Kun hiin Tvetydighed
tilslører den Vanskelighed, der kommer frem paa \emph{ethvert} Punkt,
hvor en Form, der som \emph{immateriel} intet Rumsproducerende har
\emph{i sig}, skal give Stoffet Rumsform.}, faaer sin Plads: \emph{Verden i sin Fylde kommer
til at staae som Virkelighed ved Siden af den efter
Begrebet altomfattende Virkelighed}; den \gk{Κόσμος}, i
hvilken Virkelighed evigt reproduceres af Mulighed, hævder
sig ved Siden af den rene Energie, i hvilken al Mulighed er
forsvunden.

I denne Fordobbling viser Dualismen sig som uovervindelig.
% p. 120 l. 13: "harmnoisk" AS PRINTED -- the m and the o are
% transposed, where "harmonisk" is expected.  Read the same way by both
% OCR witnesses and confirmed on the image at 1200 dpi.  NOT on the
% author's Trykfeil leaf.
I samme Øieblik som den synes forsvunden i Verdensvirkelig\-%
heden, der af Anskuelsen opfattes som harmnoisk Enhed, bryder
den frem i Fordobblingen idet \emph{Tanken} atter udsondrer Formen
i dens \emph{Reenhed} og \emph{Anskuelsen} fæstner den i Sondringen
fra \gk{Κόσμος} ved uvilkaarlig at gjøre denne til en \emph{Rums\-%
sondring}. Og det i Anskuelsen af \gk{Κόσμος} umiddelbart Forenede,
der var det oprindelig Forskjellige, viser sig atter i dualistisk
Sondring saasnart \emph{Tanken} bestemmer det ved sine Former.

c) \textbf{Det Enkelte og Erfaringen.} --- Ligesom Aristoteles, i
Bestemmelsen af Grundmodsætningernes Forhold og i Hæv\-%
delsen af det Begreb, ved hvilket Principets Udfoldelse skulde
blive mulig og Tilværelsens Concretion forklares, viste \emph{ud
over den græske Tænknings Grændser}, men uden at
kunne naae det Maal, til hvilket han stræbte hen, saaledes
vil det nu ogsaa vise sig, at han, medens han ved sin Frem\-%
hæven af det \emph{naturligt Enkelte} og af \emph{Erfaringens Be\-%
% --- p. 121 ---
% p. 121 bears no footnote.  A short modern pencil stroke stands in the
% right margin beside l. 28; not transcribed.
tydning} for Erkjendelsen antydede en \emph{senere} Tidsalders Op\-%
fattelse, heller ikke i denne Henseende naaede hvad han
tilstræbte, fordi den almeen græske Opfattelse ham ubevidst
gjorde sin Magt gjældende. Og idet han hverken kunde
gjennemføre hine Antydninger eller fastholde det Begreb, ved
hvilket Udfoldelsen af Tilværelsens Rigdom skulde forklares,
saa maatte hans \emph{almene Naturanskuelse} i samme Øieblik
som den syntes at skulle gjennembryde den græske Tænknings
Skranker, dog igjen gaae ind under dennes Grundbestem\-%
melse.

Idet Aristoteles opfattede \emph{Formen} som \emph{iboende} Form,
som det, der giver sig Virkelighed \emph{i den concrete, reale
Tilværelse}, faldt der hos ham et ganske andet Eftertryk
paa denne Tilværelse end hos \emph{Plato}: et Eftertryk, der allerede
antydes derved, at hos Aristoteles \emph{Materiens} Begreb faaer
sin Plads i \emph{Metaphysiken} (fordi Materien er væsenseet med
% p. 121: \gk{οὐσία} -- the mark over the upsilon is a comma opening to
% the LEFT, the SMOOTH breathing, positively distinct from the "C"
% opening right that stands over the eta of ἡμᾶς on p. 114.  This is one
% of the file header's calibration words for the smooth breathing.
% \gk{τόδε τι}: acute over the omicron, τι unmarked.  Neither Greek
% phrase is letterspaced, though the Danish around both is.
Formen). Det \emph{substantielle Værende}, \gk{οὐσία}, \emph{i egentlig}
Betydning, blev af Aristoteles bestemt som et \gk{τόδε τι}, som
det individuelle Subject, Naturindividualiteten i dens Concre\-%
tion, og dette concrete Enkelte maatte, som væsentligt og
substantielt, \emph{være væsentlig Gjenstand for Erkjen\-%
delsen}, ligesom den Erkjendelsens Form, der optager dette
Enkelte: \emph{Erfaringen}, maatte blive af \emph{væsentlig} Betydning.
Og idet den iboende Form saaes som en \emph{Formernes Trin\-%
række}, der strakte sig gjennem \emph{hele} den naturlige Tilværelse,
% p. 121: the ɔ: id-est mark at the end of l. 27, confirmed at 1200 dpi.
fik denne Tilværelse \emph{efter sit hele Omfang Værd} for Er\-%
kjendelsen, og den fik dette Værd \emph{som} naturlig Tilværelse ɔ:
som den Tilværelse, til hvis Fuldstændighed \emph{det Materielle},
i hvilket Formen er iboende, samt \emph{Bevægelsen}, ere \emph{nød\-%
vendige}, og som altsaa kun erkjendes \emph{i sin Virkelighed}
naar det erkjendes som denne \emph{Enhed af Form og Materie}.
Medens derfor den aristoteliske Tænknings idealistiske Grund\-%
retning førte ham til, ogsaa med Hensyn til Naturtilværelsen,
at sætte \emph{Formen} som det \emph{Væsentlige}, og at opfatte de
\emph{finale} Aarsager som de \emph{egentlige} Aarsager, de materielle
% --- p. 122 ---
% p. 122 bears no footnote.  It is the heaviest Greek page of the batch;
% every word below was read at 1200 dpi.  \gk{μορφή} acute; \gk{εῖδος}
% treated as at p. 120 (broad circumflex over the iota, breathing not
% separable, nothing supplied); \gk{φύσις} acute and LETTERSPACED;
% \gk{γένεσις} acute; \gk{ὁδὸς} rough over the first omicron and grave
% over the second; \gk{εἰς} and \gk{εἰδοποιὸς} carry a comma-shaped
% smooth hook over the iota, with a grave over the omicron of the
% latter; \gk{φυσικῶς} and \gk{λογικῶς} circumflex over the omega,
% \gk{σκοπεῖν} circumflex over the iota; \gk{οὐσία} smooth, as p. 121;
% \gk{πρώτη}, \gk{σύνολον} and \gk{διαφορά} acute.  The "=" in l. 1 is
% text, not math mode (transcription.tex §9).
kun som Betingelser, ja endog til at bruge \gk{μορφή} (= \gk{εῖδος})
som \emph{Synonym med} \emph{\gk{φύσις}}, medens \gk{γένεσις} bliver \gk{ὁδὸς εἰς
φύσιν}, saa fører hiin Stræben efter at see det Ideelle som
\emph{væsenseet med det Materielle} og som \emph{iboende i det}
ham til at lægge en Accent paa det \emph{sandselige} Phænomen,
fra hvilket alene der ad Iagttagelsens og Erfaringens Vei kan
naaes til Naturerkjendelsen, og for denne overalt at fordre en
Undersøgelse af Naturtingenes \emph{eiendommelige physiske
Grunde} og \emph{concrete Betingelser} (\gk{φυσικῶς σκοπεῖν}), i
Modsætning til et abstract-almindeligt Begrebsræsonnement
(\gk{λογικῶς σκοπεῖν}). Hans hele Retning viser sig i disse Be\-%
stemmelser som afvigende fra \emph{Platos}. Denne kunde, idet
han satte Ideen som Videns sande Indhold og som den egent\-%
lige \gk{οὐσία}, vel, gjennem Anskuelsen af \emph{Urbilledet}, naae til
en \emph{ideel Enkelthed}; men i denne afspeilede sig, ifølge
Urbilledets Natur som Hypostaseringen af Tankens Almene
for Anskuelsen, det plastiske Ideals Forhold mellem det En\-%
kelte og det Særegne. Enkeltheden \emph{som Virkelighedens
concrete Enkelthed}, som \emph{Naturindividualitetens}
Bestemthed, forsvandt som det Ligegyldige, ligesom Endelig\-%
hedens Tilfældigheder forsvinde i det plastiske Ideal, og \emph{Na\-%
turerkjendelsen} faldt \emph{udenfor} den strænge Videns Kreds,
fik kun en uvæsentlig Betydning i Sammenligning med Er\-%
kjendelsen af Ideen. \emph{Aristoteles's} Bestemmelser antyde en
\emph{Stræben} henimod den Naturerkjendelse, som først en \emph{nyere}
Tid gjennemfører, mod \emph{den principielt fordrede Erfa\-%
rings} Naturerkjendelse. Men denne Stræbens Virkeliggjørelse
hæmmes allerede ved Aristoteles's \emph{metaphysiske} Bestem\-%
melser og i sidste Instants ved \emph{den græske Naturopfat\-%
telses almindelige Charakteer}.

Naar Aristoteles, i Metaphysiken, bestemmer \emph{Individet}
som \emph{\gk{πρώτη οὐσία}}, lader han den individuelle Væsenhed con\-%
stitueres \emph{ved Formens og Stoffets Forening til et}
\emph{\gk{σύνολον}}, og i denne Forbindelse sætter han \emph{Formen}, og
% p. 122: the ɔ: id-est mark at the end of the last line, confirmed at
% 1200 dpi.  \gk{πρώτη οὐσία}, \gk{σύνολον} and \gk{εἰδοποιὸς διαφορά}
% are all LETTERSPACED -- their letters are visibly spread against the
% close-set Greek higher on the same page.
bestemtere Formen \emph{som} \emph{\gk{εἰδοποιὸς διαφορά}} ɔ: Differentsen
% --- p. 123 ---
\setcounter{footnote}{0}
(den specifiske Differents) i \emph{Enhed med} det ved den be\-%
% p. 123 l. 2: "som det Væsentlige og egentlig / Constitutive" -- ABBYY
% transposes this to "som det egentlig Væsentlige og Constitutive".
% The image and tesseract agree on the order set here.
stemte \emph{Slægtbegreb}, som \emph{det Væsentlige og egentlig
Constitutive}. Men i \gk{εῖδος}, der constituerer Væsenet, gjør
\emph{Almindelighedens} Element sig gjældende (jfr. Aristo\-%
teles's Metaphysiks Bestemmelser), og idet \emph{kun det Almene}
bliver Videns \emph{egentlige} Gjenstand, medens Sandsetingene
\emph{som saadanne} ikke ere Gjenstande for Viden, saa kan man
med Viden vel stige ned til \emph{det concreteste Artbegreb},
til det \emph{laveste Begreb}, forat bruge den \emph{formelle}
Logiks Udtryk, men \emph{ikke} gaae videre uden at komme \emph{uden\-%
for Videns Grændse}, og dette vil, ifølge det ogsaa af
Aristoteles statuerede Forhold mellem \emph{Viden og Væren},
sige, \emph{at det Almene dog hævder sig den egentlige
Væsentlighed}. Aristoteles kommer derfor til at betegne
% p. 123: \gk{πρότερον}, \gk{μᾶλλον ὄν} and \gk{πρώτη οὐσία} are
% LETTERSPACED.  μᾶλλον carries a circumflex over the alpha; ὄν a
% two-element compound over the omicron -- a hook opening LEFT (smooth)
% and a rising stroke (acute).  \gk{εῖδος} on l. 3 as at pp. 120 and 122.
\emph{Formen} som \emph{\gk{πρότερον}} og \emph{\gk{μᾶλλον ὄν}} end den enkelte
Ting, i hvilken den er iboende, ja han bruger endog ligefrem
om den Betegnelsen \emph{\gk{πρώτη οὐσία}} (Met. VII, 7).

Hans \emph{Metaphysiks} Grundbestemmelser gjorde det \emph{umu\-%
ligt} for ham at gjennemføre Enkeltvæsenets Fortrinsret til
Bestemmelsen \gk{οὐσία}. Forat dette skulde kunne skee, maatte
der i Individualiteten paavises \emph{et positivt Mere}, der kunde
adskille den fra Arten, men et Saadant kan Aristoteles ikke
hævde den. Da \emph{Formen} er sat som \emph{det Almene}, maatte
hint Plus søges i \emph{Stoffet}, der af Aristoteles deels udtrykkelig
sættes som \emph{Flerhedens og Individualitetens Grund},
deels indirecte bestemmes som saadan gjennem det Udsagn,
at der ved Individualitetsforskjelligheden \emph{intet Øiemed}\footnote{Man
erindre, hvorledes efter Aristoteles's Opfattelse Øiemedets og
Formens Bestemmelse væsentlig høre sammen.}
\emph{tilsigtes}. Men Materien som saadan er \emph{det absolut Be\-%
stemmelsesløse} og det ved sig selv \emph{absolut Virkelig\-%
hedsløse}; den kan derfor ikke tænkes at bestemme Formen
saaledes, at det Individuelle kunde faae et Virkelighedsindhold,
der var eiendommeligt for det, og Enkeltvæsenets individuelle
% --- p. 124 ---
% p. 124 bears no footnote and carries no Greek.
Substantialitet bliver altsaa et Uforklarligt. Aristoteles \emph{maa}
statuere den ifølge hans Tænknings \emph{eiendommelige} Ret\-%
ning, men han maa ved \emph{den} Tænknings Art, der bærer og
omfatter hans individuelle Tænkning, og som lader Formen,
det ideelle Almene, \emph{umiddelbart} beherske sin Modsætning,
saa at denne som en blot Mulighed forsvinder i hiin og For\-%
men ikke ved Forholdet til den kan vinde Concretion, føres
til at give Bestemmelser om de den individuelle Existents
constituerende Factorer, der ophæve det Statuerede og gjøre
dets Mulighed ubegribelig. Og idet Enkeltvæsenet atter træder
tilbage for det Almene, maa \emph{Erfaringserkjendelsen} træde
tilbage for en \emph{speculativ Naturopfattelse gjennem ab\-%
stracte Begreber}. ---

Hvad der har viist sig som umiddelbart sammenhængende
med den aristoteliske \emph{Metaphysiks} Bestemmelser, viser sig
nu tillige som dybere begrundet i den \emph{græske Verdens\-%
anskuelses} og i den ved denne bestemte \emph{Tænknings al\-%
mindelige Charakteer}. Idet Grækerne anskuede Tilværel\-%
sen som en \emph{umiddelbar Enhed af Aand og Natur med
det Aandeliges Principat}, saa maatte, ifølge Enhedens
Umiddelbarhed, det Naturlige \emph{umiddelbart} sees som bestemt
ved det Aandelige, som afspeilende dette. Naturen saaes
umiddelbart i den menneskelige Aands Lys, blev anthropomor\-%
phiseret, fik en menneskelig Charakteer, en menneskelig Væren.
Derfor kommer, i Opfattelsen af den, \emph{Erfaringen} ikke til
sin Ret. Erfaringen, i den \emph{prægnante} Betydning, i hvilken
en \emph{moderne} Naturvidenskab tager Begrebet, forudsætter en
dybere Adskillelse mellem Tanken og dens Naturgjenstand end
Grækerne kjendte; den forudsætter, at det Naturlige, Mate\-%
rielle, ligesom er blevet befriet til en relativ Selvstændighed,
at det har kunnet gjøre sine egne Love gjældende, og deri\-%
gjennem hævdet sig som \emph{eiendommelig Kilde til Er\-%
kjendelse} for Aanden. Men hos Grækerne er Materien kun
\emph{det Ikke-Værende} i Forhold til Tanken, og derfor kan den
og det, i hvilket den er \emph{charakteristisk} Bestemmelse, ikke
% --- p. 125 ---
% p. 125 bears no footnote and carries no Greek.
tilkæmpe sig en saadan Selvstændighed og en saadan Betyd\-%
ning. Grækeren betragter væsentlig Naturen, som Heelhed og
i det Enkelte, ikke med den iagttagende Naturforskers, men
med \emph{Kunstnerens} Øie. Ligesom Naturtotaliteten anskuedes
som Verdenskunstværk, saaledes anskues det Enkelte i Ana\-%
logie med det kunstneriske Ideals Enkelthed, der umiddelbart
aabenbarer Ideen og lader Endelighedens Tilfældigheder for\-%
svinde: den concrete Naturvirkeligheds Enkelthed, Naturindivi\-%
dualiteten, forvandler sig til den skjønne Individualitets Sær\-%
egenhed, og Anskuelsen, i dens umiddelbare Enhed med Tan\-%
ken, bliver den til det Særegne, ikke til det individuelt enkelte
Naturlige, sig forholdende Anskuelse. Gjennem Anskuelsen
bestemmer Tanken Naturen: Opfattelsen af det Physiske bliver,
hvor Tænkningen kommer til Klarhed over Principerne, \emph{forud
for Erfaringen} bestemt ved disse: \emph{Naturopfattelsen
bliver metaphysisk, ikke physisk}. Den græske Philo\-%
sophie kjendte derfor heller ikke den Kunst, gjennem \emph{Expe\-%
rimentet} at adspørge Naturen; thi ved en saadan Adspørgen
% p. 125 l. 19: "Uuafhængighed" AS PRINTED -- a doubled u, where
% "Uafhængighed" is expected.  Read the same way by both OCR witnesses
% and confirmed on the image at 1200 dpi.  NOT on the Trykfeil leaf.
vilde denne netop anerkjendes i sin relative Uuafhængighed.
Først sildig og i ringe Omfang kommer Experimentet frem
hos Grækerne, og da udenfor Philosophien.

Denne Naturerkjendelses Art gjenkjende vi nu ogsaa hos
\emph{Aristoteles} trods hans Stræben efter en Erfaringserkjendelse
og trods hans udstrakte Viden, om hvilken ogsaa rent empi\-%
% p. 125 ll. 25--26: the letterspacing begins INSIDE the word -- "Na-"
% at the end of l. 25 is set close and "turerkjendelse" is spaced, the
% same split within one word as at p. 114 ("Naturbestemthed") and at
% pp. 91 and 98.  "væsentligen" on l. 26 is NOT spaced.
riske Skrifter af ham gave Vidnesbyrd. Ogsaa om hans Na\-%
\emph{turerkjendelse} gjælder det, at den væsentligen har en \emph{me\-%
taphysisk, ikke en physisk} Charakteer. De \emph{metaphy\-%
siske} Grundbegreber: Form og Stof, anvendes overalt i Phy\-%
siken \emph{som} metaphysiske Begreber uden at gjøres concretere
ved Forholdet til det Physiske; Anskuelsen af Universet som
Heelhed bestemmes \emph{umiddelbart} ved dem, ikke ved en Sam\-%
menfatten af det begrebne Reale; det Materielles Grundformers
Natur bestemmes ikke gjennem Iagttagelsen, men gjennem en
\emph{apriorisk Deduction}, og Tanken drages uvilkaarlig hen
mod de Naturgebeter, hvor det \emph{metaphysiske} Grundforhold
% --- p. 126 ---
% p. 126 bears NO footnote: the text block runs to the last line and
% there is no separator rule at the foot, so nothing can run over onto
% p. 127.  Checked on the image.
faaer den \emph{letteste} Tilknytning: til det \emph{Organiskes} Gebet,
og de lavere Natursphærer sees ligesom kun i et Reflexlys fra
det Organiske, medens i selve det Organiske vel de mechaniske
Aarsager fremhæves, men uden at de hævde sig en \emph{begrebs\-%
bestemt positiv} Betydning i Forholdet til det Teleologiske.
Og det, at den aristoteliske Naturerkjendelse faaer denne Cha\-%
rakteer, kan vel tildeels være betinget ved hans Tids Mangel
paa de talrige techniske og videnskabelige Hjælpemidler til
Iagttagelse og Beregning, over hvilket en nyere Tids Natur\-%
forskning raader, og ved dens fragmentariske Kundskab til
Naturphænomenerne; men det har, hvad \emph{alle} den aristoteliske
Philosophies Hovedbestemmelser indirecte vise, sin \emph{sidste} og
\emph{væsentlige} Grund i \emph{den græske Naturopfattelses al\-%
mindelige} Charakteer, der \emph{selv} bliver medbetingende for
hine Mangler.

% p. 126 ll. 17--18: "medens den stræber ud / over den græske
% Naturanskuelses Grændser" is set CLOSE here, though the parallel
% phrase on p. 120 is letterspaced.  Checked at 1200 dpi; not
% regularised.
I det Udviklede ligger Nødvendigheden af, at \emph{Aristote\-%
les's almene Naturanskuelse}, medens den stræber ud
over den græske Naturanskuelses Grændser, dog ikke mægter,
virkelig at overskride disse, saa at dens \emph{kategoriske Be\-%
stemthed} bliver den \emph{samme} som dennes. Den aristoteliske
Naturanskuelse er Udtrykket for den græske Naturanskuelse i
dens \emph{høieste} Udvikling og ligesom i \emph{dens Overgang} til en
ny og \emph{høiere} Form. \emph{Skjønhedsanskuelsen} af Universet,
Anskuelsen af den reale Tilværelse som Verdenskunstværk, som
\gk{Κόσμος}, er bevaret hos \emph{Aristoteles}, og Formen bliver, skjøndt
den i sin Idealitet ingen Rumsbestemthed har, dog i Naturen
Rumsform. Men \emph{paa Skjønhedsanskuelsens Grundlag}
er en rigere, dybere, livfuldere Opfattelse ifærd med at arbeide
sig frem: \emph{Livets Idee} er ifærd med at afløse \emph{Skjønhedens
Idee} som den for den reale Tilværelse altbestemmende. Na\-%
turen, der af Aristoteles bestemmes som \emph{i sin Heelhed be\-%
sjælet}, og det ikke blot ved en uvilkaarlig Anthropomorphi\-%
seren, men ifølge Begrebsforholdet til den iboende virksomme
% p. 126, last line: a faint pencil tick stands between "vises" and
% "der"; both OCR witnesses pick it up ("vises 'der"). Not transcribed.
%
% SEAM p. 126 --> p. 127: p. 126 ends MID-WORD with "Be-" and p. 127
% opens "greberne om Formen som Udviklingsprincip".  The pp. 127--140
% batch must open with "greberne" and must NOT insert a paragraph break
% here.  p. 127's note area is empty at the top, so nothing runs over.
Form, er ved denne Opfattelse bragt ind under det \emph{Organi\-%
skes} Synspunkt, og til det Organiske \emph{vises} der hen ved Be\-%
% --- p. 127 ---
% SEAM CHECK, p. 126 -> p. 127: p. 126 bears NO footnote (no footnote separator
% on the page; verified by a rule sweep of the image), so nothing runs over into
% p. 127.  p. 126 ends "vises der hen ved Be-" and this page opens with
% "greberne"; the pp. 113--126 batch owns that hyphen.
% p. 127 bears no note.  Greek: δύναμις, Κόσμος.
greberne om Formen som \emph{Udviklingsprincip,} og om den
\emph{indre Teleologie i det Enkelte og i den hele Natur.}
Den \emph{organiske Naturopfattelse,} der, \emph{consequent} gjen\-%
nemført i \emph{Tankens} Form, ifølge \emph{Væsensforholdet} mellem
det Levendes \emph{Idealitetsside} og \emph{Realitetsside,} indeslutter
den sidstes \emph{positive} Betydning og dermed \emph{Realerkjen\-%
delsens Betydning,} er hos Aristoteles tilstede \emph{i en levende
Anskuelse,} og en moderne Tids Opfattelse er \emph{berørt.} Men
de ovenfor paaviste Mangler hos Aristoteles gjøre sig gjæl\-%
dende, saasnart den formbestemte Udvikling skal gjøres klar
for Tanken, og saasnart Formen, det \emph{blot} Ideelle og Evige,
der som saadant mangler Realitetsbestemmelsen, skal sees som
Naturenergie. Ligesom Bestemmelsen af den energiske Form\-%
virksomhed ikke kunde fastholdes, men uvilkaarligt gik tilbage
til Opfattelsen af Formen som den \emph{evigt færdige,} af Vorden
uberørte, \emph{hvilende Form,} der ikke \emph{væsentlig} sondrede sig
% p. 127: δύναμις -- acute over the υ, read at 1200 dpi; not letterspaced
% (the surrounding "med dets postulerede" is unspaced).
fra det \emph{platoniske} hvilende Urbillede med dets postulerede
\gk{δύναμις}, saaledes maa Anskuelsen af det \emph{levende, orga\-%
niske Universum,} naar Udviklingsprincipet tilintetgjøres af
% p. 127: Κόσμος -- capital Κ, acute over the first ο, read at 1200 dpi.  It is
% ALSO letterspaced (its letter gaps match the spaced roman around it and are
% plainly wider than the tightly-joined cursive of δύναμις eight lines above),
% so it stands inside the surrounding \emph{} run.
Tanken, atter gaae tilbage til Anskuelsen af \emph{den i Skjøn\-%
hedens Umiddelbarhed hvilende kunstnerisk for\-%
mede \gk{Κόσμος}}. Og til denne svarer hiin \emph{Naturerkjen\-%
delse,} der umiddelbart lader Stoffet gjennemlyse af Formen,
% p. 127: "æsthetisk-metaphysiske" is a compound whose own hyphen happens to
% fall at the line end; it is not a word-break hyphen and is kept as printed.
gaae op i Skjønhedsenheden med denne: \emph{den æsthetisk-metaphysiske Naturerkjendelse.}

% DIVISIONAL RULE, p. 127: short, CENTRED hairline with white space above and
% below, measured on the image at x 517--770 (300 dpi) against a text block of
% 85--1211 -- rule centre 643.5, block centre 648.0 -- and 2.15 cm long.  p. 127
% bears no footnote, so this cannot be the (left-flush) footnote separator.
\divrule

% HEAD, p. 127 (L3, transcription.tex §3): the third of the 1./2./3. heads under
% thesis II (80, 94, 127).  IT IS NOT BOLD: it is LETTERSPACED roman at body
% size, set as a justified full-measure paragraph with a first-line indent, over
% three lines -- the same treatment as the p. 94 head.  The print's line breaks
% are: 3. ... antike / Verdensanskuelse ... ubevidste / Indvirkning.
\emph{3. Nyplatonismens Stræben ud over den antike
Verdensanskuelse hæmmet ved dennes ubevidste
Indvirkning.}

I Oversigten over den græske Philosophies Udvikling er
Nyplatonismen paaviist at være denne Hovedepoches Afslut\-%
ning. Den bliver nærmere at bestemme som Slutningsleddet
i den Række af philosophiske Forsøg, der betegne den græske
Tænknings Opløsning, idet de, i deres Søgen efter Sandhed,
% --- p. 128 ---
% p. 128 bears no note, no Greek and no rule.
under Tvivlen om den ved Dualismen vildsomme menneske\-%
lige Erkjendelses Magt, flygte ud over Tankens og Videnska\-%
bens Grændser til en guddommelig Aabenbaring, og netop
fordi de forlade Tankens Sphære: den concrete Virkelighed,
som et Usandt, bestemme dette Guddommelige, der gjennem
Aabenbaringen bliver Sandhedens og Lyksalighedens Kilde,
som et absolut Hiinsidigt, der i sin rene Aandelighed er util\-%
nærmeligt for Erkjendelsen, og kun er tilgængeligt idet Aanden
gjennem den ekstatiske bevidstløse Selvhengivelse vinder
en umiddelbar Forening med Urvæsenet. Som det \emph{sidste}
Led i denne Udviklingsrække viser Nyplatonismen sig ved sin
principielle Spiritualisme og ved den Gjennemførelse af det
Guddommeliges Transscendents, ifølge hvilken det bestemmes
% p. 128: the "og" between "al Tænken" and "al Væren," is NOT letterspaced
% (its letters touch, against a clear gap in every letter of "al"), so the
% emphasis is set as two runs, as printed, and not as one.
som staaende over \emph{al Tænken} og \emph{al Væren,} og som det,
% p. 128: "Meddelelse," is NOT letterspaced -- checked at 600 dpi against the
% adjacent "dynamisk"; the space before its comma is the fount's ordinary
% comma side-bearing after a round letter, not Sperrsatz.
der gjennem \emph{en blot dynamisk} Meddelelse, i hvis \emph{quan\-%
titative} Aftagen Materien danner det sidste Led, er Kilden
til al bestemt Væren. Selv hos \emph{Philo,} den betydeligste af
de Tænkere, der banede Veien for Nyplatonismen, er der endnu
en kjendelig Afstand fra dette Maal. Han tilnærmer sig Ny\-%
platonismens Standpunkt ved sin Lære om det Guddommeliges
Transscendents, om Logos og om den ekstatiske Anskuelse af
det Guddommelige, der er hans Systems egentlige Maal, og som
er Udtrykket for den Længsel efter en umiddelbar Forening med
det Guddommelige, der var fremtvungen ved Følelsen af det
Endeliges Intethed og støttet ved Troen paa, fra den oversandselige
Verden at kunne faae denne Trang tilfredsstillet, som Menne\-%
% p. 128: ABBYY reads "den endelige. Verden" and tesseract "den endelige,
% Verden".  There IS a small round blob at the baseline after "endelige", but
% it is NOT type: its darkest pixel measures grey 90 against grey 29 for the
% adjacent letters of the same word at 1200 dpi, i.e. it carries none of the
% impression a printed sort carries.  Read as a speck, and no punctuation is
% set here; the sense ("som Menneskeaandens egen Kraft og den endelige Verden
% ikke kunne skaffe Udfyldelse") wants none.
skeaandens egen Kraft og den endelige Verden ikke kunne
skaffe Udfyldelse. Men han skiller sig afgjort fra Nyplatonis\-%
men derved, at han endnu er paavirket af den jødiske Mono\-%
theismes Forestilling om Guds Personlighed og endnu ikke vil
sætte Gud \emph{over det Værendes Begreb;} at han opfatter
Materien i en \emph{dualistisk} Modsætning til det Guddommelige,
og at Urvæsenets Meddelelse til Verden, der hos Plotinos kun
bestemmes som dynamisk, som en Kraftmeddelelse, hos Philo
% --- p. 129 ---
\setcounter{footnote}{0}
% p. 129 bears ONE note, self-contained (two lines, ending in a full stop);
% nothing runs over onto p. 130.  Greek: Νοῦς (body and note), τὸ ἁπλοῦν,
% ἐπέκεινα τῆς οὐσίας.
bliver noget Mere idet Mellemvæsenerne synes at staae i
et substantielt Enhedsforhold til Urvæsenet.

Nyplatonismens eiendommelige Charakteermærke er saa\-%
ledes en paa Anskuelsen af det \emph{uendelige} Urvæsens \emph{abso\-%
lute Transscendents} hvilende \emph{monistisk} Spiritualisme.
Spliden mellem Tilværelsens Principer, der var bleven aaben\-%
bar idet det aristoteliske System endte i Dualisme, søges hævet
ved at føre Alt tilbage til det rent Aandelige\footnote{Til Materiens Udledning af Urvæsenet danner det en Analogie, at
Plotin i \gk{Νοῦς} udleder Væren af Tanken.}: Endelighedens
Splid og Tankens Tvivl søges forsonet ved at vise hen til et
% p. 129: "al" IS letterspaced here and stands alone -- its a--l gap measures
% 36 px at 1200 dpi against 12--16 px inside "alle" nine words later on the
% same line, and against a 4--24 px range inside every other word of the line.
% spacing.py does not flag it.
uendeligt Urvæsen, ophøiet over \emph{al} Endelighed: over alle Væ\-%
% p. 129: Νοῦς -- circumflex over the υ, read at 1200 dpi.
rens Bestemmelser, over hele Tankens Verden, selv over \gk{Νοῦς},
og ved at angive Veien, ad hvilken Mennesket kan naae tilbage
til dette Guddommelige og i dets Dyb finde Hvile. Det er
den græske Tankes sidste Stræben; den stræber ikke blot ud
over sin Skranke som \emph{græsk} Tanke, idet den stræber at hæve
Materiens eiendommelige, oprindelige Væren, at opløse Skjøn\-%
hedens Enhed til Fordeel for dens ideelle Element, at opgive
Grændsen og Maalet til Fordeel for det Ubegrændsede, Be\-%
stemmelsesløse; men den stræber ud over \emph{selve Tankens}
% p. 129: τὸ carries a GRAVE (the stroke descends left-to-right, read at
% 1200 dpi); ἁπλοῦν carries a ROUGH breathing over the α (a "C" opening to the
% right, unambiguous at 1200 dpi) and a circumflex over the ου.
Sphære, idet den, søgende det Ubegrændsede, det Enkelte (\gk{τὸ
ἁπλοῦν}), der er hiinsides Væren (\gk{ἐπέκεινα τῆς οὐσίας}) og
% p. 129: ἐπέκεινα -- SMOOTH breathing over the first ε (a comma opening to the
% right) with an acute over the second; this is one of the two calibration
% words for the smooth/rough distinction (transcription.tex §2).  τῆς has a
% circumflex; οὐσίας a smooth breathing over the υ and an acute over the ι.
% Neither OCR witness has a Greek character here: ABBYY gives "{snéxsivct zrjg
% ovdiag)" and tesseract "(åréseswe vijc odoiag)".
Tanken, søger at forvandle sig selv til \emph{et Høiere end Tan\-%
ken} (jfr. Damascius).

% p. 129: "men," at the head of the next line is NOT letterspaced (m--e 16 px,
% e--n 8 px at 1200 dpi, against 20--32 px in the certainly spaced "almindelige
% Charakteer" on the line above); spacing.py's single-word flag is rejected.
I sin \emph{almindelige Charakteer} frembyder Nyplatonis\-%
men, netop i det, der adskiller den fra de andre Former af
den græske Philosophie, en væsentlig Lighed med den Verdens\-%
og Livsanskuelse, der afløste Oldtidens: med \emph{Christendom\-%
mens,} og gjennem dens Forhold til denne vil dens Forhold
til den \emph{græske} Philosophie og til den \emph{græske} Aandsretning
overhovedet blive tydelig. Nyplatonismen og Christendommen
stemme overens deri, at begges Grundanskuelse er transscendent
og spiritualistisk, og deri, at den praktiske Tendents i begge
% --- p. 130 ---
% p. 130 bears no note, no Greek and no rule.
viser hen i samme Retning. De to store historiske Phæno\-%
mener vise sig ved første Øiekast at være fremgaaede af en
og samme Retning i Tiden. Det praktiske Maal, som de begge
have, og med hvilket i Nyplatonismen det theoretiske umid\-%
delbart er forenet, er det med sig selv og med Gud splidagtige
Menneskes Forening med det Guddommelige, betinget ved Be\-%
frielsen fra Sandseligheden og dens Besmittelse, ved Livets
Fremme gjennem den høiere Indvirkning. Følelsen af at være
bleven fremmed for Gud, Utilfredsstilletheden i den givne sandselige
Virkelighed, og Længselen efter en høiere Bistand og en høiere
Aabenbaring er charakteristisk for Oldtidens Opløsningsperiode.
I den mødes den græske og den jødiske Aand, der, idet de
indgaae en Forbindelse paa Judaismens religiøse Basis, blive
Christendommens historiske Factorer. I den græske Philo\-%
sophie stod den praktiske Stræbens Eiendommelighed i nøieste
Forbindelse med Tankens Udvikling og begyndende Opløsning.
Fra det Øieblik af, at den græske Tænkning fortvivler \emph{om
Viden,} har den med det Samme opgivet det \emph{ethiske} Til\-%
hold; thi Viden var overalt den græske Ethiks Grundlag. Den
\emph{theoretiske Skepsis} maatte, hos det Folk, der satte Men\-%
neskets \emph{Væsen} i Theorien, blive til \emph{ethisk Fortvivlelse} ---
ogsaa \emph{Skepticismens} Forhold til det Praktiske er væsentlig
seet Fortvivlelse --- og Trangen til en \emph{Aabenbaring} og til
en høiere \emph{sædelig} Bistand maatte umiddelbart falde sammen.

Mod dette høiere, fra det aandelige Guddommelige Udgaa\-%
ende stræbtes der overalt hen ved den gamle Verdens Under\-%
gang. Men det Motiv, der blev denne Stræbens bevægende
Kraft, var et forskjelligt. Naar der stræbtes hen mod et
\emph{Høiere,} saa stræbte \emph{Orientaleren} derefter som efter Kilden
til en \emph{høiere Magt;} \emph{Grækeren} stræbte derefter som efter
Kilden til en \emph{høiere Viden;} den \emph{Christne} stræbte derefter
som efter Kilden til et \emph{høiere Liv.}

Heri ligger Antydningen af, at \emph{Nyplatonismen,} den
græske Tænknings sidste Repræsentant, og \emph{Christendom\-%
men,} trods den almindelige Lighed og det fælleds Udspring,
% --- p. 131 ---
% p. 131 bears no note and no rule.  Greek: καινὴ κτίσις.
dog kunde fjerne sig saameget fra hinanden, at netop de kunde
staae som \emph{stærke Modsætninger,} og at Nyplatonikerne, der
gjennem Udviklingen af en universel Philosophie vilde restau\-%
rere Oldtidens Religion som universel Religion, kunde være
Hedenskabets meest haardnakkede Forkæmpere i den Kamp
paa Liv og Død, det førte med den nye Religion. Og idet i
Nyplatonismens \emph{Modsætning} til det Christelige den \emph{græske
Aands Grundcharakteer} afspeiler sig, vil det, at denne gjør
sig gjældende trods det Fremmedartede i Nyplatonismens \emph{be\-%
vidste} Stræben, falde sammen med at den udøver en hæm\-%
mende Indvirken paa denne Stræben, ved hvilken Nyplatonis\-%
men søgte hen imod en consequent spiritualistisk Verdensan\-%
skuelse.

Vi fandt en iøinefaldende \emph{Lighed} mellem Nyplatonismen
og Christendommen i begges abstract-spiritualistiske Grund\-%
præg og i begges Retning mod Forsoningen ved det overna\-%
turlige, uendelige Guddommelige. Men ved Siden af Ligheden
træder en afgjørende og overveiende \emph{Modsætning} frem ved
\emph{Maaden,} paa hvilken denne Forsoning tilstræbes, og med det
Samme bliver en Forskjel i Spiritualismens \emph{Art} aabenbar.
\emph{Nyplatonismen} stræber at naae Forsoningen gjennem Er\-%
kjendelsen, gjennem en, om end ved en guddommelig Virkning
fremmet og ligesom op over sig selv løftet, philosophisk Spe\-%
culation, og at hæve Mennesket, der glemmer sin Endelighed,
til en overmenneskelig Enhed med Guddommen, til en, den
skjulte Egoisme tilfredsstillende, guddommelig Lyksalighed;
\emph{Christendommen} vil naae den gjennem Troen, og ved at
lade Gud stige ned i den menneskelige Svagheds Dybder forat
overtage den Skyld, fra hvilken Mennesket ikke kan flygte.
\emph{Nyplatonismen} lader Menneskenaturen, som skal guddom\-%
meliggjøres, i sit inderste Væsen være uforstyrret trods al Splid
(jfr. \emph{Plotinos's} Udsagn) og fordrer derfor ingen Selvets For\-%
vandling i dets Inderste; \emph{Christendommen} gaaer ud fra
Forudsætningen af den menneskelige Naturs Fordærvethed ved
% p. 131: καινὴ carries a GRAVE over the η and κτίσις an ACUTE over the ι,
% both read at 1200 dpi; the pair is NOT letterspaced (ordinary cursive
% spacing).  Neither witness has Greek here (ABBYY "xcarrj xtiaig",
% tesseract "xa xtiog").
Synden, og fordrer derfor en \gk{καινὴ κτίσις}. \emph{Nyplatonismen}
% p. 131: the sheet signature "9*" stands at the foot; not transcribed.
% --- p. 132 ---
% p. 132 bears no note and no rule.  Greek: Κόσμος's.
har for sin høiere Viden et Egoismens Maal, idet den, paavir\-%
ket af et østerlandsk Element, opfatter den som Kilden til en
magisk Magt, til et Herredømme over Tilværelsen (Guddom\-%
men bliver, gjennem Theurgien, Menneskets Tjener); \emph{Chri\-%
stendommen} søger gjennem Troen et høiere Liv ved Af\-%
% p. 132: Κόσμος's -- the Danish genitive apostrophe + s is set in roman after
% the Greek, as printed; the Greek is not letterspaced here.
døetheden fra Verden og fra Selvet. \emph{Nyplatonismen} drages,
i sin Flugt mod det Hiinsidige, i sin Higen mod den ekstatiske
Forening med Urgrunden i Verdens- og Selvforglemmelse,
uvilkaarligt tilbage til Verden, den priser begeistret \gk{Κόσμος}'s
Herlighed og kæmper mod Christendommens Verdensfornæg\-%
telse; \emph{Christendommen} seer i den skabte Natur det, der i
sit Væsen er et Intet og skal sættes som et Intet, og som,
naar det vil gjøre sig gjældende som et Noget, gaaer ind un\-%
der det Ondes Bestemmelse, under Bestemmelsen af det, der
ikke skal være. \emph{Nyplatonismen} seer endelig Maalet som
naaet i denne Tilværelse; \emph{Christendommen} seer det først
virkeliggjort i en tilkommende, i hvilken Alt er blevet Aand,
Naturen er forsvunden.

I den Modsætning, som disse Differentspunkter betegne,
røber sig den antike Hellenismes Indvirkning, og den røber
sig som begrændsende Nyplatonismens bevidste Stræben. Det
er en Følge af denne Indvirkning, at trods den Tvivl om Tan\-%
kens og Virkelighedens Sandhed, der bringer Aanden til at
søge et Tilhold i et spiritualistisk Guddommeligt, \emph{Menneske\-%
sjælen} dog bliver sat som det i sin Væsens Grund Uforstyr\-%
rede, altsaa bliver sat under den Bestemmelse, der følger af
Opfattelsen af Tilværelsen som et harmonisk Kunstværk, af
Ideelærens Opfattelse af Evigheden som Livets Baggrund, der
aldrig afspærres for Erindringen. Det er en Følge af den, at,
trods den Accent, der falder paa \emph{Menneskeaanden,} der kan
naae den umiddelbare Enhed med det Høieste, --- hvori det lig\-%
ger antydet, der udtales i Christendommen: at den individuelle
Sjæl med dens indre Uendelighed har større Værd end hele
Naturtilværelsen, --- Nyplatonismen dog lader Menneskesjælen
træde tilbage for de kosmiske Legemers Livsprinciper, for
% --- p. 133 ---
\setcounter{footnote}{0}
% p. 133 bears ONE note, self-contained (a single line ending in a full stop);
% recon had flagged it as a run-over candidate and cleared it, and the image
% confirms the clearance.  Greek: Κόσμος.
Stjernernes Sjæle, under hvilke den underordnes ligesom den
enkelte Personlighed i den græske Ethik underordnedes under
det upersonlige Almene. Det er en Følge af den, at, trods
den endelige Virkeligheds Splidagtighed og Uforsonethed, der
væsentlig har sin Kilde i det Materielles Modstand, dog
Anskuelsen af \gk{Κόσμος}, i hvilken Materiens harmoniske
Underordning er en væsentlig Bestemmelse, hævder sig Gyl\-%
dighed, og det skjønne Verdensalt bliver Gjenstand for tilbe\-%
dende Begeistring. Det er endelig en Følge af den, at trods
Nyplatonismens Forsøg paa at lade \emph{Materien} forsvinde som
% p. 133: the note's own text carries letterspacing ("forudsætter"), as §3 of
% BATCH-AGENT records that it does.
et Selvstændigt, dette Forsøg mislykkes, ikke blot derved, at
der ved Udledelsen af Materien\footnote{I Regelen \emph{forudsætter} Plotinos Materien.} fra det i sig afsluttede Ab\-%
solute gjennem en \emph{blot quantitativ} Aftagen af Kraftmed\-%
delelsen, indirecte tillægges den en Væsentlighed, men ogsaa
derved, at den, medens den skal udledes paa denne Maade,
bestemmes som et fra det Aandelige \emph{qualitativt} Forskjelligt,
og at der tillægges den en den tilhørende Virksomhed, ved
hvilke Bestemmelser Materien, paa den gamle græske Philoso\-%
phies Viis, sættes som et Selvstændigt og Uafledet, som For\-%
mens oprindelige, uafhængige Supplement, saa at den græske
Philosophie ved sin Slutning paany aabenbarer den \emph{Dualisme,}
der skjuler sig under dens Grundanskuelses harmoniske Enhed.

Og naar det her Udviklede viser os en hæmmende Ind\-%
virkning, ved hvilken Nyplatonismens Stræben hen imod det
\emph{Transscendente} og mod en ude over den græske Tænknings
Skranker liggende \emph{spiritualistisk Monisme} bliver forhin\-%
dret i en consequent Gjennemførelse, saa see vi en lignende
Indvirkning forhindre den fra at fastholde det \emph{ugræske Be\-%
greb} om det \emph{Absolutes Uendelighed.} Allerede i det spi\-%
ritualistiske Absolutes \emph{Begrændsning ved en uafledet
Materie} ligger Uendelighedens Ophævelse; men den viser sig
ogsaa fra en \emph{anden} Side. \emph{Uendeligheden} bliver for Ny\-%
platonismen, ifølge den græske Tænknings Charakteer, \emph{ikke}
% --- p. 134 ---
% p. 134 bears no note, no Greek and no rule.
% p. 134: the letterspaced "ikke" at the foot of p. 133 closes there; "den paa
% den" at the head of this page is unspaced, so the runs are separate.
den paa den \emph{uendelige Negativitet} hvilende \emph{concrete}
Uendelighed, men \emph{Bestemmelsesløshedens Uendelighed.}
Derfor stræber Tanken at negere \emph{enhver} Bestemmelse af det
kun i Ekstasens Berusning Tilnærmelige, men idet det skal
sees \emph{som Forklaringsgrund for Tilværelsen,} trænge
\emph{Aarsagens} og \emph{Kraftens} Bestemmelser sig ind, og de trænge
sig ind som \emph{endeliggjørende} Bestemmelser, ikke blot fordi
\emph{enhver} Bestemmelse \emph{som} Bestemmelse efter Forudsætningen
er \emph{endeliggjørende,} men fordi de sætte et \emph{Forhold} til
\emph{det, der skiller sig fra} det i sig afsluttede Absolute i
Væren.

Den modsigende Dobbelthed, vi have fundet i Nyplato\-%
nismens Væsen, og som har sin Rod i Modsætningen mellem
det bevidst Tilstræbte, med den græske Aands Væsen Uensartede
og denne Aands ubevidst indvirkende Grundeiendommelighed,
forhindrer den græske Aand fra, nogensinde virkelig at komme
til Ro i denne Tankens Form. Aandens indre Uro viser sig
ikke blot i den krampagtige Famlen, med hvilken der i Ny\-%
platonismen gribes efter Elementer fra alle philosophiske Sy\-%
stemer og fra alle Religioner, uden at Aandens Trang mere
tilfredsstilles derved, end den religiøse Aands Trang tilfreds\-%
stilledes da den ved Oldtidens Slutning fyldte et Pantheon
med døde Guder, men den finder sit afgjørende Udtryk i
\emph{Tankens,} hiin Dobbelthed afspeilende, \emph{Dobbeltbevægelse}
og i dens \emph{Selvopløsning.}

Aandens Stræben er, at naae ud \emph{over Virkeligheden}
og \emph{ud over Viden;} men idet den \emph{græske} Aand ikke kan
% p. 134: "Ro" and "Tro," are NOT letterspaced -- checked at 600 dpi against
% "Videns Modsætning," on the same line, which is.
slaae sig til Ro i en Tro, der er \emph{Videns Modsætning,} idet
den, selv hvor den higer ud over Tanken, i det hiinsides
Tanken Liggende kun faaer den ved Tanken satte yderste Ne\-%
gation af den bestemte Viden, ikke en positiv Modsætning til
Tanken, saa bliver den bunden \emph{indenfor Tankens Kreds,}
og Aanden, der ikke har en positiv Gjenstand udenfor denne,
faaer den Opgave, \emph{at ordne den af det utænkelige Ab\-%
solute udledte Tilværelse efter Tankebestemmelser.}
% --- p. 135 ---
\setcounter{footnote}{0}
% p. 135 bears ONE note only (a three-line note ending in a full stop), not the
% two or three BATCH-AGENT §3 lists for this page; nothing runs over onto
% p. 136.  Recon had flagged p. 135 as a run-over candidate and cleared it, and
% the image confirms the clearance.  No Greek.
Men i det Øieblik, Nyplatonikernes Tanke drages hen til den
\emph{concrete} Videns Sphære, er den ført tilbage til \emph{det Natur\-%
grundlag,} der er den \emph{græske} Tankes, og den er, som fastholdt
ved dette, der \emph{ikke} har sin Form fra det formløse Uendelige,
forhindret i at absorberes af det abstract Spiritualistiske. Dette
Sidste kommer frem som Tankens Grændse, som Systemets
Afslutning, medens paa alle Systemets øvrige Punkter den
\emph{græske} Tankes gamle Indhold, om end seet som i en ny
Belysning, bliver det Udfyldende, der bærer Tanken. \emph{Ver\-%
densvirkeligheden} drager Aanden til sig, men idet hint
Afsluttende viser sig \emph{bag} den, tilfredsstiller den dog ikke den
Aands \emph{høieste} Trang, der higer mod \emph{det ved selve Vir\-%
keligheden antydede Uendelige.} Aanden kaster sig da
fra Virkeligheden over i \emph{det Absolute,} men kun for atter
fra dets \emph{Tomhed} at vende tilbage til den \emph{transfigurerede
Virkelighed,} der i sin \emph{dobbelte} Væren paany vækker
Aandens Uro og kaster den svimmel tilbage i den samme
rastløse Vexlen.

Og i \emph{denne Vexlen opløser Tanken sig selv.} Idet
den stræber ud over sig selv til \emph{et Høiere end Tanken}
(Damascius), synker den i samme Øieblik ned \emph{under Tanken}
og bliver virkelighedsløs \emph{Phantasie}\footnote{Man erindre de phantastiske Elementer i Nyplatonismen: Magien,
Theurgien, Overtroen. Selv i Opfattelsen af Universets Sympathie
viser det Phantastiske sig.}. For \emph{Phantasiens
Øie} viser sig i \emph{Ekstasen} et blændende Billede, et glimrende
Lys, der, naar Ekstasens flygtige Moment forsvinder, taber sig
i Tomhedens uendelige Mørke. Idet den erkjendende Aand
vil \emph{fæstne} det i Berusningens Øieblik Skuede, nødes den til
at bruge \emph{Tankens} Former, men kun for i samme Øieblik at
% PRINTER'S DEFECT (letterspacing), p. 135: in "Tankebestemmelser" only the
% first element is letterspaced -- "Tanke" measures 10--18 px between letters at
% 600 dpi and "bestemmelser," 3--7 px, i.e. the compositor spaced half a word.
% Set as printed.  Note also that "uophørlig" is NOT spaced (only "maa" is),
% against what a reading of the 300 dpi image suggests.
\emph{negere} dem. Forat \emph{fastholde} det Absolute, \emph{maa} der uop\-%
hørlig anvendes \emph{Tanke}bestemmelser, men dets \emph{negative}
Uendelighed gjennembryder dem i \emph{samme} Øieblik, de an\-%
vendes, og Resultatet bliver \emph{det Bestemmelsesløse,} der
% --- p. 136 ---
\setcounter{footnote}{0}
% p. 136 bears ONE note (two lines, self-contained).  Greek: ὑπεράγνοια.
% NOTE ON PLACEMENT: the note's call stands ABOVE the divisional rule, the note
% itself below it at the foot of the page.
saa paany kan faae en illusorisk Væren, en bedragerisk Fylde,
for \emph{Phantasien,} men kun for paany at hjemfalde til \emph{den
negerende Tanke,} der ligesom indhenter den Phantasie,
der vil overflyve den, og forstyrrer dens Billeder. Tankens
Virksomhed bliver, idet den i Forholdet til det \emph{Absolute,}
ved at ville gaae \emph{ud over sig selv,} aabenbarer Tvivlen om
sin Evne i Forhold til det Principielle, væsentlig \emph{Skepsis.}
\emph{Uendeligheden,} der kun bliver den \emph{tomme, formløse}
Uendelighed, sees som det, der \emph{forstyrres ved Tankens
Former,} og Tanken maa derfor \emph{skeptisk} tilintetgjøre sine
egne Bestemmelser. Den \emph{Modsigelse,} at Urvæsenet \emph{kun}
bliver til for os ved de \emph{øieblikkelig tilbagetagne} Tanke\-%
bestemmelser\footnote{Jfr. Exemplerne hos Zeller: Philosophie der Griechen, III, 955 ff.
(1ste Udgave).}, at der skal udsiges Bestemmelser ved \emph{det,}
der \emph{ingen} Bestemmelser maa gives, fremtvinger saa, hos de
sidste Nyplatonikere, Klagerne over det menneskelige Sprogs
% p. 136: ὑπεράγνοια -- the mark over the υ reads at 1200 dpi as a "C" opening
% to the right, i.e. a ROUGH breathing, and there is an acute over the α; the
% breathing is a small mark and the reading is less secure than the ἁπλοῦν and
% ἐπέκεινα of p. 129, so it is recorded here.  ABBYY gives "vnsgdyvoia" and
% tesseract "ømæeodyvosn"; neither has a Greek character.
Utilstrækkelighed og Bekjendelsen af en \gk{ὑπεράγνοια} (Dama\-%
scius), af \emph{en ved Tankens Selvtilintetgjørelse poten\-%
seret Uvidenhed.}

% DIVISIONAL RULE, p. 136: short, CENTRED hairline with white space above and
% below, measured on the image at x 678--928 (300 dpi) against a text block of
% 248--1357 -- rule centre 803.0, block centre 802.5 -- and 2.12 cm long.  It is
% NOT the footnote separator, which on this page is left-flush and stands lower,
% at x 351--496.  This rule closes the three examples under thesis II; the
% paragraph that follows it is the summing-up of the whole thesis, and it is the
% only rule anywhere near the III. head of p. 137, which itself carries none
% (see the comment there).
\divrule

% p. 136: a modern PENCIL mark -- a small asterisk-like blob -- stands above
% "dens" in "medens med det Samme dens Begrændsethed"; ABBYY splices it into the
% text as "dens'’".  Not transcribed.
I de udviklede Exempler har saaledes den \emph{bevidste}
Stræben viist sig, idet den inciteres ved de af Grundtankens
indre Udfoldelse fremgaaede Problemer, at være bestemt ved dens
\emph{indre Universalitet,} der bliver \emph{det fremad Bevægende,}
medens med det Samme dens \emph{Begrændsethed} som \emph{dette
bestemte Moment} i Videns Idee bliver det, der med uimod\-%
staaelig Magt \emph{hæmmer} hiin Stræben. Ved denne dobbelte
Virkning realiserer Grundtanken sin \emph{fuldstændige} Udfoldelse
og forbereder en \emph{ny} Udviklingsrække.
% --- p. 137 ---
% p. 137 bears no note.  No Greek.  The printed folio reads 137 plainly on the
% image; ABBYY garbles it as "mmmmmmmmmmsBMmmmmmm".
%
% HEAD, p. 137 (L2, transcription.tex §3): the THIRD of the book's four theses
% (I. 50, II. 79, III. 137, IV. 170).
%
% BOLD, NOT LETTERSPACED -- this is the calibration page, and the numbers are:
% measured at 600 dpi, the inter-letter gaps in the head's first line run
% 2--6 px, exactly the 1--7 px of the ordinary body line "Vi oplyse dette
% Forhold ved følgende Exempler:" four lines below it; the letterspaced L3 head
% at the foot of the SAME page runs 8--22 px (median ~14) on its first line and
% 10--22 px on its second.  spacing.py reports the L2 head as six letterspacing
% RUNs; that is the known false positive of BATCH-AGENT §3 -- bold letters are
% wider than its per-glyph fit predicts -- and the head is therefore NOT set
% \emph{}.
%
% NOT CENTRED.  The head is set FULL MEASURE and justified over seven lines:
% every line runs x 112--1224 at 300 dpi, the same measure as the body, and
% three of them break a word with a hyphen (Sta-/diums, Eien-/dommelighed,
% In-/citament), which a centred block would not do.  REVIEW ITEM: the p. 79
% head, measured the same way, is also full measure (x 92--1211 on every line),
% but the pp. 71--84 batch set it inside \begin{center}; the two should be made
% to agree, and this one is set as printed.
%
% NO DIVISIONAL RULE.  A rule sweep of the image finds exactly one rule on
% p. 137 -- the centred one below, before the L3 head -- and none above the III.
% head; nor is there one at the foot of p. 136, whose only rules are the
% mid-page divisional rule and the left-flush footnote separator.  So thesis III
% opens without a rule, unlike thesis II on p. 79.
\noindent\textbf{\large III. Det i den historiske Udvikling tidligere Sta\-%
diums Bevidsthedsindhold hæmmer, ubevidst for de
tænkende Individer, ud fra sin oprindelige Eien\-%
dommelighed, den bevidste Stræben, der hævder
den nye Tids Princip i Kamp mod Fortidens; men
saaledes, at den hæmmende Paavirkning bliver In\-%
citament for en dybere og alsidigere Tankeudvikling.}

Vi oplyse dette Forhold ved følgende Exempler:

% p. 137: the three numbered items are ordinary roman paragraphs, unspaced and
% unbold, each with the print's own indent.  All three guillemet pairs on this
% page (items 1, the L3 head, and the body below) point INWARD, «…», read at
% 1200 dpi.
1. Det Antikes Eftervirkning i den «christelige Philoso\-%
phie», navnlig i Scholasticismen.

2. Det Middelalderliges Eftervirkning i Overgangs\-%
perioden.

3. Den middelalderlige Grundanskuelses Indvirkning som
bestemmende Ensidigheden i begge den nyere Philosophies
Udviklingsrækker.

% DIVISIONAL RULE, p. 137: short, CENTRED hairline with white space above and
% below, measured on the image at x 548--800 (300 dpi) against a text block of
% 112--1235 -- rule centre 674.0, block centre 673.5 -- and 2.14 cm long.
% p. 137 bears no footnote, so this cannot be the footnote separator.
\divrule

% HEAD, p. 137 (L3, transcription.tex §3): the first of the 1./2./3. heads under
% thesis III (137, 145, 150).  LETTERSPACED roman at body size, set as a
% justified full-measure paragraph over three lines -- gaps 8--22 px at 600 dpi
% against 2--6 px in the bold L2 head above (see the measurements there).  The
% print's line breaks are: 1. ... Eftervirkning / indenfor ... navnlig /
% indenfor Scholasticismen.
\emph{1. Den antike Verdensanskuelses Eftervirkning
indenfor den «christelige Philosophie» og navnlig
indenfor Scholasticismen.}

Trods den Opfattelse, der tidligere er gjort gjældende
med Hensyn til den saakaldte \emph{christelige Philosophies}
videnskabelige Charakteer, ifølge hvilken denne «christelige
Philosophie» \emph{ikke} danner en særegen Epoche i Philosophiens
Historie, drage vi den dog, idet den betegner \emph{en af en
eiendommelig Livsanskuelse fremgaaet Tankebe\-%
stræbelse,} ind under det angivne Hovedsynspunkt, for derved
at belyse et i \emph{almeenhistorisk} Henseende \emph{væsentligt}
Tidsafsnit.

De tidligere Udviklinger have bestemt \emph{Christendommens
Forhold til den antike Philosophie.} Det har viist sig,
at denne Philosophies Optagelse blev uundgaaelig, at kun ved
den Udviklingen og Fastsættelsen af det christelige Dogme, og
Forsvaret mod Hedenskabet blev muligt. Oldtidens Tanke\-%
% --- p. 138 ---
\setcounter{footnote}{0}
% p. 138 bears ONE note, ten lines long and self-contained (it ends
% "forraader Troen."); nothing runs over onto p. 139.  No Greek.
% p. 138: «christelige Philosophie» -- both marks checked at 1200 dpi: the
% opening one points LEFT and the closing one RIGHT, the book's normal inward
% pair.  No guillemet defect here.
former bleve saaledes i den «christelige Philosophie» i den
christelige Oldtid og Middelalder Former for det christelige
Troesindhold; den christelige Bevidstheds Fylde udfoldede sig
% p. 138: "været," is NOT letterspaced -- its letter gaps measure 3--6 px at
% 600 dpi and the 20 px before its comma is the fount's ordinary comma
% side-bearing.  spacing.py's single-word flag is rejected.
ved dem til en christelig \emph{Lærebygning.} Men Formen, der
optoges, kunde ikke saaledes sondres fra det Indhold, hvis
Form den havde været, fra den hele Livsanskuelse, der op\-%
rindelig havde udpræget sig i den, at den ikke skulde udøve
en mægtig Indvirkning paa det nye Indhold, om den end
\emph{med Bevidsthed} kun sattes i et \emph{tjenende} Forhold til
dette.

% p. 138: "to" IS letterspaced (t--o gap 18 px at 600 dpi against 4 px in "de"
% three words earlier on the same line); spacing.py does not flag it.
En saadan skjult Indvirkning viser sig da ogsaa i de \emph{to}
store Afsnit, i hvilke den saakaldte christelige Philosophie
deler sig: \emph{Kirkefædrenes} og \emph{Scholastikernes, Dogme\-%
frembringelsens} og \emph{Dogmetilegnelsens Tid.}

% p. 138: "immaculata conceptio" in the note is set in the same upright roman as
% the Danish and is not letterspaced, so it takes no markup at all
% (transcription.tex §2).
Det \emph{første} Afsnit: det med Hensyn til \emph{gjennemført}
philosophisk Tænkning mindre betydelige, omtale vi kun i
Korthed. I det udvikles, ved Optagelsen af Oldtidens Philo\-%
sophie, navnlig af saadanne Former som \emph{Philos Lære, Stoi\-%
cismen,} og \emph{Platonismen} i dens \emph{sildigere} Skikkelse, den
christelige Troes Indhold til en Kreds af Dogmer, til hvilken
intet Væsentligt senere behøvede at tilføies\footnote{Vel har endnu vor Tid tilføiet et peripherisk Dogme: Dogmet om
immaculata conceptio, men netop i sit Forhold til dette Dogme har
den røbet sin Modsætning til Christendommens classiske Tid, den
christelige Oldtid og Middelalder. Da Læren om den ubesmittede
Undfangelse i Middelalderen var Gjenstand for Strid, var \emph{Begei\-%
stringen} for Madonna, der i den Troendes Bevidsthed overstraalede
selve Christus, det bevægende Motiv for Forsvarerne. I vore Dage,
da Læren blev fastsat som Dogme, var Motivet \emph{en skjult skeptisk
Reflexion,} der, forat sikkre Christi Reenhed, gaaer et Skridt
længere tilbage, og saa standser, men netop ved dette ene Skridt
forraader Troen.}. \emph{Under} denne
Udvikling maatte der opstaae en Vanskelighed ved at bringe
de for Christendommen eiendommelige Lærdomme, som Læren
om Gudmennesket, om Skabelsen og om det Onde, uden For\-%
% --- p. 139 ---
% p. 139 bears no note and no rule.
% GREEK ON A PAGE THE RECON LIST DOES NOT NAME: μὴ ὄν, below.  The recon
% Greek-page list (127, 129, 131, 132, 133, 136 in this range) is a lower bound,
% as BATCH-AGENT §4.6 warns; p. 139 is another instance.
vanskning ind i de optagne antike Former. For den \emph{troende}
Christnes \emph{umiddelbare religiøse Bevidsthed,} for den
\emph{religiøse Følelse og Forestilling,} traadte det Hedenskes
og det Christeliges Forskjellighed \emph{umiddelbart} frem, og
man sondrede med større eller mindre Bestemthed mellem
den \emph{christelige Troes} og den \emph{hedenske Videns} Indhold.
Men dog var den gamle Kirkes Sammenhæng med Oldtiden
ved Nationalitet og Sprog saa stærk, og de Mænd, der, ikke
blot i den østerlandske og den græske, men ogsaa i den
vesterlandske Kirke, repræsenterede den christelige Oldtids
Tænkning, ved deres Dannelse, ved alle deres aandelige Livs\-%
forudsætninger, saa inderligt knyttede til det Antike, at en
klar Erkjendelse af Modsætningens \emph{Princip} ikke blev mulig,
og, som Følge deraf maatte Klarheden med Hensyn til
Grændsen mellem Gebeterne blive ufuldstændig. \emph{Uvilkaar\-%
ligt} forenedes det Heterogene; selv de eiendommeligste chri\-%
stelige Lærdomme bleve, ved at bringes ind i de antike For\-%
mer, underkastede en Forvandling, ved hvilken Troesindholdets
Eiendommelighed angrebes. I Læren om \emph{Gudmennesket}
trængte Oldtidens Logosbegreb sig ind, saaledes som det var
udviklet hos Stoikerne og Philo; i Læren om \emph{Treenigheden}
bleve psychologiske Analyser Udgangspunktet for Speculationen;
i Læren om \emph{Skabelsen} gjorde den græske Anskuelse af en
evig Verden (Ideeverden --- evig Materie) eller den orientalske
om en Emanation sig hvert Øieblik gjældende, og selv ved
Dogmet om \emph{Synden} virker, trods den christelige religiøse
Bevidstheds afgjørende Divergents fra den antike, den antike
% p. 139: μὴ carries a GRAVE over the η; over the ο of ὄν the mark separates at
% 1200 dpi into TWO elements -- a comma-shaped curl on the left and a rising
% stroke on the right -- i.e. smooth breathing plus acute, unlike the single
% undifferentiated blob of ἄπειρον (pp. 52, 81).  ABBYY reads "iV" and
% tesseract "u 6»y"; neither has a Greek character.
metaphysiske Opfattelse af det Onde som et \gk{μὴ ὄν}, som Be\-%
grændsning og Ufuldkommenhed.

Denne Beherskethed ved Oldtidens Tænken, under For\-%
søgene paa at hævde det eiendommelig Christelige, finder
ikke blot Sted hos de \emph{ældste} philosophiske Forsøg hos
\emph{Apologeterne,} hos hvem \emph{stoiske} og \emph{philoniske Be\-%
greber umiddelbart} anvendes paa det specifisk Christe\-%
lige og forvandle det til speculative Theoremer; men der kan
% --- p. 140 ---
% p. 140 bears NO note, no Greek and no rule.
% SEAM CHECK, p. 140 -> p. 141: p. 140 carries no footnote at all (confirmed by
% a rule sweep -- there is no footnote separator on the page), so nothing runs
% over onto p. 141, whose first line "Heelhed." is body text.
uden \emph{Indskrænkning} siges om \emph{hele} Patristikens Tid, at
\emph{alt} dens Tankeindhold var optaget fra Oldtidens Tænkning,
og at af den Grund ethvert Forsøg paa at bringe det christe\-%
lige Dogma ind i Tankens Form førte til \emph{en ubevidst Me\-%
tamorphose} af det Christelige, til en Omsætten af Troens
paa et overnaturligt Princip hvilende overmenneskelige Indhold
i et ved den Menneskenaturen udtrykkende Tanke bestemt
Led i en Begrebssammenhæng.

Af større Interesse i philosophisk Henseende end i Pa\-%
tristikens Tid bliver Forholdet mellem det Christelige og det
Antike i \emph{Scholasticismen,} fordi der den umiddelbare
Sammenhæng, Frændskabet med Oldtiden, var afbrudt, fordi
den principielle Modsætning var traadt klarere frem for Be\-%
vidstheden, og fordi Forsøget paa at optage Troesindholdet
uforvansket i Tankens Form gjordes med rigere Kræfter og
større Energie. Mænd som \emph{Thomas Aquinas} og \emph{Duns
Scotus} ere saa mægtige Aander, saa uforlignelige Tankebyg\-%
mestere, at deres Forsøg paa at skabe en Philosophie, der
kan fortjene Navnet: christelig, maa blive af høieste Betydning,
og dette Forsøgs Mislykken et ubedrageligt Kjendetegn paa
Umuligheden af at løse den stillede Opgave.

Tager man de store Systemer, som hine Mænd skabte i
Scholasticismens Blomstringstid, og underkaster dem en Ana\-%
lyse, saa viser \emph{Sammenføiningen af deres Bestanddele}
sig let. Der stræbes efter at give \emph{en christelig Lære\-%
bygning, der omfatter Troens og Videns Totalitet
og sammenføier Troesindholdet og Tankeindholdet
til en Heelhed.} Denne Opgaves Løsning see vi forsøgt i
den Form, at den antike Verdensanskuelse, der samlede sig i
Forestillingen om det harmonisk ordnede guddommelige Uni\-%
versum, ligesom omfattes af det Christeliges Ramme. An\-%
skuelsen af Verdensvirkeligheden, der maatte tages fra den
Philosophie, for hvilken Verden var en Virkelighed, og som
for Oldtiden var \emph{det Hele,} bliver nu, forbunden med det
% p. 140: "større" at the foot of the page is letterspaced and the phrase runs
% on into p. 141 ("Heelhed."); the pp. 141--154 batch should check whether the
% letterspacing continues there and, if it does, close a single \emph{} run.
christelige Troesindhold, til \emph{en Deel,} indføiet i en \emph{større}
% --- p. 141 ---
% SEAM p.140 -> p.141 checked (BATCH-AGENT §5): p. 140 bears NO footnote and
% no note matter stands at the head of p. 141, so nothing runs over.  The
% sentence runs across the break ("... indfoiet i en storre / Heelhed."), and
% p. 141's "Heelhed." is NOT letterspaced (measured 5 px letter gaps).
% p. 141: the id-est mark ɔ: CONFIRMED on the image ("Villiesproduct ɔ: som
% det") -- the one candidate site in this range that bears it.
% Emphasis on this page was settled by measuring letter gaps at 600 dpi, not
% by eye: unspaced words run 1--5 px between glyph bounding boxes, letterspaced
% words 9--20 px.  Three words that LOOK spaced at low magnification measure
% tight and are set plain here: "Overnaturlige" (l. 12), "Bestaaen" (l. 27),
% "Det i Naturen ... Virkende" (l. 28, where only "ene" is spaced).
Heelhed. Det fra Oldtiden optagne Indhold udfyldte kun den
\emph{midterste} Deel af et Verdensdrama, i hvilket \emph{første} og
\emph{sidste} Act tilhørte Christendommen. \emph{Forud} for det natur\-%
lige Universum sattes den naturfrie, almægtige, skabende Gud,
\emph{efter} det det christelige Himmerige, Naadens overnaturlige
Rige. \emph{Begyndelsen og Slutningen} blev, efter Scholasti\-%
cismens \emph{egen} Indrømmelse, \emph{det for Tanken Utilgænge\-%
lige, Ubegribelige}, altsaa i \emph{Væsen uensartet} med det,
der udfyldte \emph{Mellemrummet}. Den rent aandelige Guds
Skaben af en materiel Verden ved en i det aandelige Væsen
ikke begrundet, altsaa absolut vilkaarlig, Almagtens Handling;
det Naturliges Gjennembryden ved det Overnaturlige, der
sætter et absolut Nyt, ved det Foregaaende Uforklarligt, er
\emph{det Tanken Modsigende}. Og det er dette idet det er
\emph{det Naturen Negerende}. Men Begyndelsen og Slutningen
staae, efter Sagens Natur, ikke som passive og ligegyldige ved
Siden af det, der omfattes af dem; dette maa \emph{omformes}
ved Forholdet til dem. Idet Gud er sat \emph{som Skaber}, bliver
Verden, det \emph{Skabte}, sat som \emph{Villiesproduct} ɔ: som det,
der \emph{ikke har Væsenhed}, som det, der i sig er et \emph{Intet}
og kun bæres oppe af den guddommelige \emph{Almagt}. Med
andre Ord: med \emph{Skabelsen}, der er det første altomfattende
Under, er \emph{Underet} sat som Synspunktet for Opfattelsen af
Naturen. I \emph{hvert} Øieblik og paa \emph{hvert} Punkt gjennem\-%
bryder Underet Natursammenhængen og Naturloven og op\-%
hæver Loven som saadan, og selve det Naturliges uforandrede
Bestaaen er kun dets \emph{Opholdelse} ved \emph{Skabelsesunderets
Gjentagelse}. Det i Naturen \emph{ene} Virkende er \emph{for Troens
Øie} den guddommelige \emph{Almagt}, det \emph{Overnaturlige} er,
som \emph{Duns Scotus} udtrykker det, \emph{Naturens Begreb og
Grund}, Naturens egentlige Væsen. Og i denne i sig \emph{væ\-%
senløse} Natur har den \emph{christelige} Aand ingen Hvile; den
er kun i den som i \emph{et Fremmedt}, der, naar det virker
% p. 141, l. 34: "som" is letterspaced (gaps 18/12 px) but the following
% "Natur," is NOT (gaps 1--2 px).  Set as printed.  The rejected alternative
% is that the whole phrase "som Natur" is emphasised; the measurement does not
% support it.
\emph{som} Natur, sees som \emph{det Onde}, og som bestandig, blot som
tilværende i sin Intethed, føles som en \emph{hæmmende Skranke}
% --- p. 142 ---
% p. 142, l. 5: a small stray dot stands in the left margin beside
% "bryde Menneskelivets ..." (ABBYY renders it as a bullet).  Not printed
% type -- an ink speck or modern pencil -- and not transcribed.
og ønskes tilintetgjort. Aandens Retning følger den \emph{over\-%
naturlige Virkning}, der udgaaer fra den almægtige Villie
for i denne Tilværelse som \emph{Undervirkning} at gjennem\-%
bryde Naturskrankerne, for som \emph{Naadevirkning} at gjennem\-%
bryde Menneskelivets Naturbetingelser, og ved hint som ved
dette at anticipere den \emph{Fuldendelse}, der fuldbyrdes, naar
det Overnaturliges Straaler samle sig i eet Lys, i \emph{Salighe\-%
dens Rige}. I Forestillingen om dette er Naturen, hvis In\-%
tethed i det Forbigangne er sat for Heelheden ved Skabelsen,
og i det Nærværende sættes for det Enkelte ved Underet, i
en ny Form sat som det, der er \emph{et Intet for Aanden}, og
som derfor engang, naar Historien er afsluttet, Guds Raad\-%
slutning fuldbyrdet, \emph{atter skal ophøre, blive til Intet
fordi det er af Intet}. I den hiinsidige Tilværelse ere \emph{alle}
Naturens Skranker, der kun føles som hæmmende Baand,
tilintetgjorte, og den \emph{naturløse} Aand lever, i nydende Con\-%
templation af det \emph{rent aandelige} Guddommelige, dettes
uendelige \emph{salige} Liv.

Saaledes er den antike Verdensanskuelse og det Virkelig\-%
hedsindhold, der toges fra denne, forvandlet i sin Grund ved
Underordningen under det Christelige. Den hele Opfattelse
af Naturens Væsen, Aandens hele theoretiske og praktiske
Forhold til den, er forandret. Naturen har mistet sin Væ\-%
sentlighed, dens Love ere opløste af Underet; den er ikke
længere Gjenstand for Theorien, thi Aanden søger i den kun
det Overnaturlige; den er ikke Grundlag for det sædelige Liv,
thi Aanden flyer Naturen, der, naar den gjør sig gjældende
som Natur, sees som det Onde; den begrændser ikke, som et
Selvgyldigt, Aandens Vilkaarlighed, thi ved Phantasien giver
den abstracte naturløse Subjectivitet sig en Uendelighed, for
hvilken ingen Skranke bestaaer, og Aanden hersker tilsyne\-%
ladende med ubegrændset Magt over den afmægtige Natur.
Der synes da at være bragt en \emph{spiritualistisk Livsan\-%
skuelsens Enhed} tilveie, om end ikke en \emph{Tankeenhed}, ---
den vilde ialfald kun være en formel; --- det Christelige
% --- p. 143 ---
synes at have \emph{assimileret} det Optagne. Og dog viser der
sig, og \emph{maa} der vise sig en \emph{Reaction} fra \emph{det Antike},
ved hvilken Enheden bliver usikker, og Inconsequentser trænge
sig ind.

At en Tilbagevirkning af det Antike ikke kunde ude\-%
blive, ligger allerede deri, at for Mennesket, der kun er
\emph{Aandsvæsen} idet det tillige er \emph{Naturvæsen}, den \emph{supra\-%
naturale} Opfattelse af Tilværelsen altid fordrer et \emph{natura\-%
listisk} Udgangspunkt; at det \emph{Overnaturlige} kun er \emph{over\-%
naturligt} ved at hæve sig over et \emph{forudsat, ikke ved det
forklaret Naturligt}; at Mennesket kun kan leve det \emph{rent
aandelige} Liv i dets Abstraction idet det i hvert Øieblik
kommer til det ved at \emph{gaae bort} fra det \emph{verdslige}, saa
at der, hvor et spiritualistisk Standpunkt søges gjennemført,
dog \emph{indirecte} maa tillægges det Naturlige, som det Over\-%
naturliges Modsætning og Betingelse, en Selvstændighed. Vi
finde ogsaa indenfor den christelige Middelalder en Gjøren sig
gjældende af Naturgrundlaget, af hvilket den antike Livs- og
Verdensanskuelse bares, ikke blot i den almindelige Bevidsthed,
men ogsaa i den philosophiske. I hiin fremtræder den i den
dæmoniske Magt, som den fornægtede Natur, som væsentligt
Moment i det Menneskelige, vedbliver at udøve over Bevidst\-%
heden. Sandseligheden, der fornægtes, nøies ikke med at
rehabiliteres paa ideel Maade i Dyrkelsen af Qvindens gud\-%
dommeliggjorte Ideal; men den bliver en real Magt, uimod\-%
% p. 143, ll. 26 f.: the print reads "og i Don-" at the foot of l. 26 and
% "Juans Sagnet" at the head of l. 27.  The hyphen is taken here as PART OF
% THE COMPOUND (cf. "Christus-Historien" three lines below, hyphenated within
% the line), not as a line-break hyphen; the alternative reading "Don Juans
% Sagnet" cannot be excluded, since the two coincide at the line end.
staaeligt fristende, og i Sagnet om Venusbjerget, og i Don-Juans
Sagnet, der opstaaer ved Middelalderens Slutning, og
tilligemed Faustsagnet danner de profane Paralleler og Mod\-%
sætninger til Christus-Historien, fremtræder den i en Mægtig\-%
hed, som Oldtiden ikke kjendte. Og medens Naturen i sin
Fylde unddrager sig Tanken, der ikke mere har et væsentligt
Tilhold deri, bliver den \emph{Phantasiens udvalgte Gjenstand}
og behersker Aanden gjennem den. Naturen, der, seet i
Christendommens Lys, blev \emph{Underets Verden}, ikke Virke\-%
lighedens, bliver, hvor den sees i sin \emph{egen} Dæmring, en
% --- p. 144 ---
\emph{Phantasiens Verden}, i hvilken Loven er forsvunden og en
naturfremmed Phantasies magtløse Almagt er det ene Raa\-%
dende. Eventyrets Verden afløser Tankens: i digterisk Leg,
i anthropomorphistisk Forvandling, i forunderlig Blanding af
Tiltrækning og Frastøden, opfattes Naturen. Magiske Kræfter
beherske den; Sympathien sammenholder den; hemmeligheds\-%
fulde Væsener, undertiden lunefuldt venlige, hyppigere forfær\-%
dende, Afbilleder af Aandens skjulte Angest for Naturen, be\-%
folke den. Magisk raader Aanden gjennem Phantasien over
denne Eventyrets Verden, og dog bliver den netop med det
Samme gjennem Phantasien --- det efter sit Væsen Afhængige
--- magisk behersket af den uforstaaede Natur, den bliver,
beruset, overvældet, fængslet ved dens Magt.

I Middelalderens \emph{philosophiske} Bevidsthed, i Schola\-%
sticismen, træder den det Naturliges Gjøren sig gjældende, gjen\-%
nem hvilken den \emph{antike} Verdensanskuelse reagerer, frem i Læren
om \emph{Formens Realisation i Natursammenhængen}, om
\emph{Gradforskjellene i Naturen}, om \emph{det Naturliges Be\-%
tydning for Menneskelivets Udvikling i intellectuel
og ethisk Henseende}. I alle disse Bestemmelser ind\-%
rømmes der, netop ved Begrebet: \emph{Udvikling og Udvik\-%
lingstrin}, det Naturlige en \emph{Væsentlighed}; thi kun in\-%
denfor det Væsentlige kan en Gradforskjel og en Udfoldelse
gives, ikke i det Væsenløse; kun en væsensgyldig Natur kan
faae Betydning for det Aandeliges Udvikling. \emph{Det fornægtede
Naturlige tvinger sig saaledes ubevidst frem}; kun
ved det faaer Tanken Fodfæste og Indhold; derfor er det
uundgaaeligt, at det gjør sig gjældende. Og om end dets
Betydning \emph{atter} afskjæres, idet Formen sees realiseret i en
ved en \emph{Skabelse} sat Tilværelse, hvis Continuitet afbrydes af
Underet, og i hvilken Guds Almagt er det umiddelbart Vir\-%
kende og Væsentlige; og idet \emph{hele} det naturlige Liv, ogsaa
hele den naturlige \emph{sædelige} Handlen, de \emph{ethiske} Dyder,
% p. 144, ll. 34 f.: "kun" is letterspaced, "sættes" is NOT (gaps 1--5 px),
% and "som Middel" is letterspaced again.  Set as printed.
der bestemtes som Forudsætning for de theologiske, \emph{kun}
sættes \emph{som Middel} for et \emph{Overnaturligt}, der ikke rea\-%
% --- p. 145 ---
% p. 145, l. 6: "er" is letterspaced ("e r", gap 16 px) while the flanking
% "det" and "Supranaturale" are not (gaps 2--9 px).  Set as printed; the
% rejected alternative is that this is the compositor letterspacing a short
% word to justify a loose line (word spaces on this line run 37--70 px).
% p. 145, l. 7: a small isolated mark stands between "af det," and "Flugten"
% (ABBYY renders it as a raised dot).  Not printed type; not transcribed.
liserer sig gjennem den naturlige Udvikling, \emph{ikke} er tilstede
\emph{deri}, men \emph{træder ind udenfra} ved et nyt Skabelsesunder,
som en \emph{høiere, i sig afsluttet Orden}, som en fremmed
evig Verden ligeoverfor det Timelige; saa mister det kun sin
Betydning for i næste Øieblik at tiltvinge sig den igjen; thi
kun i Benægtelsen af det, i Flugten fra det, \emph{er} det Supra\-%
naturale, og \emph{Benægtelsen} af det, \emph{Flugten} fra det, der
forudsætter, at det \emph{var sat som en Virkelighed}, er kun
idet det bestandig \emph{paany} sættes som en saadan, for atter at
fornægtes, atter at bortskydes, men i \emph{evigt Kredsløb} at
komme tilbage som \emph{det i sin Oprindelse Ubegrebne}, og
\emph{derfor i sin Væren Selvstændige}.

% DIVISIONAL RULE, p. 145, standing between the end of the paragraph above
% and the L3 head "2. ...": short, CENTRED (about a third of the measure),
% with generous white space above and below.  Not the footnote separator --
% this page bears no note.
\divrule

% L3 HEAD, p. 145 (transcription.tex §3: the "2." of the sequence 1./2./3.
% under thesis III, pp. 137, 145, 150).  MEASURED at 600 dpi, as §3 requires:
% intra-word letter gaps 9--17 px (median ~12) against 1--5 px in the body of
% the same page and 2--5 px in the certainly bold L2 head on p. 79.  It is
% therefore LETTERSPACED ROMAN, at body size and body weight, set as a
% justified PARAGRAPH on the full measure with the first line indented --
% NOT bold and NOT centred.
\emph{2. Den middelalderlige Grundanskuelses ubevidst
hæmmende Indvirkning i Overgangsperioden.}

Vi have i det Foregaaende paaviist \emph{Overgangsperio\-%
dens} almindelige Charakteer, dens beherskende Grundtanker,
og den Udviklingsgang, der førte den længere og længere bort
fra det Middelalderlige. Det Bestemmende i den Bevægelse,
der mere og mere fjernede den fra dette, og som lod den
Kamp, der først kun var rettet mod \emph{Philosophien}, blive
til en Kamp ogsaa mod Middelalderens \emph{Religion} og \emph{hele
Livsanskuelse}, var \emph{det Naturliges} bestandig stærkere
Hævdelse i dets Væsentlighed. Men medens Bevægelsen gaaer
bort fra det Middelalderlige, ja endog hvor Modsætningen til
dette bliver til forbittret Kamp, see vi dog Overgangsperiodens
Tænkning ubevidst at blive paavirket af den Fortid, af hvilken
den er fremgaaet, og Consequentsens Strænghed og Klarhed
at lide derunder.

De Tanker nemlig, der ovenfor ere blevne angivne som
de centrale i Overgangsperioden: Tankerne om det guddomme\-%
lige Universum og det guddommelige, mikrokosmiske, aandelig-naturlige
Menneskevæsen, maatte, \emph{consequent} fastholdte og
% Sheet signature "10" stands at the foot of p. 145; not transcribed.
udviklede, have ført til, at see \emph{det besjælede Universum}
% --- p. 146 ---
som al Tilværelse og som evig værende \emph{ved sig selv}, og at
see det Aandelige og det Legemlige i det og i Mennesket
\emph{uopløseligt} og \emph{uadskilleligt} sammenknyttede. Disse
Consequentser saae vi ogsaa udtalte hos \emph{Giordano Bruno}.
For ham bar Materien alle Former i sig: den var selv et
Guddommeligt; nærmere bestemt var den, som Moment i Guds,
de fire aristoteliske Aarsager sammenfattende Væsen, selve
Guds almindelige Kraft eller Magt, der er Eet med Guds
frembringende Virksomhed, og Udtrykket for Guds Væsens
Nødvendighed. De mange vexlende Former, som den kraft\-%
opfyldte Materie indenfra udfolder, vise hen paa een almindelig
Form: den almindelige Forstand, Verdenssjælen, der er For\-%
mernes Kilde og som saadan i Enhed med Materien. Denne
\emph{virksomme Enhed af Materie og Form} betegnes som
\emph{Naturen}. Som \emph{Princip} er den et \emph{evigt} Væsen, den er
% p. 146: the Italian «spirito si trova in tutte le cose» is set in the SAME
% upright roman as the Danish and is NOT letterspaced, so it takes neither
% \emph{} nor \textit{}.  Both guillemets read at 1200 dpi and both point
% INWARD.
\emph{guddommelig og uendelig}, overalt besjælet («spirito si
trova in tutte le cose»). Og dette naturlige Verdensalt, der
er Eet med Verdenssjælen, sættes som \emph{al naturlig Philo\-%
% p. 146: «Gud i Tingene» re-checked at 1200 dpi.  The opening mark is «
% (points left) and the closing mark is » (points right): the book's normal
% INWARD-pointing pair, not one of the ~13 reversed openings.  The words
% inside are letterspaced, the guillemets are not, so the marks stand
% outside \emph{}.
sophies Maal og Grændse}; det kaldes «\emph{Gud i Tin\-%
gene}».

I disse Bestemmelser er en \emph{consequent Pantheismes}
Standpunkt naaet: det af den græske Philosophie i dens Cul\-%
minationstid Tilstræbte er udtalt, og den \emph{principielle Mod\-%
sætning til det Christelige} er statueret, og med disse
Bestemmelser hos \emph{Bruno} stemme Udsagn af de betydeligste
af Periodens Repræsentanter overens. Men i det Øieblik, at
Maalet synes naaet, gjør en Eftervirkning af den Grundanskuelse,
mod hvilken der kæmpedes, sig gjældende: der trænger sig
igjen en Tanke frem om \emph{en fra Universet sondret Gud},
vi finde igjen \emph{Skabelsens} Forestilling, vi see et \emph{Over\-%
naturligt} adskilt fra et Naturligt, et \emph{rent Aandeligt} fra
et Legemligt, vi see igjen Tanke og Materie i \emph{Modsætning}
til hinanden, og disse Bestemmelser finde vi kun paa en \emph{ud\-%
vortes} Maade knyttede sammen med hine.

Adskillelsen mellem \emph{Gud} og \emph{Universet} kommer hos
% --- p. 147 ---
% p. 147, l. 4 and l. 16: TWO occurrences of a lower-case `i' set from a
% MARKEDLY HEAVIER (bold) fount -- "dette som Dele i, ikke af," and "som
% værende i Tingene".  transcription.tex §2 records three such heavy `i's
% (pp. 42, 89, 95) and reads them as a WRONG SORT, transcribed as an ordinary
% `i'.  That reading does not fit here: both stand in a contrastive position
% ("i" against "af"; immanence against transcendence), and on l. 4 the very
% word contrasted with it, "af", IS LETTERSPACED (measured).  The bold `i' is
% therefore read here as the compositor's substitute for letterspacing a
% single letter, which would be invisible, and is set \emph{}.  The rejected
% alternative -- an ordinary `i' from a stray heavy sort -- is what §2 assumes
% for pp. 42, 89 and 95; those three sites should be re-examined.  REPORT.
% p. 147, l. 5 f.: "tutto infinito" is letterspaced ENTIRE on l. 5 but only in
% its first word on l. 6, where both occurrences of "infinito" measure tight
% (2--5 px) beside "totalmente" and "tutto" at 12--19 px.  Transcribed as
% printed, not regularised.
% p. 147, l. 7: "d. v. s." is printed with a word space after each point; this
% is the abbreviation's own setting, not letterspacing.
% p. 147, l. 15: a small stray dot stands between "ellers" and "af"
% (ABBYY renders it "ellers . af").  Not printed type; not transcribed.
% p. 147: the Latin (natura naturans, deus creans, monas monadum) is set in
% the SAME upright roman as the Danish; it takes \emph{} only where it is
% letterspaced and nothing where it is not.  Never \textit{}.
\emph{Bruno} frem i den Form, at han, efterat have bestemt Uni\-%
verset som \emph{uendeligt}, og efterat have sikkret denne Uende\-%
lighed ligeoverfor det Endelige, den omfatter, ved at sætte
dette som Dele \emph{i}, ikke \emph{af}, det Uendelige (nell' infinito, ikke
dell' infinito), adskiller det, der baade er \emph{tutto infinito} og
\emph{totalmente} infinito, fra det, der kun er \emph{tutto} infinito,
d. v. s. det, der ligesom er et heelt igjennem Uendeligt (et
absolut Uendeligt) fra det, der kun er som Heelhed uendeligt.
Hint er Gud, dette Universet. \emph{Gud} er \emph{absolut uendelig}
fordi han er heel i den hele Verden og i hver af dens Dele
er paa en uendelig og fuldstændig Maade. \emph{Verden} er \emph{som
Heelhed uendelig} fordi den hverken har Rum eller Grændse,
\emph{men ikke heelt igjennem uendelig} fordi hver Deel deri
er endelig og fordi i dens Dele Modsætningerne ikke ere
ubetinget forenede. Gud, der ellers af Bruno betegnes som
\emph{natura naturans}, som værende \emph{i} Tingene, bliver nu en
deus \emph{creans} og betegnes som det \emph{supernaturale} Princip,
som den \emph{supersubstantiale Substants}; han bliver den
\emph{høieste Monade} (Individualitet), Monade i \emph{ubetinget For\-%
stand} (monas monadum), den \emph{aandelige}, blivende Enhed
uden Vorden. Men denne Distinction, der fra Begyndelsen af
kun hviler \emph{paa et Postulat}, bliver \emph{uigjennemførlig} ved
Siden af de tidligere Bestemmelser af \emph{Universets Væsen}
og af dets \emph{paa ethvert Punkt værende} Uendelighed og
Udelelighed.

I hiin \emph{første} Sondring, der ligesom gjør det \emph{uendelige}
Værende \emph{splidagtigt} i sig selv, ligge alle de andre Son\-%
dringer. Sondringen mellem \emph{Gud} som et \emph{Overnaturligt}
og \emph{Verdenstilværelsen} som det \emph{Naturlige} er allerede
berørt. \emph{Skabelsesforestillingen} kommer ind idet Naturen
sees som den uendelige Aarsags \emph{Virkning}, og om end denne
Frembringen sættes som \emph{nødvendig}, altsaa indirecte betegnes
som en, Væsensenheden indesluttende, \emph{Væsensudfoldelse},
paatrænger dog Udtrykket Skabelse sig, netop fordi Aarsagen
% Sheet signature "10*" stands at the foot of p. 147; not transcribed.
\emph{sondres} fra Universet. Og idet Gud, i Sondringen fra det
% --- p. 148 ---
\setcounter{footnote}{0}
% p. 148 bears ONE note, marked *) after "Livsaanden" and before the full
% stop.  It is self-contained: it ends "...(Nerveætheren)»." above the foot
% of the page and nothing of it stands at the head of p. 149, which opens
% with its own marked note.
% p. 148, ll. 7 and 16: two small stray dots stand in the type area (after
% "dog" on l. 7, before "i den concrete" on l. 16; ABBYY renders the first
% as "dog'uvilkaarligt").  Not printed type; not transcribed.
% p. 148, note: the sign between "Ting" and "Livsaanden" is a PLUS SIGN, read
% at 1200 dpi (the OCR's "-{-" and "-+" are both misreadings of it).  It is
% set here as text, not in math mode.  Both guillemets in the note point
% inward.
Naturlige, det uendelige \emph{Legeme}, væsentlig bliver en \emph{aan\-%
delig Enhed}, saa bliver der, trods Bestræbelsen for at sætte
\emph{Materien i ham}, som Moment i hans Væsen, dog kun sat
en \emph{aandelig Kilde til den}, en Materiens «ratio», i Gud,
og i \emph{Universet}, hvor \emph{Ætheren} paa een Gang skulde være
den almene, altgjennemtrængende Materie og den levende,
virkende Kraft, \emph{skiller} dog uvilkaarligt det Legemlige og det
Aandelige sig fra hinanden. Det \emph{Legemliges} Bestemmelser
blive \emph{uforklarede} ud fra \emph{Monaderne}, der, som Afbilleder
af det Guddommeliges Væsen, blive \emph{aandelige} Kræfter, og
\emph{mellem Sjælen og Legemet} maa der derfor indskydes et
\emph{forenende Mellemled: Livsaanden}\footnote{Man sammenligne hermed Telesios Opfattelse af det enkelte levende
Væsen som «en Kjæde af Ting + Livsaanden (Nerveætheren)».}.

Denne Inconsequents, der paa den angivne Maade gjør
sig gjældende i Overgangsperiodens Opfattelse af dens Grund\-%
tanker, medfører saa atter Følger for det ved disse Betingede;
og disse Følger vise sig i den concrete Opfattelse af Natur\-%
tingene, og i Erkjendelseslæren. I Periodens Grundtanker og
Grundretning laae det begrundet, at den, \emph{consequent},
maatte stræbe efter en Erkjendelse af den concrete Natur\-%
tilværelse ud fra \emph{physiske} Principer, og at den maatte ind\-%
rømme \emph{Erfaringen} og den \emph{frie}, af fremmed Autoritet uaf\-%
hængige, \emph{Selvbevidsthed}, den høieste Betydning. Denne
Stræben og denne Indrømmelse finde vi ogsaa, navnlig i Pe\-%
riodens \emph{andet} Afsnit, hvor man ikke nøies med, gjennem
Bestemmelsen af det \emph{Materielles} Forhold til \emph{Guds Væsen},
\emph{indirecte} og i \emph{Almindelighed} at affirmere det Naturliges
Selvbegrundethed, men hvor man søger Principerne for den
concrete Naturtilværelse i det \emph{Varme} og det \emph{Kolde}, \emph{Ild} og
\emph{Vand} o. desl., og hvor man, medens man paa den ene Side
gjør Krav paa \emph{Erkjendelsens Uafhængighed af Auto\-%
riteten}, paa Tankens Frihed, paa den anden Side viser hen
\emph{til et Udgangspunkt for denne frie Erkjendelse i}
% --- p. 149 ---
\setcounter{footnote}{0}
% p. 149 bears ONE note, marked *) after "Phantasien" in l. 11.  It is
% self-contained (it ends "... paa dens bacchantiske Tog." above the foot of
% the page); p. 150 opens with running text, so nothing runs over.
% p. 149, l. 8: a short horizontal pencil stroke stands under/after
% "samvirker" (ABBYY renders it as a comma).  Not printed type; not
% transcribed.  A second stray dot stands before "det" in l. 28.
% p. 149, l. 22: "Fremskridens" MEASURES TIGHT (1--6 px) between the
% letterspaced "successive" and "Nødvendighed" (10--20 px), and is set plain
% here.  Likewise "Tanken" in l. 25, between the spaced "gjennem Naturen" and
% "Gud".  Both look spaced at low magnification and are not.
\emph{Naturen}, til \emph{Sandsernes} og \emph{Erfaringens} Vidnesbyrd, og
fordrer Erkjendelsens Vei lagt saaledes, at den fra disse fører
hen til Fornufterkjendelsen (jfr. navnlig \emph{Telesio} og \emph{Bruno}).
Men \emph{Fortidens Eftervirkning} giver ikke blot ved Over\-%
gangsperiodens Begyndelse Systemerne et theologisk Præg,
men lader \emph{aldrig} dette Præg fuldstændig forsvinde, om det
end successivt bliver svagere under Naturvidenskabernes Ind\-%
flydelse; den samvirker med Eftervirkningen af den græske
metaphysiske Naturopfattelse i at berøve de \emph{physisk be\-%
nævnede} Principer deres physiske Betydning, og i at give
\emph{Phantasien}\footnote{\emph{Phantasiens} Overmagt viser sig navnlig i den digterisk-phanta\-%
stiske Anskuelse af Naturens Livsenhed, og af dens \emph{umiddelbare},
de physiske Betingelser overspringende, Sympathie; ligeledes i Troen
paa Magien. Overalt i det Concrete bliver Phantasien raadende
Magt, og Tanken følger den ligesom beruset paa dens bacchantiske
Tog.} Overmagten over Forstanden i Opfattelsen af
Naturen, der for Aanden endnu kun er den anthropomorphi\-%
serede Natur; den begrændser den frie Erkjendelse ved \emph{Au\-%
toritetstro} (Philosopherne i første Afsnit af Perioden, \emph{Te\-%
lesio, Campanella}) og den lader den \emph{mystiske Con\-%
templation} af det \emph{overnaturlige} Guddommelige blive
det Høieste. Dette Sidste fremtræder især klart i \emph{Brunos}
udførte Erkjendelseslære. Han paaviser de Trin, ad hvilke
den menneskelige Erkjendelse hæver sig op fra Ideernes sand\-%
selige Billeder (Afbilleder) til Ideerne og til den aandelige
Enhed, der er Altets Livsgrund. Her bliver den \emph{successive}
Fremskridens \emph{Nødvendighed} betonet; det \emph{Endelige}, der
er Guddommens Virkning, sættes som det, der er \emph{uundvær\-%
ligt} til Erkjendelsen af Aarsagen: \emph{gjennem Naturen} skal
Tanken finde \emph{Gud}. Men \emph{paa een Gang} skydes dette Syns\-%
punkt tilside, og om end heri en \emph{ensidig Immanents\-%
theorie} og et \emph{ensidigt bestemt Gudsbegreb} gjør sig
gjældende, saa medvirker tillige, som det i \emph{Udførelsen} viser
sig, Eftervirkningen af det \emph{middelalderlige, supranaturale}
% --- p. 150 ---
% p. 150, l. 13: within the Italian «impossibile e nulla a chi non crede» only
% the last word is letterspaced (12--17 px against 2--7 px for the rest of the
% quotation).  Set as printed; the rejected alternative is that this is
% justification spacing.  Both guillemets point inward.
% p. 150, l. 26: GREEK.  \gk{μανία}, read off the image at 1200 dpi; the mark
% over the iota is an ACUTE (a stroke rising to the right), and there is no
% breathing on this word.  Neither OCR witness contains a Greek character
% ("fiavia" / "pavtfa").  This page is NOT on the recon list of Greek-bearing
% pages -- the list is a lower bound, as BATCH-AGENT §4.6 warns.
% p. 150, l. 26: "eroïco furore" -- the i carries a DIAERESIS, read at 1200 dpi
% (the OCR gives "eroico" / "ero'ico").  Italian, upright roman, unspaced.
% p. 150, l. 28: a small stray dot stands above the second `a' of "Maal"
% (ABBYY reads the word "Maaf").  Not printed type; not transcribed.
\emph{Transscendentsstandpunkt}. Idet Guds Iboen i os i sin
\emph{Uendelighed} accentueres, bliver den det ene Afgjørende,
ligeoverfor hvilken den \emph{endelige} Betingethed taber sin Be\-%
tydning: \emph{Erkjendelsens Grader forsvinde i det Uen\-%
delige}. I Kraft af Guds Iboen kunne vi see Alt i een An\-%
skuelse, med \emph{umiddelbar Evidents}, paa \emph{Intuitionens}
Maade. I disse Bestemmelser have vi endnu Udtryk for \emph{Pan\-%
theismen}, om end dennes \emph{eiendommelige} Form er be\-%
stemt ved Indvirkninger fra en Fortid, hvis \emph{Forhold til
Naturen} var et andet end Overgangsperiodens. Men \emph{Mid\-%
delalderens} Indvirkning røber sig saa, idet Contemplationen
udtrykkelig sættes som en \emph{overnaturlig} Act, der skal være
«impossibile e nulla a chi non \emph{crede}». Ligesom, efter \emph{Scho\-%
lasticismens} Lære, den overnaturlige Forhøielse af de men\-%
neskelige Evner vel sættes som \emph{betinget ved}, men dog \emph{ikke}
som i Sandhed staaende i \emph{teleologisk Sammenhæng} med
den \emph{naturlige} Udvikling, saaledes skal ogsaa denne \emph{Con\-%
templation} være \emph{forberedt} i os ved Tankelivets Udvikling,
men ikke fremgaae af denne; den skal komme som en \emph{plud\-%
selig} Virkning fra Gud, den «almindelige Intelligents», der,
idet den bestemmes som \emph{overnaturligt} virkende, uvilkaar\-%
ligt forvandles til det «\emph{supranaturale} Princip», til den
\emph{transscendente} Guddom, i hvis Anskuen Saligheden be\-%
staaer. I denne Anskuen er Sjælen i \emph{Enhed med Gud}, og
\emph{Sjælens Forvandling til Gud} fuldendes i den Kjærlig\-%
hedens \gk{μανία}, den eroïco furore, i hvilken Erkjendelse og
Handlen mødes, og Tankens og Villiens, Erkjendelsens og det
sædelige Livs sidste Maal er i Enhed.

% DIVISIONAL RULE, p. 150, between the end of the paragraph above and the L3
% head "3. ...": short, CENTRED, with white space above and below.  This page
% bears no note, so it is not the left-flush footnote separator.
\divrule

% L3 HEAD, p. 150 (transcription.tex §3: the "3." of the sequence 1./2./3.
% under thesis III, pp. 137, 145, 150).  MEASURED at 600 dpi: intra-word
% letter gaps 8--15 px on the first line and 8--17 px on the second, against
% 1--5 px in the body of the same page and 2--5 px in the certainly bold L2
% head on p. 79; word spaces in the head run 33--65 px.  LETTERSPACED ROMAN,
% body size and body weight, set as a justified PARAGRAPH with the first line
% indented -- NOT bold, NOT centred.
\emph{3. Ensidigheden i den nyere Philosophies to
Udviklingsrækker som bestemt ved den skjulte Ind\-%
virkning af den middelalderlige Grundanskuelse.}

Den nyere Philosophies Opgave blev i det Foregaa\-%
ende bestemt som den, ud fra Forudsætningen af det Aande\-%
% --- p. 151 ---
% p. 151, l. 7: a small stray dot stands in the left margin beside "Under den
% nyere Philosophies ..." (ABBYY renders it as a lone "*" on its own line);
% a second stands between "Maal," and "den" in l. 8, and a third between
% "sande" and "Gjenstand" in l. 33.  None is printed type; none is
% transcribed.  There is NO footnote on this page.
% p. 151, l. 32 f.: "Tankens" is letterspaced (12--16 px) but the following
% "høieste Gjenstand," measures tight (2--9 px) and is set plain, although at
% low magnification the whole phrase looks spaced.
liges og det Naturliges absolute Heterogeneitet at naae til
Modsætningernes Enhed. Der blev viist, hvorledes dens For\-%
udsætning indesluttede Nødvendigheden af, at den maatte ud\-%
vikle sig i to sideordnede Rækker, hver gaaende ud fra sit
Udgangspunkt, og der antydedes, hvorledes den bestemte den
Enheds Art, der naaedes ved disse Rækkers Afslutning.

Under den nyere Philosophies Udvikling opfatte vi altsaa
\emph{Principets Enhed} som det Maal, den \emph{med Bevidsthed}
sætter sig, ialfald saasnart den kommer til Klarhed over sig
selv, \emph{Dualismen} som det Hæmmende, der fra den tidligere
Grundanskuelse trænger ind og ikke blot som en \emph{Bevidstheds\-%
factor} gjør Enhedsbestræbelsen uklar, men som \emph{ubevidst}
indvirkende bestemmer \emph{Enheden} paa en saadan Maade, at
den bliver en \emph{skjult Form for Dualismen}.

I den \emph{idealistiske} Udviklingsrække er Forholdet let at
opfatte. \emph{Descartes} begynder i den \emph{dualistiske} Modsæt\-%
ning mellem Aand og Materie. De fæstne sig for ham til to
\emph{Substantser}, med \emph{absolut forskjellige, hinanden
udelukkende, Prædicater}. I \emph{Mennesket} mødes de, men
paa \emph{ubegribelig} Maade, og Enheden er der kun \emph{Sammen\-%
sætningens} Enhed: Mennesket «bestaaer af» Sjæl og Le\-%
geme. I \emph{Gud} mødes de atter, som i deres \emph{fælleds Grund};
men Materiens Udspring fra den \emph{kun} som Aand bestemte Gud
bliver igjen det \emph{Ubegribelige}. Gud bliver kun det \emph{anty\-%
dede} Enhedspunkt: Enheden bliver det uendeligt Fjerne, hvor
Modsætningernes parallele Linier fingeres at løbe sammen.
Først da \emph{Spinoza}, den dogmatiske Idealismes egentlige Fuld\-%
ender, gjennemtænker Descartes's halve Tanker, gjennemfører
hvad der hos ham kun var et Tilløb, faaer \emph{Enheden} en
\emph{større} Betydning, og bliver det Begreb, der behersker hele
Systemet. Den \emph{ene uendelige Substants}, hvis Væsen
omfatter \emph{hele} Tilværelsen, bliver nu, som \emph{Tankens} høieste
Gjenstand, dens eneste sande Gjenstand, og i dens Væsen gaae
\emph{Natur} og \emph{Aand}, der før bare \emph{Substantsens} Navn, ind som
Væsensbestemmelser, som \emph{Attributer}. De ere Bestemmelser
% --- p. 152 ---
% p. 152, ll. 10 f.: the word "begribes" is broken at the line end and only
% its second half is letterspaced -- "be-" at the foot of l. 10 measures tight
% (2--5 px), "gribes ved sig" at the head of l. 11 measures 9--15 px.  Set as
% printed; the compositor's own inconsistency, not regularised.
% p. 152, l. 25: a small stray dot stands between "bevare" and "Substantsens".
% Not printed type; not transcribed.
for \emph{eet og samme} Væsen; de udtrykke begge dette Væsen
\emph{heelt}; de udtrykke det i \emph{fuldkommen overensstemmende
Rækker af Modificationer}. Saaledes er \emph{Enheden sta\-%
tueret}; men indenfor den bryder \emph{Dualismen} frem saasnart
Attributernes Forhold nærmere skal fastsættes. Tanke og Ud\-%
strækning, der udtrykke Aandens og Naturens Væsen, blive
\emph{ikke} udledte af Substantsens Natur; de blive \emph{ikke} begrebne
som hver i sit Væsen indesluttende sin Modsætning; de blive
\emph{umiddelbart} satte som Substantsens uendelige Væsensbe\-%
stemmelser, men udtrykkeligt saaledes, at hver af dem be\-%
\emph{gribes ved sig}, og ikke behøver Modsætningen til at for\-%
staaes. Derfor staae de som fremmede og ligegyldige for hin\-%
anden: de uendelige Rækker af Modificationer, i hvilke de ud\-%
folde sig, blive \emph{parallele} Rækker, uden gjensidig Overgang.
Hvert Led i en af Rækkerne har vel sit tilsvarende Led i den
anden, hver Tanke (Idee) sit tilsvarende Legemlige og omvendt,
men Overensstemmelsen er ikke en Følge af et Vexelvirknings\-%
forhold mellem Rækkerne; hvert Led begribes, og begribes
heelt, af det foregaaende i \emph{samme} Række. Men saaledes blive
Natur og Aand det \emph{kun} Væsensforskjellige, deres Enhed bliver
et \emph{tomt Navn}, og det \emph{Alternativ} opstaaer: at fastholde
Attributernes Eiendommelighed, altsaa deres Forskjel, og da at
lade Substantsen, der i sit ene Væsen skal forene det \emph{kun}
Forskjellige, det ved hinanden Ubegribelige, sprænges; eller at
søge at bevare Substantsens Enhed, og da at lade Modsætnin\-%
gerne forsvinde som Modsætninger. Men i første Tilfælde er
Dualismen aabent udtalt, i sidste Tilfælde bliver Enheden kun
to tomme, indholdsløse Abstractioners Enhed, altsaa ikke En\-%
hed af Modsætningerne \emph{som} Modsætninger, men kun det Ind\-%
holdsløse, i hvilket to Indholdsløse falde sammen, og idet
Tanken nu atter vender sig mod den concrete Tilværelse, er
Modsætningen der igjen som uforsonet.

Ogsaa hos \emph{Leibnitz} bryder Dualismen frem trods den
endnu mere udprægede Stræben efter Enhed. I sit \emph{Monade\-%
begreb} har Leibnitz Begrebet om \emph{et rent Aandeligt}, et
% --- p. 153 ---
% p. 153, l. 12: only the FIRST ELEMENT of the compound is letterspaced --
% "Reflexions" measures 11--14 px between glyphs, "kategorier," 2--6 px in the
% same word.  Set as printed.  (Compare the same phenomenon at p. 152 l. 10 f.
% and p. 153 l. 15 f.)
% p. 153, ll. 15 f.: "Phænomen" is letterspaced but the following "Op-" and
% "fattedes" measure tight (Op- is borderline at 6--7 px, fattedes at 2--5),
% so "Opfattedes" is set plain.
% p. 153, l. 4: "Tilværelsen" measures tight (2--6 px) between the spaced
% "hele" and "overalt" and is set plain, although it looks spaced at low
% magnification.
virksomt Immaterielt, der i sin individuelle Substantialitet har
indre Uendelighed. Idet al Tilværelse i sin sidste Grund gaaer
tilbage til Monaderne, al Sammensætnings Elementer, er \emph{hele}
Tilværelsen idealiseret: det Aandelige er \emph{overalt}. Det \emph{Le\-%
gemlige} som saadant kommer kun frem ved Monadernes Sam\-%
mensætning, som denne Sammensætnings \emph{Phænomen}. For\-%
holdet er altsaa ikke, som hos Spinoza, et Ligeberettigelsens
Forhold; men kun det Aandelige har Væsentlighed, det Legem\-%
lige er ikke Andet end dette Væsentliges Skin; det er \emph{kun}
Phænomen, og dets Grundbestemmelser Phænomener af det
Aandeliges Forhold. Men allerede denne Bestemmelse af Grund\-%
forholdet ved \emph{Reflexions}kategorier, der forudsætte hinanden,
og paa een Gang udsige Modsætning og Enhed, og som her
skulle gjælde som \emph{definitive} Bestemmelser, antyder \emph{Dua\-%
lismens Vedvaren}. Og idet det som \emph{Phænomen} Op\-%
fattedes Grundbestemmelser \emph{ikke deduceres} af Væsenets,
men \emph{kun postuleres}, og det derfor som \emph{uafledet} hævder
sig en \emph{selvstændig} Plads ved Siden af \emph{Væsenet}, bliver den
dualistiske Adskillelse aabenbar. Paa Grund af hiin Phæno\-%
menets selvstændige Væren bliver saa hos Leibnitz det Sjælelige
og det Legemlige to \emph{sideordnede} Sphærer, der hver følger
sin Lovgivning, og \emph{Mechanismens} Bestemmelser gjøre sig
gjældende som en Magt, der ikke blot har det \emph{Legemlige}
til sit \emph{uafhængige} Omraade, men, \emph{ubegrebet}, derfra ud\-%
strækker sine Virkninger til det \emph{Aandelige}. --- Endnu mere
iøinefaldende bliver denne igjen frembrydende Dualisme i de
aandløse, uconsequente Former, i hvilke den Leibnitziske Phi\-%
losophie fortsætter sig. Idet \emph{Monaderne} miste den \emph{Leib\-%
nitziske} Grundbestemmelse: at være \emph{forestillende} Væsener,
og kun beholde \emph{Enkelthedens} Bestemmelse tilbage, saa for\-%
vandles de i Virkeligheden til \emph{Atomer}, og mellem de \emph{Mona\-%
der}, der ere \emph{det Legemliges Moleculer}, og dem, der
skulle danne Centrum i et \emph{tænkende} Væsen, og derfor gjen\-%
nem et \emph{Postulat} faae \emph{Forestillingens} Bestemmelse til\-%
bage, danner der sig en uoverskridelig Kløft.

% --- p. 154 ---
\setcounter{footnote}{0}
% p. 154 bears TWO notes, marked *) and **).  Both are self-contained --
% *) ends "Naturwissenschaft.", **) ends "satz 3, Anm." -- and p. 155 opens
% with running text ("som uudførlig for den menneskelige Erkjendelse. ---
% Ligeledes fremtræder hos Fichte ..."), so NOTHING RUNS OVER at this seam.
% p. 154, l. 5: PRINTER'S ERROR, transcribed as printed.  The page reads
% "Aarhuedrede" where "Aarhundrede" is expected; confirmed on the image at
% 1200 dpi (both OCR witnesses agree).  It is NOT one of the nine corrections
% on the author's Trykfeil leaf.
% p. 154: the German titles in the two notes are set in the SAME upright roman
% as the Danish and are not letterspaced, so they take neither \emph{} nor
% \textit{}.  "Urtheilskraft", "Anfangsgründe", "Hauptstück" all read on the
% image (ABBYY garbles them as "Uitlieilskraft", "Aufangsgründe").
% p. 154, l. 17: the guillemets round «legemlige Naturs» point inward; only
% the first word inside them is letterspaced.
% p. 154, ll. 22 f.: "Iogforsigværendes" is set solid (letterspaced), broken
% at the line end after "Iogforsig-"; the hyphen is the line-break hyphen and
% is resolved here.
Og hvad der gjælder om de Former, der fremtræde som
\emph{Fortsættelser} af Monadelæren, gjælder ligeledes om de
andre, deels \emph{uconsequent eklektiske}, deels paa den \emph{al\-%
mindelige}, i sin Grund \emph{dualistiske}, Bevidstheds Udsagn
hvilende Former af Idealismen i det 18de Aarhuedrede. I dem
alle træder Dualismen tydelig frem.

I den \emph{kritiske} Idealisme træder i Begyndelsen Problemet
om Forholdet mellem Tilværelsens Grundformer tilbage, \emph{Er\-%
kjendelsesproblemet} indtager den fornemste Plads. Og
dog røber sig hos \emph{Kant}, trods denne Stilling af Problemerne,
og trods hans kritiske Sky for at bestemme Noget om det
Objectives Væsen i og for sig, den gamle Opfattelse af det
Naturlige, Materielle, som det Aanden modsatte Passive\footnote{Jfr. Kritik der Urtheilskraft og Metaphysische Anfangsgründe der
Naturwissenschaft.}. Vel
har \emph{Naturbegrebet} hos Kant en mere udstrakt Betydning
end i den almindelige Sprogbrug og i den tidligere Philoso\-%
phie, og omfatter \emph{alle} lovmæssigt ordnede \emph{Phænomener}, ikke
blot dem, der henhøre til den «\emph{legemlige} Naturs» Sphære, saa
at altsaa Naturens Navn ikke directe udtrykker Aandens gamle
Modsætning. Men den \emph{legemlige, materielle Natur, Natu\-%
ren i snævrere} Betydning, afgiver dog Synspunktet for Opfattel\-%
sen af Naturens \emph{almindelige} Begreb i dets Modsætning til det
fra det Aandeliges (Intelligibles) Synspunkt opfattede \emph{Iogforsig\-%
værendes} Begreb, og, trods Kants \emph{dynamiske} Opfattelse af
Materien, bliver dog, ifølge den samme Inconsequents som i
Naturvidenskaben efter Newton, den gamle, paa Dualismen
hvilende, Opfattelse af det Legemlige som det Passive,
\emph{kun Bevægelige}, bevaret hos ham\footnote{Jfr. navnlig Metaph. Anfangsgr. d. Naturw., 3. Hauptstück, Lehr\-%
satz 3, Anm.}. \emph{Enheden} mellem
Naturen og Aanden, mellem deres dobbelte Lovgivning, \emph{an\-%
tyder} Kant vel i sit sidste Hovedværk, men det bliver \emph{kun}
en Antydning, der ikke udføres, og som udtrykkelig sættes
% --- p. 155 ---
\setcounter{footnote}{0}
% SEAM p.154/155 CHECKED (BATCH-AGENT §5.1).  The note area of p. 155 opens
% with a MARK -- "*) Man sammenligne Herbarts ..." -- so nothing runs over
% from p. 154, confirming the note already standing at the p. 154 marker.
% The body of p. 155 opens mid-sentence with running text continuing p. 154's
% "... men det bliver kun en Antydning, der ikke udføres, og som udtrykkelig
% sættes / som uudførlig for den menneskelige Erkjendelse."  No paragraph
% break at the seam.
% p. 155 bears ONE note, *), self-contained (ends "Monade og Atom.").
% p. 155: no Greek, no ɔ:, no guillemets, no pencil.
som \emph{uudførlig} for den \emph{menneskelige} Erkjendelse. --- Lige\-%
ledes fremtræder hos \emph{Fichte}, hos hvem Tankens Problemer
endnu mere reduceres, trods Bestræbelsen for at lade Alt frem\-%
gaae af Jegets Proces, den tidligere Philosophies \emph{dualistiske}
Modsætning som en \emph{Dualisme indenfor Jeget}, hvor
\emph{Naturskranken}, det Jegets Uendelighed paa een Gang Be\-%
tingende og Begrændsende, bliver staaende som et \emph{Uafledet}
og som et \emph{Ubegribeligt}. Hos \emph{Schelling} træder hint Pro\-%
blem om \emph{Tilværelsens Enhedsgrund} igjen frem i For\-%
grunden, og den søges \emph{i det Absolute}, der er \emph{udeelt} i sine,
kun ved en \emph{quantitativ} Differents bestemte Væreformer.
Men ligesom hos \emph{Spinoza} blive disse Former ingensteds de\-%
ducerede, og den absolute Enhed bliver det Tomme, i hvilket
Modsætningernes Virkelighed forsvinder. Hos \emph{Hegel} bliver
Enheden atter, som hos \emph{Leibnitz}, søgt gjennem en \emph{Under\-%
ordning}; men Forholdet bestemmes ved høiere Kategorier:
ved \emph{Ideebestemmelser}. Naturen bliver Moment i Ideens
evige Selvrealisationsproces; i sin Modsætning til Aanden er
den evigt ophævet i denne; Modsætningen er evigt tilbageført
til Enhed. Men, som vi allerede have seet, bliver Enheden
kun \emph{Tankens Enhed med sig selv i sin Abstraction},
ikke med \emph{Værens concrete Virkelighed}, og \emph{Dualismen}
bryder frem idet Naturphilosophien sætter Væren i dens Rea\-%
litetsbestemmelse som det med Begrebet Incongruente, og der\-%
igjennem statuerer \emph{et existerende Ufornuftigt}, der \emph{ikke}
begribes ud fra \emph{Ideens Enevirkelighed}.

Hos \emph{Herbart} endelig, der, medens han knytter sig til
% p. 155, l. 28: PRINTED "Begrunder", with a plain u and NO umlaut, read at
% 1200 dpi.  The ABBYY layer gives "Begründer" here and again at p. 156;
% tesseract gives "Begrunder" both times.  The image settles it: no umlaut.
den \emph{kritiske} Idealismes Begrunder, og ved sit Grundbegreb
bliver indenfor Idealismens Kreds, forfølger sin eiendommelige,
fra den efterkantiske Idealismes Repræsentanters afvigende, Ret\-%
ning, see vi \emph{Leibnitzes} Bestemmelser af Forholdet gjentage
sig, om end i en forandret Belysning og med Reminiscentser
af Monadelærens Fortsættelser\footnote{Man sammenligne \emph{Herbarts} Opfattelse af Begreberne Sjæl og En\-%
kelt-Reale med Begreberne: forestillende Monade og Atom.}. Den \emph{naturlige} Tilværelses
% p. 155/156: the letterspaced run ENDS at "Tilværelses"; "Grundbestemmelser"
% at the head of p. 156 is set SOLID (checked at 1200 dpi against the spaced
% "Phænomen" on the same line).  The phrase is therefore broken by the print,
% not by this transcription.
% --- p. 156 ---
Grundbestemmelser sættes ogsaa her som \emph{Phænomen} af det
virkeligt Værende, men deduceres \emph{kun tilsyneladende}, og
tiltvinge sig saaledes en Selvgyldighed og Uafhængighed.

I den \emph{realistiske Rækkes} Udvikling bliver \emph{Dualis\-%
mens} første Indvirkning utydeligere, og den har derfor und\-%
draget sig Historikernes Opmærksomhed, men ved en omhyg\-%
geligere Analyse blive dens Spor kjendelige strax fra Begyn\-%
delsen. Under Udfoldelsen af den første Paavirknings Conse\-%
quentser bliver saa Dualismens Indflydelse mere og mere iøine\-%
faldende, og idet Udviklingsrækken har naaet sin Afslutning,
viser dens endelige Resultat sig som væsentlig afhængende
af den.

\emph{Hos Baco}, den realistiske Philosophies Begrunder, træder,
efter hans Tænknings hele Retning, Problemet om \emph{Erkjen\-%
delsens Methode} frem foran Problemet om \emph{Tilværelsens
Princip}. Men idet Spørgsmaalet om Erkjendelsesprincipet
og Spørgsmaalet om Tilværelsesprincipet i Realiteten \emph{er det
samme}, kun seet fra \emph{forskjellige} Sider, og idet den \emph{de\-%
finitive} Besvarelse af hint forudsætter en Besvarelse af dette,
saa kan Baco, idet han udvikler sin \emph{Erfaringstheorie}, ikke
undgaae at berøre Problemet om \emph{Tilværelsens Grund}, og
i hans Bestemmelser derom lader et Spor af \emph{Dualismens}
Indvirkning sig paavise, der saa fra dette Punkt kan forfølges
i hans Philosophies Retning og i hans Erkjendelseslæres Ud\-%
førelse.

Naar vi ville finde Bacos \emph{sande} Mening angaaende Til\-%
værelsens Grund, maae vi ikke tage hans Yttringer om hvad
han gjør til \emph{Theologiens} Gjenstand efter Bogstaven. Disse
Yttringer ere \emph{kun accommoderende}, kun beregnede paa at
beskytte hans Philosopheren mod theologiske Indgreb. Theo\-%
logien faaer, efter disse Udsagn, \emph{Aabenbaringen og det
Overnaturlige} til sit Gebet; Philosophien opgiver villigt
dette Gebet, og undgaaer derved, indenfor sit Omraade, \emph{Dua\-%
lismen mellem Overnaturligt og Naturligt}, da den,
idet den hævder sig \emph{det Naturlige}, gjør hint til et \emph{for sig}
% p. 156/157 seam: the letterspaced run "for sig ikke Existerende" is broken
% by the page break; it is closed here and reopened at the head of p. 157.
% One emphasis in the print, two \emph{} groups here.
% --- p. 157 ---
\setcounter{footnote}{0}
% p. 157 bears ONE note, *), self-contained (ends "selvgyldige Natur.").
% p. 157: no Greek, no ɔ:, no guillemets, no pencil.
% p. 157: the LATIN is set in the same upright roman as the Danish and takes
% \emph{} only where it is letterspaced, never \textit{} (transcription.tex §2).
% Read at 1200 dpi:  "spiritus innatus" -- BOTH words letterspaced;
% "tangibile" -- letterspaced;  "spiritus" at l. 18 -- SOLID;  and in
% "spiritus aëreus aut igneus" only "aëreus" and "igneus" are letterspaced,
% while "spiritus" and "aut" are set solid.  The word is "aëreus", with a
% DIAERESIS over the e (ABBYY reads "aereus", tesseract "aéreus"; the two
% dots are unambiguous at 1200 dpi).
\emph{ikke Existerende}, til et Saadant, som den \emph{factisk} negerer.
Men \emph{i Naturen} bliver \emph{Materien} Grundbestemmelsen, saa\-%
ledes opfattet, \emph{at den sættes som oprindelig havende
Form og Bevægelse i sig}. I dette Grundbegreb er der
\emph{tilstræbt en Enhed af det i Naturen}, Philosophiens \emph{ene\-%
ste Gebet, aldrig adskilte Aandelige og Legemlige},
\emph{og i det søges en oprindelig Kilde til den Virksomhed},
ved hvilken \emph{Naturformernes og Naturtingenes Mang\-%
foldighed} udfolder sig. Men Modsætningen, der \emph{tilsyne\-%
ladende} forsonedes ved Opfattelsen af Materiens Begreb, træder
atter frem \emph{indenfor} dette. I Materien foregaaer der ligesom
% p. 157, l. 12: PRINTED AS SHOWN.  "en Deling i et" is letterspaced and
% "dobbelt Moment:" is set SOLID in the middle of the run, then the spacing
% resumes at "det formende".  Measured at 1200 dpi: the d-o gap in "dobbelt"
% is 2 px against 11--13 px in "Deling" on the same line.  The compositor
% broke the run, not the transcription.
\emph{en Deling i et} dobbelt Moment: \emph{det formende, bevæ\-%
gende Princip}, og det, der er dette \emph{Virksommes Stof (Sub\-%
strat)}. Hint betegnes som \emph{spiritus innatus}, saaledes at
Navnet viser hen paa den gamle Modsætning mellem det \emph{Aan\-%
delige som det Active og det Materielle som det Pas\-%
sive}; dette betegnes som \emph{tangibile}: det Haandgribelige, der
er Virksomhedens Materiale. Vel bestemmes igjen spiritus
\emph{som materiel}, som en spiritus \emph{aëreus} aut \emph{igneus}, saa at
altsaa Sondringen ophæves \emph{i Navnet}, men kun forat lade en
\emph{uforklaret dualistisk Væsenssondring} blive sat \emph{in\-%
denfor det Materielle}\footnote{Det er interessant at sammenligne \emph{Baco} med \emph{Giordano Bruno}
med Hensyn til Begges Opfattelse af \emph{Materien} og \emph{det Naturlige}.
\emph{Hos Bruno} finde vi en \emph{idealistisk} Opfattelse med Retning hen\-%
imod \emph{Naturconcretionen}: Monaderne, idealistiske Momenter, blive
\emph{Rumsbestemmelsernes} Grundlag. \emph{Hos Baco} finde vi et \emph{realistisk
Udgangspunkt} med \emph{uvilkaarlig} Optagen af \emph{Idealitetsbestem\-%
melsen}; det Ideelle reduceres til Bestemmelse af det Materielle.
\emph{Hos Bruno} gaaer, med Hensyn til dette Problem, Tankens Bevægelse
\emph{mellem Gud og Natur}, men overalt paa \emph{Idealitetens} Basis,
\emph{hos Baco} gaaer den mellem \emph{Materie og Aand} paa \emph{Realitetens}
Basis. \emph{For Bruno} er Naturen den \emph{anthropomorphiserede,
idealiserede} Natur, \emph{for Baco} er den den \emph{Erfaringen be\-%
stemmende, i sin Materialitet selvgyldige} Natur.}. \emph{Fra dette Materielle ude\-%
lukkes saa Aandens Væsensbestemmelse: Øiemedet},
% --- p. 158 ---
% p. 158: no notes, no Greek, no ɔ:, no pencil.
og idet Begrebet saaledes bliver \emph{snævrere end Virkelig\-%
hedens Omfang}, blive Væreformer staaende, med Hensyn
til hvilke \emph{Baco idetmindste ikke paaviser} en Mulighed til,
at de kunne føres ind under hiin Grundbestemmelse af det
Værende, f. Ex. \emph{Forstand og Villie}. Der er en \emph{Tvetydig\-%
hed} hos Baco hvor han taler om disse. Han skjelner i Men\-%
nesket mellem den \emph{fornuftige} og den \emph{sandselige} Sjæl.
Hiin skal være inspireret paa \emph{overnaturlig} Maade, og der\-%
for være Gjenstand for \emph{Theologien}; denne skal være en
\emph{legemlig} Substants, der har sin Plads i Hjernen, og kun er
usynlig paa Grund af sin Fiinhed. Dens Kræfter skulle være
% p. 158, l. 12: PRINTER'S ERROR, transcribed as printed.  The page reads
% "cg" where "og" is expected.  Read at 2400 dpi: the sort has a wide, clean
% aperture on the right at mid-height with two finished terminals -- a `c' --
% against the closed ring of the certain `o' in "Fornemmelse" 3 mm to its
% right on the same line; it is not an under-inked or worn `o' (the stroke is
% unbroken and as dark as its neighbours).  BOTH OCR witnesses independently
% read "cg".  Not one of the nine corrections on the author's Trykfeil leaf.
\emph{vilkaarlig} (spontan) \emph{Bevægelse} cg \emph{Fornemmelse}: Kræf\-%
ter, der søges \emph{overalt i det} Materielle. \emph{Villie og Forstand
sættes ikke} som dens Kræfter, men der siges \emph{ikke heller},
at de ere den \emph{fornuftige} Sjæls Kræfter, men derimod, med
% p. 158, l. 16: the guillemets round «menneskelige Sjæls» are the normal
% INWARD pair, « opening and » closing, read at 1200 dpi.  (The ABBYY layer
% gives the closing mark as «; the image does not.)  Only the words inside
% them are letterspaced.
et tvetydigt Udtryk, \emph{at de ere den} «\emph{menneskelige Sjæls}».
% p. 158, l. 18: PRINTED "Realitet".  The ABBYY layer reads "Idealitet";
% tesseract and the image both read "Realitet".
Bacos egentlige Tanke kan, idet det \emph{Overnaturlige ikke}
har Realitet for Philosophien, kun være den, at see \emph{Forstand
og Villie}, der for Bevidstheden ere givne som \emph{Kjendsgjer\-%
ninger}, som den \emph{sandselige} Sjæls Bestemmelser, men han
udtaler sig kun halvt, og \emph{paaviser ikke Væsenssammen\-%
hængen}. Vel have Forstand og Villie en \emph{Analogie} med For\-%
nemmelse og vilkaarlig Bevægelse, og med \emph{det}, der er disses
Form i de \emph{tilsyneladende livløse} Ting: Perception og Til\-%
træknings- og Frastødningskraft; men, ligesaalidt som Baco
godtgjorde \emph{Muligheden} af, at det Materielle kunde have For\-%
nemmelse og Vilkaarlighed, ligesaalidt \emph{gjennemfører} han de
aandelige Functioners Analogie med Materiens Kraftyttringer,
endsige udvikler den til \emph{Væsensidentitet}, og, trods de Præ\-%
dicater, han tillægger det Materielle, bliver der altsaa en \emph{dua\-%
listisk} Sondring tilbage.

Efterat vi saaledes hos Baco have fundet Dualismens
skjulte Indvirkning i hans Opfattelse af Tilværelsesprincipet,
søge vi den nu ogsaa i Bestemmelsen af hans \emph{Erkjendel\-%
seslæres} Centralpunkt: i Opfattelsen af \emph{Erfaringsprinci\-%
% --- p. 159 ---
\setcounter{footnote}{0}%
% (the % above eats the newline: "Erfaringsprinci-" / "pet" is a word broken
% across the PAGE break and must not acquire a space -- PLAYBOOK §6.)
% p. 159 bears ONE note, *), self-contained (ends "Charakterer).").
% p. 159 carries the ɔ: id-est mark ONCE, at "af dette ɔ: en Opfattelse";
% read on the image at 1200 dpi (both witnesses give "o:").  No Greek.
pet}, og see, hvorledes derigjennem den realistiske Rækkes
\emph{hele} Udvikling bliver betinget.

Idet \emph{Erfaringen} af Baco sættes som den sande Erkjen\-%
delses Kilde, er der hermed \emph{oprindelig} ikke sat nogen Ind\-%
skrænkning af den til et bestemt Gebet af det Værende. Baco
vil \emph{oprindelig}, hvad der stemmer overens med hans Stræben
efter at føre det Værende tilbage til en Enhed, lade den faae
\emph{alt det Værende}, det, der betegnes som det Aandelige, lige\-%
saafuldt som det, der kaldes det Legemlige, til Gjenstand (Nov.
Org. I, 127); han vilde altsaa, ligesom senere \emph{Kant}, ogsaa
kunne lade \emph{den menneskelige Aand i dens Selvvirk\-%
somhed}, med hvad der \emph{oprindelig fremgaaer af denne},
blive Object for Iagttagelsen. \emph{Men denne Stræben hæm\-%
mes øieblikkelig}: Erfaringen tager \emph{strax} sin Retning \emph{mod
det Physiske i snævrere} Forstand, imod det \emph{Ydre}, det
\emph{Sandseligt-Naturlige}, og vender sig kun imod \emph{Aanden}
forsaavidtsom den betragtes fra sin \emph{Naturside}. Aanden be\-%
tragtes kun som «\emph{sandselig Sjæl}», eller, forsaavidt \emph{For\-%
stand og Villie}, med den ovenfor berørte Tvetydighed, sæt\-%
tes som den «\emph{menneskelige Sjæls}» Evner, og der tages
Hensyn til dem, saa bliver det dog kun til de \emph{phænome\-%
nale Yttringsformer}\footnote{I Forhold til dette Gebet er Baco ikke \emph{begribende}, men \emph{beskri\-%
vende} (Lidenskaber, Charakterer).}, og hverken deres Grundformer,
deres Kilde eller deres Love undersøges. Og idet Erkjendel\-%
sens Retning saaledes væsentlig gaaer \emph{udefter} mod det givne
Naturphænomen, og Opgaven bliver den \emph{objective Opfat\-%
telse} af dette ɔ: en Opfattelse \emph{ved Objectet selv}, saa træ\-%
% p. 159, l. 30: PRINTED AS SHOWN.  In "Virksomhedsbestemmelse" only the
% first element is letterspaced: "Virksomheds" carries 9--12 px letter gaps
% at 1200 dpi and "bestemmelse" 1--3 px, the same as the certainly solid
% "tilbage" that follows it on the line.  A compound broken between a spaced
% and a solid setting inside one word.
der hos Baco den \emph{Virksomheds}bestemmelse tilbage, som
han directe og indirecte havde tillagt Aanden, og Aandens
Forhold til Erkjendelsesgjenstanden: det Erfarede, antager ube\-%
vidst Charakteren af \emph{et blot modtagende}, hvorfor ogsaa
Baco kan udtale det Haab, at hans Methode skal være istand
% --- p. 160 ---
% p. 160: no notes, no Greek, no ɔ:, no pencil.
til at udjevne Forskjellen mellem Aanderne og gjøre alle lige
skikkede til at modtage og erkjende Sandheden.

I denne \emph{Begrændsning af Erfaringen til det Na\-%
turlige}, Materielle, og i denne \emph{Opfattelse af Erkjendel\-%
sen} under \emph{Passivitetens} Bestemmelse, med Tilbagetrængen
af det Aprioriske, røber \emph{Dualismens Indvirkning} sig. Thi
ligesom hiin Begrændsning til det Naturlige, der er en Ude\-%
lukken af det Aandelige, som dog ikke kan fornægtes, hviler
paa en \emph{dualistisk} Sondring mellem de to Gebeter, saaledes
viser ogsaa Erkjendelsens Passivitet hen til det Samme. Det
at Aanden ligesom seer bort fra sig selv i sin indre Energie,
kan vel være \emph{understøttet} af Aandens Mistillid til det Ind\-%
% p. 160, l. 13: the guillemets round «udspinder af sig selv» are the normal
% inward pair; nothing inside them is letterspaced.
hold, den «udspinder af sig selv», og af dens Trang til en
ydre, objectiv Control for Tænkningen, men den er i sidste
Instants \emph{begrundet} deri, at, under Bestræbelsen for at føre
Alt tilbage til \emph{Natur}, en \emph{abstract Modsætning} først
uvilkaarlig gjør sig gjældende mellem det \emph{Aandelige} og det
\emph{Naturlige}, og saa \emph{Enhedsbestræbelsens Reaction}
ubevidst gjør \emph{det Sidstes} ensidige Bestemmelse gjældende
\emph{som almeen}. I hiin Begrændsning af Erfaringen og i hiin
Opfattelse af Erkjendelsen ligger saa Spiren til den \emph{hele} Ud\-%
vikling, gjennem hvilken \emph{Erfaringsphilosophien} omformer
sig, først til \emph{Empirisme}, og saa til \emph{Sensualisme} og \emph{Ma\-%
terialisme}.

% p. 160, l. 32: "tabula rasa" is LATIN set in the same upright roman as the
% Danish and is NOT letterspaced here; it therefore takes no markup at all
% (transcription.tex §2).  Checked at 1200 dpi against the letterspaced
% "modtagende" three words before it.
Den \emph{Tilbagetræden af Activiteten i} Erkjendelsen,
der viste sig hos \emph{Baco}, fører først til, at \emph{Erfaringen},
bestemt som det \emph{modtagende} Forhold til Erkjendelsesobjec\-%
tet, gjennem hvilket Aanden, der oprindelig sættes som en
tabula rasa, faaer sit hele Indhold, opfattes som \emph{eneste Er\-%
kjendelseskilde}. I denne Bestemmelse ligger Polemiken mod
\emph{al apriorisk Erkjendelse}, og den betegner \emph{Empiris\-%
mens} Standpunkt, der netop ved dette polemiske Forhold be\-%
grebsmæssigt adskiller sig fra \emph{Erfaringsphilosophien}, for
hvilken Erfaringen var \emph{væsentlig}, men endnu \emph{ikke ude\-%
% --- p. 161 ---
% p. 161: no notes, no Greek, no pencil.
% p. 161 carries the ɔ: id-est mark THREE times: "det Ydre ɔ: det Legemlige",
% "mechanisk bestemt Ydre ɔ: modtagende sit Indhold", and "uden de sandsede
% ɔ: materielle, Objecter".  All three read on the image at 1200 dpi; ABBYY
% gives "o:" twice and drops the third, tesseract gives "o;", "o:" and "»:".
% The sheet signature "11" stands at the foot of p. 161; not transcribed.
lukkende} Erkjendelseskilde. Men i den \emph{første} Consequents
ligge \emph{nye} Consequentser. Erfaringen, der \emph{blot er modta\-%
gende}, forholder sig til en Gjenstand, der kun er bestemt
som \emph{Bevidsthedens Andet}, idet selv den indre Erfarings
Gjenstand: \emph{Bevidsthedstilstanden}, opfattes som et \emph{udenfra
Bestemt og udenfra Bevirket}. Erkjendelsens \emph{primitive Form}
bliver da den \emph{sandselige Fornemmelse}, ved hvis \emph{uvil\-%
kaarlige} Omformning de andre Bevidsthedsformer blive til,
og Erkjendelsens \emph{primitive og virkelige Indhold} bliver
% p. 161, ll. 10 f.: "Begrændsede" straddles the line break and the print
% straddles the emphasis with it -- "er det Be-" is set SOLID at the crowded
% line end (1--2 px letter gaps) while "grændsede," at the head of the next
% line is letterspaced (9--11 px), as are "Enkelte," beside it.  The word is
% set here as one \emph{} group, to keep the word whole; the split is
% recorded rather than reproduced.
det \emph{Ydre} ɔ: \emph{det Legemlige}, der, som Legemligt, er det \emph{Be\-%
grændsede}, \emph{Enkelte}, og nærmere, \emph{det mechanisk Bevæ\-%
gede}. Det \emph{aandelige} Element, der hos \emph{Baco} blev tilbage
i det rigtignok \emph{kun postulerede} Begreb: den \emph{besjælede
Materie}, maa nemlig forsvinde fra det, der bliver det \emph{egent\-%
lige} Materielle, idet Dualismen gjør sig gjældende indenfor
Materien, og, idet det forsvinder, kan dette Materielle kun
sættes som \emph{det mechanisk Bevægede}, der \emph{ikke} har Virk\-%
somhedens (Bevægelsens) Udgangspunkt i sig. Ved hiin Op\-%
fattelse af den sandselige Fornemmelses Forhold til de andre
Bevidsthedsformer, er Empirismen udviklet til \emph{Sensualisme};
men Sensualismen gjemmer atter i sig \emph{Materialismens}
Consequents. For Sensualismen var Aanden \emph{kun receptiv},
nærmere: receptiv i Forhold til et \emph{mechanisk bestemt
Ydre} ɔ: modtagende sit Indhold, sine Bestemmelser, gjennem
\emph{Indtryk paa Sandserne, Receptivitetens Organer}.
For denne \emph{dogmatiske} Sensualisme bliver der Intet tilbage
uden \emph{de sandsede} ɔ: materielle, Objecter og det \emph{sandsende}
Subject, der, i sin blotte Receptivitet, forholder sig til Gjen\-%
standen \emph{som Ting til Ting}, og bliver et \emph{afhængigt Led
i Tingenes Kjæde} idet det \emph{passivt} modtager Indtrykkene
af de mechanisk forbundne endelige Sandsegjenstande, hvilke
Indtryk saa i Bevidstheden, som i en Beholder, bundfældes som
Begreber, i hvilke der ikke er et Selvstændigt, der viser hen
paa en ideel Region, men i hvilke kun det \emph{Sandselige} er
% --- p. 162 ---
\setcounter{footnote}{0}
% p. 162 bears ONE note, *), self-contained (ends "Causalsammenhængen.").
% p. 162: no Greek, no ɔ:, no guillemets, no pencil.
som det \emph{Reale}. Men idet Subjectet bliver mechanisk bestemt
\emph{Led i Tingenes Række} maa det gaae ind under Rækkens
\emph{Fælledsbestemmelse}; det bliver \emph{sandseligt, materielt
Subject}, og \emph{Alt i Tilværelsen}: det Subjective og det Objective,
er nu først tilbage til den \emph{samme Enhed}: \emph{til Natur, be\-%
stemt som passiv Materie}, og dette Begreb af Materien er
\emph{blevet det altomfattende, der i sig absorberer det
Aandelige}.

Dette er den Udviklingsgang, den realistiske Rækkes Ud\-%
gangspunkt indeslutter. Den sidste Consequents kan kun und\-%
gaaes, naar der fra Empirismens Standpunkt gaaes i \emph{for\-%
mel-idealistisk} eller \emph{skeptisk} Retning; naar der enten,
idet alle vore Ideer sættes som Ideer om \emph{enkelte Ting}, idet
Tingene opløses i \emph{Egenskaberne}, og disse i \emph{Perceptioner},
kun statueres perciperende Aander med deres Bevidsthedsind\-%
hold: Ideerne, uden et tilsvarende Legemligt; --- eller naar den
som \emph{enegyldig} satte Erfarings Begrændsning til det \emph{Enkelte} og
\emph{Phænomenale} sees som udelukkende Erkjendelsens \emph{Almeen\-%
gyldighed}, \emph{Nødvendighed} og \emph{Objectivitet}, saa at \emph{Sand\-%
synligheden} afløser \emph{Sandheden} og Tanken holdes svævende i
en \emph{Tvivl}, der \emph{Intet} bestemmer om det \emph{Værendes Væsen}. Disse
Consequentser, som Sensualismen og Materialismen slippe forbi,
idet de med Forsæt lukke Øinene for deres Motiver, bestemme
Leddene i den Forgrening med \emph{negativt} Slutningspunkt, der
\emph{fra Lockes Empirisme} skiller sig fra den \emph{sensualistisk-%
materialistiske} Forgrening.

\emph{Schematisk} kan Udviklingsgangen i den \emph{nyere rea\-%
listiske Philosophie}\footnote{Vi forbigaae i denne Oversigt \emph{Hobbes}, trods hans Betydning som
Tænker, og trods de vidtgaaende Consequentser, han allerede har
draget, fordi hans Indvirkning paa de efterfølgende Systemer kun
har været underordnet, og fordi vi her nærmest have Blikket fæstet
paa Causalsammenhængen.} angives saaledes:
% --- p. 163 ---
% =====================================================================
% SCHEMA-TODO p.163 -- see RECON.md §3.  Words captured; the construction
% (the vertical stems, the crossbar, the column positions) is set in a
% later pass.  THIS MUST NOT BE LEFT AS RUNNING TEXT.
%
% A GENEALOGICAL TREE BUILT OUT OF RULES, standing at the head of p. 163
% between the "angives saaledes:" of p. 162 and the paragraph "Medens vi
% i den Hume'ske ..." below.  It is set to the ORDINARY body measure
% (x 85--1203 at 300 dpi, centre 644), not to a wider one.
%
% ALL SIX LABELS ARE ORDINARY UPRIGHT ROMAN at body size -- NOT bold, NOT
% letterspaced, NOT small caps.  Checked at 1200 dpi: "Empirismen (Locke)."
% has 1--2 px letter gaps and the same stem weight as the body text on the
% same page, i.e. it is NOT a display head.  (The recon pass's text-layer
% sweep listed it as a heading; that is WRONG.)  "Système" carries a grave
% accent; it is French set in the same roman as the Danish and, being
% unspaced, takes no markup (transcription.tex §2).
%
% What is printed, exactly, and where.  Measurements are in pixels of the
% 300 ppi scan (1 px = 0.085 mm); x is measured from the left paper edge.
%
%   line 1  "Erfaringsphilosophien (Baco)."
%           x 373--873, centre 623 -- centred on the measure.
%   rule A  a VERTICAL hairline, x 631--632, y 318--366, i.e. 49 px
%           (4.1 mm) long, centred under line 1 and hanging free: it
%           touches nothing at either end.
%   line 2  "Empirismen (Locke)."
%           x 446--802, centre 624 -- centred on the measure, directly
%           under line 1.
%   rule B  a VERTICAL hairline, x 632--633, y 432--481, 50 px (4.2 mm),
%           centred under line 2.  Its FOOT MEETS THE CROSSBAR (rule C):
%           the two are separated by a single blank pixel row.
%   rule C  the CROSSBAR: one long HORIZONTAL hairline, y 483--488 (its
%           densest rows 486--487), x 290--897, i.e. 607 px = 5.1 cm.
%           It has NO turned-down ticks at its ends; it simply stops.
%           The stem (rule B) meets it at x 632, which is 36 px RIGHT of
%           the crossbar's own midpoint (x 594) -- the stem is centred on
%           the head above it, not on the bar.  The bar sits very slightly
%           askew in the forme: its left terminus is at y 488 and its
%           right at y 482.
%   line 3  TWO labels side by side, on one line, y 498--543:
%             left   "Sensualismen (Condillac)."          x  87-- 523
%             right  "Den formelle Idealisme (Berkeley)." x 595--1204
%           The crossbar's LEFT terminus (x 290) stands over the left
%           label's centre (x 305) and its RIGHT terminus (x 897) over the
%           right label's centre (x 899).  There is no rule dropping from
%           the crossbar to this line.
%   rule D  a VERTICAL hairline, x 290--291, y 557--606, 50 px (4.2 mm),
%           below the LEFT label and above "Materialismen".
%   rule E  a VERTICAL hairline, x 897--898, y 553--601, 49 px (4.1 mm),
%           below the RIGHT label and above "Skepticismen".
%           D and E hang free: neither touches the crossbar above it.
%   line 4  TWO labels side by side, y 613--659:
%             left   "Materialismen (Système"             x  84-- 496
%             right  "Skepticismen (Hume)."               x 711--1089
%   line 5  the TURNOVER of the left label only, y 670--712:
%             "de la nature)."                            x 289-- 528
%           It is NOT centred under its first line (its centre is x 408
%           against x 290 for "Materialismen (Système"); it is indented
%           about 205 px (1.7 cm) from the left label's left edge.
%
% In reading order the six labels are therefore:
%   "Erfaringsphilosophien (Baco)."
%   "Empirismen (Locke)."
%   "Sensualismen (Condillac)."   "Den formelle Idealisme (Berkeley)."
%   "Materialismen (Système de la nature)."   "Skepticismen (Hume)."
%
% NOT PRINTED TYPE, therefore not transcribed: a short thin stroke stands
% immediately after the full stop of "de la nature)." at x 529--552,
% y ~697, on the baseline.  It is modern pencil, not an em-dash: its
% darkest pixel is grey 152 against a paper of 203 (49 levels below),
% where the full stop 4 px to its left reaches grey 31 (172 levels below)
% and the type on the line 31--56.  It tapers and is not straight.
% tesseract reports it as "-"; ABBYY ignores it.  (BATCH-AGENT §7.)
% =====================================================================
% p. 163: no notes, no Greek, no ɔ:, no guillemets.
% The sheet signature "11*" stands at the foot of p. 163; not transcribed.

\emph{Medens vi i den Hume'ske Skepticismes} Opløsning
af Erkjendelsen see den \emph{realistiske} Philosophies \emph{Ensidig\-%
hed} træde frem i en Consequents, der har viist sig bestemt
ved \emph{Dualismens Indvirkning} paa Opfattelsen af \emph{Erkjendel\-%
sesprincipet}, see vi i \emph{Materialismens} Grundbegreb \emph{Ensi\-%
digheden} træde frem i den ved en Indvirkning fra samme Kilde
betingede \emph{Opfattelse af Tilværelsens Princip}, og idet
vi nærmere analysere Bestemmelsen af Grundbegrebet, vil
det eiendommelige Forhold mellem den \emph{bevidst} tilstræbte
\emph{Enhed} og den \emph{ubevidst} indtrængende \emph{Dualisme} blive
tydelig.

Det Begreb, til hvilket i Materialismen \emph{Alt} førtes tilbage,
var \emph{Naturens Begreb}, saaledes opfattet, at det blev ligt med
\emph{den uendelige Kjæde af det mechanisk bevægede
Materielle}. \emph{Den evige og evigt bevægede Materie er
det Eneværende}: \emph{Aanden} er, som \emph{individuel}, kun en
Existentsform \emph{indenfor} Materiens Grændser, kun en \emph{bestemt
Art af Moleculernes} Bevægelse; som \emph{guddommelig} Aand
\emph{er den en tom og skadelig Illusion}. \emph{Naturens Bevæ\-%
gelseslove} beherske \emph{Alt}, ogsaa det, vi pleie at føre ind un\-%
der ganske andre Bestemmelser: \emph{Ethiken er kun anvendt
Physik}, Lægekunsten det vigtigste ethiske Middel.

Her er den skarpeste Modsætning naaet til \emph{Middelal\-%
derens} Livsanskuelse: en Modsætning, der røber sig i Mate\-%
rialismens lidenskabelige Had til den christelige Religion. Og
dog viser sig i selve dens Grundbegreb, trods Bestræbelsen for
% --- p. 164 ---
% p. 164: no notes, no Greek, no ɔ:, no guillemets.
% The printed folio reads 164; the ABBYY layer garbles it as "1G4"
% (pagemap.py gives PDF 173).
% p. 164: "Système de la nature" IS letterspaced here (broken "Sy-/stème"
% across the line), and therefore takes \emph{}.  On pp. 165 and 168 the
% same phrase is set SOLID and takes nothing.  All three read at 1200 dpi.
at hævde dets Gyldighed som \emph{altomfattende}, Eftervirknin\-%
gen af hiin Livsanskuelse gjennem den \emph{Dualisme}, der var
charakteristisk for den. Denne Dualisme fandt sit \emph{begrebs\-%
mæssige Udtryk} ved den \emph{nyere idealistiske Philoso\-%
phies Begyndelse}, idet Aand og Materie bestemtes ved \emph{ab\-%
solut modsatte}, hinanden \emph{udelukkende}, Prædicater, og
\emph{Materiens Væsen} udtømtes ved \emph{Udstrækning og Pas\-%
sivitet}. \emph{Denne Opfattelse gjør sig i Materialismen
endnu gjældende efterat Materiens Modsætning er
fornægtet}. \emph{Man lægger vel de Bestemmelser, som
Dualismen tillagde Aanden, og som Kjendsgjernin\-%
ger forbøde at benægte, fra Aanden over i Mate\-%
rien; men de blive staaende som fremmede og he\-%
terogene, og i Gjennemførelsen træde de tilbage
for den gamle Opfattelse}. Det, at Aanden gaaer op i
Materien, gjør ikke Materien aandig og activ; den bliver hvad
den var: \emph{den Materie}, hvis Bestemmelse netop var given ved
\emph{Modsætningen til Aanden}, ligesom \emph{Aandens Affald}.
I selve den yderste Modsætning til Middelalderens i Spiritua\-%
lismen grundende \emph{Dualisme} røber sig \emph{Paavirkningen} af denne.

Dette viser sig let naar vi undersøge Materialismens Be\-%
greb om \emph{Materien}, saaledes som dette er bestemt i \emph{Sy\-%
stème de la nature}, den materialistiske Retnings classiske
Værk.

Ved første Øiekast seer det ud, som om Materiens Be\-%
greb er blevet et andet end det var i Idealismens dualistiske
Modsætning: Materien sættes som \emph{i sit Væsen uadskille\-%
lig fra Bevægelsen}; der protesteres altsaa herved mod Fore\-%
stillingen om en \emph{død, kun passiv}, Materie. Spørgsmaalet:
hvorfra kommer Bevægelsen? besvares med et: hvorfra skulde
den vel komme uden fra Materien selv, det Ene-Existerende?
Men dette er foreløbigt ikke Andet end et i Formen positivt
Udtryk for det Negative: den kommer ikke, som Cartesianerne
statuerede, fra Gud, den active Aand; thi Gud er jo kun en
Fiction. Forat være Mere, maatte Svaret være begrundet
% --- p. 165 ---
% p. 165: no notes, no Greek, no ɔ:, no pencil.
% p. 165, l. 4: "Système de la nature" is set SOLID here and takes no
% markup (contrast p. 164, where the same phrase is letterspaced).
% p. 165: the two guillemet pairs -- «primære Qualiteter» and «Sæde» --
% are both the normal INWARD pair, « opening and » closing, read at 1200 dpi.
gjennem en Analyse af Materiens Begreb, der viste hiin For\-%
bindelses Mulighed; men dette er ikke skeet. I Materiens
\emph{oprindelige Egenskaber} («primære Qualiteter»), saaledes
som Système de la nature opregner dem, ligger \emph{intet Til\-%
knytningspunkt} for Bevægelsen, Intet, der kan gjøre den for\-%
klarlig, og idet det udtrykkelig siges, at Bevægelsen paa \emph{et\-%
hvert Punkt i Naturen kun} er tilstede og \emph{kun} kan blive
\emph{virkelig som meddelt Bevægelse}, saa er hermed sagt, at
\emph{Bevægelsen ingensteds} er paa \emph{oprindelig Maade} tilstede
i Materien, at den altsaa \emph{overalt er fremmed for denne},
\emph{og ikke} kan begribes ved dens Natur. Den postulerede, uad\-%
skillelige Sammenhæng opløses altsaa atter, Materien beholder
kun de Bestemmelser tilbage som den mechaniske Opfattelse
tillagde den: den bliver passiv, død, \emph{udenfra} mechanisk be\-%
væget Materie.

Naar nu ikke engang \emph{Bevægelsens} physiske Begreb, der
i Idealismen kun ved et Postulat sattes i Forbindelse med
Aanden, kunde forklares ud fra dette Begreb om Materien,
saa vil det allerede forud være givet, at \emph{Aandens eiendom\-%
melige Bestemmelser} ikke ville kunne det, at den Ma\-%
terie, der skulde omfatte \emph{hele} Tilværelsen, maa lade en \emph{Ho\-%
vedform} af Tilværelsen staae som \emph{ubegribelig}: den Ho\-%
vedform, hvis Væsensbestemmelser netop oprindelig ere \emph{ude\-%
lukkede} fra dens Begreb.

Dette viser sig, idet Materialismen skal forklare den men\-%
neskelige Aands Grundbestemmelse: \emph{Bevidstheden}. Forat
gjøre dette reducerer Materialismen først, paa samme Maade
som allerede Sensualismen havde gjort, alle Bevidsthedens
Former til \emph{Fornemmelse}, og, idet den saa beraaber sig paa,
hvad den kalder en Kjendsgjerning, at Fornemmelsen \emph{bevir\-%
kes ved Bevægelser i Nervesystemet}, der forplantes til Hjer\-%
nen, og paa Iagttagelsen af, at \emph{Hjernen} er «\emph{Sæde}» for
Fornemmelsen, saa mener den at have ført \emph{Beviset} for, at
et \emph{materielt} Organ \emph{factisk} har Evne til at kunne fornemme,
overhovedet til at udøve Bevidsthedsacter, og slutter saa, at
% --- p. 166 ---
% p. 166: no notes, no Greek, no pencil.
% p. 166 carries the ɔ: id-est mark ONCE, in the last line, at "gjensidig
% indvirke paa hverandre ɔ: med-"; read on the image at 1200 dpi (ABBYY
% "o:", tesseract "9:").
% p. 166: the guillemets round «død» and «levende» are both the normal
% inward pair; nothing inside them is letterspaced.
% p. 166, l. 12: "vis inertiæ" is Latin set in the same upright roman as the
% Danish and is NOT letterspaced; it takes no markup.
det maa have den fordi Fornemmelse er en \emph{Egenskab ved
Materien}, enten ved \emph{Materien i Almindelighed}, der
skal kunne have den som «død» eller «levende» Fornemmelse,
eller ved bestemte \emph{Combinationer af det Materielle}.
Men, hvilket af disse Alternativer man vælger, saa faaer man
\emph{ingen} Forklaring af Fornemmelsen, men kommer kun til den
ved et tomt Postulat. Tage vi det \emph{første} Alternativ: at Ma\-%
terien \emph{som saadan} skulde være fornemmende, saa tillægger
man den Materie, hvilken man ikke engang kunde hævde Be\-%
% p. 166, ll. 10 f.: "Materiens" straddles the line break and the print
% straddles the emphasis with it -- "har i Ma-" is set SOLID at the crowded
% line end while "teriens primære Qualiteter:" at the head of the next line
% is letterspaced (read at 1200 dpi).  The word is set here as one \emph{}
% group, to keep it whole; the split is recorded rather than reproduced.
% (Same case as p. 161 "Begrændsede".)
vægelsen, et Prædicat, der \emph{intet} Tilknytningspunkt har i \emph{Ma\-%
teriens primære Qualiteter}: Udstrækning, Delelighed, Be\-%
vægelighed, Uigjennemtrængelighed og vis inertiæ, og ligefrem
strider mod det saaledes bestemte Materielles Væsen. For\-%
nemmelsen, i hvilken den legemlige Affection activt forvand\-%
les, forudsætter et \emph{Selv}, der fornemmer \emph{sig}, en activ Sub\-%
jectivitet, der i sin Fornemmen ligesom \emph{fordobbler} sig paa
ideel Maade. Men det Materielle er det i \emph{Objectivitetens}
Form værende Passive, \emph{Selvløse}, det atomistisk Enkelte, i
hvilket en saadan ideel Fordobbling, en saadan Enhed af Sub\-%
jectivitet og Objectivitet, ikke kan tænkes. For at faae et
Tilknytningspunkt for Forestillingen om en saadan Fordobbling
og Enhed maatte man ialfald gaae ud over det enkelte Mate\-%
rielle og sætte et \emph{Forhold} mellem \emph{forskjellige} Dele af
Materien. Dette vilde saa føre os til det \emph{andet} Alternativ:
at sætte Fornemmelsen \emph{som Resultat af Combinationen
af Materiens Dele}. Men ogsaa denne Udvei afskjærer Ma\-%
terialismen sig selv. \emph{Materien}, hvis Navn giver \emph{Skinnet}
af en Enhed, er, strængt taget, kun en \emph{Fælledsbetegnelse},
altsaa en \emph{subjectiv} Abstraction, og det \emph{Reale}, der sam\-%
menfattes under dette Fælledsnavn, er de utilblevne og ufor\-%
anderlige \emph{Moleculer}, af hvis forskjellige Combinationer Na\-%
turens Former skulle dannes. Men disse Former i deres til\-%
syneladende Enhed ere kun \emph{Aggregater}, Forbindelser af i
og for sig \emph{selvstændige} mindste Dele. Disse Dele kunne,
som aggregerede, \emph{gjensidig} indvirke paa hverandre ɔ: med\-%
% --- p. 167 ---
% p. 167: no notes, no Greek, no ɔ:, no pencil.
% The printed folio reads 167; the ABBYY layer garbles it as "1G7"
% (pagemap.py gives PDF 176).
% p. 167, ll. 7 f.: the guillemets round «se replier sur soi-même» are the
% normal INWARD pair, « opening and » closing, read at 1200 dpi.  The French
% inside them is set in the same upright roman as the Danish and is NOT
% letterspaced, so it takes no markup; the hyphen in "soi-même" is a printed
% hyphen, not a line-break hyphen.
dele hverandre deres meddelte Bevægelse; men Bevægelsens
Strøm gaaer fra hver Deel bestandig \emph{udefter}, al Virksomhed
bliver, som \emph{Meddelen} af Bevægelse, en Virken \emph{paa et An\-%
det, et Forhold til et Andet}; et Forhold til \emph{sig selv}, en
\emph{Reflexion i sig selv} i Virksomheden, der skulde svare til hiin
Selvfordobbling, bliver \emph{udelukket} ved Forudsætningerne. Vel
tillægger Materialismen Hjernen den Egenskab, at kunne «se
replier sur soi-même», at kunne iagttage sine egne Forandrin\-%
ger. Men i dette Udtryk skjuler sig en Tvetydighed. Strængt
taget kan der kun forstaaes derved, at en af \emph{Hjernens
Moleculer} skal kunne indvirke paa \emph{en anden}, og denne
altsaa bestemmes i sin Væren ved hiin, og kun idet man for
\emph{Aggregatet af selvstændige} Moleculer, hvis \emph{gjensidige}
Paavirkning dog aldrig kan blive et \emph{Selvhedsforhold}, sub\-%
stituerer \emph{Heelhedens} og den \emph{subjective Enheds Begreb},
tillister sig Anskuelsen af den \emph{organiske Totalitet}, og
derved ubevidst gaaer ud \emph{over} Materialismens Grændser, til\-%
veiebringes \emph{Skinnet} af hiin Mulighed.

End ikke Selvets \emph{elementære} Form: \emph{Fornemmelsen},
formaaer saaledes Materialismen at forklare, fordi den til dens
Forklaring kun har det, hvis Begreb netop er dannet \emph{ved at
udelukke Selvets Bestemmelser}. Og Forklaringen af
denne elementære, allerede i Dyrelivet fremtrædende Form for
Selvets Virken vilde dog være den nødvendige, om end i og
for sig utilstrækkelige Betingelse for Forklaringen af de \emph{spe\-%
cifisk menneskelige} Aandsfunctioner: \emph{Bevidsthedsfunc\-%
tionerne}. Men naar nu \emph{Bevidstheden} staaer som en uaf\-%
viselig \emph{Kjendsgjerning}, der ikke uden \emph{umiddelbar Selv\-%
modsigelse} kan benægtes, saa bliver dens \emph{Mulighed} et
\emph{uafviseligt Problem}, og \emph{Umuligheden} af at kunne løse
det bliver \emph{afgjørende} for Dommen over Materialismens Prin\-%
cip, der viser sig som \emph{utilstrækkeligt og ensidigt}. Og
idet det \emph{Aandelige}, der \emph{ikke} kan skaffes tilside, staaer som
\emph{uforklaret}, saa fremkalder det uundgaaeligt \emph{Inconsequent\-%
ser i Systemet}, ved hvilke der uvilkaarligt indrømmes det en
% --- p. 168 ---
% SEAM p.168/169 CHECKED (BATCH-AGENT §5.2).  p. 168 bears NO notes at all
% -- there is no footnote separator and no note matter below the last body
% line -- so nothing can run over onto p. 169.  p. 169 opens with running
% text ("Betingelse, som et Incitament for den nye Tids Tanke ..."),
% continuing p. 168's last words without a paragraph break.
% NOTE FOR THE NEXT BATCH: the letterspaced run at the foot of p. 168,
% "en fremskridende Udviklings", CONTINUES onto p. 169, whose first word
% "Betingelse," is also letterspaced (read at 1200 dpi; "som et" after it is
% solid, "Incitament" spaced again).  The \emph{} group is closed here at the
% page break and must be reopened at "Betingelse," on p. 169.
% p. 168: no Greek, no ɔ:, no guillemets.
% p. 168: "Système de la nature" is set SOLID here and takes no markup.
Gyldighed, uafhængig af det Materielle. Dette viser sig i Sy\-%
stème de la nature, naar det benytter Begreber, der ikke kunne
være givne ved Erfaringsforholdet til Naturtingene, og altsaa
vise hen til et Apriorisk; det viser sig i dets Opfattelse af
Historien, idet der tillægges den en fra Naturtautologien for\-%
skjellig Udvikling; det viser sig i dets Ethik, idet det, i Strid
med sit Princip, gjør Frihed, Dyd og Menneskelighed til sit
Løsen.

\emph{Materialismens Grundbegreb} viser sig saaledes som
mangelfuldt og uigjennemførligt, men i det er den \emph{ensidige}
Enhed bleven aabenbar, hvis Spire blev paaviist i \emph{den nyere
Philosophies Udgangspunkt}. Det, der bliver tilbage
\emph{som Heelheden}, er her, ligesom ved den \emph{idealistiske}
Rækkes Slutningspunkt, \emph{netop kun det, der oprindelig
var sat som Modsætningens ene Led}. \emph{Det andet Mod\-%
sætningsled} er ikke forsonet med dette, saa at det som Heel\-%
hed bestemte er blevet indholdsrigere, men er \emph{skudt tilside},
kastet bort, og idet det \emph{ikke} kan tilintetgjøres for Bevidsthe\-%
% p. 168, ll. 21 f.: NOT a paragraph break.  "deelt." and "I denne nye
% Sondring" stand on the SAME printed line; the ABBYY layer inserts a break
% and an indent here that the page does not have.
den, tiltvinger det sig paany en Plads og gjør Tilværelsen
deelt. I denne \emph{nye Sondring og i Grundbegrebets Be\-%
stemmelse}, der netop er den \emph{dualistiske} Modsætnings, er
% p. 168, l. 24: a small faint mark stands between "Eftervirkningen" and
% "af"; a second stands before "Materialismen" on the line below.  Both are
% far lighter than the type and are modern pencil or dirt, not printed
% sorts; not transcribed.  (The ABBYY layer reads the first as a "t" --
% "Eftervirkningen taf den" -- and the second as an apostrophe.)
det altsaa, at \emph{Eftervirkningen} af den Livsanskuelse, mod
hvilken Materialismen \emph{bevidst kæmpede}, endnu \emph{ubevidst}
gjør sig gjældende, som hæmmende Aandens bevidste Stræben,
og som bringende en Dobbelthed ind i dens Retning.

% p. 168: a SHORT CENTRED DIVISIONAL RULE stands here, between the two
% paragraphs, with white space above and below.  Measured on the image at
% 300 dpi: x 685--934, i.e. 249 px = 2.1 cm, its centre at x 810 against
% the measure's centre at x 807 -- centred, therefore the book's divisional
% rule and NOT the left-flush footnote separator (transcription.tex §2).
% It closes Afsnit III.
\divrule

I det her afsluttede Afsnit (III) er, gjennem de valgte
Exempler, den ubevidste hæmmende Indvirkning af det For\-%
tidens Princip, mod hvilket der bevidst kæmpedes, bleven paa\-%
viist i \emph{hvert} af de Tidsafsnit, der følge efter Oldtiden. Det
bliver nu at paavise, hvad der tidligere (S. 49) blev antydet,
at denne ubevidste Hæmmen, idet Fortidens Princip opfatte\-%
des som Moment i et mere Omfattende, fra det Universelles
Synspunkt maatte sees som en \emph{fremskridende Udviklings}
% --- p. 169 ---
% SEAM 168->169.  p. 168 ends mid-sentence with "...maatte sees som en
% \emph{fremskridende Udviklings}" and p. 169 opens "Betingelse, som et
% \emph{Incitament}...".  THERE IS NO PARAGRAPH BREAK HERE, and the letterspaced
% run that begins on p. 168 ("fremskridende Udviklings") STOPS before
% "Betingelse," -- measured on the image: "Betingelse," has ordinary body gaps
% while "Incitament" three words later is spaced.  The pp. 155--168 batch must
% end its fragment with "...en \emph{fremskridende Udviklings}" and NO blank
% line; run joints.py after the splice.
% p. 169 bears no note and no Greek.  ABBYY splices no pencil into this page.
Betingelse, som et \emph{Incitament} for den nye Tids Tanke, der
directe eller indirecte fremkalder en rigere og alsidigere Tanke\-%
udvikling. Dette Forhold, der umiddelbart er givet naar Videns
Idee i dens Concretion sees som det den menneskelige Videns
Udvikling Beherskende, fremtræder i en særegen Form i hvert
af de ovenfor behandlede Tidsafsnit. Idet den \emph{abstracte
spiritualistiske} Livs- og Verdensanskuelse søger at udfolde
sig i et Tankeheles Form, trænger det Element sig frem, der,
skjøndt væsentligt, søgtes fortrængt, og idet Tanken stræber
at drage det ind under Aanden, tvinges, trods det Naturliges
Forvandling, Aandelighedens Abstraction til at \emph{articulere}
sig, til at forme sig \emph{henimod Virkeligheden.} Idet \emph{Na\-%
turelementet} nu er kommet frem som anerkjendt af \emph{Over\-%
gangsperiodens Bevidsthed,} viser det sig først ligesom
\emph{igjennem det Ideelle,} overstraalet af dettes Lys, og selv
idet det bestemtere kommer tilsyne i sin \emph{egen,} i Naturens
Form, faaer det dog den samme, ved Phantasien satte, \emph{umid\-%
delbare Enhed med det Aandelige} som i den græske
Philosophie. Idet den \emph{middelalderlige Dualisme} gjør sig
gjældende, hævdende det overnaturlige Guddommeliges rene
Aandelighed, træde Aand og Natur ud fra hinanden i en \emph{Son\-%
dring,} der, som betingende det Naturliges Opfattelse i dets
\emph{principielle Eiendommelighed,} er Forudsætningen for
Erkjendelsen af den \emph{concrete} Enhed af Natur og Aand, be\-%
stemte efter deres \emph{Væsensforskjel:} for Erkjendelsen af en
\emph{sand Forskjelsenhed.} Idet endelig den \emph{nyere Philoso\-%
phie} paa Basis af Dualismen stræber efter Enheden, giver
\emph{Dualismens Indvirkning} denne Enhed en \emph{ensidig} Form;
men herved \emph{forberedes,} paa den ene Side, gjennem den
\emph{isolerede} Udvikling af de enkelte Momenter, der overhovedet
er mulig, Modsætningernes \emph{fyldige} Enhed, og, paa den anden
Side, fremtvinge de ifølge Enhedens Ensidighed uundgaaelige
\emph{Inconsequentser} Forsøget paa at tilveiebringe en saadan
Enhed, gjøre denne Opgave til en \emph{uafviselig Nødvendig\-%
hed for Tanken.}

% DIVISIONAL RULE at the FOOT of p. 169, below the last line of text and with
% white space above and below -- the pp. 15 and 79 pattern (rule at the foot,
% head at the top of the next page).  Measured at 300 dpi: x 547--799, i.e.
% 252 px = 2.13 cm, centre 673.0 against a text-block centre of 675.5 (block
% x 120--1231); a plain hairline, one to two pixels thick.  It is NOT the
% footnote separator (p. 169 has no note at all), and it is the rule that
% belongs to the IV. head on p. 170, which itself carries none.
\divrule

% --- p. 170 ---
% p. 170 bears no note and no Greek.
%
% HEAD, p. 170 (L2, transcription.tex §3): the FOURTH and last of the book's
% theses (I. 50, II. 79, III. 137, IV. 170).
%
% BOLD, NOT LETTERSPACED.  Measured at 600 dpi on this page, three lines
% compared: the head's first line has inter-letter gaps of 1--6 px (median 4);
% the ordinary body line "Forat belyse en saadan gjensidig Indvirkning,
% gjennem" four lines below runs 1--11 px (median 6); the LETTERSPACED L3 head
% lower on the same page ("Idealismens og Realismens Indvirkning paa hin-/
% anden i den nyere Philosophie.") runs 8--21 px on its first line (median 15)
% and 8--17 px on its second.  So the head is TIGHTER than the body, not wider:
% it is bold, and it is NOT set \emph{}.  (spacing.py reports two RUNs inside
% it; that is the known bold false positive of BATCH-AGENT §3.)
%
% NOT CENTRED -- it runs the FULL MEASURE, like p. 137 and unlike the p. 79
% head, which the pp. 71--84 batch set inside \begin{center}.  Measured at
% 300 dpi: head lines 1--3 run x 263--1371, x 263--1370, x 192--1369, against
% body lines at x 260--1364, x 259--1367; identical measure, and the head is
% justified.  THE OPEN REVIEW ITEM (p. 79 centred vs p. 137 full measure) is
% therefore settled two-to-one in favour of full measure; p. 79 should be
% brought into line.  ONE FURTHER DATUM for that review: the head's SHORT LAST
% LINE ("Bestemmelse og Udfoldelse.") is CENTRED on the measure, not flush
% left -- x 528--1103, centre 815.5, against a head-block centre of 817.  That
% is not reproduced below, because p. 137 was set as a plain justified block
% and the two heads must agree; it is recorded here instead.
%
% NO DIVISIONAL RULE stands above this head.  Its rule is at the FOOT OF
% p. 169 (see the comment there) -- the pp. 15/79 pattern.
\noindent\textbf{\large IV. To samtidige, coordinerede Udviklingsformers
modsatte Principer indvirke, under den bevidste
Modsætning, ubevidst for Tænkerne, paa hinandens
Bestemmelse og Udfoldelse.}

Forat belyse en saadan gjensidig Indvirkning, gjennem
hvilken ét mere Omfattende, skjult for de tænkende Individer,
forbereder Fremtidsudviklingen, udvikle vi:

% L3 head, p. 170 (transcription.tex §3): LETTERSPACED roman at body size, set
% as an ordinary indented, justified paragraph of two lines -- not bold, not
% centred.  Gaps measured at 600 dpi: 8--21 px against 1--11 px in the body
% line immediately below it and 1--6 px in the bold L2 head above.
\emph{Idealismens og Realismens Indvirkning paa hin\-%
anden i den nyere Philosophie.}

Den gjensidige Indvirkning af den nyere Philosophies to
Hovedretninger see vi under Synspunktet af en gjensidig,
ubevidst Paavirken af deres dualistisk sondrede Principer.
Principernes gjensidige Bestemmen viser sig i \emph{negativ} Form
idet hvert oprindelig hævder sig sine Prædicater \emph{exclusivt,} og
gjennem denne exclusive Tilegnen af dem indirecte bestemmer
Modsætningens. Men fra denne Bestemmen, der falder sammen
med den i det Foregaaende paaviste Virkning af \emph{Dualismen,}
see vi i denne Sammenhæng bort, og betragte her den ubevidste
gjensidige Paavirkning, som de to Principer udøve paa hinan\-%
den under den nærmere Bestemmelse af selve Principet og
under Udfoldelsen af det af Principbestemmelsen Afhængige, og
som fremtræder deri, \emph{at de her overføre deres Bestem\-%
melser paa Modsætningen,} hvor de vise sig som \emph{In\-%
consequentser.} Denne ubevidste Overføren af Bestemmelser
paa Modsætningen opfatte vi som nødvendig paa Grund af
\emph{Ensidigheden} i det enkelte Princip, der i denne Ensidig\-%
hed skal udtrykke \emph{den omfattende Totalitet,} og derfor
ligesom kræver en Fuldstændiggjørelse. Og gjennem den \emph{ud\-%
vortes,} inconsequente Optagen af Modsætningens Bestemmel\-%
ser see vi \emph{den sande concrete Enhed} antydet, som Frem\-%
tidsudviklingen skal virkeliggjøre.

De to Principer, der bestemme den nyere Philosophies
Hovedformer, betegne vi, naar vi gaae ud fra Ideeforhol\-%
% --- p. 171 ---
% p. 171 bears no note and no Greek.
dene, som Andetheden sat \emph{som Modsætning,} og Subjectivi\-%
teten, sat \emph{som Modsætning.} \emph{Realismens Princip var Na\-%
turen,} sat i Andethedens Form, som Materie = Ikke-Aand;
% p.171: the two "=" here and below are the TEXT sort (a double stroke),
% not math mode -- BATCH-AGENT §9.
\emph{Idealismens Princip var Aanden,} sat i Modsætning til
Andetheden, reflecteret i sig som Aand = Ikke-Natur. Af
disse Principer har Idealismens Princip allerede viist sig som
det overgribende idet det oprindelig bestemmer Opfattelsen af
% p.171: a faint hooked speck stands between `det' and `altsaa' in the line
% below; ABBYY reads it as an apostrophe ("det'altsaa").  At 1200 dpi it is a
% thin light mark well above the x-height, unlike this page's certain
% punctuation, and it is not transcribed (BATCH-AGENT §7).
Realismens Princip, saaledes at det altsaa, idet dette virker
bestemmende tilbage paa det, bestemmes af det, som det selv
først har bestemt.

Søge vi nu først den Paavirkning, \emph{Realismen} har mod\-%
taget fra sin Modsætning, saa kunne vi, med Hensyn til dens
Princips \emph{oprindelige} Bestemthed, vise tilbage til hvad der
allerede i det foregaaende Afsnit blev udviklet: at Opfattelsen
af den realistiske Philosophies Grundbegreb er paa \emph{negativ}
Maade bestemt ud fra et Aandsbegreb, hvis dualistiske Mod\-%
sætning til Naturen var fastslaaet af Idealismen. Denne gav
Aanden Activitetens Bestemmelse og reducerede Aandens ab\-%
stracte Modsætning: \emph{Naturen,} til det \emph{Passive} og kun ved
\emph{Udstrækningens} Abstraction Bestemte. Og trods Erfarings\-%
philosophiens Stræben efter at bestemme Naturbegrebet saa\-%
ledes, at det Materielle i sig fik Form og Bevægelse, og Na\-%
turbegrebet altsaa blev det \emph{omfattende Concrete;} saa ved\-%
blev dog det Materielle, under Paavirkning af hiin af Idealis\-%
men givne Bestemmelse, at være det Uvirksomme, mechanisk
Bevægede, og Materialismen, der opfattede Aanden som et
Uselvstændigt og overførte Materiens Bestemmelse paa den,
kunde dog ikke frigjøre denne, tilsyneladende ikke blot til
Selvstændighed, men til Enegyldighed naaede Materie fra \emph{den}
Opfattelse, som den havde faaet fra den Philosophie, der satte
% p.171: PRINTER'S ERROR, TRANSCRIBED AS PRINTED -- `sem' where the sense (and
% the parallel `som' eight words earlier) wants `som'.  Read on the image at
% 1200 dpi: the second sort is unambiguously `e', not `o'.  Both OCR witnesses
% give `sem'.  Not on the author's Trykfeil leaf.
Aanden som det \emph{kun} Active og sem det \emph{eneste} Active.

Men indenfor den Tænkningsform, der er bestemt ved det
% p.171: only `Natur' is letterspaced in `Naturbegreb' -- measured at 600 dpi,
% gaps 13--18 px through `Natur' and 2--7 px through `begreb,'.  This is a
% recurring habit of this compositor; see pp. 174, 175, 177, 178, 179.
\emph{ensidigt} opfattede \emph{Natur}begreb, finde vi saa en ubevidst
\emph{Overføren} af Modsætningens Bestemmelser. Den viser sig i
Opfattelsen af \emph{Principet,} ved Forsøget paa, at tillægge dette
% --- p. 172 ---
% p. 172 bears no note and no Greek.
Activitetens og Fornemmelsens Bestemmelser, der vare ude\-%
lukkede ved dets Begreb, og som derfor kun gjennem Postu\-%
% p.172: the letterspacing begins IN THE MIDDLE OF A BROKEN WORD.  `lige-' at
% the end of the line below has ordinary body gaps (2--7 px at 600 dpi) while
% `ledes, under Udviklingen af Systemerne,' on the next line runs 13--19 px.
% Transcribed as printed, so the \emph{} opens at `ledes'.
later inconsequent kunde tillægges det. Den viser sig lige\-%
\emph{ledes, under Udviklingen af Systemerne,} gjennem
de Inconsequentser, der deels fremtvinges ved Forsøget paa at
forklare ubenægtelige, til et andet Princip henvisende Kjends\-%
gjerninger, deels ligesom spontant fremkomme idet den om\-%
fattende Heelhed, som det enkelte Princip ensidigt udtrykker,
uvilkaarlig fører Tanken til Modsætningens Bestemmelser. Saa\-%
danne Inconsequentser finde vi i Realismens \emph{Erkjendelses\-%
lære,} hvor ikke blot paa Udviklingens fremskridende Stadier
et \emph{Apriorisk} deelviis anerkjendtes, men hvor det ved Udvik\-%
lingens Afslutning trods den bevidste Fornægtelse ubevidst
anerkjendtes gjennem den factiske Anvendelse af aprioriske
% p.172: `Système de la nature' is set in the SAME upright roman as the Danish
% and is NOT letterspaced -- checked at 600 dpi, grave accent plainly resolved.
% No \textit{} (transcription.tex §2).
Bestemmelser (Hume; Système de la nature). Og ligesom i
Erkjendelseslæren, gjennem denne ubevidste Anerkjendelse, en
anden Grundanskuelse af Aandens Virken trænger ind, saaledes
trænger den ind i Opfattelsen af \emph{Menneskelivet,} naar \emph{Per\-%
fectibilitetens} Begreb afgiver Synspunktet for \emph{Historiens}
Udvikling, der danner en Modsætning til \emph{Naturens} ensfor\-%
mige Gang; og i Opfattelsen af den \emph{Natur,} der bestemmes
som omfattende al Tilværelse, gjør, gjennem en analog In\-%
consequents, en Indvirkning fra Modsætningen sig gjældende,
idet et Naturligt bestemmes som \emph{Heelhed,} ikke blot som Ag\-%
gregat; thi i Heelhedens Bestemmelse er Idealismens Grund\-%
bestemmelse tilstede.

Mere skjult end Idealismens Indvirkning paa \emph{Realismen,}
er den realistiske Retnings Grundbegrebs Indvirkning paa \emph{Idea\-%
lismen;} men Indvirkningen er derfor ikke mindre væsentlig.

Vi opdage den i den nyere Idealismes første Form: den
\emph{dogmatiske} Idealisme, idet vi følge dens Grundbegreb til
dets consequente Udfoldelse. Det Begreb, om hvilket \emph{hele}
den dogmatiske Idealismes Tænkning samler sig, er \emph{Sub\-%
stantsens} Begreb: Begrebet om det væsentlig Værende i dets
Selvstændighed og Nødvendighed. Den \emph{dogmatiske} Idealisme,
% --- p. 173 ---
% p. 173 bears no note and no Greek.
der bærer dette Navn fordi den \emph{forudsætter} Erkjendelsens
Evne til at opfatte det Objective, vender sig \emph{umiddelbart} mod
Erkjendelsen af \emph{Tingenes Væsen.} Men dette Tingenes Væ\-%
sen er just \emph{Substantsen,} hvis Begreb, idet det ogsaa gjæl\-%
% p.173: `res cogitans' is upright roman AND letterspaced -- not italic.
der for Aanden, der som \emph{res cogitans} gaaer ind under
det \emph{Værendes} Begreb, kommer til at omfatte baade den
% p.173: a small pale speck follows the full stop after `Tilværelse' in the line
% below; both witnesses read it as a second period ("Tilværelse..").  At 1200 dpi
% it is markedly smaller and lighter than this page's certain full stops, so it
% is read as a speck and a single full stop is transcribed (PLAYBOOK §2 rule 3).
aandelige og legemlige Tilværelse. \emph{Dette} Substantsbegrebs
\emph{Indhold} udvikler sig nu, medens \emph{Navnet} bevares, i den
dogmatiske Idealisme, og netop denne Udvikling udgjør denne
Tænkningsforms Historie.

\emph{Descartes} bestemmer, idet han sætter Aand og Materie
som absolute, hinanden udelukkende, altsaa i Forhold til hin\-%
anden selvstændige, Modsætninger, begge som \emph{Substantser:}
som \emph{tænkende Substants} og som \emph{udstrakt Substants.}
Substantsen er hvad der i Retning af Væren og Tænken har
% p.173: both guillemet pairs on this page point INWARD, «…», read at 1200 dpi.
Selvstændighed, hvad der «er \emph{ved sig selv} og begribes ved sig
selv.» Men denne \emph{Væren ved sig selv} faaer hos Descartes
strax en Begrændsning. Være ved sig selv kan kun Eet: det
% p.173: only `Absolute,' is letterspaced; `Guddommelige;' beside it has
% ordinary body gaps (measured at 600 dpi, 14--19 px against 4--7 px).
\emph{Absolute,} Guddommelige; kun paa dette passer Definitionen
fuldstændigt, paa det endelige Tænkende og det endelige Le\-%
gemlige kun med den Modification, at det til sin Existents
ikke behøver noget Andet end Gud. Der kommer saaledes,
naar Substantsbegrebet alligevel bibeholdes som Bestemmelse
for det Endelige, en \emph{Tvetydighed} ind deri. Der tales om
den \emph{uendelige} Substants og om de \emph{endelige} Substantser,
og om disse atter snart saaledes, at det Aandelige og det Le\-%
gemlige som Heelhed betegnes ved dette Navn, snart saa\-%
ledes, at det \emph{individuelle} Aandelige og Legemlige tilegner
sig det.

Denne Vaklen hæver \emph{Spinoza} idet han udskiller de for\-%
% p.173: a pale speck stands after `skjellige' in the line below; both witnesses
% read it as a full stop.  It is much smaller and lighter than the certain full
% stop after `eet' on the same line, and is not transcribed.
skjellige Begreber, Descartes sammenblandede i eet. \emph{Sub\-%
stantsen er een: den uendelige;} Tanke og Udstrækning
ere ikke \emph{bestemte Substantser,} men \emph{Substantsens Be\-%
stemmelser, dens Attributer,} det Uendeliges uendelige
Væreformer; det Individuelle er Attributernes \emph{Modificatio\-%
% --- p. 174 ---
% p. 174 bears no note and no Greek.  The word `Modificationer,' is broken
% across the page break and the letterspacing runs across it.
ner,} den uendelige, substantielle Aarsags uendelige Virkninger.
Disse spinozistiske Bestemmelser ere correcte Consequentser af
de cartesiske Forudsætninger, og Spinozas Opfattelse af den
% p.174: `Ene-Værende' is a printed (lexical) hyphen that happens to fall at a
% line end -- the second element is capitalised.  Set as an ordinary hyphen,
% not \-.
uendelige Substants som det \emph{abstract-uendelige Ene-Værende}
er begrundet i selve Substantsbegrebets Grundbe\-%
tydning. Substantsen var det \emph{væsentlige Værende i dets
Nødvendighed.} I denne Bestemmelse bliver \emph{Værens} Be\-%
greb det Tilgrundliggende og Fremherskende, ligesom Subjec\-%
tet; \emph{Activiteten, Tanken,} forholder sig til dette Sub\-%
ject kun som et bagefter tilføiet \emph{Prædicat,} og derfor
bliver \emph{Uendeligheden,} der er hiin Værens Væsen, ikke
bestemt, men bliver \emph{abstract, hvilende Uendelighed.}
Den er ikke Uendelighed ved den Activitet, ved hvilken den
sætter og i sin Enhed ophæver sine Bestemmelser: den er
Uendelighed \emph{forud for} Activiteten, og derfor sætter denne
kun et Forsvindende. Til dette \emph{Spinozistiske} Substantsbegreb,
% p.174: PRINTER'S ERROR, TRANSCRIBED AS PRINTED -- `udsrykker' for
% `udtrykker'.  Read at 600 dpi: `uds-' plainly, with `r' following.  Both OCR
% witnesses give `udsrykker'.  Not on the author's Trykfeil leaf.
der \emph{consequent} udsrykker den dogmatiske Idealismes Grund\-%
tanke, er det \emph{Leibnitziske} med dets Individualitetsbestem\-%
melse (Monaden) kun det postulerede, om end ifølge selve Spi\-%
nozismens Mangler uundgaaelige, Supplement, der staaer paa
Grændsen mellem Dogmatismen og Kriticismen, og som, naar
det gjennemtænktes, gjennem Selvbevidsthedsbegrebet vilde føre
til den sidste.

I det \emph{Spinozistiske} Grundbegreb gjør nu ikke blot den
tidligere paaviste \emph{almindelige} dualistiske Sondring af Aand
og Natur, og den derved betingede Opfattelse af Modsætnings\-%
% p.174: only `Natur' is letterspaced in `Naturbegrebets' -- gaps 13--18 px
% against 2--7 px through `begrebets', measured at 600 dpi.
leddenes Forhold sig gjældende; men \emph{Natur}begrebets Indvirk\-%
ning gjør sig indenfor det \emph{idealistiske} System tillige gjæl\-%
dende paa \emph{eiendommelig} Maade. Til sit \emph{Substantsbegreb}
var Spinoza oprindelig naaet ud fra \emph{Jeget,} fra \emph{Selvbevidst\-%
heden} som Udgangspunkt; Substantsen var oprindelig \emph{Tan\-%
kens} uendelige Væsen. Men med dette Tankens uendelige
Væsen er \emph{Naturens,} Materiens, uendelige Væsen sat i Enhed.
Substantsen er ikke blot en res cogitans, men en res extensa,
% p.174: `Fælledsbestemmelse' is one word in the print; only its first element
% is letterspaced (gaps 13--17 px through `Fælleds', 3--6 px through
% `bestemmelse').  ABBYY's space in "Fæ lieds bestemmelse" is an OCR artefact.
og i denne \emph{Fælleds}bestemmelse af \emph{Ting,} i hvilken Over\-%
% --- p. 175 ---
\setcounter{footnote}{0}
% p. 175 carries ONE note, *), of three lines, self-contained -- it does not run
% over onto p. 176, whose note area is empty.  No Greek.
% p.175: `natura naturata' is upright roman throughout, letterspaced in the body
% (line 14) and NOT letterspaced inside the note.  No \textit{}.
vægten af \emph{Værens} (Objectivitetens) Begreb antydes, antydes
allerede \emph{Natur}begrebets bestemmende Indvirkning. Men denne
bliver aabenbar, idet Attributernes Forbindelse i Substantsens
Enhed lader dem gjensidig overføre Bestemmelser paa hinan\-%
den. Medens saaledes \emph{Tankens} Attribut bestemmer Mod\-%
sætningens Begreb derved, at det Naturlige reduceres til en
abstract Tankebestemmelse med Udelukkelse af de sandselige
Concretionsbestemmelser, og tillige derved, at Udstrækningen
faaer Prædicaterne Udelelighed og Activitet, altsaa Prædicater,
der ere tagne fra det Aandelige, og kun have Mening der, ikke i
det til Udstrækning reducerede Naturlige; saa see vi omvendt
\emph{Natur}begrebet bestemme Aandsbegrebet idet den uendelige
% p.175: the footnote mark *) is set on `Nødvendighedens', inside the
% letterspaced run.
Tænknings Activitet opfattes under \emph{Nødvendighedens}\footnote{Med Hensyn til \emph{Spinozas} Opfattelse af Nødvendighedsbegrebet i
dets Anvendelse paa Substantsen og paa natura naturata henviser
jeg til mit Skrift om Spinoza, S. 66 ff.}, ikke
% p.175: ABBYY prints "jeg*" in the note above.  There is NO asterisk after
% `jeg' on the image at 1200 dpi -- the mark is an OCR artefact -- so the
% self-reference reads simply "henviser jeg til mit Skrift om Spinoza, S. 66 ff."
% (Brøchner's own *Benedict Spinoza*, also in this repo.)
under \emph{Øiemedets} Synspunkt, og idet den ideelle \emph{natura na\-%
turata,} der er sat ved denne Activitet, opfattes i Analogie
med det mechanisk sammenkjædede bevægede Materielle.
\emph{Causalitetens Nødvendighedsforhold} bestemme de \emph{me\-%
chanisk forbundne} Ideers Sammenhæng idet de individuelle
Tankemodificationer gaae ind under Tingens Begreb, og i deres
Følge og \emph{Sammensætning} anskues i Analogie med Udstræk\-%
ningsmodificationerne. Og Naturbegrebets Indvirkning faaer
sit \emph{definitive} Udtryk deri, at Substantsbegrebet, der i sin
Sandhed er Begrebet om \emph{Gud,} faaer \emph{Naturens Navn,} saa
at \emph{Deus} og \emph{Natura} blive \emph{Synonymer.} Under Attributernes
gjensidige Bestemmen er Naturbegrebets Indvirkning mægtig
% p.175: only the second element of `Tilværelsesheelheden,' is letterspaced
% (gaps 3--6 px through `Tilværelses', 13--18 px through `heelheden,') -- the
% mirror image of the `Natur'/`begrebet' cases above.  Transcribed as measured.
nok til at bestemme den Opfattelse af Tilværelses\emph{heelheden,}
der fører denne ind under Naturnødvendighedens Lov. Men i
denne Indvirkning af Naturbegrebet gjør, som allerede antydet,
ifølge det idealistiske Moments Grundforhold, en eiendommelig
Kredsbevægelse sig gjældende. Det Naturbegreb, der bestemmer
hiin Opfattelse af den \emph{tænkende} Substantses Activitet og dermed
% --- p. 176 ---
% p. 176 bears no note (its note area is empty, so p. 175's note does not run
% over) and no Greek.
Opfattelsen af Substantsens Activitet overhovedet, er det, der selv
er bestemt ved den Maade, paa hvilken Aanden, under Dualismens
Paavirkning, opfattede sit eget Begreb. Aandsbegrebet, der
bestemte Naturbegrebet ved at udelukke det abstract fra sig,
bestemmer altsaa i hiint Vexelvirkningsforhold \emph{indirecte} sig
selv gjennem det saaledes Bestemtes Tilbagevirkning. Men
dog maa det ikke oversees, at dette Tilbagevirkende udtrykker
det, der fra Begyndelsen af som et \emph{Selvstændigt og Væ\-%
sentligt} gjør sig gjældende for den Aand, der søger Væren
og Vished i sig selv, og idet det tvinger den til at anerkjende
sig, trods det Aandeliges Overgriben, hævder sig som en \emph{væ\-%
sentlig} Væreform, der, selv under Bestemmetheden ved det
Aandelige, bevarer Muligheden til en Virken ud fra sin Eien\-%
dommelighed, til en Bestemmen af den Aand, der fra Begyn\-%
delsen af stræber at drage Tilværelsens Heelhed ind under sig.

Idet \emph{Spinozas} Substantsbegreb bliver det consequente Ud\-%
tryk for den \emph{dogmatiske} Idealismes Grundtanke, er dets
Bestemmethed ved hiin Indvirkning afgjørende for den dog\-%
matiske Idealisme, naar den tages \emph{som Heelhed,} og vi vilde
kunne indskrænke os til at paavise Indvirkningen paa dette
Punkt. Men det har Interesse at see, hvorledes ogsaa i de
andre Systemer, der hvile paa samme Grundbegreb: Begrebet
om \emph{Substantsen,} en lignende Paavirkning bliver kjendelig.
Hos \emph{Descartes} træder den væsentligst frem i den Forvandling,
ved hvilken Jeget, hvis \emph{Væsen} var bestemt som \emph{activ}
Tænkning, bliver til en \emph{res cogitans,} altsaa kun faaer Tænk\-%
ningen til \emph{Værens} Attribut. I denne Forvandling, der gaaer
forud for Enhedstankens Gjennemførelse hos Spinoza, gjør
nemlig en Reaction sig gjældende fra det, der hos Descartes,
trods hans Udgangspunkt, giver Aanden dens væsentlige Er\-%
kjendelsesindhold, fra Naturen, hvis Værens Bestemmelse uvil\-%
kaarlig overføres paa den i \emph{umiddelbart} Forhold til den
staaende Aand. Hos \emph{Leibnitz,} der, medens han vil forvandle
Alt til Aand, dog uvilkaarligen kommer til at anerkjende det
Naturliges Væsentlighed, træder Paavirkningen frem idet Me\-%
% --- p. 177 ---
\setcounter{footnote}{0}
% p. 177 carries ONE note, *), of thirteen lines, filling the lower third of the
% page.  It is self-contained: the top of p. 178 is ordinary body text with no
% unmarked matter, so the note does not run over.  No Greek.
% p.177 also carries the sheet signature "12" at the foot (p. 179 carries "12*").
% Signatures are not transcribed (BATCH-AGENT §6).
chanismen trænger ind i Forholdet mellem Substantserne (Mo\-%
naderne) indbyrdes, og i Substantsens (forvirrede) Forestillinger.
Hos \emph{Wolf} træder den frem idet Monadebegrebet, som han
% p.177: only `Atom' is letterspaced in `Atombegreb'.
trivialiserer, deelviis ombyttes med et \emph{Atom}begreb. Og vil
man tage Hensyn til den hensygnende dogmatiske Idealisme
i anden Halvdeel af det 18de Aarhundrede, saa vil man finde,
at den modsatte Tankeretnings og Livsanskuelses Overmagt
overalt gjør sig gjældende og gjør Idealismen eklektisk og
inconsequent.

% p.177: a thin hooked ink or pencil stroke stands immediately after `Overmagt'
% in the line below, at x-height; ABBYY reads it "Overmagt^", tesseract
% "Overmagt«".  Examined at 1200 dpi it is a hairline curve, not a printed sort
% (no comma, no guillemet), and is not transcribed.  p. 177 is one of the known
% pencil-pollution sites of BATCH-AGENT §7.
Ved denne Retnings Overmagt bliver Idealismen i det
18de Aarhundrede overvældet; men den samler ny Kraft ved
paany at gaae tilbage til sit første Udgangspunkt: til \emph{Jeget,
Selvbevidstheden,} for ud fra det atter at stræbe hen
mod det Maal, til hvilket den nyere Idealisme som Heelhed
sigtede hen, men som den \emph{dogmatiske} Idealisme, ifølge sit
% p.177: only `Værens' is letterspaced in `Værensbegrebets'.
Grundbegrebs Mangel paa indre Bevægelighed, --- en Mangel,
der var uadskillelig fra \emph{Værens}begrebets Overvægt --- ikke
kunde naae, skjøndt det blev \emph{antydet} ved selve dens Ud\-%
gangspunkt, saasnart den individuelle Tanke tænkte sit Væ\-%
sen: mod Begrebet om \emph{den absolute Aand}\footnote{Allerede hos \emph{Descartes} antydes det Begreb af \emph{den absolute
Aand,} med hvilket den nyere Idealisme afslutter. Dette Begreb
er antydet, idet Descartes opfatter den guddommelige Aand som
\emph{den active Selvbevidstheds uendelige Væsen;} idet han
viser hen til, at det \emph{Uendelige} er vor \emph{første} Tanke, fra hvilken
der først naaes til det Endelige gjennem en Begrændsning; idet
han lader \emph{Gud være virksom i Gudsideen,} der opfattes som
hans \emph{Væren i os} og som given med \emph{selve vor Selvbevidsthed,}
og idet han sætter \emph{Gud,} der er \emph{Naturens Ophav,} som i en
\emph{høiere} Form indesluttende Naturens Væsen i sig, i Analogie med
den Maade, paa hvilken vor Bevidsthed indeslutter Forestillingen
om det Legemlige.}, der i sin
Enhed skulde omfatte \emph{Naturens Væsen.} Men paa dette nye
Udviklingsstadium, der, ved at begynde med at gjøre sig det
tidligere oversprungne Problem: \emph{Muligheden af en Er\-%
% --- p. 178 ---
% p. 178 bears no note and no Greek.  The letterspaced phrase
% `Muligheden af en Erkjenden af det Objective,' runs across the page break.
kjenden af det Objective,} til Gjenstand, faaer Charakteren
af \emph{kritisk} Idealisme, viser Paavirkningen sig paany. Vi
kunne see bort fra den \emph{negative} Form af den, der er iden\-%
tisk med \emph{Dualismens Indvirkning,} og som, naar vi tage
Stadiet som Heelhed, viser sig deri, at i det afsluttende Sy\-%
stem den Enhed, der naaes, trods Bestræbelsen for at aner\-%
kjende Naturens Væsentlighed, dog bliver en \emph{ensidig} Enhed,
idet Aandsbegrebet, der skulde omfatte Naturens Væsen i sig,
ifølge den dualistiske Modsætning, ikke kan tilegne sig det
Naturliges eiendommelige Bestemmelser, men maa skyde dem
fra sig som usande. Men vi fremhæve en \emph{positiv} Paavirk\-%
ning gjennem Overførelser fra Modsætningen, saavel i Be\-%
stemmelsen af Principet, som i det af Principet Afledede.
Hvad Bestemmelsen af \emph{Principet} angaaer, saa forbigaae vi
\emph{Kant} fordi den Maade, paa hvilken han unddrager det Ab\-%
solute fra Erkjendelsen, idethøieste tillader at søge en Paa\-%
virkning gjennem svage Antydninger. Hos \emph{Fichte} finde vi
Paavirkningen, i hans \emph{ældre} Philosophie, idet en \emph{uforklaret
Natur}bestemmelse, en \emph{Naturskranke,} trænger ind i \emph{Jegets}
Activitet, og i hans Philosophies \emph{sildigere} Form idet \emph{Væ\-%
rens} Begreb afløser \emph{Jegets} som det høieste. Hos \emph{Schel\-%
ling} finde vi den i Identitetsphilosophiens Grundbegreb, der,
trods Fordringen af Subjectivitetens Overmagt, dog gaaer til\-%
bage til den spinozistiske \emph{uendelige Væren.} Hos \emph{Hegel}
endelig finde vi den i den tvetydige Sammenblanding af Sub\-%
jectiviteten og Objectiviteten, der bringer en Værens Reflex
ind i den logiske Idee og ligesom indsmugler Naturen i et
Moment af dens Proces: i Processens Sammenslutning i Umid\-%
delbarhed. --- \emph{I det af Principet Afledede} viser Paavirk\-%
ningen sig hos \emph{Kant} naar Naturcausaliteten bliver Syns\-%
punktet for Opfattelsen af Aandsphænomenerne, idet \emph{det af}
Aanden, der bliver Gjenstand for \emph{Erkjendelse,} gaaer ind
under denne Causalitet, og den viser sig --- hvad der umid\-%
delbart hænger sammen hermed --- i den Bestemmelse af
% --- p. 179 ---
% p. 179 bears no note and no Greek.  It carries the sheet signature "12*" at
% the foot; not transcribed.
\emph{Mechanismens Forhold til Teleologien,} der, trods For\-%
dringen af denne, og trods Udtalelsen af hiins \emph{principielle
Afhængighed,} dog lader Mechanismen, i dens Uensartethed
% p.179: the letterspacing here is DISCONTINUOUS as printed -- `Selvstændighed'
% and `hele' are spaced (gaps 13--20 px at 600 dpi) but `over' between them and
% `vor Erkjendelses Gebet.' after them have ordinary body gaps (3--8 px).
% Transcribed as measured, not regularised.
med det \emph{Teleologiske,} hævde sig i \emph{Selvstændighed} over
\emph{hele} vor Erkjendelses Gebet. Den viser sig hos \emph{Schelling}
i den Reaction mod Videnskabslærens Standpunkt, der lader
% p.179: only `Natur' is letterspaced in `Naturphilosophie,' and in
% `Naturbestemthed'.
Identitetsphilosophien væsentlig blive \emph{Natur}philosophie, og
som fører til at lade Aandens \emph{Natur}bestemthed i \emph{Geniet}
blive det \emph{Høieste.} Den viser sig endelig hos \emph{Hegel} idet den
paaviste Tvetydighed ved det Absolute er gjennemgaaende i
den logiske Tankes hele Udfoldelse, og bestemmer dens Be\-%
vægelse.

Gjennem den \emph{ubevidste} Vexelvirkning mellem de to
Udviklingsformers med Bevidsthed mod hinanden hævdede
Grundtanker virker den begge omfattende Heelhed, den \emph{con\-%
crete Enhed,} der omslutter de \emph{som} Modsætninger satte
Principer, og har det Forhold, i hvilket Modsætningerne gjøre
sig gjældende \emph{som saadanne,} i sig \emph{som Moment;} og ved
dens Virken forberedes, skjult for Individerne, en \emph{høiere}
Tænkningsform, gjennem hvilken den sande concrete, ind\-%
holdsrige Enhed erkjendes. Denne høiere Form forberedes
baade naar de enkelte Modsætningsled ved den gjensidige
\emph{negative} Indvirkning udvikle sig i en \emph{Ensidighed,} der
ender i \emph{Modsigelser,} og naar de, ved paa \emph{udvortes} Maade
at overføre Bestemmelser paa hinanden, bringe en Splid til\-%
veie i Tanken, gjennem hvilken den sættes i Bevægelse hen\-%
imod den sande, \emph{indenfra} udfoldede Enhed. ---

% p.179: a faint speck stands between `det' and `allerede' in the line below;
% ABBYY reads "det-".  Not printed type; not transcribed.
Vi tilføie her, i Tilknytning til det allerede Udviklede,
en schematisk Oversigt over den nyere Idealismes Grundtankes
Genesis, saaledes som den i Formen bestemmes ved den
paaviste Indvirkning.
% --- p. 180 ---
% =====================================================================
% SCHEMA-TODO p.180 -- see RECON.md §3 and BATCH-AGENT §6.  Words captured;
% the construction (the four typographic braces, the eleven vertical connector
% rules, the two-column arrangement under head B) is set in a later pass.
% THIS MUST NOT BE LEFT AS RUNNING TEXT.
%
% p. 180 carries NOTHING but the schema: no running text, no note, no Greek,
% no folio words.  All measurements below are at 600 dpi on a page image
% 2966 x 4757 px; the folio "180" is centred at x 1621.
%
% ---- the two display heads ----
% Both are BOLD roman, centred, and NOT letterspaced.  Measured gaps: head A
% 1--10 px (median 5), head B 3--9 px (median 6), against 1--11 px in an
% ordinary body line elsewhere in the batch and 12--20 px in the letterspaced
% lines of this very schema.  Head A x 1208--2040 (centre 1624); head B
% x 1232--1963 (centre 1597.5).
%
%   head A:  "A.  Dogmatisk Idealisme."
%   head B:  "B.  Kritisk Idealisme."
%
% THE SECOND HEAD READS "B.", NOT "II.".  Read at 1200 dpi: a plain roman
% capital B with two bowls.  The ABBYY layer's "11." / "II." is an OCR error,
% and tesseract independently reads "B.".  The recon's suspicion is confirmed.
%
% ---- section A, in reading order ----
% Four author labels stand at the left, each LETTERSPACED and followed by a
% colon.  Descartes and Spinoza each have a TYPOGRAPHIC BRACE opening to the
% right ("{") between the label and its matter; Leibnitz and Wolf have none.
%
%   "Descartes:"   letterspaced label, brace spanning 2 lines
%                  (brace y 703--924, i.e. 221 px = 0.94 cm, at x about 990--1005)
%     line 1  "Jeg = activ Selvbevidsthed."
%             -- "Jeg" and "activ" letterspaced; "Selvbevidsthed." NOT.
%                The "=" is the TEXT sort (a double stroke), not math mode.
%     [connector rule]
%     line 2  "Res cogitans (substantia cog.)"
%             -- ENTIRELY UNSPACED, and "Res" is CAPITALISED here, unlike the
%                lower-case letterspaced "res" in Spinoza's and Wolf's lines.
%   [connector rule]
%   "Spinoza:"     letterspaced label, brace spanning 3 lines
%                  (brace y 1035--1369, i.e. 334 px = 1.41 cm, at x about 990--1005)
%     lines 1--2  "Substantsen som res cogitans infinita"
%                 "     og res extensa infinita."
%             -- BOTH LINES ENTIRELY LETTERSPACED, the Latin included; the
%                second line is centred under the first.
%     [connector rule]
%     line 3  "Substantsen = Deus = Natura."
%             -- entirely letterspaced; both "=" are the text sort.
%   [connector rule]
%   "Leibnitz:"    letterspaced label, NO brace
%     "Substantserne (Monaderne,"
%     "     aandeligt bestemte.)"
%             -- "Substantserne" and "aandeligt" letterspaced (gaps 12--19 px);
%                "(Monaderne," and "bestemte.)" NOT (gaps 1--7 px).  The second
%                line is centred under the first.
%   [connector rule]
%   "Wolf:"        letterspaced label, NO brace
%     "Monaden som res (substantia) simplex"
%     "     (berørende Atomets Begreb)."
%             -- "res" and "simplex" letterspaced, "Monaden som" and
%                "(substantia)" NOT; "Atomets" letterspaced, "(berørende" and
%                "Begreb)." NOT.  Second line centred under the first.
%
% NO connector rule stands between the end of section A and head B.
%
% ---- section B, in reading order ----
% A single tall brace opening to the right stands after the label "Kant:",
% spanning SEVEN lines (y 2171--3075, i.e. 904 px = 3.83 cm, at x about
% 755--800); the label sits beside the middle of the brace, level with
% "overhovedet.".  A SECOND brace, opening to the LEFT ("}"), stands to the
% right of the first FOUR of those lines (y 2184--2631, i.e. 447 px = 1.89 cm,
% at x about 1900--1950), and the Fichte column stands to the right of it.
% So the foot of the page is a TWO-COLUMN arrangement: the Kant column at the
% left, the Fichte column at the right, joined by the left-opening brace.
%
%   "Kant:"   letterspaced label, 7-line brace
%     line 1  "Jeget = Selvbevidstheden."
%             -- entirely letterspaced; the "=" is set tight against "Jeget".
%     [connector rule]
%     lines 2--4  "Selvbevidsthedens rene"
%                 "Form (Selvbevidstheden"
%                 "overhovedet."
%             -- "Selvbevidsthedens rene" letterspaced; "Form (Selvbevidstheden"
%                and "overhovedet." NOT.  THE OPENING PARENTHESIS IS NEVER
%                CLOSED: the print reads "(Selvbevidstheden overhovedet." with
%                a full stop and no ")".  Transcribed as printed; the expected
%                reading is "...overhovedet)."  Both OCR witnesses agree that
%                there is no closing parenthesis.
%     [connector rule]
%     line 5  "Homo noumenon."
%             -- entirely letterspaced, upright roman, not italic.
%     [connector rule]
%     lines 6--7  "Det fælleds intelligible Substrat"
%                 "     for Naturen og Aanden."
%             -- both lines entirely letterspaced; the second centred under
%                the first.  (End of the Kant brace.)
%
%   right column, after the left-opening brace:
%     "Fichte: Jeget som"
%     "     absolut"
%     [connector rule]
%     "     (Væren)."
%             -- "Fichte: Jeget som" and "absolut" letterspaced; "(Væren)." NOT.
%                The three lines are centred on one another.
%
%   [connector rule, in the main column, below "for Naturen og Aanden."]
%   "Schelling: Det Absolute (Fornuften, Subject-"
%   "          Objectiviteten)."
%             -- "Schelling: Det Absolute" letterspaced; "(Fornuften, Subject-"
%                and "Objectiviteten)." NOT.  Second line centred under the first.
%   [connector rule]
%   "Hegel: Den absolute Aand (den logiske Idee)."
%             -- "Hegel: Den absolute Aand" letterspaced (gaps 13--20 px),
%                "(den logiske Idee)." NOT (gaps 2--6 px).
%
% ---- the connector rules ----
% ELEVEN in all: five in section A, five in the main column of section B, and
% one in the Fichte column between "absolut" and "(Væren).".  Each is a short
% THIN VERTICAL hairline standing in the white space between two successive
% items, 4--6 px wide and 45--90 px tall at 600 dpi (i.e. about 0.02 x 0.3 cm),
% and each is set at the horizontal centre of its own column -- x 1464--1484 in
% section A, x 1418--1431 in the main column of section B.  They are plainly
% rules, not hyphens or dashes: they are vertical, they sit in the interline
% space, and they carry no serifs.  ABBYY renders them "I" or "!" and tesseract
% drops most of them.
%
% ---- the rule at the foot ----
% A SWELLED (spindle-shaped) rule closes the schema at the foot of the page:
% x 1313--1790, i.e. 477 px = 2.02 cm, hairline-thin at both ends and swelling
% to a heavy lobe at the centre.  IT IS NOT THE SAME SORT AS THE BOOK'S
% ORDINARY DIVISIONAL RULE, which is a plain even hairline (compare the rule at
% the foot of p. 169, rendered at 600 dpi for the comparison).  It is centred
% on the measure and has white space above and below, so it is transcribed
% \divrule below like the other divisional rules -- but the difference of sort
% is a REVIEW ITEM: the preamble has only one rule macro, and if the swelled
% rule occurs elsewhere in the book a second macro is wanted.
% =====================================================================
\divrule

% --- p. 181 ---
% p. 181 bears no note and no Greek.  spacing.py finds no letterspacing at all
% on this page, and the image confirms it: there is none.
%
% LARGE OPENING INITIAL, NOT IN THE HEADING LIST OF transcription.tex §3.
% p. 181 opens the concluding part of the book with a RAISED CAPITAL "V" about
% twice the cap-height of the surrounding roman (measured at 300 dpi: the V is
% about 60 px tall against 30 px for the "F" of "Foregaaende" on the same
% line), standing on the same baseline and rising above the line -- a raised
% initial, not a drop cap.  The first line is set FLUSH LEFT with no paragraph
% indent.  There is no rule and no display head above it; the division is
% marked by the initial alone (and by the swelled rule at the foot of p. 180).
% Not reproduced typographically below; recorded here and reported.
Ved de i det Foregaaende givne Oversigter over Udvik\-%
lingsgangen i Philosophiens Historie, og ved de historiske
Exempler paa de Grundforhold, der følge af den teleologisk
bestemte historiske Tankeudviklings Natur, vil den i Indled\-%
ningen formulerede Opgave for Opfattelsen af denne Tanke\-%
udvikling være bleven anskueligere, og det vil være blevet
klart, hvilke Fordringer til Fremstillingen af Philosophiens
Historie den af os hævdede Opfattelse indeslutter. Fordringen
til Fremstillingen maa først blive den, at den skal anlægges
saaledes, at de mange Systemer, der komme frem indenfor
de enkelte Hovedafsnit, saaledes samle sig om den ene Grund\-%
tanke, at vi see denne vorde, udfolde sit Indhold, og saa,
naar det er udfoldet, og den Tænkningsorganisme, som den
besjælede, opløses, at vige for en høiere Tanke, og, idet den
% p.181: a faint speck stands between `et' and `mere' in the line below; both
% witnesses read it as a full stop ("i et. mere").  At 600 dpi it is much
% lighter and smaller than this page's certain full stops.  Not transcribed.
gjør det, at vise sig som Led i et mere Omfattendes Proces:
i den gjennem alle de historiske Udviklingsstadier sig articu\-%
lerende totale Ideebevægelse, der, medens den fuldbyrder sig
gjennem den af Grundtanken bestemte Udvikling, optager
denne i sin større Heelhed. Men forat en Periodes Grund\-%
tankes Væsensforhold til de Momenter, den omslutter i sig,
og til de andre Perioders Grundtanker, og derigjennem de
begrændsede Principers levende Enhed i det omfattende Totale
skal kunne vise sig i det fulde Lys, fordres saa dernæst
netop Tydeliggjørelsen af deres i det Foregaaende paaviste
bevidst-ubevidste Spænding og Kamp i de tænkende Individer;
der fordres, at de Magter skulle blive synlige, der have virket
ligesom bagved Individernes Bevidsthed; der fordres, at Hi\-%
storiefremstilleren skal paavise de skjulte Strømninger i Tan\-%
kens Liv længe før de fra Dybet naae op til Overfladen forat
blive til Strømninger i Tankens bevidste Liv, og længe efter
at de, som et for Bevidstheden Forbigangent, paany ligesom
have sænket sig i Dybet; der fordres, at han skal forstaae at
belyse den dobbelte Betydning: den bevidste og den ubevidste,
% --- p. 182 ---
% p. 182 bears no note and no Greek.  The sentence runs on from p. 181 with no
% paragraph break.
i Tænkernes Ord, saa at den sidste ligesom skinner igjennem den
% p.182: tesseract inserts an em-dash before `Og forat' below.  There is none
% on the image -- only the compositor's wide sentence space.  ABBYY agrees.
første. Og forat kunne aabenbare de skjulte Bevægelser,
for overalt at kunne drage det Tankens hemmelighedsfulde
% p.182: a faint speck stands over the `e' of `de' in the line below; both
% witnesses read it as an acute ("dé").  At 600 dpi it is a pale round speck,
% not a printed accent, and the book uses no accented `e' in Danish.  Not
% transcribed.
Liv, i hvilket de universelle Magter raade, frem for Dagen,
maa han, hvor det er nødvendigt forat kaste Lys over den
philosophiske Tankeudviklings ofte skjulte Veie og utydelige,
næsten usynlige, Spor, benytte al anden Aandsudviklings
Historie, der altid staaer i indre Sammenhæng med hiin.

Men naar nu en historisk Fremstilling følger denne Me\-%
thode; naar den saaledes samler Alt under een Tanke; naar
den lader en heel historisk Udfoldelse vise sig som denne
ene Tankes Historie og lader Forholdet til den bestemme
hvad der er væsentligt og hvad der er uvæsentligt; naar den
forklarer Alt ud fra denne Tankes Raaden, og, idet den be\-%
standig lader det Universelles Magt blive synlig, hyppigt lige\-%
som underlægger Tænkerne en anden Tanke end den, de be\-%
vidst forholde sig til, og lader en dobbelt Magt virke i deres
Tænken, --- saa udsætter den sig for den Fare, \emph{at blive
vilkaarlig,} eller idetmindste for at paadrage sig \emph{Skinnet}
af en Vilkaarlighed, der gjør en subjectiv Opfattelse gjældende
paa den objectives Bekostning, og ikke respecterer den histo\-%
riske Kjendsgjerning i dens Givethed og Concretion. Opgaven
bliver da, baade i Virkeligheden at undgaae en saadan Vil\-%
kaarlighed og at frigjøre sig for Skinnet af den. Forat Vil\-%
kaarligheden kan undgaaes, fordres der foruden de Aandsevner,
der overhovedet ere Betingelsen for en klar Tilegnelse af de
historiske Kjendsgjerninger og en dybere Indtrængen i deres
Sammenhæng, et Studium af de philosophiske Systemer, der
paa een Gang er omfattende, universelt nok til at vise Tan\-%
kens store Bevægelser, ligesom Acterne i Tankedramaet, og
med det Samme fører saaledes ind i de enkelte store Tænkeres
Eiendommelighed, giver en saadan Fortrolighed med dem, at
man faaer Gehør for deres Tankebevægelses Rhythmus og
derved ledes til den inderlige Forstaaelse. Til dette Øiemed
% SEAM 182->183.  THE LAST WORD OF THIS FRAGMENT IS BROKEN ACROSS THE BATCH
% BOUNDARY: p. 182 ends "...en Tænkers Hoved-" and p. 183 opens "værker,
% eller ialfald at have en Oversigt...".  The discretionary hyphen and the
% comment-eaten newline below are what join them, so the pp. 183--196 fragment
% MUST begin its first text line with "værker," and MUST NOT open with a blank
% line.  Neither page's note area carries anything, so no footnote runs over
% this seam.  Run joints.py after the splice.
er det ikke tilstrækkeligt, blot at kjende en Tænkers Hoved\-%
% --- p. 183 ---
% SEAM p. 182 -> p. 183 checked (BATCH-AGENT §5).  p. 182 bears NO footnote --
% there is no footnote separator and no note matter anywhere on it -- and
% p. 183 has no note area at all, so NOTHING runs over into this batch.
% BUT THE SEAM IS A WORD BREAK: p. 182's last line ends "en Tænkers Hoved-"
% and p. 183 opens "værker,".  The pp. 169--182 batch must end its fragment
% "Hoved\-%" with no newline of its own, and joints.py must be run after the
% splice: the marker in transcription.tex has a blank line above it, so the
% joined file would otherwise acquire a paragraph break the book does not have
% AND a dangling hyphen (TRANSCRIPTION-PLAYBOOK §6).  This fragment therefore
% begins with running text on the line after this comment block.
%
% p. 183 bears no note.  No Greek.  No ɔ: mark.  No guillemets.
% Letterspacing on this page is TWO words only, both confirmed at 1000 dpi:
% "Baco" at l. 5 (first occurrence) and "Skin" at l. 24.  The two LATER
% occurrences of Baco (ll. 12 and 16) are PLAIN -- letter gaps 1--4 px at
% 1000 dpi against 14--22 px in the spaced one -- and are set plain here.
% "conditio sine qua non" is upright roman, unspaced, and takes no markup
% (transcription.tex §2: italic is reserved for Greek).
værker, eller ialfald at have en Oversigt over de mindre be\-%
tydelige. Der gives philosophiske Former, hvor den egentlige
Nøgle til Forstaaelsen af Systemet i dets dybere, omfattende
Sammenhæng maa søges, ligesom afsides, i Udsagn, der have
Tilfældighedens Præg. Et Exempel herpaa var \emph{Baco}. Hans
Udgangspunkt er Erkjendelseslæren, og dens Problemer den
Gjenstand, der væsentligst beskjæftiger ham. Man har derfor
ogsaa, i Fremstillingerne af hans Lære, saagodtsom udeluk\-%
kende holdt sig til disse, uden at betænke, at Erkjendelses\-%
lærens Problemer staae i uopløselig Sammenhæng med Pro\-%
blemet om Tilværelsens Principer, og at hines eiendommelige
Løsning hos Baco, der indvirker afgjørende paa hele den
realistiske Philosophies Udvikling, derfor kun kan forstaaes
fuldstændig gjennem Besvarelsen af dette. Men Bestemmelserne
om hvad vi vilde kalde de metaphysiske Grundbegreber, men
som Baco vilde kalde de physiske, finder man hos ham ad\-%
spredte og henkastede, og kun idet man ligesom forlænger de
Linier, der bestemme Erkjendelseslærens Grundform, seer man,
at de løbe hen til hine Punkter, og samle sig i dem.

% p. 183, l. 21: a small raised tick stands between "for," and "at" (ABBYY
% reads it as an apostrophe, "for,' at").  It is above the x-height, not on
% the baseline, and is not printed type -- an ink speck or modern pencil,
% BATCH-AGENT §7.  Deleted.
Naar saaledes det er angivet, der maa være conditio sine
qua non for, at hiin Methode kan anvendes uden at det Hi\-%
storiskes Objectivitet virkelig krænkes ved en vilkaarlig Be\-%
handling, saa bliver det nu at betragte, hvorledes Vilkaarlig\-%
hedens \emph{Skin} kan undgaaes ved Fremstillingens Art. Vil\-%
kaarligheden kunde først træde frem i den vilkaarlige Vælgen
mellem det historisk Givne, i en Optagen og Udskillen, be\-%
stemt ved et ensidigt Synspunkt. Mod en saadan Vilkaarlighed
er en Garantie given naar de Phænomener, der synes at und\-%
% p. 183, l. 30: the print reads "som en heel Udviklings\-forms," -- read at
% 700 dpi on both halves.  It is the genitive "Udviklingsforms" (the principle
% posited as that of a whole developmental form), not a compositor's error,
% and is not emended.
drage sig fra det Princip, der er sat som en heel Udviklings\-%
forms, klart sættes i Forbindelse med det derved at det for\-%
følges ud i den tilsyneladende Modsætning. Saaledes ville de
enkelte store Udviklingsstadier, med alle de tilsyneladende
Inconsequentser og Modsigelser --- i Udviklingsgangen over\-%
hovedet og i de enkelte Systemer, --- blive klare og conse\-%
quente ved Opfattelsen ud fra det Centrale, og det historiske
% --- p. 184 ---
% p. 184 bears no note.  No Greek.  No ɔ: mark.  No guillemets.  The whole
% page is ONE paragraph, continuing p. 183 and running onto p. 185.
% Letterspacing: "Væsentlige" (broken Væsent-/lige) and nothing else.  Three
% places that look spaced at low magnification measure tight at 1000 dpi and
% are set plain: "central Enhed, der samler dette Adskilte" (ll. 14--15),
% "samler det, ideelt behersker det" (l. 16) -- both merely loose
% justification -- and "lige" in "veiledende" (l. 10), a spacing.py flag.
% ABBYY and tesseract both insert em-dashes around "Umiddelbart ... Centrale"
% (ll. 16--17); there are NO dashes there on the image at 700 dpi.
Stof, saaledes som det er givet i Overleveringen, vil uden
vilkaarlig Udsondring af et Væsentligt, det vil her sige: af
Noget, der factisk har havt historisk Betydning, bringes ind
under det Princip, der er bestemt som det ledende. Naar
den antydede Opfattelse viser sig at kunne bestaae denne
Prøve ligeoverfor det tilsyneladende med Principet Uensartede,
saa gjør den, idet den vælger og grupperer, ordner og ud\-%
skiller, idet den angiver det Centrum, om hvilket det givne Stof
organisk kan forme sig, idet den fortolker det Givne ved den
veiledende Tanke, kun i det Større, ligeoverfor hele Tidsrum,
hvad den historiske Fremstilling af en enkelt begrændset Be\-%
givenhed, af en enkelt Persons Liv gjør og maa gjøre i det
Mindre. Ud fra det Givne i dets Adskilthed og i dets
uundgaaelige Ufuldstændighed skal der stræbes hen imod en
central Enhed, der samler dette Adskilte, og idet den
samler det, ideelt behersker det. Umiddelbart er dette
Centrale ikke tilstede; det er kun for den Opfattelse, der,
gjennem Aandens totaliserende Magt, forvandler den reale,
fragmentariske Mangfoldighed til ideel, principbestemt Enhed.
En Forvandling finder Sted med det Givne, idet det samles
om sit Centrum, men kun den Forvandling, der er Sandheds\-%
erkjendelsens: Forvandlingen, gjennem hvilken det \emph{Væsent\-%
lige} hæves frem i Belysning af den besjælende Enhed. Som
Kunstneren forvandler den Skikkelse, hvis Afbillede han giver,
netop ved kunstnerisk at gjengive den, ved at lade den gaae
gjennem Aanden, gjennem Phantasiens ideelle Anskuelse, ved
at lade det i Phænomenet fordunklede Ideelle træde frem som
det Besjælende og Forenende, saaledes forvandler den begri\-%
bende Erkjenden og den paa denne hvilende Fremstilling af
det historisk Givne bestandig sin Gjenstand, og maa forvandle
den, fordi den ideelle Enhed, der er Phænomenets Sandhed,
Kjendsgjerningens Sjæl, kun er for Aanden, der opfatter det
Givne, i sin Idealitet kun bringes frem af den, for saa, idet
den aabenbarer sig for Aanden, at blive ledende Norm under
dens Formen af det Givne. Men i denne Forvandling bliver
% --- p. 185 ---
% p. 185 bears no note.  No Greek.  No ɔ: mark.  No guillemets.
% Emphasis settled by measurement at 600 dpi, not by eye: letterspaced words
% run 14--25 px between glyph bounding boxes, unspaced words 1--11 px.
% Spaced: "efter sin Sandhed", "ideelle", "anden", "Bevidste", "Ubevidste,",
% "bagved", "er", "umiddelbare", "centraliserende".  Set PLAIN, though they
% adjoin a spaced word: "efter sin" in l. 2 (the second occurrence),
% "Retning" (l. 3), "og" between Bevidste and Ubevidste (l. 4),
% "Bevidstheden," (l. 13).
% "Netop i denne umiddelbare Berøring" (l. 23) is FLUSH LEFT on the image --
% no indent -- so it continues the paragraph; both witnesses put a paragraph
% break there and both are wrong.  "Vi forsøge her," (l. 30) IS indented.
det Objective uforvansket; det bliver sig selv \emph{efter sin Sand\-%
hed}, efter sin \emph{ideelle} Natur.

Vilkaarligheden kunde i en \emph{anden} Retning træde frem
i Sondringen mellem det \emph{Bevidste} og det \emph{Ubevidste}, i
Adskillelsen mellem hvad der i stræng Forstand hører Tæn\-%
keren til, og hvad der tilhører det gjennemvirkende Univer\-%
selle. Der kunde da enten vilkaarligt underskydes Tænkeren
en for ham fremmed Tanke, eller vilkaarligt benægtes, at han
med klar Bevidsthed forholdt sig til en Tankebevægelse, der
hæmmede hans Tænknings oprindelige Stræben. Mod en
saadan Vilkaarlighed giver Fremstillingen en Garantie, naar
den, hvor en fremskyndende eller en hæmmende Virken af
det Universelle, der foregaaer ligesom \emph{bagved} Bevidstheden,
statueres, gjør det anskueligt, hvorledes den høiere Tanke
allerede \emph{er} udtalt i Former, der umiddelbart føre til videre\-%
gaaende Consequentser, og dog i samme Øieblik kaldes tilbage
ved den bevidste Tankebevægelse, der ikke naaer op til den;
eller hvorledes en bevidst Stræben hen mod et Maal, der
allerede er berørt, idet alle Præmisserne til Conclusionen
ere opstillede, pludselig hæmmes, idet det allerede Opstillede
omstyrtes, ved en Magt, der ligesom overvælder Individet,
som kæmpede imod den, og syntes at have frigjort sig for den.
Netop i denne \emph{umiddelbare} Berøring af det Modstridende
ligger Garantien for det Ubevidste i Tankebevægelsen, og hvad
der gjennem den anskuelige Fremstilling ligesom paanøder sig
som Vished, faaer for Reflexionen sin definitive Betryggelse gjen\-%
nem Erkjendelsen af den teleologiske Udviklings Grundforhold.

Vi forsøge her, forat godtgjøre Rigtigheden af vor \emph{cen\-%
traliserende} Methodes Anvendelse i det tidligere Udviklede,
en Prøve af samme Art som den, den sammenhængende
Fremstilling i sin sluttede Enhed skulde give, idet vi fra den
ene af Philosophiens Hovedepocher fremdrage Phænomener,
der synes at staae som Instantser mod vor Bestemmelse af
Grundbegrebet. Prøven paa vor Methodes Rigtighed skal da
bestaae i, at ogsaa disse Phænomener blive forstaaede, og
% --- p. 186 ---
% p. 186 bears no note.  Greek: Κόσμος (one word, l. 9; capital kappa,
% acute over the omicron, final sigma; no breathing, read at 1200 dpi).
% No ɔ: mark.  No guillemets.
% Letterspaced, all confirmed at 1000 dpi: "dybere", "Naturopfattelses",
% "metaphysisk og teleologisk" (broken metaphy-/sisk), "umiddelbart",
% "Atomlæren".  "Charakteer", "Verdensheelheden", "Saadan er den" and the
% three enumerated items below are PLAIN.
% Punctuation checked at 1000 dpi: "Philosophies Historie;" is a SEMICOLON
% (tesseract reads a colon).
netop i \emph{dybere} Betydning forstaaede, ud fra hint Grund\-%
begreb og fra den Grundanskuelse, det betinger.

% THE FORECAST LIST, p. 186.  The batch brief called these "run-in
% enumerators inside ordinary paragraphs"; THEY ARE NOT.  On the image each
% item stands ALONE ON ITS OWN LINE, beginning at the ordinary paragraph
% indent, and each line ends short of the measure (the three items end at
% x 0.42, 0.68 and 0.44 of the text block against a full measure of 1.00).
% The leading is unchanged from the body -- baselines at y 0.1544, 0.1767,
% 0.1989, 0.2215 of the page, a uniform 0.0222 -- so there is NO extra space
% between the items and none between "Prøve:" and item 1.  The enumerators
% are "1)" "2)" "3)", not "1." "2." "3.".  Set with \\ plus an explicit
% \hspace*{\parindent} so that the print's indent and its uniform leading are
% both preserved; nothing here is bold and nothing is letterspaced.
Vi vælge til denne Prøve:\\
\hspace*{\parindent}1) Den antike Atomlære,\\
\hspace*{\parindent}2) Den aristippiske og epikureiske Lystlære, og\\
\hspace*{\parindent}3) Det sokratiske Standpunkt.

% DIVISIONAL RULE, p. 186 (BATCH-AGENT §6): a short CENTRED hairline with
% white space above and below, measured on the image at x 681--930 at 300 dpi
% (span 248 px = 2.10 cm, centre 805) against a text block of x 232--1332
% (centre 782).  p. 186 bears no footnote, so this is not the footnote
% separator.
\divrule

% HEAD, p. 186.  transcription.tex §3 files "1. Den græske Atomlære. (186)"
% under L3 -- "LETTERSPACED ROMAN at body size, set as a justified PARAGRAPH,
% not bold and not centred".  THAT DESCRIPTION DOES NOT FIT THIS HEAD, and the
% measurements are:
%   * BOLD.  Median stem width 9 px at 600 dpi against 5--6 px in the body
%     lines immediately above and below it on the same page (p. 186 l. 13
%     median 6, l. 21 median 5).
%   * NOT LETTERSPACED.  Inter-letter gaps within its words run 1--11 px
%     (median ~7) at 600 dpi -- indistinguishable from ordinary roman on the
%     same page ("Saadan er den i de Systemer, der", gaps 2--11) -- against
%     14--23 px in the certainly letterspaced "Naturopfattelses" nine lines
%     below and 17--25 px in "metaphysisk og teleologisk".
%   * CENTRED and SHORT, not full measure: x 584--1020 at 300 dpi (span
%     436 px = 3.69 cm, centre 802) against a body measure of ~1100 px.
%   * Slightly SMALLER than the body: x-height band 37 px against 41--42 px
%     in the body lines above and below, at 600 dpi.
% So it is a short bold centred head under a divisional rule -- typographically
% the L2 pattern, not the L3 pattern of pp. 137/145/150.  Set \textbf{} inside
% \begin{center}; NOT \emph{}.  REVIEW ITEM for the final pass: the §3 list
% should either move the 186/192/201 heads out of L3 or record that L3 has two
% distinct settings.  (The p. 192 head, measured the same way, agrees.)
\begin{center}\textbf{1. Den græske Atomlære.}\end{center}

Ved at sætte Anskuelsen af Verdensheelheden som \gk{Κόσμος}
som den charakteristisk græske, førtes vi til at bestemme
den græske \emph{Naturopfattelses} Charakteer som \emph{metaphy\-%
sisk og teleologisk}. Saadan er den i de Systemer, der
betegne den græske Philosophies Høidepunkt; saadan maatte
den være hvor de metaphysiske Grundtankers Begrebsforhold
\emph{umiddelbart} overførtes paa det Physiske, og hvor den arti\-%
culerede Forms Realisation gjennem Stoffets gradvise Beher\-%
skethed afgav Synspunktet for Opfattelsen af det concrete
Naturlige. Men ligeoverfor den platonisk-aristoteliske Natur\-%
opfattelse, der var forberedt ved Pythagoræerne og bevaredes
i den efteraristoteliske Philosophies Hovedformer, staaer indenfor
den græske Philosophies Kreds en Opfattelse, der minder om
den moderne Tids physiske og mechaniske Naturopfattelse, og
ikke synes at ville gaae ind under samme Princip og samme
Grundanskuelse som hiin. Det er \emph{Atomlæren}, hvilken vi
to Gange see fremtræde i den græske Philosophies Historie;
som oprindelig og selvstændig Theorie i den første Periode,
som Gjentagelse og Efterligning i den efteraristoteliske Tid.
I denne sidste Fremtræden er det ikke saa meget Theorien
selv med de ubetydelige og vilkaarlige Modificationer, den
anbringer i den ældre, der har Interesse, som selve Optagel\-%
sen: det, at en Theorie af denne Art har kunnet optages
efterat den ideelle Naturopfattelse i dens Gjennemførthed
havde gjort sig gjældende.
% --- p. 187 ---
\setcounter{footnote}{0}
% p. 187 bears ONE note, *) called at "moderne" in l. 2, five lines long and
% SELF-CONTAINED (it ends "fandt Tiltro."); nothing runs over onto p. 188.
% No Greek.  No ɔ: mark.  No guillemets.
% Letterspaced in the body: "mechaniske", "væsentlig Forskjel",
% "eiendommelig Udformning".  "moderne" itself is PLAIN (letter gaps 2--5 px
% at 1000 dpi).
% Letterspaced in the note: "Demokrit" (broken De-/mokrit), "Baco",
% "Descartes".  "Aristoteles's" in the note is PLAIN.
% APOSTROPHE, p. 187: the note reads "Aristoteles's hadefulde", checked at
% 800 dpi -- the book's ordinary 's form (cf. Sokrates's p. 193, Kosmos's
% pp. 191/192, Aristoteles's p. 190).  Set as a plain ' per BATCH-AGENT §9.
Den græske Atomlæres Berøringspunkter med den mo\-%
derne\footnote{Det er betegnende, at Atomlærens væsentligste Repræsentant: \emph{De\-%
mokrit}, ved den nyere Philosophies Begyndelse blev Gjenstand
for Lovprisninger af de betydeligste Tænkere (f. Ex. hos \emph{Baco} og
\emph{Descartes}), og at taabelige Fabler om Aristoteles's hadefulde
Misundelse mod ham fandt Tiltro.} \emph{mechaniske} Naturopfattelse ligge klart for Dagen.
Det Naturliges Forklaring søges i selve det Naturlige; dette
bliver den virkende Aarsag, den eneste Aarsag, som det me\-%
chanisk bevægede Materielle. Materien bliver ikke et Tje\-%
nende, der bøier sig for Formen og som bevægelig almeen
Mulighed faaer sin Existentsform bestemt ved denne, den
bliver tværtimod det Uforanderlige, der, forat kunne bruges
som Forklaringsgrund for Tilværelsens Mangfoldighed, maa
bestemmes som et oprindelig Mangfoldigt. Det, der bliver
til af de materielle Elementer, bliver Aggregatet, i hvilket
Elementerne vedblive at være det Selvstændige. Det For\-%
bindende og Adskillende er den Tilfældighed, der tiltager sig
Nødvendighedens Navn.

Men ved Siden af disse Berøringspunkter mellem det
Antike og Moderne træder der saa en \emph{væsentlig Forskjel}
frem, der viser den antike Atomlæres Sammenhæng med den
græske Tænknings Princip og den græske Grundanskuelse.

Naar Atomlæren gjør sig gjældende i \emph{eiendommelig
Udformning} og consequent Gjennemførelse i den nyere
Videnskab, saa har den til sine Forudsætninger den dualistiske
Sondring af Aand og Natur, og en Erkjendelseslære, der lader
det gjennem Erfaringen Modtagne være Erkjendelsens eneste
Indhold, og Naturen eneste Gjenstand for Erfaringen. Idet
Tænkningen, paa Grundlag af disse Forudsætninger, stræber
at føre Alt ind under den fra Aanden væsensadskilte Naturs
Synspunkt, og at finde et Middel til en exact Bestemmen af
Erfaringens concrete Gjenstand, kommer den til Atomlæren
som en Hypothese, der er tjenlig til Forklaringen af den ydre
% --- p. 188 ---
\setcounter{footnote}{0}
% p. 188 bears ONE note, *) called at "omsættes" in l. 30, three lines long
% and SELF-CONTAINED (it ends "Iklædning."); nothing runs over onto p. 189.
% No Greek.  No ɔ: mark.  No guillemets.
% "Système de la nature" (l. 4) is set in the SAME UPRIGHT ROMAN as the
% Danish and is not letterspaced -- checked at 900 dpi, and the accent over
% the e is a GRAVE.  Per transcription.tex §2 it therefore takes no markup at
% all; ABBYY reads "Systeme", tesseract "Systéme", and both are wrong.
% Letterspaced in the body: "gamle", "ved Speculationen," (the comma is
% inside the spaced run), "almindelige Form" (broken alminde-/lige),
% "ad Speculationens Vei alene,", "Atomerne og det tomme Rum",
% "physisk benævnede", "physisk" (l. 25), "metaphysiske" (l. 28),
% "Værens", "Ikke-Værens Begreber,", "physiske" (l. 31).
% Set PLAIN although adjacent: "og" between Værens and Ikke-Værens (its o and
% g touch at 1200 dpi, against 17--23 px gaps in the spaced words on the same
% line); "(det Udfyldte og det Tomme)"; "ere vel"; "Bestemmelser, og";
% "løses Problemet"; "omsættes" IN THE BODY (l. 30) -- it is spaced only in
% the note.
% p. 188, l. 4: a small raised tick follows "nature)." at the right margin
% (tesseract reads it as an apostrophe).  Not printed type; deleted.
Erfarings concrete Phænomener, der vise hen paa det Mate\-%
rielles Deelthed i Elementer af ubestemmelig Lidenhed, og
som kan udvides til en indre Erfarings usynlige Objecter (jfr.
Forklaringen af Aandsfunctionerne i Système de la nature).
Denne Hypothese gjøres saa, idet Erfaringsprincipet glemmes,
til Dogma, og Dogmet forkynder Atomernes Væren som evige,
uforanderlige, i al Sammensætning selvstændige Grunddele,
der ved deres nødvendighedsbestemte Sammentræden consti\-%
tuere alt realt Værende.

Idet Atomlæren kommer frem i den \emph{gamle} Philosophie,
er Sondringen mellem Aand og Natur endnu ikke klaret for
Tanken, og det er ikke ved Erfaringen og ved Analysen af Na\-%
turens concrete Phænomener, men \emph{ved Speculationen,} at
Philosophien føres til at opstille Læren om Atomerne. Kun
det ganske almindelige Problem: at forklare Muligheden af
Vorden og en Tilværelsens Mangfoldighed, har Anskuelsen af
den reale Tilværelse bragt frem, men i denne \emph{alminde\-%
lige Form} løses Problemet \emph{ad Speculationens Vei alene,}
og netop Problemets Almindelighed maa føre ind paa denne
Vei. \emph{Atomerne og det tomme Rum} (det Udfyldte og det
Tomme) ere vel \emph{physisk benævnede} Bestemmelser, og
Forklaringen af Tingenes Tilblivelse ved Atomernes Bevægelse
hviler paa en \emph{physisk} Anskuelse: Anskuelsen af Rummet og
af det i Rummet værende Bevægede; men bag de physiske
Navne skjule sig kun abstracte, \emph{metaphysiske} Begreber: de
Begreber, hvis Forhold allerede tidligere Tænkere havde grundet
over. Det er \emph{Værens} og \emph{Ikke-Værens Begreber,} der om\-%
sættes\footnote{Hos Atomisterne \emph{omsættes} i Virkeligheden Begreber i Anskuelser.
Hos \emph{Heraklit} er Begrebet \emph{oprindelig} iført Anskuelsens Iklæd\-%
ning.} i \emph{physiske} Bestemmelser idet de anskues som Ato\-%
merne og det tomme Rum, som det værende Udfyldte og det
værende Tomme; men de beholde endnu i denne Omsæt\-%
ning det tidligere Indhold. Det, at Atomerne og det Tomme
% --- p. 189 ---
% p. 189 bears no note.  No Greek.  No ɔ: mark.  The whole page is ONE
% paragraph, continuing p. 188 and running onto p. 190.
% GUILLEMETS, p. 189: the one quotation, «Det Værende er slet ikke mere end
% det Ikke-Værende», was read at 1200 dpi -- the opening mark points LEFT and
% the closing mark points RIGHT, the book's normal INWARD pair.  No defect.
% The quoted matter is NOT letterspaced: the "er" that opens l. 3 measures an
% e--r gap of 8 px at 1200 dpi against 96 px for the following word space,
% i.e. the fount's ordinary sidebearing.
% Letterspaced: "Begrebet" x3 (the first broken Be-/grebet), "det Værendes og
% det Ikke-Værendes" (here the "og" IS spaced, unlike p. 188 l. 30),
% "metaphysiske" (broken metaphy-/siske, l. 9), "materielle", "physiske"
% (l. 13), "nødvendig", "Anaxagoras", "Atomisterne,", "Eleaterne", "rumlige",
% "teleologiske" (broken tele-/ologiske), "fornægtes,", "metaphysiske"
% (broken metaphy-/siske, l. 28), "physiske" (l. 30).
% Set PLAIN although adjacent: "Atomer," at the head of l. 6 (the spaced run
% is "Begrebet" alone, the noun after it is plain -- three times over);
% "Charakteer"; "Værende"; "altomfattende"; "Modsætninger"; "af metaphysiske
% Begreber i physiske Anskuelser" in l. 20, where the same two words that are
% spaced in l. 13 are plain.
% p. 189, l. 15: a small speck stands before "græske" (ABBYY reads ".græske").
% Not printed type; deleted.
er forklarende Værens-Grund for det Tilværende, udtrykker
netop det Samme som Atomisternes Paradox: «Det Værende
er slet ikke mere end det Ikke-Værende». Tilværelsen for\-%
klares ikke physisk, men metaphysisk, ved Atomlæren; \emph{Be\-%
grebet} Mangfoldighed (Flerhed) forklares ved \emph{Begrebet}
Atomer, \emph{Begrebet} Vorden (Bevægelse) forklares ved \emph{det
Værendes og det Ikke-Værendes} Forhold; og om en
exact Bestemmen af Phænomenerne er der ikke Tale. Saa\-%
ledes har Naturopfattelsen i Atomlæren samme \emph{metaphy\-%
siske} Charakteer som den græske Naturopfattelse overhovedet.
Og at denne Lære gjør \emph{materielle} Tilværelsesprinciper gjæl\-%
dende og omsætter de allerede klarede metaphysiske Begreber
i \emph{physiske} Anskuelser, forklares ved den som \emph{nødvendig}
paaviste Udviklingsgang i den græske Philosophies ældste
Periode. Fra Stoffet gik den Tanke ud, der gjennem Stoffet
skulde finde Ideen; kun i en ufuldkommen og uconsequent
Antydning naaede den hos \emph{Anaxagoras} til Aandens Begreb,
og hvor den tidligere berørte det Ideelle, forvandlede det sig
strax for den til et legemligt Værende. Til den Omsætning
af metaphysiske Begreber i physiske Anskuelser, der finder
Sted hos \emph{Atomisterne,} svarer hos \emph{Eleaterne} den For\-%
vandling, ved hvilken Væren, trods de mod det Ideelle sig\-%
tende Bestemmelser, der tillægges den, bliver til det anskuede
\emph{rumlige} Værende. Og at den græske Naturopfattelses \emph{tele\-%
ologiske} Charakteer ikke blot, som i ældre philosophiske
Forsøg, ikke kommer tilsyne i Atomlæren, men udtrykkelig
\emph{fornægtes,} bliver en uundgaaelig Følge af, at de \emph{metaphy\-%
siske} altomfattende Modsætninger, idet de antage physisk
Form, som disse \emph{physiske} Modsætninger, der fastholdes i
Anskuelsens Sondring, maae blive Principer for al Væren, og
ikke give nogen Plads for det af Anaxagoras statuerede, paa
udvortes Maade ordnende, aandelige Bevægelsesprincip. Den
gamle Atomlære, der giver Principerne physisk Form og, med
Udelukkelse af det Teleologiske og ubevidst Anthropomorphi\-%
stiske, gjør Tingenes Heelhed til Product af de materielle
% --- p. 190 ---
% p. 190 bears no note.  No ɔ: mark.  No guillemets.
% Greek (two words, both read at 1200--1600 dpi): \gk{Νοῦς} -- capital nu,
% omicron, upsilon carrying a CIRCUMFLEX (a tilde, resolved positively; the
% same reading as p. 88), final sigma; and \gk{χωριστόν} -- acute over the
% omicron, no breathing on either word since both open with a consonant.
% The fount's cursive rho is typed as plain ρ per BATCH-AGENT §4 rule 4.
% Letterspaced: "anden", "Epikur," (l. 7), "gamle", "uvæsentlige og
% vilkaarlige" (the "og" IS inside the run), "metaphysiske" (l. 13), "ny"
% (l. 14), "praktiske", "physiske" (l. 17), "enkelte", "materialistisk",
% "første", "ved Forholdet mellem Aristoteles og Plato." (broken Aristo-/teles),
% "oversandselige Væren" (broken oversandse-/lige), "Entelechie;" (broken
% Entele-/chie), "naturlige Universum", "Aristoteles's Skole", "Universets
% Form,", "levende Universum seer det Guddommelige selv." (broken Gud-/dommelige),
% "Stoikerne", "Epikur" (l. 35).
% Set PLAIN although adjacent: "Form" after "første" (l. 25); "Belysning,"
% after "ny"; "Modifica-/tioner"; "altsaa i det".
% "ny" was measured because it is only two letters: n--y gap 20 px at
% 1200 dpi against 4 px in the "en" immediately before it, and 14--25 px in
% the certainly spaced words on the same page.
Grunddeles tilfældig-nødvendige Sammenhobning, viser sig
saaledes som en Consequents af et nødvendigt Moment i den
ideelle Verdensanskuelses Genesis, og bag dens Forklaring af
Tilværelsen skjuler sig en metaphysisk Forklaring, der viser
tilbage til den Aands Grundretning, af hvilken den er frem\-%
gaaet, og antyder det, hvortil denne Retning maatte føre.

Hvor Atomlæren \emph{anden} Gang kommer frem: hos \emph{Epikur,}
finde vi, i Overensstemmelse med den efteraristoteliske, væ\-%
sentlig kun paa det Ethiskes Gebet productive, Philosophies
Charakteer, paa Physikens Gebet den \emph{gamle} Atomlære gjen\-%
optagen, kun med \emph{uvæsentlige og vilkaarlige} Modifica\-%
tioner, der ikke berøre Theorien som Heelhed eller tilintetgjøre
dens Grundpræg som \emph{metaphysiske} Theorie. Kun faaer Theo\-%
riens metaphysiske Charakteer hos Epikur en \emph{ny} Belysning,
idet det bliver \emph{praktiske} Motiver, Hensyn til det Individs
Interesse, der er sig selv det sidste Øiemed, der bestemmer
Valget af denne \emph{physiske} Theorie; og idet Ligegyldigheden
af en Erkjendelse af det \emph{enkelte} Physiske udtrykkelig ud\-%
tales. Optagelsen af en \emph{materialistisk} Theorie bliver for\-%
klaret ved hvad der i Oversigten over den græske Philosophies
Perioder er udviklet angaaende Nødvendigheden af en mate\-%
rialistisk Enhedslære som Gjennemgangsled i den sidste Pe\-%
riodes Udfoldelse. Idet Dualismen ogsaa hos Aristoteles er
uovervunden, gaaer Tankens Retning hen efter en Enhed, der
trænger den ene Modsætning tilside. Denne Enheds \emph{første}
Form er allerede antydet \emph{ved Forholdet mellem Aristo\-%
teles og Plato.} Den platoniske Idee i dens \emph{oversandse\-%
lige Væren} føres af Aristoteles ind \textbf{i} Stoffet som \emph{Entele\-%
% p. 190, l. 30: the `i' of "ind i Stoffet" is set from a MARKEDLY HEAVIER
% sort than its neighbours.  Measurement (BATCH-AGENT §3): at 1200 dpi its
% median stem width is 20 px; every other glyph stem on the same line
% measures 8--12 px (median 8), including the `i' of "ind" three glyphs
% earlier and the `l' of "lige" at the line head.  It is the widest stem on
% the line by a factor of ~1.7--2.5.  Set \textbf{} so the fact is greppable.
% NO CLAIM is made here as to whether this is a foul case or an emphasis; it
% joins the sites at pp. 42, 89, 95, 100, 119 and 147 for the single
% adjudication pass at the end.
chie;} det \emph{naturlige Universum} kommer eet Skridt nær\-%
mere henimod at opfattes som den høieste Forms Realisation.
I \emph{Aristoteles's Skole} drages denne Consequents, der lader
\gk{Νοῦς} ophøre at være et \gk{χωριστόν} forat blive \emph{Universets
Form,} altsaa i det \emph{levende Universum seer det Gud\-%
dommelige selv.} Hos \emph{Stoikerne} fremtræder Pantheismen
i Enhed med Materialismen, og \emph{Epikur} føier kun, ved at
% --- p. 191 ---
% p. 191 bears no note.  No ɔ: mark.  The whole page is ONE paragraph,
% continuing p. 190 and running onto p. 192.
% NO GUILLEMETS ANYWHERE ON THIS PAGE.  The recon sweep suspected a quotation
% at "mangfoldige Gjentagelser i Tiden"; the marks there are EM DASHES, read
% at 650 dpi -- "den historiske Udviklings --- ligesom det aandelige
% Universums --- mangfoldige Gjentagelser i Tiden".  Every page of this batch
% was read for guillemets and only pp. 189, 193 and 194 carry any.
% GREEK, p. 191: \gk{απειρον} in "det tomme Rums" (l. 33).  THE ALPHA IS SET
% UNACCENTED, per BATCH-AGENT §4 rule 2 and the precedent at pp. 51/52/81.
% The mark over it is the breathing-plus-acute compound printed as one sort.
% At 2400 dpi (8x interpolation over the 300 ppi original) it does separate
% into two lobes with one light pixel between them, but neither lobe is
% ~4 px wide and neither shows a resolvable shape: the curl's opening cannot
% be assigned a direction, so ἄ cannot be separated from ἅ.  NOTHING IS
% SUPPLIED.  (At p. 86 the same compound WAS resolvable -- two components with
% a clear gap, the left curl opening LEFT = smooth -- so ἄπειρον is the likely
% reading; it is not the reading this site establishes.)  The cursive rho is
% typed as plain ρ.
% Letterspaced: "Materialismen Benægtelsen af det Teleologiske" (broken
% det/Teleologiske across the line), "atomistiske Materialisme", "Enhed"
% (after Kosmos's), "mangfoldige," (l. 17), "Verdener" (l. 19), "græske
% Grundanskuelse" (broken Grundan-/skuelse), "resulterende
% Hensigtsmæssighed.", "uendelig mange", "historiske Udviklings" (broken
% hi-/storiske), "aandelige Universums" (broken Univer-/sums),
% "Gjentagelser i Tiden".
% Set PLAIN although adjacent: "sidestillede, mod hverandre ligegyldige";
% "faae Plads i det tomme Rum."; "af Tilværelsen."; "tilfældig"; "Verdener,"
% at l. 34 (the SECOND occurrence, unlike l. 19); "mangfoldige" at ll. 25 and
% 36 (both plain, unlike l. 17); "ligesom det"; "en Opfattelse,".
% p. 191, l. 10: a small raised tick stands between "høre" and "de" (both
% witnesses read "høre'de"), and a second one above the l of "til" on the
% same line.  Neither is printed type -- both sit above the x-height, off the
% baseline -- and the printed reading is "der høre de to Sphærer til:".
% Deleted and logged, BATCH-AGENT §7.
optage Atomlæren, til \emph{Materialismen Benægtelsen af det
Teleologiske}. Den \emph{atomistiske Materialisme} er con\-%
sequentere end den pantheistiske, fordi, skjøndt hverken Ato\-%
met eller den almene Materie kan blive Sandsningens Gjen\-%
stand, Atomet ialfald har det Sandseliges Enkelthedsbestem\-%
melse. Ogsaa Benægtelsen af det Teleologiske er consequent.
Modsætningen mellem Aand og Natur er paa dette Standpunkt
traadt frem, og med den Modsætningen mellem de Aarsager,
der høre de to Sphærer til: Forsøget paa, at forklare Tingenes
Orden ved den ene Art Aarsager med Udelukkelse af den
anden, er en uundgaaelig Consequents hvor Principernes
Modsætning viser sig som uforsonlig. Epikur har derfor den
logiske Tankes Medhold, om end praktiske Motiver have været
væsentlig afgjørende for ham i hans Benægtelse af det Teleo\-%
logiske. Med Forkastelsen af dette og af det ideelle Princip
hænger det saa sammen, at Anskuelsen af Kosmos's \emph{Enhed}
opgives, at \emph{mangfoldige,} sidestillede, mod hverandre lige\-%
gyldige \emph{Verdener} faae Plads i det tomme Rum. Det Sam\-%
menholdende er opgivet med det Ideelle, Formen, hvis Enhed
betingede Universets Enhed. Men ligesom hele den epikurei\-%
ske Philosophie allerede har viist sig som et Gjennemgangsled,
gjennem hvilket den græske Tænkning, der ikke kan gjennem\-%
føre Materialismen, vender tilbage til den idealistiske Verdens\-%
anskuelse, saaledes viser sig indenfor dens Opfattelse af de
mangfoldige Verdener Virkninger af den \emph{græske Grundan\-%
skuelse} af Tilværelsen. En Ordning kan Anskuelsen heller
ikke undvære hos Epikur: derfor bliver hver Verden sat som
en ved Nødvendighedsforhold behersket Enhed, med en ligesom
tilfældig \emph{resulterende Hensigtsmæssighed.} Og naar
der i det tomme Rums \gk{απειρον} bliver Plads for \emph{uendelig
mange} Verdener, saa finder disses Mangfoldighed i Rummet,
der ophæver Tilværelsens Enhed, dog en Parallel hos Græ\-%
kerne, i den for dem charakteristiske Opfattelse af den \emph{hi\-%
storiske Udviklings} --- ligesom det \emph{aandelige Univer\-%
sums} --- mangfoldige \emph{Gjentagelser i Tiden}: en Opfattelse,
% --- p. 192 ---
\setcounter{footnote}{0}
% p. 192 bears ONE note, *) called at "Indhold" in the last line before the
% head, eight lines long and SELF-CONTAINED (it ends "følgende Perioder.").
% No ɔ: mark.  No guillemets.
% Greek (one phrase, l. 12, read at 1400 dpi): \gk{νόησις τῆς νοήσεως} --
% acute over the omicron of νόησις, CIRCUMFLEX over the eta of τῆς, acute
% over the eta of νοήσεως.  No breathings: all three words open with a
% consonant.  The comma after it is roman, outside the \gk{}.
% Letterspaced in the body: "Plato og Aristoteles.", "Grændsen",
% "Fuldkomne.", "Tiden er den vordende Uendelighed" (broken Ti-/den),
% "uendelig Tidsudvikling", "aldrig tilfredsstillet", "ikke" (l. 8),
% "denne" (l. 9), "evigt Kredsløb" (broken Kreds-/løb), "afsluttet".
% Set PLAIN although they sit between two spaced runs: "som det" between
% Grændsen and Fuldkomne (measured at 1400 dpi -- s-o-m and d-e-t touch,
% against 15--25 px gaps in Grændsen and Fuldkomne); "(Aristoteles), men en";
% "for Menneskeaanden,"; "Stræben henimod et uendeligt Maal, kunde";
% "gjentagende det samme Indhold" (ABBYY's "gj en tagende" is OCR noise).
% Letterspaced in the note: "Uendeligheden i Rum og Tid.", "Rummets",
% "i Tiden".
der ogsaa findes hos \emph{Plato og Aristoteles.} Denne Opfat\-%
telse af det Historiske hang i den græske Bevidsthed sammen
med Anskuelsen af \emph{Grændsen} som det \emph{Fuldkomne}. \emph{Ti\-%
den er den vordende Uendelighed} (Aristoteles), men en
\emph{uendelig Tidsudvikling} for Menneskeaanden, en \emph{aldrig
tilfredsstillet} Stræben henimod et uendeligt Maal, kunde
Grækeren \emph{ikke} tænke sig. Maalet maatte kunne naaes og
naaes i \emph{denne} Tilværelse. Men naar det var naaet, saa
maatte Tiden dog gaae sin uendelige Gang, og Aanden, der
ikke kunde hvile og ikke gaae op i den Guddommen forbe\-%
holdte evig \gk{νόησις τῆς νοήσεως}, maatte, forat faae en Be\-%
vægelse, gaae tilbage til Begyndelsen og i et \emph{evigt Kreds\-%
løb} gjentage sin Udviklingshistorie. Ligesom derfor, efter
Atomlærens Opfattelse, Verden i det Uendelige mangfoldiggjør
sig i Rummet, saaledes mangfoldiggjør, efter den almindelige
græske Opfattelse af Historien, Menneskeslægtens Udvikling
sig i Tiden, idet store Perioder, hver omfattende en \emph{afsluttet}
Menneskeslægtens Historie, afløse hinanden, gjentagende det
samme Indhold\footnote{Det er interessant at følge den græske Tænknings frugtesløse Kamp
mod \emph{Uendeligheden i Rum og Tid.} Mod \emph{Rummets} Uende\-%
lighed sættes Kosmos's Begrændsethed, men udenfor det begrændsede
Universum kan Rummet ikke fjernes for Anskuelsen, men kun benæg\-%
tes ved en vilkaarlig, Anskuelsen modsigende, Begrebsbestemmelse.
Og \emph{i Tiden} sættes vel Menneskehistorien som afsluttet i de be\-%
grændsede Perioder, men der bliver Plads til en Uendelighed af
saadanne paa hinanden følgende Perioder.}.

% DIVISIONAL RULE, p. 192: short, CENTRED hairline, x 672--924 at 300 dpi
% (span 252 px = 2.13 cm, centre 798) against a text block centre of 788,
% with white space above and below.  It stands ABOVE the note area, well
% clear of the footnote separator, and is not that rule.
\divrule

% HEAD, p. 192, measured exactly as the p. 186 head was and agreeing with it:
% BOLD -- median stem 9 px at 600 dpi against 6 px in the body line below;
% NOT LETTERSPACED -- intra-word gaps 1--15 px (median ~9), the same as
% ordinary roman on this page, against 15--25 px in the certainly
% letterspaced "uendelig Tidsudvikling" eleven lines above;
% CENTRED and SHORT -- x 406--1178 at 300 dpi (span 6.54 cm, centre 792)
% against a text-block centre of 788 and a full measure of ~1090 px.
% So \textbf{} inside \begin{center}, not \emph{}.  See the p. 186 head note.
%
% PRINTER'S DEFECT, p. 192 head: the `g' of "og" is UNDER-INKED -- the bowl
% below the baseline prints solidly, the upper bowl only faintly (checked at
% 1600 dpi).  The word is "og"; the witnesses read "o;", "o:" and "«g" off
% this damage.  Transcribed as the word it is, with the defect logged here
% rather than emended into or out of the text.
\begin{center}\textbf{2. Den Aristippiske og Epikureiske Lystlære.}\end{center}

En Grundanskuelse, der i Formens og Stoffets umiddel\-%
bare, harmoniske Enhed seer Formen som det Væsentlige; der
i Menneskevæsenet, hvis harmoniske Udfoldelse er Livsopgaven,
seer det Ideelle som det i den umiddelbare Enhed Constitutive,
% --- p. 193 ---
% p. 193 bears no note.  No ɔ: mark.
% GREEK, p. 193, both read at 1400--2400 dpi:
%   \gk{ἡδονή} -- the ROUGH breathing over the eta is positive here, a "C"
%   opening to the RIGHT, and the acute over the final eta likewise.  This is
%   one of the edition's calibration words (transcription.tex §GREEK) and it
%   was read FIRST in this batch, before the rest of the page's Greek, to set
%   the eye.
%   \gk{ἔχω οὐκ ἔχομαι} -- inside guillemets.  Over the epsilon of both ἔχω
%   and ἔχομαι stands the breathing-plus-acute compound; unlike the p. 191
%   ἄπειρον it SEPARATES at 2000 dpi into two components with a light gap
%   between them, the left a curl whose opening faces LEFT (= SMOOTH on this
%   edition's calibration, the same glyph as on οὐσία p. 97 and ἄπειρον p. 86,
%   NOT the rough of ὕλη p. 35 or of ἡδονή above it on this very page) and the
%   right a stroke running down to the left (= acute).  Over the upsilon of
%   οὐκ stands a single curl, opening LEFT = smooth.  The book's kappa is the
%   cursive ϰ; typed as plain κ, on the same reasoning as θ and ρ in
%   BATCH-AGENT §4 rule 4.
% GUILLEMETS, p. 193: the one pair, around the Greek, points INWARD «…» --
% read at 1400 dpi.  No defect.
% Letterspaced: "sandselige Bestemthed," (broken Be-/stemthed), "Lysten"
% (l. 7), "atomistisk", "kyrenaiske" (LOWER-CASE k, checked at 1400 dpi;
% ABBYY reads "Ryrenaiske"), "Aristippos,", "Lysten" (l. 13), "positiv",
% "enkelte", "legemlig", "Indsigten." (the FIRST of the two on l. 20),
% "Nydelsens Kunst,", "Lysten" (l. 33), "Indsigten" (l. 33), "Philosophien".
% Set PLAIN although adjacent: "(ἡδονή)," -- the parentheses and comma are
% roman and unspaced; "Sokrates's Tilhører og Sophistikens Discipel.";
% the SECOND "Indsigten" on l. 20; "som Maalet og"; "som den væsentlige
% Betingelse"; "der gjør os det muligt".
% p. 193, l. 20: a small raised tick stands over "om" in "derigjennem om det
% alment Græske" (ABBYY reads "o"m"), and a second over the k of
% "kyrenaiske" in l. 12.  Neither is printed type; deleted.
% The sheet signature `13' stands at the foot of p. 193; not transcribed.
og som derfor bestemmer Dyden som Viden, maa føre til en
Ethik med ideel Charakteer, hvilende paa Erkjendelsen, med
den Enkeltes Underordnen under det Almene, med hvilket
han er harmonisk sammenknyttet (jfr. de tidligere Udviklinger).
Men indenfor den græske Philosophie finde vi dog en ethisk
Anskuelse, der opfatter Individet efter dets \emph{sandselige Be\-%
stemthed,} gjør \emph{Lysten} til Maalet, og ved dette Maal lader
Individet i \emph{atomistisk} Sondring løsne sig fra det Almene.
Denne ethiske Anskuelse bliver det da Opgaven at forklare
ud fra det alment Græske.

Vi møde den først i den \emph{kyrenaiske} Philosophie hos
dennes Grundlægger \emph{Aristippos,} Sokrates's Tilhører og So\-%
phistikens Discipel. Som Maalet for Menneskets Stræben be\-%
stemmer han \emph{Lysten} (\gk{ἡδονή}), som \emph{positiv} Lyst, ikke som
Frihed for Smerte; som den \emph{enkelte} Lyst, det nærværende
Øiebliks Lyst, ikke som almeen Livs-Tilstand, som Lyksa\-%
lighed; og endelig: som \emph{legemlig} Lyst. Men til Virkelig\-%
gjørelse af dette Lystens Maal, og som bestemmende Nydelsens
Form, fordrer han, hvad der minder om det Sokratiske, og
derigjennem om det alment Græske: \emph{Indsigten.} Indsigten
skal af ethvert Livsforhold udvinde den Lyst, det indeslutter;
den skal, ved Bevidstheden om Lystens Følger, kunne vælge
ret mellem de forskjellige i og for sig lige Lyster, der tilbyde
sig; den skal frigjøre for falske Forestillinger, for Fordomme
og Indbildninger; den skal fjerne Længselen efter det For\-%
svundne og Attraaen efter det Tilkommende; den skal skaffe
Frihed i selve Nydelsen («\gk{ἔχω οὐκ ἔχομαι}») og den Bevidst\-%
hedens Uafhængighed, der i hver Situation giver os Tilfredshed
i det Nærværende. Den skal med eet Ord muliggjøre en
\emph{Nydelsens Kunst,} der gjør os det muligt, med Bevidst\-%
hedens Frihed og Sikkerhed at hengive os til Øieblikkets
Nydelse, og saaledes at udvinde den størst mulige Lyst af
Livet. \emph{Lysten} som Maalet og \emph{Indsigten} som den væsent\-%
lige Betingelse blive saaledes satte i den nøieste Forbindelse,
og \emph{Philosophien} betegnes som den Vei, der fører til et i
% --- p. 194 ---
% p. 194 bears no note.  The whole page is ONE paragraph, continuing p. 193
% and running onto p. 195.
% THE ɔ: MARK IS PRESENT on p. 194 and is CONFIRMED on the image at 1400 dpi:
% "en Efterklang af den egentlige ɔ: legemlige Lyst" -- a reversed c followed
% by a colon, at the end of the line after "egentlige".  ABBYY reads "o:",
% tesseract a bare ":".  This is the one candidate site in the batch that
% bears it; every other page was read for it and none has one.
% GREEK, p. 194 (one phrase, l. 4, read at 1400 dpi): \gk{λεία κίνησις} --
% acute over the iota of λεία and over the iota of κίνησις; no breathings
% (both words open with a consonant).  The cursive kappa is typed plain κ.
% GUILLEMETS, p. 194: «sjælelig Lyst» -- INWARD pair, read at 1400 dpi.
% Letterspaced: "Øieblikkets" (l. 3), "Øieblikkets Sag," (broken Øieblik-/kets),
% "leve i Øieblikket, concentrere sig deri;", "Indsigtens Forhold til Lysten.",
% "Paa positiv", "ikke" (l. 18), "at afværge det Forstyrrende,", "enhver",
% "som" (l. 35).
% TWO PLACES WHERE THE SPACED RUN SKIPS A WORD, both measured at 1400 dpi and
% both set as printed:
%   l. 20 "Paa positiv Maade" -- "Paa" and "positiv" are spaced, "Maade" is
%   tight;
%   l. 35 "enhver Lyst som Lyst et Gode:" -- "enhver" and "som" are spaced,
%   both occurrences of "Lyst" are tight.
% Set PLAIN: "Lyst" after "Øieblikkets" in l. 3; "og netop heri ligger".
Sandhed menneskeligt Liv. Men dog bliver Lystens og Ind\-%
sigtens Forhold strax tvivlsomt, saasnart det erindres, at Ly\-%
sten udtrykkelig er sat som \emph{Øieblikkets} Lyst, fordi kun det
Nærværende tilhører os, og fordi Lysten som den \gk{λεία κίνησις},
der finder Sted i den Fornemmende som sandselig bestemt, i
sin Opstaaen er et Forsvindende, der maa gribes i Øieblikket,
og vel kan bevares som en «sjælelig Lyst», men saa kun som
en mattere Erindring, som en Efterklang af den egentlige ɔ:
legemlige Lyst. Naar nemlig Lysten, Livets Maal, er \emph{Øieblik\-%
kets Sag,} saa maa den Nydende, der forholder sig til dette
Maal, ogsaa \emph{leve i Øieblikket, concentrere sig deri;}
og netop heri ligger Vanskeligheden for Bestemmelsen af
\emph{Indsigtens Forhold til Lysten.} \emph{Paa positiv} Maade
skaffe Lyst tilveie, vil Indsigten \emph{ikke} kunne, thi saaledes
maatte der gaaes ud over Livet i Øieblikket og Standpunktet
vilde være opgivet. Opgaven skulde saa nærmest blive den,
\emph{at afværge det Forstyrrende,} og saa at lade den Lyst,
som Livsgangen fører med sig, uhæmmet udfylde Øieblikket.
Under dette Synspunkt maatte da Indsigtens Function, saa\-%
ledes som den ovenfor er angivet, lade sig opfatte. Men
ogsaa herved opstaae uløselige Vanskeligheder. Det kan ikke
gjøres klart, hvorledes Indsigten, hvis Indhold kun kunde
blive Bevidstheden om Lysten som Livets Maal og som væ\-%
rende i Øieblikket, skulde kunne give Bevidstheden Uafhæn\-%
gighed af Øieblikkets tilfældige Udfyldelse, der ogsaa kan være
en Udfyldelse blot med Smerte; Smerten maa være forstyrrende
for Lykken, da Bevidstheden ikke maa abstrahere fra Øie\-%
blikket, og ikke har et høiere Livsindhold nærværende i sig.
Og netop fordi Bevidstheden ikke har et Indhold, der er
høiere end Nydelsen, kan det ikke heller gjøres klart, hvor\-%
ledes Indsigten skulde kunne give en Frihed i den Nydelse,
der skal udfylde den til Øieblikket bundne Menneskeaand.
Det kan ikke heller blive klart, hvorledes Indsigten kan vælge
mellem Lysterne efter Hensynet til deres Følger, thi deels er
\emph{enhver} Lyst \emph{som} Lyst et Gode: der er ingen Artforskjel
% --- p. 195 ---
% p. 195 bears no note.  No ɔ: mark.  No guillemets.
% Greek (one word, last line, read at 1400 dpi): \gk{Κόσμος} -- capital
% kappa, acute over the omicron, final sigma; no breathing.
% Letterspaced: "Indsigtens Uddannelse,", "Umiddelbarheden", "Skjønheden",
% "sandselig".  Set PLAIN: "som" between Skjønheden and sandselig;
% "Skjønhed" after it; "Vilkaarligheden"; "Aristips Opfattelse".
% PRINTER'S DEFECT, p. 195, l. 21: the print reads "ikke kan opfyldes, De to
% Momenter" -- a COMMA where a full stop is required, followed by the
% fount's wide sentence space and a capital D (read at 1200 dpi; the comma's
% tail is unmistakable and there is no period beside it).  Transcribed AS
% PRINTED per BATCH-AGENT §2; not on the author's Trykfeil leaf.
% This drives check.py's punctuation counts, not its quote balance.
% The sheet signature `13*' stands at the foot of p. 195; not transcribed.
mellem dem, og selv Gradforskjellen forsvinder naar Livets
Opløsen i discrete Øieblikke tager Sammenligningens Maale\-%
stok bort; deels viser Hensynet til Følgerne ud over det Nær\-%
værende, hvad der strider mod Forudsætningerne. Der blev
da kun tilbage, at skyde alt det tilside, der som forstyrrende
Forestillinger, som Længsel og Attraa, trængte sig ind paa
Øieblikkets Nydelse, og at lade Indsigten føre til, at Concen\-%
trationen i Øieblikket blev fuldstændig, at Fortid og Fremtid
ophørte at være for den Nydendes Bevidsthed. Men denne
yderligste Consequents, der vilde umuliggjøre \emph{Indsigtens
Uddannelse,} idet der ingen Plads blev til en saadan i det
til Nydelse bestemte Liv, vilde directe modsige Indsigtens Væ\-%
sen som Bevidsthed; thi Bevidsthedens Sphære er det Almene;
den kan ikke bindes til Øieblikket og til Øieblikkets Lyst, men
dens Væsen viser bestandig ud derover. Hiin Fordring paa en
Sammenknytten af Lysten, i Aristips Opfattelse, og Indsigten
viser sig da vel som en Fordring, der ikke kan undgaaes fordi
den Nydende er Menneske, men tillige som en Fordring, der
ikke kan opfyldes, De to Momenter, der i deres Væsen ere
heterogene, maae igjen skille sig fra hinanden, og Standpunktet
maa gjennemløbe en Udvikling, i hvilken den indre Modsigelse
aabenbarer sig, og den Aand gjør sig gjældende, af hvilken
dette ethiske Standpunkt, der vil gribe Mennesket midt i Til\-%
værelsen, som sandseligt og bevidst nydende denne Sandselig\-%
hed, og som frigjort i den ved Indsigten, er fremgaaet, som
en ensidig, og dog relativ berettiget Form.

For nu at see, hvorledes dette Standpunkt, der opløser
Livet i Nydelsens Skjønhedsmomenter, har sin Rod i den græske
Bevidsthed, og forat forstaae dets Udvikling og Opløsning ud
fra denne, bringe vi det paa følgende Maade i Sammen\-%
hæng med det Centrale i den græske Livsanskuelse. Idet vi
mindes, hvorledes for Grækerne \emph{Umiddelbarheden} var Be\-%
stemmelsen for Enheden mellem Tilværelsens Grundmomenter,
hvorledes \emph{Skjønheden} som \emph{sandselig} Skjønhed blev En\-%
hedens Anskuelsesform i \gk{Κόσμος} og i den individuelle Tilvæ\-%
% --- p. 196 ---
% p. 196 bears no note.  No Greek.  No ɔ: mark.  No guillemets.
% The whole page is ONE paragraph, continuing p. 195 and running onto p. 197.
% Letterspaced: "æsthetisk", "Individet som levende Kunstværk.", "ideale,"
% (l. 25), "Vorden,", "Væsens" (the first element only of "Væsensbestemmelse"
% -- "bestemmelse" is tight, checked at 1200 dpi), "Tiden gjør sin Magt",
% "nydende Livsfølelses Indhold".
% Set PLAIN, both spacing.py single-word flags refuted on the image at
% 1200 dpi: "Punkt," (l. 18, gaps 2--5 px) and "Umiddelbare" (l. 3, where
% only "æsthetisk" before it is spaced).  Also plain: "harmoniske
% Tilværelse" after "ideale,"; "gjældende." after "Magt"; "vilde" at the
% foot of the page.
%
% SEAM p. 196 -> p. 197 checked (BATCH-AGENT §5).  p. 196 bears NO footnote at
% all, so nothing can run over; p. 197's note area is likewise empty at its
% head.  The last word of p. 196 is "vilde", COMPLETE -- no dangling hyphen --
% but the sentence continues on p. 197 ("nu som før blive Øieblikkets
% Livsfylde", read at 600 dpi), so this fragment ends WITHOUT a trailing blank
% line and joints.py must be run after the splice: the marker in
% transcription.tex has a blank line below it and the joined file would
% otherwise acquire a paragraph break the book does not have.
relse, og hvorledes Opfattelsen af Personligheden og af Menne\-%
skelivets ethiske Udvikling bevarede det \emph{æsthetisk} Umiddel\-%
bare, søge vi at forklare Lystlæren ved at fastholde Anskuel\-%
sen af \emph{Individet som levende Kunstværk.} Vi kunne da
gaae ud fra Forestillingen om en plastisk Fremstilling af en
ideal Menneskeskikkelse, som vi ved en Fiction lade, ligesom
Pygmalions Billedstøtte, paa een Gang faae Liv og Bevidsthed.
Denne Bevidstheds Indhold maa da blive selve Skikkelsens har\-%
moniske, ideale, i sig tilfredsstillede Liv. Der bliver ingen
Stræben, ingen Attraa: i plastisk Ro hviler Livet i sig selv, i
rige tidløse Øieblikke. Bevidstheden er som det Lys, der fal\-%
der over Livets Harmonie, og i hvilket denne seer sig selv:
den er hvilende som Livet. Lade vi nu denne ideale Enkelt\-%
skikkelse blive sat ind i en ligeledes ideal, af Harmonie gjen\-%
nemtrængt Heelhed, af hvilken den udgjør et Led, saa bliver
kun Livsfølelsen, Nydelsen, mere omfattende; det er Heelhe\-%
dens Liv, der individualiserer sig paa dette bestemte Punkt, i
den i Heelheden harmonisk indordnede Individualitet. Der er
en Strømning af Liv, men Harmonien er i ethvert Øieblik til\-%
stede: Tilværelsen bliver en Række af ideale Øieblikke; hvert
enkelt afspeiler sig i Bevidstheden i plastisk Ro. Bevidstheden
bliver heller ikke her en Søgen, en Stræben: den er kun som
Harmoniens Gjenklang i sig selv. Men ombytter man nu den
\emph{ideale,} harmoniske Tilværelse med en Tilværelse, i hvilken
Alt er i \emph{Vorden,} i Forandring; i hvilken uløst Modsætning,
Tilbliven og Undergang faae Plads, og lader det ideale Individ,
uden at dets \emph{Væsens}bestemmelse opgives, blive et Led i denne
Tilværelse, saa forandrer Forholdet sig. Harmonie og Dishar\-%
monie vexle nu med hinanden, det Modstridende og det Over\-%
ensstemmende berører Individet, og Bevidstheden bliver ikke
længer blot en Afspeiling af Individets og af Universets i In\-%
dividet centraliserede Harmonie; men den maa optage et Frem\-%
med, Uharmonisk i sig, og derved drages den ud over Indi\-%
videt og bort fra Evighedsøieblikket: \emph{Tiden gjør sin Magt}
gjældende. Men den \emph{nydende Livsfølelses Indhold} vilde
% =====================================================================
% BATCH pp. 197--210.  Page map: printed 197--209 = PDF +9; printed 210
% = PDF +11 (PDF 219/220 are a SECOND EXPOSURE of printed 208/209).
% Every page was obtained through pagemap.py; no offset was added by hand.
%
% SEAM, HEAD: p. 197 begins IN MID-SENTENCE.  p. 196 ends "...Men den
% nydende Livsfølelses Indhold vilde" and p. 197 opens "nu som før blive".
% There must be NO paragraph break at the splice point (joints.py).
% p. 196 bears no note, and p. 197's note area is empty: nothing runs over
% into this batch.
%
% SEAM, FOOT: p. 210 ends "...at udtrykke dets Væsen; deels" and p. 211
% opens "viser det sig gjennem de Udsagn af Sokrates om Guderne" --- again
% NO paragraph break at the splice point.  p. 210 bears no note, so
% nothing runs over out of this batch.  (p. 211 opens a note of its own,
% "*) Man sammenligne med Stederne hos Plato...", which belongs to that
% page and is that batch's business.)
% =====================================================================
% --- p. 197 ---
% No note, no Greek on this page.  Guillemets «Indsigten» point INWARD.
nu som før blive \emph{Øieblikkets Livsfylde}, det er, paa dette
Standpunkt: \emph{den harmoniske Bevægethed ved et An\-%
det.} Bevidsthedens Opgave maatte altsaa blive: at \emph{beskytte}
% "at" is NOT letterspaced; only "beskytte" is (checked at 350 dpi, the
% following "Nydelsens Øieblik" is plainly normal set).
Nydelsens Øieblik mod det Uharmoniske, Forstyrrende, og dette
Forstyrrende maatte da deels tænkes som de forstyrrende Fore\-%
stillinger, deels som den Stræben, den Attraa og det Savn,
der gaae ud over Øieblikket. Kun en Beskytten af Øieblikket
kunde blive Opgaven: skulde Bevidstheden («Indsigten») paa
directe Maade tilveiebringe Lyst, vilde dette netop vise ud over
Øieblikket.

Det Forhold, vi sidst have skildret, er netop det af \emph{Ari\-%
stip} satte. For ham staaer Individet, som et levende Kunst\-%
værk, midt i Sandselighedens Verden, i dens Bevægelse og
Strømning: Opgaven er: \emph{kunstnerisk at nyde Øieblikket,}
at lade Bevidstheden holde det Forstyrrende borte, og saa lade
sig gjennemstrømme af den harmoniske Livsfølelse, der faaer
% "Lyst." on the next line is NOT letterspaced -- checked at 600 dpi
% against the certainly unspaced "Lyst," nine lines above.
sin Bestemthed ved \emph{Øieblikkets} Lyst. Men saaledes kom\-%
mer hiin Modsigelse tilsyne: at Tanken skal fastholde Øieblik\-%
ket, concentrere sig i Øieblikket, medens dens Natur er, ikke
at være i Øieblikket, og netop idet den bevidst sænker sig
deri, at række ud derover. Og Øieblikkets Indhold bliver be\-%
standig det Tilfældige, det for Tanken Fremmede. I den fin\-%
gerede, ideale, harmoniske, tidfrie Tilstand vare de to Elemen\-%
ter: Bevidsthed og Lyst, sammensmeltede til een uadskillelig
harmonisk Enhed; nu er, ved Disharmoniens Indtrængen, Ad\-%
skillelsen gjort, og Bevidstheden, Tanken, har sondret sig fra
det sandselige Nydelsesindhold, der bliver tilfældigt for den.
Men Tanken er det Mægtigste, den ideale Menneskeskikkelses
egentlige Væsen; den maa derfor, naar Adskillelsen er gjort,
gjøre sin Magt gjældende ved at fortrænge det andet Element.
Dette, der allerede er antydet derigjennem, at Fordringen af
\emph{Aandsfrihed i Nydelsen} i Virkeligheden sætter Nydelsen
som det Aandslivet Underordnede, skeer, under den kyrenaiske
Skoles Udvikling, hos \emph{Theodoros og Hegesias}, positivt:
gjennem Tankemomentets Udvikling, negativt: gjennem Lyst\-%
% --- p. 198 ---
\setcounter{footnote}{0}%
% GREEK, all read off the image at 2400 dpi.  This page carries the
% edition's cleanest breathing calibration and it was used for the whole
% batch: the rough on the η of ἡδονή (line 1) is a "C" opening to the
% RIGHT, the smooth on the α of ἀφροσύνη (line 3) is the mirror image,
% a "ɔ" opening to the LEFT.  The two are positively distinguishable at
% this magnification.
% χαρὰ (twice) carries a GRAVE, not an acute: its mark falls to the right
% while the acute on the υ of λύπη in the same line rises to the right.
% Set as printed, though the lexical form is χαρά.
elementets Forsvinden. \emph{Theodorcs} sætter \gk{ἡδονή} og \gk{πόνος}
% PRINTER'S ERROR, transcribed as printed: the page reads "Theodorcs",
% with a `c' where the `o' of "Theodoros" belongs.  The glyph is open on
% the right at 2400 dpi against the two closed `o' bowls earlier in the
% SAME word, and both OCR witnesses read `c' here while both read
% "Theodoros" on p. 197.  Expected: Theodoros.
som ligegyldige; Betydning faae kun \gk{χαρὰ} og \gk{λύπη}, der op\-%
staae af \gk{φρόνησις} og \gk{ἀφροσύνη}. Men Indholdet af \gk{φρόνησις}
bliver kun Bevidstheden om Lystens og Ulystens Ligegyldig\-%
hed: Indholdet af \gk{χαρὰ} bliver den Vises tomme, men netop
derfor ubegrændsede Selvfølelse i den Uafhængighed, denne
Bevidsthed giver. Den Vise er \gk{αὐτάρκης}; alle objective For\-%
hold tabe deres Betydning for ham. --- Hos \emph{Hegesias} fører
Overbevisningen om Lystens Tilfældighed og Relativitet, og om
Umuligheden af at realisere Aristips Princip, til at sætte
\emph{Smerteløshed} som Maalet og at bestemme det \emph{Negative:}
at afværge Onder, som det \emph{Eneste}, der staaer i vor Magt;
og Fortvivlelsen om at give Livet et Nydelsesindhold driver
ham til at sætte \emph{Livet} som det \emph{Ligegyldige}, \emph{Døden} som
et \emph{Tillokkende}.

Ved at sees fra det angivne Synspunkt viser den kyre\-%
naiske Lystlære, i dens Fremtræden og Udvikling, sig som et
\emph{eiendommelig græsk} Phænomen og forstaaes ud fra den
græske Grundanskuelse. Det er \emph{Skjønhedens} Moral, Aristip
lærer, idet han fra den \emph{sandselige Skjønheds} Synspunkt
seer Individets Liv som Summen af de harmoniske Øieblikke,
i hvert af hvilke Sandseligheden er harmonisk afspeilet i Be\-%
vidsthedens Idealitet. Det \emph{Æsthetiske}, der hos \emph{Plato} og
endnu hos \emph{Aristoteles er Moment i den ethiske} Handlen,
er her \emph{Alt}, og \emph{Virkelighedens Verden} er ved en Illusion
forvandlet til \emph{Skjønhedens tidløse Verden}, der bliver In\-%
dividets Skueplads indtil Tiden og Virkelighedens Modsætninger
bringe Illusionen til at briste. Og idet der nu søges en Norm
for Livet, gaaer Udviklingen ilsomt bort fra hint Udgangspunkt,
og ender i dets skarpeste Contrast: i \emph{Døden}. Ogsaa \emph{Plato}
fordrede en Afdøen\footnote{Ved denne Møden vises der hen til, at
\emph{Døden} er Skjønhedslivets Maal, hvad enten det opfattes
sandseligt eller ideelt.}, men gjennem Ideen reddede han Tanken
% PRINTER'S ERROR in the note, transcribed as printed: "Ved denne Møden",
% capital M, read at 1200 dpi (the two stems with the V between are
% unmistakable) and so read by both witnesses.  The text calls for
% "Døden" / "Afdøen"; "Møden" is not the word the sense wants.
ud af det Foranderliges Verden og lod kun Sjælen afdøe fra
% --- p. 199 ---
% GREEK: ἡδὺς βίος is the book's grave/acute calibration pair -- rough
% breathing on the η, a GRAVE on the υ (falling to the right), and an
% ACUTE on the ι of βίος (rising to the right), all three positively
% seen at 2400 dpi.  ὅταν carries a compound rough-plus-acute over the
% ο; the breathing element reads as a "C" opening right, but a compound
% over a round vowel is at the resolution limit -- see the note there.
dette forat finde sit Liv igjen i det evige Værende. Den \emph{ky\-%
renaiske} Anskuelse, der vil holde sig i \emph{Sandseverdenens}
Vorden, tvinges ved dennes Dialektik ud i Døden, og i den
Død, der ikke er Ideens Begyndelse, men det Vordendes Op\-%
hør. Den Lære, hvis Princip er Nydelsen, ender med at tale
saa tillokkende om det, der staaer i hemmelighedsfuld Sam\-%
menhæng med Nydelsen: om Døden, at de, som hørte det,
frivilligt valgte den, og \emph{Hegesias}, den kyrenaiske Skoles
egentlige Afslutter, fik Navnet \emph{Peisithanatos}.

Hvor Lystlæren i den sildigere græske Philosophie frem\-%
træder i en fornyet Skikkelse, hos \emph{Epikur}, viser sig det samme
Vexelforhold mellem de to Elementer; men saaledes, at Tanke\-%
elementet fra Begyndelsen af er stærkere betonet og dets Over\-%
vægt afgjørende. Idet \emph{Epikur} sætter det \emph{høieste Gode i
Lyksaligheden}, og lader denne bestaae i \emph{Lyst}, saa at alt\-%
saa Lysten er det egentlig \gk{τέλος}, bestemmer han, til Forskjel
fra Aristippos, Lysten som \emph{varig Tilstand}, og sætter den
\emph{sjælelige} Lyst som den høieste og varigste, om end al Lyst
i sidste Instants gaaer tilbage til et Legemligt. Den \emph{ethiske
Opgave} formuleres saaledes: ved den fornuftige Indsigt (\gk{φρόνη\-%
σις}) at træffe et saadant Valg mellem de i Grad forskjellige
Lyster, at der deraf resulterer \gk{ἡδὺς βίος}. Men ogsaa her maa
Lysten vise sig som det, der i sin Tilbliven er \emph{uafhængigt
af Individet}, og som Betingelsen for en \emph{varig} Nydelse
sættes derfor \emph{Nøisomhed og Tilfredshed med det Givne,
Begrændsningen af Attraaer og Lidenskaber}. Det
Væsentlige bliver \gk{ἡ τοῦ σώματος ὑγίεια καὶ ἡ τῆς ψυχῆς
ἀταραξία}, og denne Ataraxie tilveiebringes idet Forstanden
tilintetgjør Fordomme. Idet den \emph{positive} Lyst forudsætter
en Trang, altsaa en Smerte, bliver Lystens egentlige Maal og
Væsen at bestemme \emph{negativt}, som \emph{Smertefrihed:} \gk{ὅταν
% p. 199: over the ο of ὅταν the scan shows a COMPOUND mark, breathing
% plus accent, as one cluster.  The accent element rises to the right
% (acute) and the breathing element reads as a "C" opening right, i.e.
% ROUGH, when set beside the certain rough of ἡδονή and the certain
% smooth of ἀφροσύνη on p. 198.  Less certain than an isolated breathing;
% recorded here so it can be re-checked in the final Greek pass.
μὴ ἀλγῶμεν οὐκέτι τῆς ἡδονῆς δεόμεθα}. Den Vise har, ved
% δεόμεθα is printed with the cursive variant ϑ; per the header, plain θ.
den fornuftige Indsigt, i \emph{ethvert} Livsforhold Sindsro og Til\-%
fredshed.

Her er altsaa Nydelsen saaledes aandeliggjort, og saaledes
% --- p. 200 ---
% No note, no Greek.  NOTE THE COMPOUND EMPHASIS PATTERN that recurs on
% this page and on p. 203: the compositor letterspaces only the FIRST
% element of a compound noun and sets the rest normally
% ("A a n d s bestemmelse", "S a n d s e l i g h e d s- og
% I d e a l i t e t s momentet").  It is one word in each case, and it is
% set here as \emph{Aands}bestemmelse etc., not as two words.
draget ud af Afhængighedsforholdet til det Ydre, at kun \emph{Ly\-%
stens Navn} i Virkeligheden er blevet tilbage. Det, der stræ\-%
bes henimod, er \emph{det individuelle Livs smertefrie Ufor\-%
styrrethed i Aandsdannelsens Bevidsthed om Uaf\-%
hængigheden af det Ydre.} Saaledes er i Nydelsesbegrebet,
i hvilket oprindelig \emph{det individuelle Kunstværks Sand\-%
selighedsside} afspeilede sig, gjennem det ideelle Moment den
Bestemthed trængt ind, der berøver Nydelsen dens sandselige
Indhold og gjør den til en \emph{Aands}bestemmelse, til \emph{den uen\-%
delige frie Selvbevidstheds Hvile i sig selv.} Ikke
sandselig Nydelse, men aandelig Bevidsthed, ikke det Givne,
men det Erhvervede, bliver det Høieste. Her er saaledes denne
ethiske Retning, der tager sit Udgangspunkt fra Lysten, ved
Hævdelsen af det Moment, der oprindelig skjultes, ført hen til
det Punkt, hvor den mødes med den græske Ethiks Ho\-%
vedstrøm, og, idet den forener sig med denne, viser det fæl\-%
leds Udspring.

Og hvad der gjælder om Forholdet mellem \emph{Sandselig\-%
heds- og Idealitets}momentet, gjælder ogsaa om Forholdet
mellem det \emph{enkelte} Individ i dets Isolation og det \emph{Almene}.

I \emph{Epikurs} Lære staaer Individet som et \emph{aandeligt
Atom}, hvilende i sin abstracte Selvbevidsthed, afsluttet i sig
ligeoverfor det objective Almene. I sin Selvbevidstheds Frihed veed
Individet sig uafhængigt af Alt: alt det Ydre, og selve det, i
hvilket Grækeren tidligere fandt sit Livs substantielle Indhold:
Familie- og Statslivet, bliver fremmedt for det og ligegyldigt.
Subjectet bliver sin egen Verden, gjennem Autarkien og Ata\-%
raxien løsnet fra den ydre. Den ethiske Opgave bliver ikke:
at realisere sædelige Ideer i Verden, men at skaffe det tæn\-%
kende Individ Uafhængigheden af Verden. Den abstract uen\-%
delige Subjectivitet, der som saadan er sit eget Maal, hævder sig
selv i Sondringen fra Naturen og det menneskelige Samfund
som det ubetinget Høieste, og hviler tilfredsstillet i sig. Men
under denne Individets atomistiske Frigjørelse gjør Forholdet
til det Almene sig dog gjældende. Individet, der søgte sig
% --- p. 201 ---
selv som nydende, fandt sig i sin, ved Aandsdannelsen fri\-%
gjorte Aand, i sin Selvbevidstheds om end abstracte Uende\-%
% The letterspacing on the next line is SPLIT, and it was checked twice
% at 2400 dpi: "a l m e n e" is spaced, "Bestemthed" is NOT, "s o m" is
% spaced again, "Menneske:" is not.  Two separate emphases, not one.
lighed, væsentlig efter sin \emph{almene} Bestemthed \emph{som} Menneske:
i dets Bevidsthed begynder derfor \emph{Humanitetens} almene
Begreb at arbeide sig frem medens Statsborgerlighedens og Na\-%
tionalitetens fordunkledes. Og netop idet Individet, der er
ligesom et aandeligt Atom, i sin indre Frihed løftes op over
Individualitetens Begrændsning og Afhængighed, mødes det i
Aandens Universalitet med \emph{de andre frie Individualite\-%
ter} og sluttes sammen med dem i \emph{et frit Humanitetens
Samfund}. Vel aabenbarer nu den i Bevidstheden indtrængte
Dualisme, som hele denne Udvikling har til Forudsætning, sig
atter her i den humane Bevidstheds abstracte og negative
Charakteer, i de frie, ligesom sublimerede, Subjectivite\-%
ters Abstraction, men under Abstractionen gjør det Almenes
Form sig gjældende som Formen for \emph{en aandelig, alt det
Menneskelige i sig til Enhed forbindende Humani\-%
tetens} \gk{Κόσμος}.

% DIVISIONAL RULE, p. 201: short, CENTRED, generous white above and
% below, standing between the end of the paragraph and the head.  Not the
% footnote separator -- this page bears no note.
\divrule

% L3 HEAD, p. 201 ("3. Sokrates.", the third of the 186/192/201 sequence
% in transcription.tex §3).  MEASURED at 600 dpi, as §3 requires, and the
% answer is NOT the answer the other L3 heads gave:
%
%                       x-height   median stroke   intra-word letter gaps
%   p. 201 head            41 px       10 px            14--16 px
%   p. 201 body            43 px        6 px            2--10 px (med. 5)
%   p.  79 L2 head         48 px       13 px            1--7 px  (med. 3)
%   (certainly bold)
%
% So this head is at BODY SIZE (41 vs 43 px x-height) but at 1.7x the
% body's stroke weight, and stroke/x-height puts it at 0.244 against
% 0.271 for the certainly bold p. 79 head and 0.140 for the body.  It is
% BOLD.  It is ALSO letterspaced -- 14--16 px gaps against 2--10 px in
% the body of the same page -- and it is CENTRED and alone on its line.
% It is therefore NOT the letterspaced-roman justified paragraph that
% pp. 80 and 145 turned out to be; it is a bold letterspaced centred head.
% Per BATCH-AGENT §3 a bold head must not be set \emph{}, so it is set
% \textbf{} and the letterspacing is recorded here rather than in markup.
% pp. 186 and 192 (same sequence) should be measured the same way and
% transcription.tex §3 corrected once the pattern is known.
\begin{center}
\textbf{3. Sokrates.}
\end{center}

Blandt alle den græske Oldtids Skikkelser er der ingen,
der viser sig saa eiendommelig, saa afstikkende, tilsyneladende
saa helt uensartet med alle de andre og med den græske Aands
almindelige Charakteer, som \emph{Sokrates}. Ugræsk som han var
i sit Ydre, synes han ogsaa at være i sit Indre. Midt iblandt
de Vidende, midt i det Folk, hvis Tænkere gjorde Viden til
Menneskets Væsen, staaer han som den, der kun veed sig som
\emph{uvidende}, og med denne \emph{Uvidenhedens Vished} ikke som
Skeptikeren behøver al sin Kraft til at holde sig svævende i
en Bevidsthed, der aldrig tør fæstne sig til Viden, men med
Uvidenheden til Udgangspunkt kan finde Sikkerhed for sin
Bevidsthed i ethisk Handlen og i det religiøse Forhold. Mel\-%
lem dem, der søge Vished i Tankens Klarhed og dem, der
% --- p. 202 ---
% PENCIL, p. 202: two short vertical strokes in the LEFT margin beside
% the lines "Personlighedens Uendelighed er hos ham ifærd med at / komme
% til Gjennembrud".  Modern pencil, not type; not transcribed.  p. 202 is
% not on the header's list of twelve known pollution sites.
søge den i det ydre Orakels Dunkelhed, staaer han som den,
der søger den i det Orakel, han bærer i sit Indre, og som
lyder til ham som den iboende Guddoms advarende Stemme.
\emph{Personlighedens Uendelighed} er hos ham ifærd med at
komme til Gjennembrud, en ny Tids ethisk-religiøse Princip
forkynder sig i hans Bevidsthed, om end Uendelighedens Syn
kun viser sig for ham som et blændende Lynglimt, men ikke
vinder Skikkelse og Aabenbaring for Tanken, og kun holdes
fast i Ekstasens bevidstløse Stirren.

Denne for sin Samtid og for Eftertiden gaadefulde Skik\-%
kelse, i hvis Indre man anede det Høieste: «kostelige Gude\-%
billeder, gjemte i Silenen», denne \gk{ἄτοπος}, der i sine Tilhø\-%
% p. 202: over the α of ἄτοπος stands a COMPOUND breathing-plus-accent
% mark.  The accent element rises to the right (acute); the breathing
% element is a "ɔ" opening LEFT, matching the certain SMOOTH on the α of
% ἀφροσύνη (p. 198) and not the certain ROUGH on the η of ἡδονή on the
% same page -- the two calibration marks were rendered at 2400 dpi and
% set beside this one before the call was made.  The header names the
% compound over a low vowel (ἄπειρον, pp. 52, 81) as the one thing the
% scan does not resolve; this instance WAS resolved by that comparison,
% but the basis is recorded here so the final Greek pass can revisit it.
reres Bevidsthed proteusagtigt forvandlede sig i saa mangfol\-%
dige Former, og som gjennem sin Ikke-Viden bliver Anledning
til al den sildigere græske Videns Udvikling, har dog, trods
sin tilsyneladende Fremmedhed, sin dybeste Rod i den græske
Aands Væsen; han --- og han alene --- repræsenterer heelt og
gjennemført et Aandsstandpunkt, der ligger som Consequents
i den græske Grundanskuelse og fuldstændig forklares ved
denne. Derfor skinner hos ham den græske Baggrund overalt
igjennem, selv hvor han synes at fjerne sig meest afgjørende
fra sin Tid og sit Folk.

Gaae vi for et Øieblik tilbage til det, der for den græske
Tanke blev \emph{Grundforholdet i Tilværelsen}: Forholdet mel\-%
lem Verdenskunstværkets to Momenter, Stoffet og Formen,
saa see vi en \emph{dobbelt} Consequents at være indesluttet deri
med Hensyn til den menneskelige Tanke. Idet hint Forhold
paa den ene Side viser sig saaledes for os, at Formen, der
i Stoffet realiserer \emph{sig selv}, er det sande og ene Værende,
saa sees fra dette Synspunkt Formen som det ubetinget Her\-%
skende, som det, der ene bestemmer det Formedes Væsen,
og deri ligger, ifølge Formens Slægtskab med \gk{Νοῦς}, den Con\-%
sequents, \emph{at Tanken (Begrebet) er det Værendes Væ\-%
sen}, og med Hensyn til den \emph{menneskelige} Tanke, at denne,
som det i Mennesket Herskende, Menneskevæsenet Bestem\-%
% --- p. 203 ---
\setcounter{footnote}{0}%
% PENCIL, p. 203: a vertical stroke in the RIGHT margin at the head of
% the page and a second beside "I det her Fremsatte".  Not transcribed.
% The ɔ: id-est mark stands in line 16 of this page, confirmed at 1200
% dpi: a reversed c opening LEFT, followed by a colon.
mende, i Kraft af sin Enhed med den guddommelige \gk{Νοῦς},
har Muligheden til, \emph{som vidende at gjennemtrænge
Objectiviteten}, hvis Væsen selve Tanken er. --- Men idet
Forholdet paa den anden Side viser sig saaledes, at Stoffet sees,
ligesom under en anden Belysning, som den af Formen \emph{ikke}
satte Betingelse for Formens Realisation; saa faaer Stoffet en Uaf\-%
hængighed af Formen, en Art af selvstændig Virken. Og dette
Uafhængighedsforhold fører saa indenfor det \emph{Relatives} Sphære\footnote{Jfr.
Platos og Aristoteles's Physik.}
--- thi i Totaliteten er Forholdet uforstyrret, og derfor er den
\emph{guddommelige} Viden uforstyrret Viden\footnote{Jfr. Xenoph. Mem.
I, 1, 8, IV, 6, 7. Plato, Apol. 20. 23.} --- til et \emph{nega\-%
tivt} Forhold, til et \emph{Modsætningsforhold} mellem Stof og
Form, og den \emph{menneskelige} Aand sees, i Kraft deraf,
skjøndt ved Formen \emph{væsentlig} bestemt som \emph{Viden}, idet
\emph{den ikke behersker det fremmedartede Stof}, men \emph{for\-%
dunkles} ved det, som den, hvis \emph{væsentlige} Viden, under Indi\-%
videts Existents ɔ: under Formens Sammenknytning med det
Fremmedartede, faaer \emph{Uvidenhedens} Form. Netop i disse to
Consequentsers Forskjel er Forskjellen given mellem \emph{Sokra\-%
tes} paa den ene Side, og \emph{Plato} og \emph{den hele græske} Specu\-%
lation paa den anden. \emph{Sokrates} standser ved \emph{Uvidenhe\-%
den}, opgiver Speculationen; men gjør, idet det Ideelle, For\-%
men, har Realitet for ham som Væsensbestemmelse og Maal,
det \emph{Ethisk-Religiøse} til sin Interesse; \emph{Plato} fastholder,
om og det andet Forhold hos ham, ligesom ogsaa hos Aristo\-%
teles, kommer tilsyne, skjøndt kun for let at affærdiges, \emph{Væ\-%
sens}bestemmelsen Viden som \emph{Virkeligheds}bestemmelse, og
% Again the split compound: "Væsens" and "Virkeligheds" are letterspaced,
% "bestemmelsen"/"bestemmelse" are not, and each is one printed word.
Viden bliver for ham til \emph{speculativ Viden}, hvis Gjenstand
er \emph{Ideen}; hans Interesse bliver \emph{Philosophien}.

I det her Fremsatte er den sokratiske \emph{Uvidenhed} taget
i \emph{egentlig} Forstand, --- ikke ironisk --- kun med den nær\-%
mere Bestemmelse, at Uvidenheden, medens den sees som en
\emph{hæmmet Viden}, selv bliver et \emph{negativt Videns Ind\-%
% --- p. 204 ---
% The letterspaced run "negativt Videns Ind-hold." RUNS ACROSS THE PAGE
% BREAK; \emph{} opens on p. 203 and closes on the first word here.
% PENCIL, p. 204: a cursive word in the LEFT margin beside "Disse Udsagn
% ere ligesom Summen", and a vertical stroke lower down beside "bestandig
% med sine Omværende".  Not transcribed.
hold.} I egentlig Forstand maa den opfattes, naar der ikke
skal gjøres Vold paa det historisk Overleverede. Vel viser det
om Sokrates gjennem de forskjellige Kilder Overleverede, navn\-%
lig ved første Øiekast, kun ringe indbyrdes Overensstemmelse,
men der er, i de fra hinanden afvigende, trivialiserende og
idealiserende Opfattelser af Sokrates visse Punkter, i hvilke Op\-%
fattelserne berøre hverandre, og ud fra hvilke det bliver muligt,
at bringe den virkelige historiske Sokrates's Skikkelse frem, at
klare hans Standpunkt og at finde Sammenhængen i den
ethiske Anskuelse, der svarer til det, og som for en nærmere
Betragtning viser sig i sin eiendommelig græske Charakteer.

Der berettes om Sokrates, at han, idet han ligesom førte
Philosophien fra Himlen ned til Jorden, d. v. s. fra Natur\-%
speculationen førte den tilbage til Selverkjendelsen, \emph{udeluk\-%
kende beskæftigede sig med det Ethiske}; --- der be\-%
rettes, at han \emph{bestemte Dyden som Viden}; --- der beret\-%
tes, at han \emph{var den Første, der gjennem Inductionen
bragte de almindelige Begreber frem}; --- og der be\-%
rettes endelig, at han \emph{erklærede sig for uvidende}.

Disse Udsagn ere ligesom Summen af det, der af de for\-%
skjellige Beretningers Berøringer foreløbig fremgaaer som Holde\-%
punkt for Opfattelsen; men charakteristisk er det, at det, hvori
de forskjellige Beretninger væsentlig enes, umiddelbart taget,
er et i sig Uoverensstemmende, et Modsigende. Modsigelsen
viser sig strax naar man sammenstiller Udsagnene; men netop
gjennem Løsningen af Modsigelsen naae vi til Opfattelsen af
Sokrates og til Forstaaelsen af Misforstaaelserne.

Sammenstille vi nemlig først de to Udsagn om Sokrates's
\emph{Inductionsmethode} og om hans \emph{udelukkende Beskæf\-%
tigelse med det Ethiske}, saa kommer der strax en Mod\-%
sigelse ind; thi naar Sokrates, som \emph{Xenophon} udtrykker det,
bestandig med sine Omværende undersøgte \gk{τὶ ἕκαστον εἴη τῶν
% p. 204, GREEK: three compound breathing-plus-accent marks in one
% phrase, all read at 2400 dpi against the p. 198 calibration pair.
% ἕκαστον -- breathing element a "C" opening RIGHT: rough.
% εἴη and ὄντων -- breathing element a "ɔ" opening LEFT: smooth.
% τὶ carries a grave (falling right), τῶν a circumflex.
ὄντων}, saa bliver Interessen en \emph{logisk}, ikke en ethisk. Men
ved nærmere Undersøgelse af, hvorledes det forholder sig med
Resultatet af Sokrates's Søgen af Tingenes Begreb og med
% --- p. 205 ---
% PENCIL, p. 205: several vertical strokes in the RIGHT margin, and a
% cursive marginal word beside "hvis han var naaet til et positivt
% Resultat".  Not transcribed.
hans Definitioner, viser det sig, at hans inductive Methode ---
bestemtere: hans begrebsdannende Abstractions Methode ---
enten \emph{ligefrem} fører til et \emph{negativt} Resultat, idet de for\-%
søgte almene Bestemmelser af Sagen atter omstødes ved mod\-%
sigende Tilfælde, eller at den \emph{kun tilsyneladende} fører til
et \emph{positivt}, nemlig til et saadant, der allerede er viist som
uholdbart ved de forudgaaende dialektiske Udviklinger (f. Ex.
hvor det Gode bestemmes som det Nyttige; jfr. Undersøgelsen
i Platos Laches). Og netop dette synes at antyde, at Sokra\-%
tes's Induction og Definition, overhovedet hans logiske Tanke\-%
virksomhed, maa opfattes, \emph{ikke} som et \emph{Øiemed}, men som
Noget, der paa en eller anden Maade tjener til \emph{Redskab for
et Andet}, og dette Andet vilde da naturligst være at søge i
det \emph{Ethiske} som det, paa hvilket Sokrates's Interesse jo be\-%
standig siges, væsentligen at have hvilet.

Det, at det definitive Tankeresultat af Sokrates's Induc\-%
tioner bestandig var \emph{negativt}, og at derfor det sidste Maal
for den Tanke, der uafladelig vendte tilbage til den samme
Bevægelse, ikke kunde være det logiske, et Videns Maal, finder
sin Bestyrkelse i det Berettede, at Sokrates erklærede sig for
\emph{uvidende}. Thi naar han, som Traditionen udsiger, satte
\emph{Begrebet} (\gk{τί ἐστι}) som det Væsentliges sande Udtryk, og i
% τί acute (rising right); ἐστι smooth, unaccented -- both at 2400 dpi.
Definitionen vilde give Begrebets Bestemmelse, saa kunde han,
hvis han var naaet til et positivt Resultat, om han end nok
saa meget «saae Begrebet i dets Relation til det subjective Liv»,
ikke uden Selvmodsigelse kalde sig uvidende; thi Definitionens
Indhold vilde jo netop være en virkelig Viden.

Men naar vi nu ville fastholde den Antydning, at Sokra\-%
tes's \emph{logiske} Interesse tjente den \emph{ethiske}, saa kommer der
atter en Vanskelighed frem, idet han, den \emph{Uvidende}, be\-%
stemte \emph{Dyden som Viden}. At Dyden for Sokrates maatte
være en Virkelighed, vil Ingen nægte; thi det Modsatte vilde
ikke blot forudsætte et helt andet Standpunkt end det græske,
men stride mod alle Beretninger om hans personlige Liv. Men
var \emph{Dyden} en Virkelighed, maatte ogsaa dens synonyme Be\-%
% --- p. 206 ---
stemmelse: \emph{Viden}, være det. Viden maa være en Realitet;
men det bliver den kun \emph{som Viden om Noget}; og hvorfra
skulde Sokrates's Viden, der er = Uvidenhed, faae dette No\-%
% The "=" between "er" and "Uvidenhed" is printed type (two short
% parallel strokes), not an editorial addition; set as text, not math.
get? \emph{Hvorom} er Dyden en Viden? Naar der svares: om
\emph{det Gode}, men det Gode som det \emph{ethisk} Gode saa atter
kun kan bestemmes ved Viden, saa bevæge vi os kun i en
Cirkel, som vi ikke kunne komme ud af. Men Sokrates maatte,
hvis han var Ethiker, ud deraf; thi forat der kan naaes til
Handling, maa Viden have et Indhold, dens Gjenstand maa
have Bestemthed: Form, som en Græker vilde sige. De hi\-%
storiske Efterretninger lade nu Sokrates bestemme det Godes
Indhold som \emph{det med Lovene Overensstemmende} (lige\-%
som han fordrer, at man skal dyrke Guderne \gk{νόμῳ πόλεως}),
% νόμῳ carries an iota subscript under the ω, seen at 2400 dpi.
og hos \emph{Xenophon} fremstilles han som bestemmende det Gode
ved \emph{det Nyttige}, altsaa eudaimonistisk. Men Dialektiken i
\emph{det Nyttige} gjør sig strax gjældende: Begrebet viser sin Re\-%
lativitet, og Spørgsmaalet: hvad er det Nyttige? kan ikke be\-%
svares med Sikkerhed i \emph{positiv} Form af den, hvis Viden er
Uvidenhed, og en Besvarelse i denne Form er nødvendig naar
der skal handles, og Livet ikke gaae op i en Askese. Skal
der blive noget Sikkert tilbage i Bestemmelsen, maa der ial\-%
fald fordres en Mellembestemmelse. Det Samme er Tilfældet
hvor det Gode bestemmes som \emph{det med Lovene Overens\-%
stemmende}; thi Modsætningen mellem Lovene indbyrdes og
mellem Lovene og en Viden, der, i negativ Form i Menne\-%
sket, i positiv i det Guddommelige, var høiere end Lovene,
var Sokrates klart bevidst. Over Staternes modsigende Love
staae \gk{νόμοι ἄγραφοι}, men hvem skal læse deres uskrevne
% ἄγραφοι: compound over the α, breathing element opening LEFT (smooth),
% accent rising right, read as at p. 202.
Skrift?

De Vanskeligheder og Modsigelser, der have viist sig, ville
neppe ad anden Vei kunne fjernes end ved at opfatte Sokrates
\emph{som ethisk-religiøs Personlighed}, og ved i det \emph{Ethisk-%
% The hyphen of "Ethisk-Religiøse" falls at the compositor's line break
% and is a REAL hyphen, not a word-break hyphen: kept, with % eating the
% newline.
Religiøse} at see den Gjenstand, til hvilken hans Interesse
\emph{med Bevidsthed} begrændsedes, og som hans Dialektik,
der maa opfattes som \emph{rent negativ}, tjente.
% --- p. 207 ---
% PENCIL, p. 207 (a known pollution site): vertical strokes in the RIGHT
% margin in both halves of the page; a small blob between "thi" and "vi"
% in "thi vi have seet" (both witnesses read it as an apostrophe,
% "thi'vi"); a thin diagonal stroke through the comma after "paa dette
% Standpunkt," (ABBYY reads "Standpunkt^"); and a second such stroke
% through the comma after "til Negativitetens Uendelighed," (ABBYY drops
% the comma).  Both commas ARE printed -- confirmed at 1800 dpi.  None of
% the pencil is transcribed.
Det blev ovenfor berørt, at den logiske Interesse ikke var
det hos ham, der var det i sidste Instants Tilskyndende i hans
Dialektik. Uagtet han var den Første, der bragte den Me\-%
thode frem, som man har kaldt Inductionens, saa var hans
Interesse ikke Methoden som saadan, saaledes som f. Ex.
\emph{Plato} og endnu mere \emph{Aristoteles} gjøre den til Gjenstand
for Undersøgelsen; ikke heller var det et positivt Resultat af
Methoden: Begrebsbestemmelsen; thi vi have seet, at et saa\-%
dant Resultat var illusorisk, og Sokrates var sig bevidst, at
det \emph{maatte} være det. Men opfatte vi hans Dialektik i dens
Eiendommelighed som et blot \emph{tjenende} Middel, saa bliver
den ligesom den Magt, der gjennem \emph{alt} det Endeliges For\-%
virring bringer \emph{Uendeligheden} frem, og denne Uendelighed,
som Sokrates vilde bringe tilsyne, var \emph{Selvbevidsthedens,
foreløbig endnu negative Uendelighed, Udgangs\-%
punktet for det Ethiske}, for Selvbestemmelsen i Kraft af
Frihed og i Kraft af Tanke. Netop gjennem \emph{Uvidenhedens
Bevidsthed} skal Individet føres til det \emph{Ethiske}, og Dia\-%
lektiken, i hvilken Subjectet gjør sin Uendelighed gjældende i
Negativitetens Form, holder bestandig hiin Bevidsthed vaagen,
samtidig med, at den holder Uendelighedsenergien levende i
Selvet. Men naar \emph{den ironisk opløsende Dialektik} kunde
vise Veien tilbage til hint det Ethiskes Udgangspunkt, saa for\-%
dredes der, paa dette Standpunkt, for atter at komme ud fra
dette usynlige Punkt, hvor Selvbevidstheden slutter sig sammen
med sig selv idet den befrier sig fra al Endelighed, idet den
bruger sin negative Uendelighed til denne Befrielse, og for
atter at give denne Selvbevidsthed, der endnu kun var naaet
til Negativitetens Uendelighed, et Videns Indhold, et nyt Mo\-%
ment: det \emph{religiøse} Forhold. Sokrates kunde ikke ad Tan\-%
kens Vei, ved Tanken selv, give Selvbevidsthedens Uendelig\-%
hed Form. Han maatte søge Uendelighedens Udfyldelse gjen\-%
nem et Andet end sin subjective Tanke. Dette Andet kunde
kun være den \emph{guddommelige} Tanke. Denne er i Enhed
med den menneskelige, er dennes Væsen, men i denne Enhed
% --- p. 208 ---
\setcounter{footnote}{0}%
% p. 208 is one of the two pages in the whole book with a SECOND
% EXPOSURE (pagemap.py --alt 208 = PDF 219); both were consulted.
% RUN-OVER FOOTNOTE: note **) below begins here and CONTINUES ON P. 209,
% where the note area opens with "rethed. I Bestemmelsen af det Nyttige
% mødes saaledes..." and carries NO mark.  The whole note is set here, in
% one \footnote{} at its call, per transcription.tex §4; a % comment at
% the p. 209 marker records where the print breaks it.
bliver den dog, skjøndt et Ideelt i Subjectet, som en \emph{An\-%
skuelsens}\footnote{Vi minde atter her om Sokrates's ekstatiske
Bevidstløshed, der er Formen for hans Fortaben af sig i en Anskuen af
det Uendelige, i hvilken den endelige Viden negeres.} Gjenstand, et fra Subjectets Tanke og Person\-%
lighed Sondret og taler som et guddommeligt \emph{Orakel i Sub\-%
jectet}. Og idet nu \emph{Sokrates}, gjennem det religiøse Forhold
til det Guddommelige, kan anskue \emph{det Givne} som det, hvori
der er en \emph{guddommelig Ordning} (jfr. hans teleologiske Opfat\-%
telse), saa har han, naar Oraklet i hans Indre taler i Uendelig\-%
hedens \emph{negative} Form, i dets negative, \emph{advarende} Tilsagn et
\emph{Correctiv} i Forholdet til dette Givne. Naar ikke dette reli\-%
giøse Forhold bliver \emph{Mellembestemmelse}, er det uforstaae\-%
ligt, at Sokrates, trods sin opløsende Dialektik, atter som
Norm for den personlige Handlen kan sætte \emph{det Lovmæs\-%
sige}, der i sin Enkelthed allerede var blevet gjort vaklende
ved Dialektiken, og det \emph{Nyttige}, der var underkastet en for\-%
tærende Skepsis. Kun naar Forholdet til det Guddommelige
drages med ind, kan Loven, som et guddommeligt Anordnet,
atter faae Fasthed, og det Nyttige faae en Sikkerhed midt i
Usikkerheden, blive en Gjenstand for en positiv Viden midt i
Uvidenheden\footnote{Under Bestemmelsen af det \emph{Nyttige} kan
\emph{Sokrates}, naar Begrebet skal have et \emph{positivt} Indhold,
sammenfatte de Forhold, som den almindelige \emph{ethiske} Bevidsthed
sanctionerer, og som han kan opfatte under samme Synspunkt som det
Lovmæssige, nemlig som \emph{anordnede ved Guddommen}. Fra et
\emph{andet} Synspunkt gaaer den \emph{ethiske Selverkjendelse} som
Forudsætning for den ethiske Handling, ind derunder, ligeledes
\emph{Autarkien}, den ethiske Friheds \emph{negative} Udtryk, og hvad
Sokrates opfatter som Midler til dens Erhvervelse (Afholdenhed, Omsorg
for Legemets Sundhed, maaskee ogsaa Erhvervelse af Kundskaber; jfr.
Xenophons Memorab.). Men saavel hiin Selverkjendelse som denne Autarkie,
i dens væsentlige Forskjellighed fra det Asketiske, vise hen til et
positivt Handlingens Indhold; medens, omvendt, de ethiske Forhold igjen
kunne sees fra den Side, at de, som f. Ex. Venskabet, efter Sokrates's
Opfattelse, blive Midler til ethisk Selverkjendelse, eller fremme Livets
Uforstyr\-%
rethed. I Bestemmelsen af det Nyttige mødes saaledes den negative og
den positive Side i den sokratiske Ethik: den Frigjørelse, der udtrykkes
gjennem Autarkien, og den concrete positive Bestemmen af Handlingen
gjennem Forholdet til det Guddommelige. Som Udtryk for det Negative, for
\emph{Autarkien}, erkjendes det Nyttige \emph{ved Viden}, ligesom
\emph{Uvidenheden} vides, og ligesom \emph{Grundforholdet} --- i dets
\emph{negative} Bestemthed --- vides.}. Og til det \emph{Guddommelige} viser Sokrates
% --- p. 209 ---
% p. 209 also has a SECOND EXPOSURE (pagemap.py --alt 209 = PDF 220);
% every doubtful reading below was checked on both.
% In the print the foot of this page carries the CONTINUATION of the
% note **) called on p. 208 -- it opens "rethed. I Bestemmelsen af det
% Nyttige..." with no mark of its own.  It is set in full inside that
% \footnote{} above; nothing is dropped and no new note is opened here.
% SHEET SIGNATURE "14" at the foot: not transcribed.
% PENCIL, p. 209: a tick and two vertical bracket strokes in the RIGHT
% margin, beside three passages in the lower half.  Not transcribed.
% LETTERSPACING: this is the most heavily spaced page in the book.
% ELEVEN spaced phrases in the body, not twelve -- the recon's
% "a l m e n e  V æ s e n" is wrong.  Checked at 1800 dpi on BOTH
% exposures: "almene" is spaced, "Væsen," is plainly normal set.
% The footnote adds five more: Autarkien, ved Viden, Uvidenheden,
% Grundforholdet, negative.  spacing.py recovers only nine of the body's
% eleven; the page was read by eye, not from its output.
overalt tilbage, naar han taler om sin \emph{Livsopgave} og om
hvad der definitivt bestemmer hans Gjerning.

Dette sokratiske Standpunkts Væsen og Betydning bliver
klarere ved en Sammenstilling med dets \emph{historiske Forud\-%
sætninger}. Hos \emph{Sophisterne} havde Subjectiviteten gjort
sin Negativitet gjældende gjennem den opløsende Dialektik;
men for dem var endnu \emph{den tilfældige Individualitet}
bleven det Afgjørende: Tingenes Maal; deres Dialektik vendte
sig opløsende mod Objectiviteten, men ikke uendeliggjørende
mod Tanken og Subjectet selv. Sophisterne vare ikke for
meget Dialektikere, men for lidet: de standsede ved det En\-%
delige, ved et illusorisk Positivt, og bragte ikke den Uendelig\-%
hed tilsyne, ud fra hvilken et Almeengjældende, i Sandhed
Positivt, skulde begrundes. Netop dette var \emph{Sokrates's} Op\-%
% A short low mark stands between "Netop" and "dette" -- present in BOTH
% exposures at 1800 dpi, so it is on the leaf and not a scanning
% artefact, but it sits near the baseline, well below hyphen height, and
% there is a full word space on either side of it.  Read as a printer's
% speck (a quad or bearer rising to ink), not as a hyphen; not
% transcribed.  Both witnesses splice it in as "Netop-dette".
gave. Baaren af \emph{Ironiens} uendelige Negativitet forstyrrer hans
Dialektik \emph{alt Endeligt}, al endelig Positivitet, ogsaa den af
selve Individet \emph{satte}, og aabenbarer i denne gjennemgaaende
Opløsen \emph{Selvbevidsthedens indre Uendelighed}. Dia\-%
lektikens Maal er for ham ikke at nedbryde og opløse: gjen\-%
nem Opløsningen af de endelige Bestemmelser, gjennem For\-%
virringen af Meningens, den indbildte Videns Indhold, lader
han Tanken søge et Fast og Blivende, det \emph{almene} Væsen,
udtrykt i \emph{Begrebet}: det \emph{formede Uendelige}. Men naar
nu atter Begrebet flygter for Tanken, der kun besidder Nega\-%
tivitetens Uendelighed uden at den ud fra denne ved sig selv
kan vinde et nyt Indhold: et Begrebsindhold, saa tvinges den
Tænkende, der ikke kan komme til Ro i det Endelige og Til\-%
% --- p. 210 ---
% THE PAGE MAP STEPS HERE: printed 210 = PDF 221, not 219.  Obtained
% from pagemap.py, as every page in this batch was.
% The ɔ: id-est mark stands in line 23 of this page, confirmed at 1200
% dpi (reversed c opening LEFT, then a colon).  No note on this page.
% PENCIL, p. 210: a vertical stroke and a small mark in the LEFT margin
% beside "bestemmer Dyd som Viden".  Not transcribed.
fældige, over i Erkjendelsen af sin \emph{Uvidenhed}, og med denne
Erkjendelse atter over i det \emph{Ethisk-Religiøse}, hvor Subjec\-%
tet, med Uvidenhedens Bevidsthed til Forudsætning, finder et
\emph{nyt} Indhold for sin Viden, betrygget ved Tilliden til det Gud\-%
dommelige og ved det Correctiv, som den \emph{dæmoniske}
Stemme i det Indre anbringer, og som, \emph{psychologisk} seet,
er den i sit Udspring ubevidste Virkning af den gjennem Iro\-%
niens negative Uendelighed frigjorte ethiske Følelses uvilkaar\-%
lige Reaction mod den enkelte usande Handling.

Men at nu Sokrates, skjøndt Viden for ham ender i
\emph{Uvidenhed}, dog kan sætte en \emph{Viden om Tingenes Be\-%
greb som et Maal for Tanken}, og bestemme \emph{Dyden som
Viden}, bliver forklarligt ved en \emph{anden} af hans historiske
Forudsætninger og ved et Henblik til den græske Grundan\-%
skuelse.

Ved den Sokratiske Bestemmelse af Dyden som Viden
mindes vi nærmest om \emph{Anaxagoras's} Princip, der vistnok
har havt Betydning for Sokrates. For \emph{Anaxagoras} var
\gk{Νοῦς}, Tanken, den høieste Virkelighedsmagt; for \emph{Sokrates} er
ligeledes, skjøndt denne Tanke ikke udtales directe eller bliver
Udgangspunktet for en dogmatisk Lære, den \emph{guddommelige
Tanke} (Fornuft), der i \emph{Menneskeaanden væsentlig
er som Viden} ɔ: som en Erkjenden af det Blivende og
% "(Fornuft), der i" is NOT letterspaced -- checked at 350 dpi; the
% spacing resumes at "Menneskeaanden".
Almene i Tingen, Principet, og denne Opfattelse danner
ligesom den gjennemskinnende Baggrund for hans hele An\-%
skuelse. Dette bliver saaledes at begrunde. Idet Sokrates
bestemmer \emph{Dyd som Viden}, saa ligger heri først, idet Dyd,
efter den græske Opfattelse, er Væsenets Virken efter dets
\emph{eiendommelige}, uforstyrrede Natur, at Viden er \emph{Menne\-%
skets Væsen}. Men den \emph{vidende Tanke}, der er Menne\-%
skets Væsen, bliver tillige det \emph{Objectives Væsen og det i
Tilværelsen Herskende}. Dette viser sig deels derved, at
Tanken hos Sokrates ikke er en fra Objectiviteten løsreven Tanke,
men en Tanke, der sætter det Objective som Videns Gjenstand for i
det af Tanken formede Begreb at udtrykke dets \emph{Væsen}; deels

% [text to be added: pp. 211--224]

% [text to be added: pp. 225--238]

% [text to be added: pp. 239--252]

% [text to be added: pp. 253--266]

% [text to be added: pp. 267--280]

% [text to be added: pp. 281--288]

% =====================================================================
%  REGISTER, printed pp. 289--294 = PDF 300--305.
%  Set here as a single-column list.  In the print it is TWO columns
%  separated by a thin vertical rule, each column carrying its own small
%  bold flush-right head "Side."; leader dots run from a headword to a
%  right-aligned figure when the entry has one reference and are absent
%  when it has several; each alphabet group opens with an oversized
%  two-line initial capital.  All of that is page-making furniture and is
%  not reproduced.  What IS kept: the bold page figures, which mark the
%  principal treatment of a topic; the — separators between run-on
%  sub-entries; and f., ff., Anm. (= the reference is to a footnote on
%  that page) and S. exactly as printed.
% =====================================================================

% [text to be added: pp. 289--294]

% =====================================================================
%  TRYKFEIL OG RETTELSER, the unnumbered leaf at PDF 306.
%  The author's own errata.  NOT applied to the body: every one of the
%  nine printed readings stands as printed, with a % comment at its site.
% =====================================================================

% [text to be added: pp. 295--295]

\end{document}
